\documentclass[12pt]{article}
\usepackage[ukrainian]{babel}
\usepackage{fontspec}
% --- НАЛАШТУВАННЯ ШРИФТІВ ---
\setmainfont{Schoolbook-Regular.otf}[
    Path = ../,
    BoldFont = Schoolbook-Bold.otf,
    ItalicFont = Schoolbook-Italic.otf,
    BoldItalicFont = Schoolbook-BoldItalic.otf
]
\usepackage[italic]{mathastext}
\usepackage[a4paper,margin=1.8cm,bottom=2cm,top=2cm]{geometry}
\usepackage{amsmath,amssymb}
\usepackage{tikz}
\usepackage{pgfplots}
\pgfplotsset{compat=1.16}
\usetikzlibrary{calc,patterns,angles,quotes}
\usepackage{xcolor,array,fancyhdr,enumitem,multicol}
\usepackage{colortbl}
\usepackage{tcolorbox}
\tcbuselibrary{skins}
\usepackage[hidelinks]{hyperref}
\usepackage[firstpage=false, text=@pvtr2525, scale=0.6, color=gray!12]{draftwatermark}
\definecolor{mainGreen}{RGB}{34, 120, 64}
\definecolor{yearOrange}{RGB}{220, 100, 30}
\definecolor{yearcolor}{RGB}{220, 100, 30}
\definecolor{titleBg}{RGB}{225, 240, 220}
\definecolor{instrBg}{RGB}{255, 240, 220}
\definecolor{instrBorder}{RGB}{220, 150, 80}
\pagestyle{fancy}
\fancyhf{}
\renewcommand{\headrulewidth}{0.4pt}
\renewcommand{\footrulewidth}{0.4pt}
\fancyhead[C]{\small\color{gray!80} Збірник НМТ \textendash{} сесії 23.05--5.06.2026}
\fancyhead[R]{\small\color{mainGreen}\textbf{\href{https://t.me/pvtr2525}{@pvtr2525}}}
\fancyfoot[L]{\small\color{gray!80}\href{https://t.me/pvtr2525}{ТГ: @pvtr2525}}
\fancyfoot[C]{\small\color{gray!80}\thepage}
\fancyfoot[R]{\small\color{gray!80}НМТ 2026}
% --- ТАБЛИЦІ ВІДПОВІДЕЙ ---
\newcommand{\answerTable}[5]{%
\par\vspace{0.15cm}
\noindent
\begin{tabular}{|*{5}{>{\centering\arraybackslash}m{3cm}|}}
\hline
\rule[-0.15cm]{0pt}{0.6cm}\textbf{А} & \rule[-0.15cm]{0pt}{0.6cm}\textbf{Б} & \rule[-0.15cm]{0pt}{0.6cm}\textbf{В} & \rule[-0.15cm]{0pt}{0.6cm}\textbf{Г} & \rule[-0.15cm]{0pt}{0.6cm}\textbf{Д} \\
\hline
\rule[-0.2cm]{0pt}{0.75cm}#1 & \rule[-0.2cm]{0pt}{0.75cm}#2 & \rule[-0.2cm]{0pt}{0.75cm}#3 & \rule[-0.2cm]{0pt}{0.75cm}#4 & \rule[-0.2cm]{0pt}{0.75cm}#5 \\
\hline
\end{tabular}
\par\vspace{0.25cm}
}
\newcommand{\answerTableTall}[5]{%
\par\vspace{0.15cm}
\noindent
\begin{tabular}{|*{5}{>{\centering\arraybackslash}m{3cm}|}}
\hline
\rule[-0.2cm]{0pt}{0.7cm}\textbf{А} & \rule[-0.2cm]{0pt}{0.7cm}\textbf{Б} & \rule[-0.2cm]{0pt}{0.7cm}\textbf{В} & \rule[-0.2cm]{0pt}{0.7cm}\textbf{Г} & \rule[-0.2cm]{0pt}{0.7cm}\textbf{Д} \\
\hline
\rule[-0.45cm]{0pt}{1.2cm}#1 & \rule[-0.45cm]{0pt}{1.2cm}#2 & \rule[-0.45cm]{0pt}{1.2cm}#3 & \rule[-0.45cm]{0pt}{1.2cm}#4 & \rule[-0.45cm]{0pt}{1.2cm}#5 \\
\hline
\end{tabular}
\par\vspace{0.25cm}
}
\newcommand{\instructionBox}[1]{%
\par\vspace{0.3cm}
\begin{tcolorbox}[colback=instrBg, colframe=instrBorder, boxrule=0.6pt, arc=2pt,
    left=8pt, right=8pt, top=4pt, bottom=4pt]
\centering\small\bfseries #1
\end{tcolorbox}
\par\vspace{0.15cm}
}
\newcommand{\sectionTitle}[1]{%
\par\vspace{0.4cm}
\begin{tcolorbox}[colback=titleBg, colframe=titleBg, boxrule=0pt, arc=8pt,
    halign=center, fontupper=\Large\bfseries\color{mainGreen}]
#1
\end{tcolorbox}
\par\vspace{0.3cm}
}
\newcommand{\task}[2]{\noindent\textbf{#1.}\ \ #2\par\vspace{0.2cm}}
% --- НМТ-СТИЛЬ ПОЛЕ ВІДПОВІДІ ---
\newcommand{\nmtAnswerBox}{%
\par\vspace{0.2cm}\noindent
Відповідь:\ \
\begin{tabular}{@{}*{4}{|>{\centering\arraybackslash}p{0.8cm}}|@{\hspace{0.1cm}\textbf{,}\hspace{0.1cm}}*{3}{|>{\centering\arraybackslash}p{0.8cm}}|@{}}
\hline
\rule[-0.2cm]{0pt}{0.7cm} & & & & & & \\
\hline
\end{tabular}
\par\vspace{0.3cm}
}
% --- СІТКА ДЛЯ МАТЧИНГУ ---
\newcommand{\matchingGrid}{%
    \begingroup
    \renewcommand{\arraystretch}{1.15}
    \setlength{\tabcolsep}{4pt}
    \footnotesize
    \begin{tabular}{r|c|c|c|c|c|}
         \multicolumn{1}{c}{} & \multicolumn{1}{c}{\textbf{А}} & \multicolumn{1}{c}{\textbf{Б}} & \multicolumn{1}{c}{\textbf{В}} & \multicolumn{1}{c}{\textbf{Г}} & \multicolumn{1}{c}{\textbf{Д}} \\ \cline{2-6}
         \textbf{1} & \phantom{X} & \phantom{X} & \phantom{X} & \phantom{X} & \phantom{X} \\ \cline{2-6}
         \textbf{2} & & & & & \\ \cline{2-6}
         \textbf{3} & & & & & \\ \cline{2-6}
    \end{tabular}
    \endgroup
}
% --- ДОДАТКОВІ МАКРОСИ ЗБІРНИКА ---
\newcommand{\taskBlock}[1]{%
\par\vspace{0.45cm}
\noindent{\color{mainGreen}\rule{\linewidth}{0.8pt}}\par\vspace{0.12cm}
\noindent{\small\bfseries\color{mainGreen}#1}\par\vspace{0.25cm}
}
\newcounter{zad}
\newcommand{\zadtask}[1]{\stepcounter{zad}\noindent\textbf{\thezad.}\ \ #1\par\vspace{0.2cm}}
\newcommand{\zadnum}{\stepcounter{zad}\textbf{\thezad.}\ \ }
\newcounter{ans}
\newcommand{\ansitem}[1]{\stepcounter{ans}\noindent\textbf{\theans.}\ #1\par\vspace{0.07cm}}
\newcommand{\nmtyear}[1]{\hfill{\small\color{yearOrange}(НМТ #1)}}
% --- БІБЛІОТЕКА ДЛЯ РОЗРИВНИХ БОКСІВ ---
\tcbuselibrary{breakable}
% --- МАКРОС КОРОТКОГО РОЗВ'ЯЗКУ ОРИГІНАЛУ ---
\newcommand{\solution}[1]{%
\par\vspace{0.12cm}
\begin{tcolorbox}[colback=titleBg!45, colframe=mainGreen!55, boxrule=0.4pt, arc=2pt,
    left=6pt, right=6pt, top=3pt, bottom=3pt, breakable]
\footnotesize\textbf{Короткий розв'язок (оригінал).}\ #1
\end{tcolorbox}
\par\vspace{0.12cm}
}
% --- ЗАГОЛОВКИ З ПУНКТАМИ КЛІКАБЕЛЬНОГО ЗМІСТУ ---
\setcounter{tocdepth}{2}
\newcommand{\chapterTitle}[1]{%
\clearpage\phantomsection\addcontentsline{toc}{section}{#1}%
\sectionTitle{#1}}
\newcommand{\typeTitle}[1]{%
\phantomsection\addcontentsline{toc}{subsection}{\quad #1}%
\par\vspace{0.3cm}
\noindent{\large\bfseries\color{yearOrange}#1}\par\vspace{0.08cm}
\noindent{\color{yearOrange}\rule{\linewidth}{0.4pt}}\par\vspace{0.2cm}}
\newcommand{\ansTheme}[1]{\par\vspace{0.25cm}\noindent{\bfseries\color{mainGreen}#1}\par\vspace{0.12cm}}
\newcommand{\ansType}[1]{\par\vspace{0.1cm}\noindent{\itshape\color{yearOrange}#1}\par\vspace{0.08cm}}

\begin{document}

% ===================== ТИТУЛ =====================
\thispagestyle{empty}
\vspace*{1.5cm}
\begin{center}
{\Huge\bfseries\color{mainGreen} ЗБІРНИК ЗАВДАНЬ НМТ-2026}\\[0.5cm]
{\Large\bfseries з математики}\\[0.8cm]
{\large Оригінали та авторські аналоги за сесіями}\\[0.2cm]
{\large \textbf{23.05 -- 5.06.2026}}\\[1.2cm]
\begin{tcolorbox}[colback=titleBg, colframe=mainGreen!60, boxrule=0.8pt, arc=6pt, width=0.8\textwidth, halign=center]
\large Згруповано за \textbf{6 розділами} та \textbf{типами завдань}.\\
До кожного оригіналу --- короткий розв'язок. Відповіді в кінці.
\end{tcolorbox}
\vfill
{\large\color{mainGreen}\textbf{\href{https://t.me/pvtr2525}{Телеграм: @pvtr2525}}}\\[0.3cm]
{\small\color{gray!70} НМТ 2026}
\end{center}
\clearpage

% ===================== ПОЯСНЮВАЛЬНА ЗАПИСКА =====================
\phantomsection\addcontentsline{toc}{section}{Пояснювальна записка}
\sectionTitle{ПОЯСНЮВАЛЬНА ЗАПИСКА}

Цей збірник укладено для підготовки до національного мультипредметного тесту (НМТ) з математики 2026 року. Матеріал охоплює завдання, що пропонувались учасникам на сесіях основного періоду з 23 травня до 5 червня 2026 року, разом із авторськими аналогами, складеними для відпрацювання кожного типу задач.

Збірник побудовано за \textbf{шістьма тематичними розділами} відповідно до чинної програми: числа і вирази; рівняння і нерівності; функції (їхні властивості, графіки, прогресії, похідна та первісна); планіметрія; стереометрія; комбінаторика, теорія ймовірностей і статистика. Усередині кожного розділу завдання впорядковано за \textbf{типом}: спершу завдання з вибором однієї правильної відповіді (А--Д), далі завдання на встановлення відповідності, наприкінці --- завдання з короткою числовою відповіддю.

Матеріал згруповано в \textbf{блоки}. Кожен блок починається з \textit{оригінального} завдання сесії, після якого вміщено \textbf{короткий розв'язок} цього оригіналу. Далі йдуть \textbf{два авторські аналоги}: перший повторює тип і спосіб міркування оригіналу (для закріплення прийому), другий є \textit{варіацією} --- зі зміненим знаком, формулою або формулюванням, щоб та сама ідея перевірялася під іншим кутом. Розв'язки до аналогів свідомо не наводяться: їх призначено для самостійної роботи.

Усі квадратні рівняння в розв'язках доведено через \textbf{дискримінант} або \textbf{теорему Вієта}. Структура кожного завдання відповідає реальному формату НМT: 15 завдань з вибором відповіді, 3 на встановлення відповідності та 4 з короткою відповіддю на повний варіант.

\textbf{Як користуватися.} Скористайтеся клікабельним змістом, щоб перейти до потрібного розділу чи типу завдань. Спочатку спробуйте розв'язати оригінал самостійно, потім звіртеся з коротким розв'язком; після цього відпрацюйте обидва аналоги. Правильні відповіді до всіх завдань (оригіналів і аналогів) зібрано в розділі \textbf{«Відповіді»} наприкінці збірника.

\textit{Зауваження.} У кількох випадках виявлено та виправлено помилки у відкритих ключах сесій; такі місця позначено коментарями у вихідному коді. Якщо помітите неточність --- напишіть на \href{https://t.me/pvtr2525}{@pvtr2525}.

% ===================== ЗМІСТ =====================
\clearpage
\phantomsection\addcontentsline{toc}{section}{Зміст}
\sectionTitle{ЗМІСТ}
{\hypersetup{linkcolor=black}\renewcommand{\contentsname}{}\vspace{-1cm}\tableofcontents}

\chapterTitle{РОЗДІЛ 1. ЧИСЛА І ВИРАЗИ}
\typeTitle{Завдання з вибором однієї відповіді (А--Д)}
\taskBlock{Спрощення раціонального виразу --- сесія 23.05, завдання №7}
\zadtask{Спростіть вираз $\dfrac{(2x - 3)^2 - 9}{x}$. \nmtyear{2026}}
\answerTable{$4x - 12$}{$4x - 6$}{$4x$}{$2x - 6$}{$2x - 12$}
\solution{За формулою $(a-b)^2$: $(2x-3)^2-9=4x^2-12x+9-9=4x^2-12x=4x(x-3)$. Поділивши на $x$, дістаємо $4x-12$. Відповідь: \textbf{А}.}
\zadtask{Спростіть вираз $\dfrac{(3x - 2)^2 - 4}{x}$.}
\answerTable{$9x - 4$}{$9x - 12$}{$3x - 4$}{$9x$}{$3x - 12$}
\zadtask{Спростіть вираз $\dfrac{9 - (3 - 2x)^2}{x}$.}
\answerTable{$4x - 12$}{$12 - 4x$}{$12 + 4x$}{$-4x$}{$12 - 2x$}
\taskBlock{Обчислення логарифма --- сесія 23.05, завдання №9}
\zadtask{Обчисліть $\log_{0{,}5} 8$. \nmtyear{2026}}
\answerTable{$-2$}{$-4$}{$-3$}{$\dfrac{16}{3}$}{$3$}
\solution{Оскільки $0{,}5=2^{-1}$, а $8=2^3$, то $\log_{0{,}5}8=\dfrac{\log_2 8}{\log_2 0{,}5}=\dfrac{3}{-1}=-3$. Відповідь: \textbf{В}.}
\zadtask{Обчисліть $\log_{0{,}2} 25$.}
\answerTable{$-2$}{$2$}{$-5$}{$5$}{$\dfrac{1}{2}$}
\zadtask{Обчисліть $\log_{3} \dfrac{1}{81}$.}
\answerTable{$-4$}{$4$}{$-3$}{$\dfrac{1}{4}$}{$3$}
\taskBlock{Ділення з остачею (порції) --- сесія 30.05, завдання №5}
\zadtask{Для приготування однієї порції коктейлю потрібно $40$ г морозива. Яку \textit{найбільшу} кількість таких порцій коктейлю можна приготувати з $0{,}5$ кг морозива?}
\answerTable{$10$}{$11$}{$12$}{$13$}{$20$}
\solution{$0{,}5$ кг $=500$ г. Кількість порцій: $500:40=12{,}5$, тобто ціла частина $12$. Відповідь: \textbf{В} ($12$).}
\zadtask{Для приготування однієї порції смузі потрібно $60$ г ягід. Яку \textit{найбільшу} кількість таких порцій смузі можна приготувати з $0{,}8$ кг ягід?}
\answerTable{$12$}{$13$}{$14$}{$15$}{$20$}
\zadtask{Для випікання одного кексу потрібно $35$ г родзинок. Пекар випік \textit{найбільшу} можливу кількість таких кексів, маючи $0{,}6$ кг родзинок. Скільки \textit{грамів} родзинок у нього залишилося?}
\answerTable{$0$}{$5$}{$10$}{$15$}{$25$}
\taskBlock{Спрощення раціонального виразу --- сесія 30.05, завдання №6}
\zadtask{Спростіть вираз $\dfrac{9 - x^2}{x^2 + 6x + 9}$.}
\answerTableTall{$\dfrac{3-x}{x+3}$}{$\dfrac{x-3}{x+3}$}{$\dfrac{3+x}{x-3}$}{$-\dfrac{x+3}{x-3}$}{$\dfrac{1}{x+3}$}
\solution{Чисельник $9-x^2=(3-x)(3+x)$, знаменник $x^2+6x+9=(x+3)^2$. Тоді $\dfrac{(3-x)(3+x)}{(x+3)^2}=\dfrac{3-x}{x+3}$. Відповідь: \textbf{А}.}
\zadtask{Спростіть вираз $\dfrac{16 - x^2}{x^2 + 8x + 16}$.}
\answerTableTall{$\dfrac{x-4}{x+4}$}{$\dfrac{4-x}{x+4}$}{$\dfrac{4+x}{x-4}$}{$-\dfrac{x+4}{x-4}$}{$\dfrac{1}{x+4}$}
\zadtask{Спростіть вираз $\dfrac{x^2 - 10x + 25}{x^2 - 25}$.}
\answerTableTall{$\dfrac{x-5}{x+5}$}{$\dfrac{x+5}{x-5}$}{$\dfrac{5-x}{x+5}$}{$-\dfrac{x+5}{x-5}$}{$\dfrac{1}{x-5}$}
\taskBlock{Спрощення тригонометричного виразу --- сесія 30.05, завдання №13}
\zadtask{Спростіть вираз $\mathrm{tg}^2 x \cdot \cos^2 x - 1$.}
\answerTableTall{$-\sin^2 x$}{$-\cos^2 x$}{$-1$}{$\cos^2 x$}{$\sin^2 x$}
\solution{$\mathrm{tg}^2 x\cdot\cos^2 x=\dfrac{\sin^2 x}{\cos^2 x}\cdot\cos^2 x=\sin^2 x$. Тоді $\sin^2 x-1=-(1-\sin^2 x)=-\cos^2 x$. Відповідь: \textbf{Б}.}
\zadtask{Спростіть вираз $\mathrm{ctg}^2 x \cdot \sin^2 x - 1$.}
\answerTableTall{$-\sin^2 x$}{$-\cos^2 x$}{$-1$}{$\cos^2 x$}{$\sin^2 x$}
\zadtask{Спростіть вираз $(1 - \cos x)(1 + \cos x) - 1$.}
\answerTableTall{$-\sin^2 x$}{$-\cos^2 x$}{$\sin^2 x$}{$\cos^2 x$}{$-1$}
\taskBlock{Спрощення виразу зі степенями --- сесія 1.06, завдання №4}
\zadtask{Спростіть вираз $\dfrac{6x^3 + 3x^3}{3}$.}
\answerTable{$3x^6$}{$6x^6$}{$5x^3$}{$3x^3$}{$7x^3$}
\solution{$\dfrac{6x^3+3x^3}{3}=\dfrac{9x^3}{3}=3x^3$. Відповідь: \textbf{Г}.}
\zadtask{Спростіть вираз $\dfrac{10x^5 + 5x^5}{5}$.}
\answerTable{$3x^{10}$}{$15x^5$}{$3x^5$}{$5x^5$}{$2x^5$}
\zadtask{Спростіть вираз $\dfrac{4x^2 \cdot 6x^3}{8x}$.}
\answerTable{$3x^4$}{$3x^5$}{$2x^4$}{$24x^5$}{$3x^6$}
\taskBlock{Сюжетна задача на спільну роботу --- сесія 1.06, завдання №11}
\zadtask{У цеху на лінії виробництва пляшкової води працюють два автомати для клейового маркування. Перший маркує $14\,000$ пляшок за годину, а другий~--- $16\,000$ пляшок за годину. Обидва автомати одночасно працювали $12$ хвилин. Скільки всього пляшок було промарковано за цей час?}
\answerTable{$2000$}{$3000$}{$5000$}{$6000$}{$7500$}
\solution{За годину разом маркують $14\,000+16\,000=30\,000$ пляшок. $12$ хв $=\frac{12}{60}=0{,}2$ год, тож $30\,000\cdot 0{,}2=6000$. Відповідь: \textbf{Г}.}
\zadtask{На пакувальній лінії працюють два автомати. Перший пакує $12\,000$ коробок за годину, а другий~--- $18\,000$ коробок за годину. Обидва автомати одночасно працювали $18$ хвилин. Скільки всього коробок було запаковано за цей час?}
\answerTable{$5400$}{$7500$}{$9000$}{$10\,000$}{$30\,000$}
\zadtask{У друкарні працюють два верстати. Перший друкує $9000$ аркушів за годину, а другий~--- $15\,000$ аркушів за годину. Обидва верстати одночасно працювали $20$ хвилин. На скільки \textit{більше} аркушів надрукував за цей час другий верстат, ніж перший?}
\answerTable{$1500$}{$2000$}{$3000$}{$4000$}{$6000$}
\taskBlock{Знаки тригонометричних функцій --- сесія 1.06, завдання №13}
\zadtask{Відомо, що $\cos 210^\circ \cdot \sin\alpha > 0$. Якого з наведених значень може набувати кут $\alpha$? \nmtyear{2026}}
\answerTable{$0^\circ$}{$30^\circ$}{$90^\circ$}{$150^\circ$}{$300^\circ$}
\solution{$\cos 210^\circ<0$, тому для $\cos 210^\circ\cdot\sin\alpha>0$ потрібно $\sin\alpha<0$. Серед варіантів лише $\sin 300^\circ<0$. Відповідь: \textbf{Д}.}
\zadtask{Відомо, що $\cos 120^\circ \cdot \sin\alpha > 0$. Якого з наведених значень може набувати кут $\alpha$?}
\answerTable{$0^\circ$}{$60^\circ$}{$90^\circ$}{$120^\circ$}{$210^\circ$}
\zadtask{Відомо, що $\cos 300^\circ \cdot \sin\alpha < 0$. Якого з наведених значень може набувати кут $\alpha$?}
\answerTable{$30^\circ$}{$90^\circ$}{$150^\circ$}{$180^\circ$}{$210^\circ$}
\taskBlock{Розкриття дужок і спрощення виразу --- сесія 2.06, завдання №2}
\zadtask{Спростіть вираз $15 - 4(2x^2 + 3)$.}
\answerTable{$3 - 8x^2$}{$18 - 8x^2$}{$27 - 8x^2$}{$22x^2 + 33$}{$22x^2 + 3$}
\solution{Розкриваємо дужки: $15-4(2x^2+3)=15-8x^2-12=3-8x^2$. Відповідь: \textbf{А}.}
\zadtask{Спростіть вираз $12 - 5(3x^2 + 2)$.}
\answerTable{$22 - 15x^2$}{$2 - 15x^2$}{$10 - 15x^2$}{$2 + 15x^2$}{$14 - 15x^2$}
\zadtask{Спростіть вираз $3(4x^2 + 5) - 20$.}
\answerTable{$12x^2 - 5$}{$12x^2 + 5$}{$12x^2 - 35$}{$5 - 12x^2$}{$12x^2 - 15$}
\taskBlock{Спрощення раціонального виразу --- сесія 2.06, завдання №15}
\zadtask{Спростіть вираз $\dfrac{x^2 - 4xy + 4y^2}{x - 2y} : (2y - x)$.}
\answerTable{$-1$}{$-x^2 + 4xy - y^2$}{$2y - x$}{$1$}{$x^2 - 4xy + y^2$}
\solution{Чисельник $x^2-4xy+4y^2=(x-2y)^2$, тож дріб дорівнює $x-2y$. Ділення на $(2y-x)=-(x-2y)$ дає $\dfrac{x-2y}{-(x-2y)}=-1$. Відповідь: \textbf{А}.}
\zadtask{Спростіть вираз $\dfrac{x^2 + 6xy + 9y^2}{x + 3y} : (x + 3y)$.}
\answerTable{$-1$}{$1$}{$x + 3y$}{$-(x + 3y)$}{$(x + 3y)^2$}
\zadtask{Спростіть вираз $\dfrac{x - 5y}{x^2 - 10xy + 25y^2} \cdot (5y - x)$.}
\answerTable{$1$}{$x - 5y$}{$-1$}{$5y - x$}{$\dfrac{1}{x - 5y}$}
\taskBlock{Пропорційні величини (швидкість відтворення) --- сесія 3.06, завдання №3}
\zadtask{Вадим планує переглянути відеоролик тривалістю $60$ хвилин. Перед переглядом він збільшив швидкість відтворення в $1{,}5$ раза. Скільки хвилин тепер триватиме перегляд?}
\answerTable{$40$}{$90$}{$45$}{$30$}{$35$}
\solution{Збільшення швидкості в $1{,}5$ раза зменшує час у стільки ж разів: $60 : 1{,}5 = 40$ (хв). Відповідь: \textbf{А} ($40$).}
\zadtask{Оксана планує переглянути відеолекцію тривалістю $80$ хвилин. Перед переглядом вона збільшила швидкість відтворення в $1{,}25$ раза. Скільки хвилин тепер триватиме перегляд?}
\answerTable{$64$}{$100$}{$60$}{$55$}{$75$}
\zadtask{Тарас слухає аудіокнигу тривалістю $90$ хвилин. Перед прослуховуванням він \textit{зменшив} швидкість відтворення до $0{,}75$ від звичайної. Скільки хвилин тепер триватиме прослуховування?}
\answerTable{$67{,}5$}{$105$}{$162$}{$120$}{$115$}
\taskBlock{Дії з мішаними дробами --- сесія 3.06, завдання №7}
\zadtask{Обчисліть $5\dfrac{1}{3} : 1\dfrac{3}{5}$.}
\answerTableTall{$\dfrac{3}{10}$}{$\dfrac{128}{15}$}{$\dfrac{10}{3}$}{$\dfrac{5}{3}$}{$\dfrac{15}{128}$}
\solution{$5\dfrac13 : 1\dfrac35 = \dfrac{16}{3} : \dfrac{8}{5} = \dfrac{16}{3}\cdot\dfrac{5}{8} = \dfrac{10}{3}$. Відповідь: \textbf{В} ($\dfrac{10}{3}$).}
\zadtask{Обчисліть $4\dfrac{2}{3} : 1\dfrac{1}{6}$.}
\answerTableTall{$\dfrac{1}{4}$}{$4$}{$\dfrac{49}{9}$}{$\dfrac{9}{49}$}{$\dfrac{7}{2}$}
\zadtask{Обчисліть $3\dfrac{3}{4} \cdot 2\dfrac{2}{5}$.}
\answerTableTall{$\dfrac{75}{8}$}{$9$}{$\dfrac{25}{16}$}{$6$}{$\dfrac{15}{2}$}
\taskBlock{Обчислення тригонометричного виразу --- сесія 3.06, завдання №13}
\zadtask{Обчисліть значення виразу $2\sin\left(-\dfrac{\pi}{2}\right) + \cos 3\pi$.}
\answerTable{$-3$}{$-2$}{$-1$}{$2$}{$3$}
\solution{$\sin\left(-\dfrac{\pi}{2}\right) = -1$, $\cos 3\pi = -1$, тому $2\cdot(-1) + (-1) = -3$. Відповідь: \textbf{А} ($-3$).}
\zadtask{Обчисліть значення виразу $4\sin\dfrac{3\pi}{2} + \cos 2\pi$.}
\answerTable{$-5$}{$-3$}{$-4$}{$3$}{$5$}
\zadtask{Обчисліть значення виразу $\operatorname{tg}\dfrac{3\pi}{4} + 2\sin\dfrac{\pi}{2}$.}
\answerTable{$-3$}{$1$}{$-1$}{$3$}{$0$}
\taskBlock{Стандартний вигляд числа --- сесія 4.06, завдання №1}
\zadtask{В астрономії для визначення відстані використовують позасистемну одиницю довжини~--- астрономічна одиниця (а.\,о.), що дорівнює середній відстані від Землі до Сонця: $1$ а.\,о. $\approx 149{,}6$~млн~\textit{км}. Запишіть це значення (у \textit{км}) у стандартному вигляді.}
\answerTable{$149{,}6 \cdot 10^6$}{$1{,}496 \cdot 10^8$}{$0{,}1496 \cdot 10^9$}{$1496 \cdot 10^5$}{$14{,}96 \cdot 10^7$}
\solution{$149{,}6$~млн $=149\,600\,000=1{,}496\cdot10^8$. У стандартному вигляді множник має задовольняти $1\leqslant a<10$, тож $1{,}496\cdot10^8$. Відповідь: \textbf{Б}.}
\zadtask{В астрономії використовують одиницю довжини парсек (пк), що приблизно дорівнює $30{,}86$~трлн~\textit{км}. Запишіть це значення (у \textit{км}) у стандартному вигляді.}
\answerTable{$30{,}86 \cdot 10^{12}$}{$3{,}086 \cdot 10^{13}$}{$0{,}3086 \cdot 10^{14}$}{$3086 \cdot 10^{10}$}{$3{,}086 \cdot 10^{11}$}
\zadtask{Діаметр еритроцита людини приблизно дорівнює $0{,}0000075$~\textit{м}. Запишіть це значення (у \textit{м}) у стандартному вигляді.}
\answerTable{$7{,}5 \cdot 10^{-6}$}{$7{,}5 \cdot 10^{6}$}{$0{,}75 \cdot 10^{-5}$}{$75 \cdot 10^{-7}$}{$7{,}5 \cdot 10^{-5}$}
\taskBlock{Спрощення раціонального виразу --- сесія 4.06, завдання №6}
\zadtask{$\dfrac{a^2 - 25b^2}{10b + 2a} =$}
\answerTable{$\dfrac{a - 5b}{2}$}{$2a - 10b$}{$\dfrac{a + 5b}{2}$}{$2a + 10b$}{$\dfrac{a - 5b}{4}$}
\solution{Чисельник $a^2-25b^2=(a-5b)(a+5b)$, знаменник $10b+2a=2(a+5b)$. Скорочуємо на $(a+5b)$: $\dfrac{(a-5b)(a+5b)}{2(a+5b)}=\dfrac{a-5b}{2}$. Відповідь: \textbf{А}.}
\zadtask{$\dfrac{a^2 - 9b^2}{6b + 2a} =$}
\answerTable{$2a - 6b$}{$\dfrac{a + 3b}{2}$}{$\dfrac{a - 3b}{2}$}{$2a + 6b$}{$\dfrac{a - 3b}{4}$}
\zadtask{$\dfrac{a^2 - 16b^2}{4b - a} =$}
\answerTable{$a + 4b$}{$\dfrac{a + 4b}{4}$}{$-a - 4b$}{$a - 4b$}{$\dfrac{a - 4b}{4}$}
\taskBlock{Тригонометричний вираз (формула подвійного кута) --- сесія 4.06, завдання №13}
\zadtask{Якому проміжку належить значення виразу $6 \sin 15^\circ \cos 15^\circ$?}
\answerTable{$(0;\, 1)$}{$[1;\, 2)$}{$[2;\, 3)$}{$[3;\, 5)$}{$[5;\, 8)$}
\solution{За формулою $2\sin\alpha\cos\alpha=\sin2\alpha$: $6\sin15^\circ\cos15^\circ=3\sin30^\circ=3\cdot\dfrac{1}{2}=1{,}5\in[1;\,2)$. Відповідь: \textbf{Б}.}
\zadtask{Якому проміжку належить значення виразу $8 \sin 15^\circ \cos 15^\circ$?}
\answerTable{$(0;\, 1)$}{$[1;\, 2)$}{$[2;\, 3)$}{$[3;\, 5)$}{$[5;\, 8)$}
\zadtask{Якому проміжку належить значення виразу $6\left(\cos^2 15^\circ - \sin^2 15^\circ\right)$?}
\answerTable{$(0;\, 1)$}{$[1;\, 2)$}{$[2;\, 3)$}{$[3;\, 5)$}{$[5;\, 8)$}
\taskBlock{Пропорційне обчислення (швидкість роботи) --- сесія 5.06, завдання №1}
\zadtask{Настінний принтер друкує зображення зі сталою швидкістю $3$ м$^2$ за одну годину. Скільки \textit{хвилин} принтер друкував на стіні будинку зображення площею $2{,}4$ м$^2$?}
\answerTable{$36$}{$40$}{$51$}{$48$}{$45$}
\solution{Час друку: $\dfrac{2{,}4}{3}=0{,}8$ год $=0{,}8\cdot 60=48$ хв. Відповідь: \textbf{Г}.}
\zadtask{Насос перекачує воду зі сталою швидкістю $5$ м$^3$ за одну годину. Скільки \textit{хвилин} насос наповнював резервуар об'ємом $3{,}5$ м$^3$?}
\answerTable{$35$}{$40$}{$42$}{$45$}{$48$}
\zadtask{Верстат обробляє поверхню зі сталою швидкістю. За $36$ \textit{хвилин} він обробив панель площею $2{,}4$ м$^2$. Скільки квадратних метрів поверхні верстат обробляє за одну \textit{годину}?}
\answerTable{$3{,}6$}{$4$}{$4{,}8$}{$4{,}4$}{$5$}
\taskBlock{Спрощення виразу (формули скороченого множення) --- сесія 5.06, завдання №6}
\zadtask{Спростіть вираз $(a+6)(a-6) - (a-4)^2$.}
\answerTable{$8a - 52$}{$-8a - 52$}{$8a - 20$}{$-52$}{$8a + 52$}
\solution{$(a+6)(a-6)-(a-4)^2 = (a^2-36)-(a^2-8a+16)=8a-52$. Відповідь: \textbf{А}.}
\zadtask{Спростіть вираз $(b+5)(b-5) - (b-3)^2$.}
\answerTable{$6b - 34$}{$-6b - 34$}{$6b - 16$}{$-34$}{$6b + 34$}
\zadtask{Спростіть вираз $(c-2)^2 - (c+7)(c-7)$.}
\answerTable{$-4c + 53$}{$4c - 53$}{$-4c - 45$}{$53$}{$-4c - 53$}
\taskBlock{Спрощення тригонометричного виразу --- сесія 5.06, завдання №10}
\zadtask{Спростіть вираз $\mathrm{tg}^2\alpha - \mathrm{tg}^2\alpha \cdot \sin^2\alpha$.}
\answerTableTall{$\sin^2\alpha$}{$\cos^2\alpha$}{$\mathrm{tg}^2\alpha$}{$1$}{$\sin^2\alpha\cos^2\alpha$}
\solution{$\mathrm{tg}^2\alpha-\mathrm{tg}^2\alpha\sin^2\alpha=\mathrm{tg}^2\alpha(1-\sin^2\alpha)=\dfrac{\sin^2\alpha}{\cos^2\alpha}\cdot\cos^2\alpha=\sin^2\alpha$. Відповідь: \textbf{А}.}
\zadtask{Спростіть вираз $\mathrm{ctg}^2\alpha - \mathrm{ctg}^2\alpha \cdot \cos^2\alpha$.}
\answerTableTall{$\sin^2\alpha$}{$\cos^2\alpha$}{$\mathrm{ctg}^2\alpha$}{$1$}{$\sin^2\alpha\cos^2\alpha$}
\zadtask{Спростіть вираз $\dfrac{1 - \sin^2\alpha}{\mathrm{ctg}^2\alpha}$.}
\answerTableTall{$\sin^2\alpha$}{$\cos^2\alpha$}{$\mathrm{ctg}^2\alpha$}{$1$}{$\mathrm{tg}^2\alpha$}
\typeTitle{Завдання на встановлення відповідності}
\taskBlock{Відповідність: вирази та їх значення --- сесія 23.05, завдання №17}
\zadtask{Установіть відповідність між виразом (1--3) та його значенням (А--Д). \nmtyear{2026}

\vspace{0.3cm}

\noindent
\begin{minipage}[t]{0.52\textwidth}
\textit{Вираз} \par \vspace{0.2cm}
\textbf{1} \ \ $\left|\dfrac{5}{3} - 2\right|$\\[0.35cm]
\textbf{2} \ \ $3^{-2} \cdot \left(\dfrac{1}{3}\right)^{-3}$\\[0.35cm]
\textbf{3} \ \ $\sqrt{3} \cdot \sin\dfrac{\pi}{3}$
\end{minipage}
\begin{minipage}[t]{0.24\textwidth}
\textit{Значення} \par \vspace{0.2cm}
\textbf{А} \ \ $\dfrac{3}{2}$\\[0.15cm]
\textbf{Б} \ \ $3$\\[0.15cm]
\textbf{В} \ \ $\dfrac{1}{3}$\\[0.15cm]
\textbf{Г} \ \ $-\dfrac{1}{3}$\\[0.15cm]
\textbf{Д} \ \ $\dfrac{1}{2}$
\end{minipage}
\hfill
\begin{minipage}[t]{0.22\textwidth}
\vspace{-0.2cm}
\begin{flushright}\matchingGrid\end{flushright}
\end{minipage}
}
\solution{$1)\ \left|\dfrac{5}{3}-2\right|=\left|-\dfrac{1}{3}\right|=\dfrac{1}{3}$~(В); $2)\ 3^{-2}\cdot 3^{3}=3^{1}=3$~(Б); $3)\ \sqrt{3}\cdot\sin\dfrac{\pi}{3}=\sqrt{3}\cdot\dfrac{\sqrt{3}}{2}=\dfrac{3}{2}$~(А). Відповідь: \textbf{1--В, 2--Б, 3--А}.}
\zadtask{Установіть відповідність між виразом (1--3) та його значенням (А--Д).

\vspace{0.3cm}

\noindent
\begin{minipage}[t]{0.52\textwidth}
\textit{Вираз} \par \vspace{0.2cm}
\textbf{1} \ \ $\left|\dfrac{7}{4} - 3\right|$\\[0.35cm]
\textbf{2} \ \ $2^{-3} \cdot \left(\dfrac{1}{2}\right)^{-5}$\\[0.35cm]
\textbf{3} \ \ $\sqrt{2} \cdot \cos\dfrac{\pi}{4}$
\end{minipage}
\begin{minipage}[t]{0.24\textwidth}
\textit{Значення} \par \vspace{0.2cm}
\textbf{А} \ \ $1$\\[0.15cm]
\textbf{Б} \ \ $\dfrac{5}{4}$\\[0.15cm]
\textbf{В} \ \ $\dfrac{1}{2}$\\[0.15cm]
\textbf{Г} \ \ $4$\\[0.15cm]
\textbf{Д} \ \ $-\dfrac{5}{4}$
\end{minipage}
\hfill
\begin{minipage}[t]{0.22\textwidth}
\vspace{-0.2cm}
\begin{flushright}\matchingGrid\end{flushright}
\end{minipage}
}
\zadtask{Установіть відповідність між виразом (1--3) та його значенням (А--Д).

\vspace{0.3cm}

\noindent
\begin{minipage}[t]{0.52\textwidth}
\textit{Вираз} \par \vspace{0.2cm}
\textbf{1} \ \ $\log_2 \dfrac{1}{4}$\\[0.35cm]
\textbf{2} \ \ $\left(\dfrac{2}{3}\right)^{-1}$\\[0.35cm]
\textbf{3} \ \ $2 \cdot \cos\dfrac{\pi}{3}$
\end{minipage}
\begin{minipage}[t]{0.24\textwidth}
\textit{Значення} \par \vspace{0.2cm}
\textbf{А} \ \ $1$\\[0.15cm]
\textbf{Б} \ \ $-2$\\[0.15cm]
\textbf{В} \ \ $\dfrac{3}{2}$\\[0.15cm]
\textbf{Г} \ \ $-\dfrac{1}{2}$\\[0.15cm]
\textbf{Д} \ \ $2$
\end{minipage}
\hfill
\begin{minipage}[t]{0.22\textwidth}
\vspace{-0.2cm}
\begin{flushright}\matchingGrid\end{flushright}
\end{minipage}
}
\taskBlock{Обчислення значень виразів (відповідність) --- сесія 30.05, завдання №18}
\zadtask{Установіть відповідність між виразом (1--3) і його значенням (А--Д), якщо $a = \dfrac{3}{8}$.

\vspace{0.2cm}
\noindent
\begin{minipage}[t]{0.45\textwidth}
\textit{Вираз}\par\vspace{0.15cm}
\textbf{1} \quad $\sqrt[3]{\dfrac{a}{3}}$\\[0.35cm]
\textbf{2} \quad $\log_2 a + \log_2 \dfrac{16}{3}$\\[0.35cm]
\textbf{3} \quad $\sin(4\pi a)$
\end{minipage}
\hfill
\begin{minipage}[t]{0.3\textwidth}
\textit{Значення}\par\vspace{0.15cm}
\textbf{А} \quad $2$\\[0.15cm]
\textbf{Б} \quad $1$\\[0.15cm]
\textbf{В} \quad $\dfrac{1}{2}$\\[0.2cm]
\textbf{Г} \quad $0$\\[0.15cm]
\textbf{Д} \quad $-1$
\end{minipage}
\hfill
\begin{minipage}[t]{0.2\textwidth}
\vspace{0.4cm}
\begin{flushright}\matchingGrid\end{flushright}
\end{minipage}
}
\solution{$a=\dfrac38$. \textbf{1}: $\sqrt[3]{\dfrac{a}{3}}=\sqrt[3]{\dfrac18}=\dfrac12$ (В). \textbf{2}: $\log_2\!\left(a\cdot\dfrac{16}{3}\right)=\log_2 2=1$ (Б). \textbf{3}: $\sin\!\left(4\pi\cdot\dfrac38\right)=\sin\dfrac{3\pi}{2}=-1$ (Д). Відповідь: \textbf{1~--~В, 2~--~Б, 3~--~Д}.}
\zadtask{Установіть відповідність між виразом (1--3) і його значенням (А--Д), якщо $a = \dfrac{5}{8}$.

\vspace{0.2cm}
\noindent
\begin{minipage}[t]{0.45\textwidth}
\textit{Вираз}\par\vspace{0.15cm}
\textbf{1} \quad $\sqrt[3]{\dfrac{a}{5}}$\\[0.35cm]
\textbf{2} \quad $\log_2 a + \log_2 \dfrac{32}{5}$\\[0.35cm]
\textbf{3} \quad $\sin(4\pi a)$
\end{minipage}
\hfill
\begin{minipage}[t]{0.3\textwidth}
\textit{Значення}\par\vspace{0.15cm}
\textbf{А} \quad $2$\\[0.15cm]
\textbf{Б} \quad $1$\\[0.15cm]
\textbf{В} \quad $\dfrac{1}{2}$\\[0.2cm]
\textbf{Г} \quad $0$\\[0.15cm]
\textbf{Д} \quad $-1$
\end{minipage}
\hfill
\begin{minipage}[t]{0.2\textwidth}
\vspace{0.4cm}
\begin{flushright}\matchingGrid\end{flushright}
\end{minipage}
}
\zadtask{Установіть відповідність між виразом (1--3) і його значенням (А--Д), якщо $a = \dfrac{1}{4}$.

\vspace{0.2cm}
\noindent
\begin{minipage}[t]{0.45\textwidth}
\textit{Вираз}\par\vspace{0.15cm}
\textbf{1} \quad $\sqrt{a}$\\[0.35cm]
\textbf{2} \quad $\log_2 \dfrac{1}{2a}$\\[0.35cm]
\textbf{3} \quad $\cos(4\pi a)$
\end{minipage}
\hfill
\begin{minipage}[t]{0.3\textwidth}
\textit{Значення}\par\vspace{0.15cm}
\textbf{А} \quad $2$\\[0.15cm]
\textbf{Б} \quad $1$\\[0.15cm]
\textbf{В} \quad $\dfrac{1}{2}$\\[0.2cm]
\textbf{Г} \quad $0$\\[0.15cm]
\textbf{Д} \quad $-1$
\end{minipage}
\hfill
\begin{minipage}[t]{0.2\textwidth}
\vspace{0.4cm}
\begin{flushright}\matchingGrid\end{flushright}
\end{minipage}
}
\taskBlock{Відповідність: значення виразів (степені, логарифми) --- сесія 1.06, завдання №16}
\zadtask{Узгодьте вираз (1--3) із його значенням (А--Д), якщо $a = 4$.

\vspace{0.3cm}
\noindent
\begin{minipage}[t]{0.45\textwidth}
\textit{Вираз}\par\vspace{0.2cm}
\textbf{1} \quad $\dfrac{2a^2 - 12a}{a^2 - 36}$\\[0.5cm]
\textbf{2} \quad $0{,}04^{\,\log_5 a}$\\[0.45cm]
\textbf{3} \quad $a^0 - a^{\frac{3}{2}}$
\end{minipage}
\hfill
\begin{minipage}[t]{0.3\textwidth}
\textit{Значення}\par\vspace{0.2cm}
\textbf{А} \quad $\dfrac{1}{16}$\\[0.25cm]
\textbf{Б} \quad $-8$\\[0.2cm]
\textbf{В} \quad $-7$\\[0.2cm]
\textbf{Г} \quad $0{,}8$\\[0.2cm]
\textbf{Д} \quad $-\dfrac{1}{16}$
\end{minipage}
\hfill
\begin{minipage}[t]{0.2\textwidth}
\vspace{0.5cm}
\begin{flushright}\matchingGrid\end{flushright}
\end{minipage}
}
\solution{При $a=4$: \textbf{1)} $\dfrac{2a^2-12a}{a^2-36}=\dfrac{2a(a-6)}{(a-6)(a+6)}=\dfrac{2a}{a+6}=\dfrac{8}{10}=0{,}8$ (Г); \textbf{2)} $0{,}04^{\log_5 4}=(5^{-2})^{\log_5 4}=4^{-2}=\dfrac{1}{16}$ (А); \textbf{3)} $a^0-a^{3/2}=1-8=-7$ (В). Відповідь: \textbf{1\,--\,Г, 2\,--\,А, 3\,--\,В}.}
\zadtask{Узгодьте вираз (1--3) із його значенням (А--Д), якщо $a = 9$.

\vspace{0.3cm}
\noindent
\begin{minipage}[t]{0.45\textwidth}
\textit{Вираз}\par\vspace{0.2cm}
\textbf{1} \quad $\dfrac{a^2 - 3a}{a^2 - 9}$\\[0.5cm]
\textbf{2} \quad $0{,}25^{\,\log_2 a}$\\[0.45cm]
\textbf{3} \quad $a^0 - a^{\frac{3}{2}}$
\end{minipage}
\hfill
\begin{minipage}[t]{0.3\textwidth}
\textit{Значення}\par\vspace{0.2cm}
\textbf{А} \quad $-26$\\[0.25cm]
\textbf{Б} \quad $\dfrac{1}{81}$\\[0.2cm]
\textbf{В} \quad $0{,}75$\\[0.2cm]
\textbf{Г} \quad $-\dfrac{1}{81}$\\[0.2cm]
\textbf{Д} \quad $-28$
\end{minipage}
\hfill
\begin{minipage}[t]{0.2\textwidth}
\vspace{0.5cm}
\begin{flushright}\matchingGrid\end{flushright}
\end{minipage}
}
\zadtask{Узгодьте вираз (1--3) із його значенням (А--Д), якщо $a = 16$.

\vspace{0.3cm}
\noindent
\begin{minipage}[t]{0.45\textwidth}
\textit{Вираз}\par\vspace{0.2cm}
\textbf{1} \quad $\dfrac{a - 9}{\sqrt{a} - 3}$\\[0.5cm]
\textbf{2} \quad $\log_2 a$\\[0.45cm]
\textbf{3} \quad $a^{-\frac{1}{2}}$
\end{minipage}
\hfill
\begin{minipage}[t]{0.3\textwidth}
\textit{Значення}\par\vspace{0.2cm}
\textbf{А} \quad $0{,}25$\\[0.25cm]
\textbf{Б} \quad $4$\\[0.2cm]
\textbf{В} \quad $7$\\[0.2cm]
\textbf{Г} \quad $16$\\[0.2cm]
\textbf{Д} \quad $2$
\end{minipage}
\hfill
\begin{minipage}[t]{0.2\textwidth}
\vspace{0.5cm}
\begin{flushright}\matchingGrid\end{flushright}
\end{minipage}
}
\taskBlock{Відповідність: значення виразу та проміжок --- сесія 2.06, завдання №16}
\zadtask{Узгодьте вираз (1--3) із проміжком (А--Д), якому належить значення цього виразу.

\vspace{0.3cm}
\noindent
\begin{minipage}[t]{0.4\textwidth}
\textit{Вираз}\par\vspace{0.2cm}
\textbf{1} \quad $\dfrac{7}{5} - \dfrac{5}{3}$\\[0.5cm]
\textbf{2} \quad $\sqrt{3} \cdot \sqrt{5}$\\[0.4cm]
\textbf{3} \quad $\log_4 8$
\end{minipage}
\hfill
\begin{minipage}[t]{0.32\textwidth}
\textit{Проміжок}\par\vspace{0.2cm}
\textbf{А} \quad $[-1;\, 0)$\\[0.15cm]
\textbf{Б} \quad $[0;\, 1)$\\[0.15cm]
\textbf{В} \quad $[1;\, 2)$\\[0.15cm]
\textbf{Г} \quad $[2;\, 4)$\\[0.15cm]
\textbf{Д} \quad $[4;\, +\infty)$
\end{minipage}
\hfill
\begin{minipage}[t]{0.2\textwidth}
\vspace{0.5cm}
\begin{flushright}\matchingGrid\end{flushright}
\end{minipage}
}
\solution{$\dfrac{7}{5}-\dfrac{5}{3}=\dfrac{21-25}{15}=-\dfrac{4}{15}\in[-1;0)$ -- \textbf{А}; $\sqrt{3}\cdot\sqrt{5}=\sqrt{15}\approx3{,}87\in[2;4)$ -- \textbf{Г}; $\log_4 8=\dfrac{3}{2}=1{,}5\in[1;2)$ -- \textbf{В}. Відповідь: \textbf{1--А, 2--Г, 3--В}.}
\zadtask{Узгодьте вираз (1--3) із проміжком (А--Д), якому належить значення цього виразу.

\vspace{0.3cm}
\noindent
\begin{minipage}[t]{0.4\textwidth}
\textit{Вираз}\par\vspace{0.2cm}
\textbf{1} \quad $\dfrac{9}{4} - \dfrac{5}{3}$\\[0.5cm]
\textbf{2} \quad $\sqrt{2} \cdot \sqrt{7}$\\[0.4cm]
\textbf{3} \quad $\log_9 27$
\end{minipage}
\hfill
\begin{minipage}[t]{0.32\textwidth}
\textit{Проміжок}\par\vspace{0.2cm}
\textbf{А} \quad $[-1;\, 0)$\\[0.15cm]
\textbf{Б} \quad $[0;\, 1)$\\[0.15cm]
\textbf{В} \quad $[1;\, 2)$\\[0.15cm]
\textbf{Г} \quad $[2;\, 4)$\\[0.15cm]
\textbf{Д} \quad $[4;\, +\infty)$
\end{minipage}
\hfill
\begin{minipage}[t]{0.2\textwidth}
\vspace{0.5cm}
\begin{flushright}\matchingGrid\end{flushright}
\end{minipage}
}
\zadtask{Узгодьте вираз (1--3) із проміжком (А--Д), якому належить значення цього виразу.

\vspace{0.3cm}
\noindent
\begin{minipage}[t]{0.4\textwidth}
\textit{Вираз}\par\vspace{0.2cm}
\textbf{1} \quad $\cos 60^\circ - 1$\\[0.5cm]
\textbf{2} \quad $\sqrt{50} - \sqrt{2}$\\[0.4cm]
\textbf{3} \quad $\log_2 12 - \log_2 3$
\end{minipage}
\hfill
\begin{minipage}[t]{0.32\textwidth}
\textit{Проміжок}\par\vspace{0.2cm}
\textbf{А} \quad $[-1;\, 0)$\\[0.15cm]
\textbf{Б} \quad $[0;\, 1)$\\[0.15cm]
\textbf{В} \quad $[1;\, 2)$\\[0.15cm]
\textbf{Г} \quad $[2;\, 4)$\\[0.15cm]
\textbf{Д} \quad $[4;\, +\infty)$
\end{minipage}
\hfill
\begin{minipage}[t]{0.2\textwidth}
\vspace{0.5cm}
\begin{flushright}\matchingGrid\end{flushright}
\end{minipage}
}
\taskBlock{Тотожні перетворення виразів (відповідність) --- сесія 3.06, завдання №16}
\noindent\zadnum Узгодьте вираз (1--3) з тотожно рівним йому виразом (А--Д), якщо $a \neq 0$.

\vspace{0.3cm}
\noindent
\begin{minipage}[t]{0.42\textwidth}
\textit{Вираз}\par\vspace{0.2cm}
\textbf{1} \quad $\sqrt[3]{\dfrac{27}{a^3}}$\\[0.55cm]
\textbf{2} \quad $a^{-2}\cdot a^3 \cdot \log_{27} 3$\\[0.4cm]
\textbf{3} \quad $\dfrac{(a-3)^2 - a^2 - 9}{2}$
\end{minipage}
\hfill
\begin{minipage}[t]{0.3\textwidth}
\textit{Тотожно рівний вираз}\par\vspace{0.2cm}
\textbf{А} \quad $-3a$\\[0.2cm]
\textbf{Б} \quad $\dfrac{3}{a}$\\[0.25cm]
\textbf{В} \quad $0$\\[0.2cm]
\textbf{Г} \quad $\dfrac{a}{3}$\\[0.25cm]
\textbf{Д} \quad $3a$
\end{minipage}
\hfill
\begin{minipage}[t]{0.2\textwidth}
\vspace{0.5cm}
\begin{flushright}\matchingGrid\end{flushright}
\end{minipage}

\solution{\textbf{1}: $\sqrt[3]{\dfrac{27}{a^3}} = \dfrac{3}{a}$~--- Б. \textbf{2}: $a^{-2}\cdot a^3\cdot\log_{27}3 = a\cdot\dfrac13 = \dfrac{a}{3}$~--- Г. \textbf{3}: $\dfrac{(a-3)^2 - a^2 - 9}{2} = \dfrac{-6a}{2} = -3a$~--- А. Відповідь: \textbf{1~--~Б, 2~--~Г, 3~--~А}.}

\vspace{0.4cm}
\noindent\zadnum Узгодьте вираз (1--3) з тотожно рівним йому виразом (А--Д), якщо $b \neq 0$.

\vspace{0.3cm}
\noindent
\begin{minipage}[t]{0.42\textwidth}
\textit{Вираз}\par\vspace{0.2cm}
\textbf{1} \quad $\sqrt[3]{\dfrac{64}{b^3}}$\\[0.55cm]
\textbf{2} \quad $b^{-3}\cdot b^4 \cdot \log_{16} 2$\\[0.4cm]
\textbf{3} \quad $\dfrac{(b-4)^2 - b^2 - 16}{2}$
\end{minipage}
\hfill
\begin{minipage}[t]{0.3\textwidth}
\textit{Тотожно рівний вираз}\par\vspace{0.2cm}
\textbf{А} \quad $\dfrac{4}{b}$\\[0.25cm]
\textbf{Б} \quad $4b$\\[0.2cm]
\textbf{В} \quad $\dfrac{b}{4}$\\[0.25cm]
\textbf{Г} \quad $-4b$\\[0.2cm]
\textbf{Д} \quad $0$
\end{minipage}
\hfill
\begin{minipage}[t]{0.2\textwidth}
\vspace{0.5cm}
\begin{flushright}\matchingGrid\end{flushright}
\end{minipage}

\vspace{0.4cm}
\noindent\zadnum Узгодьте вираз (1--3) з тотожно рівним йому виразом (А--Д), якщо $c \neq 0$.

\vspace{0.3cm}
\noindent
\begin{minipage}[t]{0.42\textwidth}
\textit{Вираз}\par\vspace{0.2cm}
\textbf{1} \quad $\left(\dfrac{c}{5}\right)^{-1}$\\[0.55cm]
\textbf{2} \quad $\sqrt[3]{-125c^3}$\\[0.4cm]
\textbf{3} \quad $\dfrac{(c+5)^2 - c^2 - 25}{2}$
\end{minipage}
\hfill
\begin{minipage}[t]{0.3\textwidth}
\textit{Тотожно рівний вираз}\par\vspace{0.2cm}
\textbf{А} \quad $5c$\\[0.2cm]
\textbf{Б} \quad $-5c$\\[0.2cm]
\textbf{В} \quad $\dfrac{5}{c}$\\[0.25cm]
\textbf{Г} \quad $0$\\[0.2cm]
\textbf{Д} \quad $\dfrac{c}{5}$
\end{minipage}
\hfill
\begin{minipage}[t]{0.2\textwidth}
\vspace{0.5cm}
\begin{flushright}\matchingGrid\end{flushright}
\end{minipage}
\taskBlock{Обчислення значень виразів (відповідність) --- сесія 4.06, завдання №16}
\zadtask{Узгодьте вираз (1--3) із його значенням (А--Д), якщо $a = \dfrac{2}{3}$.

\vspace{0.3cm}
\noindent
\begin{minipage}[t]{0.4\textwidth}
\textit{Вираз}\par\vspace{0.2cm}
\textbf{1} \quad $6a$\\[0.4cm]
\textbf{2} \quad $\sqrt{\left(2^a\right)^3}$\\[0.4cm]
\textbf{3} \quad $\log_3 (1 - a)$
\end{minipage}
\hfill
\begin{minipage}[t]{0.32\textwidth}
\textit{Значення виразу}\par\vspace{0.2cm}
\textbf{А} \quad $4$\\[0.25cm]
\textbf{Б} \quad $-1$\\[0.25cm]
\textbf{В} \quad $1$\\[0.25cm]
\textbf{Г} \quad $\dfrac{20}{3}$\\[0.25cm]
\textbf{Д} \quad $2$
\end{minipage}
\hfill
\begin{minipage}[t]{0.2\textwidth}
\vspace{0.4cm}
\begin{flushright}\matchingGrid\end{flushright}
\end{minipage}
}
\solution{При $a=\dfrac{2}{3}$: $6a=4$ (А); $\sqrt{(2^a)^3}=\sqrt{2^{3a}}=\sqrt{2^2}=2$ (Д); $\log_3(1-a)=\log_3\dfrac{1}{3}=-1$ (Б). Відповідь: \textbf{1~--~А, 2~--~Д, 3~--~Б}.}
\zadtask{Узгодьте вираз (1--3) із його значенням (А--Д), якщо $a = \dfrac{3}{2}$.

\vspace{0.3cm}
\noindent
\begin{minipage}[t]{0.4\textwidth}
\textit{Вираз}\par\vspace{0.2cm}
\textbf{1} \quad $4a$\\[0.4cm]
\textbf{2} \quad $\sqrt{2^{4a}}$\\[0.4cm]
\textbf{3} \quad $\log_2 \left(a + \dfrac{1}{2}\right)$
\end{minipage}
\hfill
\begin{minipage}[t]{0.32\textwidth}
\textit{Значення виразу}\par\vspace{0.2cm}
\textbf{А} \quad $1$\\[0.25cm]
\textbf{Б} \quad $6$\\[0.25cm]
\textbf{В} \quad $8$\\[0.25cm]
\textbf{Г} \quad $3$\\[0.25cm]
\textbf{Д} \quad $-1$
\end{minipage}
\hfill
\begin{minipage}[t]{0.2\textwidth}
\vspace{0.4cm}
\begin{flushright}\matchingGrid\end{flushright}
\end{minipage}
}
\zadtask{Узгодьте вираз (1--3) із його значенням (А--Д), якщо $a = -2$.

\vspace{0.3cm}
\noindent
\begin{minipage}[t]{0.4\textwidth}
\textit{Вираз}\par\vspace{0.2cm}
\textbf{1} \quad $a^3$\\[0.4cm]
\textbf{2} \quad $|a| + a^2$\\[0.4cm]
\textbf{3} \quad $\left(\dfrac{1}{2}\right)^{a}$
\end{minipage}
\hfill
\begin{minipage}[t]{0.32\textwidth}
\textit{Значення виразу}\par\vspace{0.2cm}
\textbf{А} \quad $6$\\[0.25cm]
\textbf{Б} \quad $-8$\\[0.25cm]
\textbf{В} \quad $4$\\[0.25cm]
\textbf{Г} \quad $-2$\\[0.25cm]
\textbf{Д} \quad $2$
\end{minipage}
\hfill
\begin{minipage}[t]{0.2\textwidth}
\vspace{0.4cm}
\begin{flushright}\matchingGrid\end{flushright}
\end{minipage}
}
\taskBlock{Обчислення значень виразів (відповідність) --- сесія 5.06, завдання №16}
\zadtask{Узгодьте вираз (1--3) із його значенням (А--Д).

\vspace{0.3cm}
\noindent
\begin{minipage}[t]{0.4\textwidth}
\textit{Вираз}\par\vspace{0.2cm}
\textbf{1} \quad $\left|\,-6 + \cos 4\pi\,\right|$\\[0.4cm]
\textbf{2} \quad $\sqrt{2}\cdot\sqrt{18}$\\[0.4cm]
\textbf{3} \quad $\log_{1/2} 32$
\end{minipage}
\hfill
\begin{minipage}[t]{0.3\textwidth}
\textit{Значення}\par\vspace{0.2cm}
\textbf{А} \quad $-6$\\[0.15cm]
\textbf{Б} \quad $-5$\\[0.15cm]
\textbf{В} \quad $5$\\[0.15cm]
\textbf{Г} \quad $6$\\[0.15cm]
\textbf{Д} \quad $7$
\end{minipage}
\hfill
\begin{minipage}[t]{0.2\textwidth}
\vspace{0.5cm}
\begin{flushright}\matchingGrid\end{flushright}
\end{minipage}}
\solution{$1$: $|-6+\cos4\pi|=|-6+1|=5$ (В); $2$: $\sqrt{2}\cdot\sqrt{18}=\sqrt{36}=6$ (Г); $3$: $\log_{1/2}32=\log_{1/2}2^5=-5$ (Б). Відповідь: \textbf{1~--~В, 2~--~Г, 3~--~Б}.}
\zadtask{Узгодьте вираз (1--3) із його значенням (А--Д).

\vspace{0.3cm}
\noindent
\begin{minipage}[t]{0.4\textwidth}
\textit{Вираз}\par\vspace{0.2cm}
\textbf{1} \quad $\left|\,-8 + \cos 2\pi\,\right|$\\[0.4cm]
\textbf{2} \quad $\sqrt{3}\cdot\sqrt{27}$\\[0.4cm]
\textbf{3} \quad $\log_{1/3} 81$
\end{minipage}
\hfill
\begin{minipage}[t]{0.3\textwidth}
\textit{Значення}\par\vspace{0.2cm}
\textbf{А} \quad $-7$\\[0.15cm]
\textbf{Б} \quad $-4$\\[0.15cm]
\textbf{В} \quad $7$\\[0.15cm]
\textbf{Г} \quad $9$\\[0.15cm]
\textbf{Д} \quad $-9$
\end{minipage}
\hfill
\begin{minipage}[t]{0.2\textwidth}
\vspace{0.5cm}
\begin{flushright}\matchingGrid\end{flushright}
\end{minipage}}
\zadtask{Узгодьте вираз (1--3) із його значенням (А--Д).

\vspace{0.3cm}
\noindent
\begin{minipage}[t]{0.4\textwidth}
\textit{Вираз}\par\vspace{0.2cm}
\textbf{1} \quad $2^{-3} \cdot 2^{5}$\\[0.4cm]
\textbf{2} \quad $\sqrt[3]{-64}$\\[0.4cm]
\textbf{3} \quad $6\cos\pi$
\end{minipage}
\hfill
\begin{minipage}[t]{0.3\textwidth}
\textit{Значення}\par\vspace{0.2cm}
\textbf{А} \quad $-6$\\[0.15cm]
\textbf{Б} \quad $-4$\\[0.15cm]
\textbf{В} \quad $4$\\[0.15cm]
\textbf{Г} \quad $10$\\[0.15cm]
\textbf{Д} \quad $2$
\end{minipage}
\hfill
\begin{minipage}[t]{0.2\textwidth}
\vspace{0.5cm}
\begin{flushright}\matchingGrid\end{flushright}
\end{minipage}}
\typeTitle{Завдання з короткою відповіддю}
\taskBlock{Сюжетна задача на відсотки і частини --- сесія 23.05, завдання №21}
\zadtask{У магазині є лише музичні диски, диски з науково-популярними фільмами та диски з художніми фільмами. Кількість музичних дисків утричі більша за кількість дисків із науково-популярними фільмами й становить $40\%$ від кількості дисків з художніми фільмами. Загальна кількість дисків $345$. Визначте кількість музичних дисків у магазині. \nmtyear{2026}}
\nmtAnswerBox
\solution{Нехай науково-популярних~--- $x$, тоді музичних~--- $3x$, а художніх~--- $\dfrac{3x}{0{,}4}=7{,}5x$. Сума: $x+3x+7{,}5x=11{,}5x=345$, звідки $x=30$. Музичних: $3x=90$. Відповідь: \textbf{90}.}
\zadtask{У бібліотеці є лише довідники, словники та енциклопедії. Кількість довідників удвічі більша за кількість словників і становить $25\%$ від кількості енциклопедій. Загальна кількість книжок $220$. Визначте кількість довідників у бібліотеці.}
\nmtAnswerBox
\zadtask{У кошику є лише яблука, груші та сливи. Кількість яблук удвічі менша за кількість груш і становить $20\%$ від кількості слив. У кошику $12$ яблук. Визначте загальну кількість фруктів у кошику.}
\nmtAnswerBox
\taskBlock{Задача на відсотки (суміш) --- сесія 30.05, завдання №20}
\zadtask{Для приготування трав'яного чаю змішали м'яту й чебрець. Маса суміші становить $208$ г, причому м'яти в суміші на $60\,\%$ більше, ніж чебрецю. Скільки грамів м'яти у цій суміші?
\nmtAnswerBox}
\solution{Нехай чебрецю $c$ г, тоді м'яти $1{,}6c$ г. Маємо $c+1{,}6c=208$, звідки $2{,}6c=208$, $c=80$. М'яти: $1{,}6\cdot80=128$ г. Відповідь: $128$.}
\zadtask{Для приготування трав'яного збору змішали ромашку й липу. Маса збору становить $252$ г, причому ромашки у зборі на $80\,\%$ більше, ніж липи. Скільки грамів ромашки у цьому зборі?
\nmtAnswerBox}
\zadtask{Для приготування мюслі змішали вівсяні пластівці й горіхи. Маса суміші становить $170$ г, причому горіхів у суміші на $30\,\%$ менше, ніж пластівців. Скільки грамів горіхів у цій суміші?
\nmtAnswerBox}
\taskBlock{Текстова задача на відсотки --- сесія 1.06, завдання №20}
\zadtask{У будинку $162$ квартири~--- одно-, дво- і трикімнатні. Двокімнатних квартир на $15\,\%$ менше, ніж однокімнатних, а трикімнатних~--- стільки ж, скільки двокімнатних. Скільки в будинку \textit{двокімнатних} квартир?}
\nmtAnswerBox
\solution{Нехай однокімнатних $x$, тоді дво- і трикімнатних по $0{,}85x$. Усього $x+0{,}85x+0{,}85x=2{,}7x=162$, звідки $x=60$. Двокімнатних: $0{,}85\cdot 60=51$. Відповідь: \textbf{$51$}.}
\zadtask{У будинку $150$ квартир~--- одно-, дво- і трикімнатні. Двокімнатних квартир на $25\,\%$ менше, ніж однокімнатних, а трикімнатних~--- стільки ж, скільки двокімнатних. Скільки в будинку \textit{двокімнатних} квартир?}
\nmtAnswerBox
\zadtask{У гаражному кооперативі $240$ боксів~--- малі, середні та великі. Малих боксів на $20\,\%$ \textit{більше}, ніж великих, а середніх~--- стільки ж, скільки великих. Скільки в кооперативі \textit{великих} боксів?}
\nmtAnswerBox
\taskBlock{Сюжетна задача на відсотки (вартість таксі) --- сесія 2.06, завдання №20}
\zadtask{Перед подорожжю турист проаналізував інфографіку й дізнався вартість (у грн) поїздки від аеропорту до центру міста (або у зворотному напрямку) у деяких світових столицях. На час туристичного сезону вартість проїзду в таксі в Осло збільшується на $20\,\%$, а в Рейк'явіку~--- на $30\,\%$ від фіксованої вартості проїзду. Визначте сумарну кількість коштів, яку витратить турист на таксі від аеропорту до центру міста й назад у цих містах, якщо відвідає Осло та Рейк'явік у туристичний сезон.

\vspace{0.3cm}
\begin{center}
\begin{tikzpicture}[scale=1.0]
    % фон-рамка
    \draw[gray!40, rounded corners=3pt] (-0.3,-0.4) rectangle (9.2,5.6);
    % дані: місто / значення / довжина бара (масштаб: 5561 -> 4.2)
    \def\scale{0.000755}
    \foreach \i/\city/\val in {0/{аеропорт Лондона}/2502, 1/{аеропорт Монако}/2502, 2/{аеропорт Рейк'явіка}/2780, 3/{аеропорт Осло}/2780, 4/{аеропорт Токіо}/5561} {
        \pgfmathsetmacro\y{\i*1.05}
        \pgfmathsetmacro\barlen{\val*\scale}
        % прапорець-заглушка
        \filldraw[fill=gray!25, draw=gray!50] (0,\y+0.32) rectangle (0.45,\y+0.62);
        % назва
        \node[right, font=\small] at (0.55,\y+0.47) {\city};
        % бар
        \filldraw[fill=cyan!45, draw=cyan!60!black] (0.55,\y) rectangle (0.55+\barlen,\y+0.28);
        \node[right, font=\scriptsize\bfseries, text=white] at (0.7,\y+0.14) {\val};
    }
    % легенда
    \filldraw[fill=cyan!45, draw=cyan!60!black] (6.0,4.55) rectangle (6.35,4.85);
    \node[right, font=\scriptsize] at (6.4,4.7) {Фіксована вартість};
    \node[right, font=\scriptsize] at (6.4,4.4) {проїзду, \textit{грн}};
\end{tikzpicture}
\end{center}
\nmtAnswerBox
}
\solution{Фіксована вартість в Осло й Рейк'явіку по $2780$ грн. У сезон: Осло $2780\cdot1{,}2$, Рейк'явік $2780\cdot1{,}3$, а туди й назад -- подвоюємо: $2\cdot(2780\cdot1{,}2+2780\cdot1{,}3)=2\cdot2780\cdot2{,}5=13900$ грн.}
\zadtask{Перед подорожжю турист проаналізував інфографіку й дізнався вартість (у грн) поїздки від аеропорту до центру міста (або у зворотному напрямку) у деяких європейських столицях. На час туристичного сезону вартість проїзду в таксі в Парижі збільшується на $25\,\%$, а в Римі~--- на $10\,\%$ від фіксованої вартості проїзду. Визначте сумарну кількість коштів, яку витратить турист на таксі від аеропорту до центру міста й назад у цих містах, якщо відвідає Париж та Рим у туристичний сезон.

\vspace{0.3cm}
\begin{center}
\begin{tikzpicture}[scale=1.0]
    % фон-рамка
    \draw[gray!40, rounded corners=3pt] (-0.3,-0.4) rectangle (9.2,5.6);
    \def\scale{0.0016}
    \foreach \i/\city/\val in {0/{аеропорт Праги}/1500, 1/{аеропорт Відня}/1800, 2/{аеропорт Мадрида}/2000, 3/{аеропорт Парижа}/2400, 4/{аеропорт Рима}/2600} {
        \pgfmathsetmacro\y{\i*1.05}
        \pgfmathsetmacro\barlen{\val*\scale}
        \filldraw[fill=gray!25, draw=gray!50] (0,\y+0.32) rectangle (0.45,\y+0.62);
        \node[right, font=\small] at (0.55,\y+0.47) {\city};
        \filldraw[fill=cyan!45, draw=cyan!60!black] (0.55,\y) rectangle (0.55+\barlen,\y+0.28);
        \node[right, font=\scriptsize\bfseries, text=white] at (0.7,\y+0.14) {\val};
    }
    \filldraw[fill=cyan!45, draw=cyan!60!black] (6.0,4.55) rectangle (6.35,4.85);
    \node[right, font=\scriptsize] at (6.4,4.7) {Фіксована вартість};
    \node[right, font=\scriptsize] at (6.4,4.4) {проїзду, \textit{грн}};
\end{tikzpicture}
\end{center}
\nmtAnswerBox
}
\zadtask{Перед подорожжю турист проаналізував інфографіку й дізнався вартість (у грн) поїздки від аеропорту до центру міста (або у зворотному напрямку) у деяких європейських столицях. У міжсезоння вартість проїзду в таксі у Відні \textit{зменшується} на $10\,\%$, а в Мадриді~--- на $20\,\%$ від фіксованої вартості проїзду. Визначте сумарну кількість коштів, яку витратить турист на таксі від аеропорту до центру міста й назад у цих містах, якщо відвідає Відень та Мадрид у міжсезоння.

\vspace{0.3cm}
\begin{center}
\begin{tikzpicture}[scale=1.0]
    % фон-рамка
    \draw[gray!40, rounded corners=3pt] (-0.3,-0.4) rectangle (9.2,5.6);
    \def\scale{0.0016}
    \foreach \i/\city/\val in {0/{аеропорт Праги}/1500, 1/{аеропорт Відня}/1800, 2/{аеропорт Мадрида}/2000, 3/{аеропорт Парижа}/2400, 4/{аеропорт Рима}/2600} {
        \pgfmathsetmacro\y{\i*1.05}
        \pgfmathsetmacro\barlen{\val*\scale}
        \filldraw[fill=gray!25, draw=gray!50] (0,\y+0.32) rectangle (0.45,\y+0.62);
        \node[right, font=\small] at (0.55,\y+0.47) {\city};
        \filldraw[fill=cyan!45, draw=cyan!60!black] (0.55,\y) rectangle (0.55+\barlen,\y+0.28);
        \node[right, font=\scriptsize\bfseries, text=white] at (0.7,\y+0.14) {\val};
    }
    \filldraw[fill=cyan!45, draw=cyan!60!black] (6.0,4.55) rectangle (6.35,4.85);
    \node[right, font=\scriptsize] at (6.4,4.7) {Фіксована вартість};
    \node[right, font=\scriptsize] at (6.4,4.4) {проїзду, \textit{грн}};
\end{tikzpicture}
\end{center}
\nmtAnswerBox
}
\taskBlock{Відсотки і знижки (сюжетна задача) --- сесія 3.06, завдання №20}
\zadtask{Сім'я з трьох осіб~--- батька, мами й сина~--- запланувала сплав на байдарці. Вартість сплаву~--- $800$ грн з особи. Мамі та сину надають знижки: $20\,\%$ і $25\,\%$ відповідно (батько сплачує повну вартість). Скільки гривень заплатить сім'я за сплав?}
\nmtAnswerBox
\solution{Батько платить повну вартість $800$ грн, мама $800\cdot0{,}8 = 640$ грн, син $800\cdot0{,}75 = 600$ грн. Усього $800 + 640 + 600 = 2040$ грн. Відповідь: \textbf{$2040$}.}
\zadtask{Сім'я з трьох осіб~--- батька, мами й доньки~--- запланувала екскурсію до печер. Вартість екскурсії~--- $600$ грн з особи. Мамі та доньці надають знижки: $10\,\%$ і $50\,\%$ відповідно (батько сплачує повну вартість). Скільки гривень заплатить сім'я за екскурсію?}
\nmtAnswerBox
\zadtask{Сім'я з трьох осіб~--- батька, мами й сина~--- купує квитки на виставу. Вартість квитка~--- $900$ грн. Мамі та сину надають знижки: $20\,\%$ і $30\,\%$ відповідно (батько сплачує повну вартість квитка). \textit{На скільки гривень менше} заплатить сім'я порівняно з повною вартістю трьох квитків?}
\nmtAnswerBox
\taskBlock{Текстова задача на відсотки --- сесія 4.06, завдання №19}
\zadtask{Валерія за два дні прочитала $132$ вірші. За перший день вона прочитала на $20\,\%$ віршів більше, ніж за другий. Скільки віршів прочитала Валерія за перший день?
\nmtAnswerBox}
\solution{Нехай за другий день прочитано $x$ віршів, тоді за перший $1{,}2x$. Сума: $x+1{,}2x=132$, $2{,}2x=132$, $x=60$. За перший день $1{,}2\cdot60=72$. Відповідь: \textbf{72}.}
\zadtask{Остап за два дні прочитав $180$ сторінок книжки. За перший день він прочитав на $25\,\%$ сторінок більше, ніж за другий. Скільки сторінок прочитав Остап за перший день?
\nmtAnswerBox}
\zadtask{Соломія за два дні розв'язала $119$ задач. За перший день вона розв'язала на $30\,\%$ задач менше, ніж за другий. Скільки задач розв'язала Соломія за перший день?
\nmtAnswerBox}
\taskBlock{Сюжетна задача на відсотки (акційна покупка) --- сесія 5.06, завдання №19}
\zadtask{Чоловік купив три футболки за $1056$ грн. У магазині діяла акція: першу футболку продавали за повну ціну, другу~--- на $20\,\%$ дешевше, а третю~--- на $40\,\%$ дешевше за повну ціну. Скільки гривень коштувала одна футболка \textit{без} акції?}
\nmtAnswerBox
\solution{Нехай повна ціна футболки $x$ грн. Тоді $x+0{,}8x+0{,}6x=2{,}4x=1056$, звідки $x=440$. Відповідь: \textbf{$440$}.}
\zadtask{Жінка купила три однакові чашки за $891$ грн. У магазині діяла акція: першу чашку продавали за повну ціну, другу~--- на $30\,\%$ дешевше, а третю~--- на $50\,\%$ дешевше за повну ціну. Скільки гривень коштувала одна чашка \textit{без} акції?}
\nmtAnswerBox
\zadtask{Покупець придбав три однакові рюкзаки й заплатив за них $1840$ грн. У магазині діяла акція: перший рюкзак продавали за повну ціну, другий~--- на $25\,\%$ дешевше, а третій~--- на $45\,\%$ дешевше за повну ціну. Скільки гривень \textit{заощадив} покупець завдяки акції порівняно з купівлею трьох рюкзаків за повною ціною?}
\nmtAnswerBox

\chapterTitle{РОЗДІЛ 2. РІВНЯННЯ І НЕРІВНОСТІ}
\typeTitle{Завдання з вибором однієї відповіді (А--Д)}
\taskBlock{Найпростіше лінійне рівняння --- сесія 23.05, завдання №1}
\zadtask{Розв'яжіть рівняння $\dfrac{x}{7} = 9$. \nmtyear{2026}}
\answerTable{$2$}{$63$}{$-63$}{$16$}{$\dfrac{9}{7}$}
\solution{Помноживши обидві частини на $7$: $x=9\cdot 7=63$. Відповідь: \textbf{Б}.}
\zadtask{Розв'яжіть рівняння $\dfrac{x}{6} = 8$.}
\answerTable{$48$}{$14$}{$-48$}{$2$}{$\dfrac{4}{3}$}
\zadtask{Розв'яжіть рівняння $\dfrac{12}{x} = 4$.}
\answerTable{$3$}{$48$}{$\dfrac{1}{3}$}{$-3$}{$8$}
\taskBlock{Текстова задача на складання рівняння --- сесія 23.05, завдання №5}
\zadtask{У коробці $30$ тістечок двох видів: бісквітні й безе. Скільки всього бісквітних тістечок, якщо їх у $5$ разів більше, ніж безе? \nmtyear{2026}}
\answerTable{$24$}{$25$}{$6$}{$15$}{$5$}
\solution{Нехай безе~--- $x$, тоді бісквітних~--- $5x$. Усього $x+5x=6x=30$, звідки $x=5$. Бісквітних: $5x=25$. Відповідь: \textbf{Б}.}
\zadtask{У вазі $36$ цукерок двох видів: шоколадні й карамельні. Скільки всього шоколадних цукерок, якщо їх у $3$ рази більше, ніж карамельних?}
\answerTable{$24$}{$12$}{$27$}{$9$}{$18$}
\zadtask{У гаражі $42$ велосипеди двох видів: гірські й шосейні. Скільки всього шосейних велосипедів, якщо їх на $14$ менше, ніж гірських?}
\answerTable{$28$}{$21$}{$14$}{$7$}{$18$}
\taskBlock{Квадратна нерівність --- сесія 23.05, завдання №11}
\zadtask{Знайдіть найбільший цілий розв'язок нерівності $(x - 4)(4 + x) \leq 8x - 7$. \nmtyear{2026}}
\answerTable{$0$}{$9$}{$-1$}{$10$}{$8$}
\solution{$(x-4)(x+4)\le 8x-7 \Rightarrow x^2-16\le 8x-7 \Rightarrow x^2-8x-9\le 0$. Дискримінант $D=64+36=100$, корені $x=\dfrac{8\pm 10}{2}$, тобто $-1$ і $9$. Отже $-1\le x\le 9$, найбільший цілий~--- $9$. Відповідь: \textbf{Б}.}
\zadtask{Знайдіть найбільший цілий розв'язок нерівності $(x - 3)(3 + x) \leq 4x + 12$.}
\answerTable{$7$}{$-3$}{$8$}{$6$}{$21$}
\zadtask{Знайдіть кількість цілих розв'язків нерівності $(x - 5)(5 + x) < 6x - 9$.}
\answerTable{$8$}{$10$}{$7$}{$9$}{$11$}
\taskBlock{Система рівнянь (показникове + дробове) --- сесія 23.05, завдання №15}
\zadtask{Розв'яжіть систему рівнянь $\begin{cases} 5^x = \sqrt{5}, \\ \dfrac{y + 1}{2x + 3} = \dfrac{1 - 4x}{2}. \end{cases}$ Якщо $(x_0;\, y_0)$~--- розв'язок системи, то $x_0 + y_0 = $ \nmtyear{2026}}
\answerTable{$-2{,}5$}{$-1{,}5$}{$2{,}5$}{$-3{,}5$}{$1{,}5$}
\solution{З першого рівняння $5^x=5^{1/2}$, тож $x_0=0{,}5$. Підставимо: $\dfrac{y+1}{2\cdot 0{,}5+3}=\dfrac{1-4\cdot 0{,}5}{2}$, тобто $\dfrac{y+1}{4}=-\dfrac{1}{2}$, звідки $y+1=-2$, $y_0=-3$. Тоді $x_0+y_0=0{,}5-3=-2{,}5$. Відповідь: \textbf{А}.}
\zadtask{Розв'яжіть систему рівнянь $\begin{cases} 3^x = \sqrt{3}, \\ \dfrac{y - 2}{2x + 1} = \dfrac{3 - 2x}{2}. \end{cases}$ Якщо $(x_0;\, y_0)$~--- розв'язок системи, то $x_0 + y_0 = $}
\answerTable{$3{,}5$}{$-4{,}5$}{$4{,}5$}{$5{,}5$}{$2{,}5$}
\zadtask{Розв'яжіть систему рівнянь $\begin{cases} \log_2 x = 3, \\ \dfrac{y - 1}{x + 2} = \dfrac{x - 6}{2}. \end{cases}$ Якщо $(x_0;\, y_0)$~--- розв'язок системи, то $x_0 - y_0 = $}
\answerTable{$3$}{$-19$}{$19$}{$-3$}{$9$}
\taskBlock{Система рівнянь --- сесія 30.05, завдання №1}
\zadtask{Розв'яжіть систему рівнянь $\begin{cases} \dfrac{2}{x} = -0{,}5 \\[0.2cm] \dfrac{y}{3} = 2 \end{cases}$}
\answerTableTall{$(4;\, 6)$}{$(-4;\, 6)$}{$(-4;\, -6)$}{$(4;\, -6)$}{$\left(-\dfrac{1}{4};\, 6\right)$}
\solution{З першого рівняння $\dfrac{2}{x}=-0{,}5\Rightarrow x=-4$; з другого $\dfrac{y}{3}=2\Rightarrow y=6$. Розв'язок $(-4;\,6)$. Відповідь: \textbf{Б}.}
\zadtask{Розв'яжіть систему рівнянь $\begin{cases} \dfrac{3}{x} = -0{,}25 \\[0.2cm] \dfrac{y}{5} = 2 \end{cases}$}
\answerTableTall{$(12;\, 10)$}{$(-12;\, -10)$}{$(-12;\, 10)$}{$(12;\, -10)$}{$\left(-\dfrac{1}{12};\, 10\right)$}
\zadtask{Розв'яжіть систему рівнянь $\begin{cases} \dfrac{x}{4} = -0{,}5 \\[0.2cm] \dfrac{6}{y} = 2 \end{cases}$}
\answerTableTall{$(2;\, 3)$}{$(-2;\, 3)$}{$(-2;\, -3)$}{$(2;\, -3)$}{$\left(-\dfrac{1}{2};\, 3\right)$}
\taskBlock{Логарифмічна нерівність --- сесія 30.05, завдання №12}
\zadtask{Розв'яжіть нерівність $\log_5(4x - 7) < 2$.}
\answerTableTall
{$\left(\dfrac{7}{4};\, 8\right)$}
{$(8;\, +\infty)$}
{$\left(\dfrac{7}{4};\, +\infty\right)$}
{$(-\infty;\, 8)$}
{$\left(-\infty;\, \dfrac{7}{4}\right)$}
\solution{ОДЗ: $4x-7>0$. Оскільки $2=\log_5 25$, нерівність $4x-7<25$. Разом: $0<4x-7<25\Rightarrow \dfrac74<x<8$. Відповідь: \textbf{А} $\left(\dfrac74;\,8\right)$.}
\zadtask{Розв'яжіть нерівність $\log_3(2x - 5) < 2$.}
\answerTableTall
{$(7;\, +\infty)$}
{$\left(\dfrac{5}{2};\, 7\right)$}
{$\left(\dfrac{5}{2};\, +\infty\right)$}
{$(-\infty;\, 7)$}
{$\left(-\infty;\, \dfrac{5}{2}\right)$}
\zadtask{Розв'яжіть нерівність $\log_{\frac{1}{3}}(2x - 1) \ge -2$.}
\answerTableTall
{$\left(\dfrac{1}{2};\, 5\right]$}
{$[5;\, +\infty)$}
{$\left(\dfrac{1}{2};\, +\infty\right)$}
{$(-\infty;\, 5]$}
{$(5;\, +\infty)$}
\taskBlock{Ірраціональне рівняння --- сесія 30.05, завдання №14}
\zadtask{Розв'яжіть рівняння $\sqrt{12 - 2x} = 6$.}
\answerTable{$-12$}{$-3$}{$3$}{$12$}{$24$}
\solution{Піднесемо до квадрата: $12-2x=36\Rightarrow -2x=24\Rightarrow x=-12$. Перевірка: $\sqrt{12-2\cdot(-12)}=\sqrt{36}=6$. Відповідь: \textbf{А} ($-12$).}
\zadtask{Розв'яжіть рівняння $\sqrt{18 - 3x} = 6$.}
\answerTable{$-18$}{$-6$}{$2$}{$6$}{$18$}
\zadtask{Розв'яжіть рівняння $\sqrt{x + 7} = x - 5$.}
\answerTable{$-2$}{$2$}{$9$}{$2$ і $9$}{$16$}
\taskBlock{Показникове рівняння --- сесія 1.06, завдання №1}
\zadtask{Позначте число, що є коренем рівняння $3 \cdot 3^x = 27$.}
\answerTable{$2$}{$3$}{$4$}{$9$}{$27$}
\solution{$3\cdot 3^x=27 \Rightarrow 3^{x+1}=3^3 \Rightarrow x+1=3 \Rightarrow x=2$. Відповідь: \textbf{А}.}
\zadtask{Позначте число, що є коренем рівняння $2 \cdot 2^x = 32$.}
\answerTable{$3$}{$4$}{$5$}{$16$}{$30$}
\zadtask{Позначте число, що є коренем рівняння $2^{x-1} = \dfrac{1}{8}$.}
\answerTable{$-3$}{$-2$}{$2$}{$3$}{$4$}
\taskBlock{Система лінійних нерівностей --- сесія 1.06, завдання №12}
\zadtask{Розв'яжіть систему нерівностей $\begin{cases} 9 + 3x < 18, \\ 2 - 3x > -10. \end{cases}$}
\answerTable{$(-\infty;\, 3)$}{$(3;\, 4)$}{$(4;\, +\infty)$}{$(3;\, +\infty)$}{$(-\infty;\, 4)$}
\solution{Перша: $3x<9 \Rightarrow x<3$. Друга: $-3x>-12 \Rightarrow x<4$. Перетин: $x<3$, тобто $(-\infty;\,3)$. Відповідь: \textbf{А}.}
\zadtask{Розв'яжіть систему нерівностей $\begin{cases} 7 + 2x < 15, \\ 3 - 2x > -7. \end{cases}$}
\answerTable{$(-\infty;\, 5)$}{$(4;\, 5)$}{$(-\infty;\, 4)$}{$(4;\, +\infty)$}{$(5;\, +\infty)$}
\zadtask{Розв'яжіть систему нерівностей $\begin{cases} 3x - 6 < 6, \\ 2x + 1 > -5. \end{cases}$}
\answerTable{$(-3;\, 4)$}{$(-\infty;\, 4)$}{$(4;\, +\infty)$}{$(-3;\, +\infty)$}{$(-\infty;\, -3)$}
\taskBlock{Ірраціональне рівняння --- сесія 1.06, завдання №14}
\zadtask{Розв'яжіть рівняння $\sqrt{\dfrac{25+2x}{9}} = \dfrac{1}{3}$.}
\answerTable{$-12$}{$-11$}{$-8$}{$0$}{$12$}
\solution{Піднесемо до квадрата: $\dfrac{25+2x}{9}=\dfrac{1}{9}$, звідки $25+2x=1$, $2x=-24$, $x=-12$. Відповідь: \textbf{А}.}
\zadtask{Розв'яжіть рівняння $\sqrt{\dfrac{16+3x}{4}} = \dfrac{1}{2}$.}
\answerTable{$5$}{$0$}{$-4$}{$-5$}{$-6$}
\zadtask{Розв'яжіть рівняння $\sqrt{2x + 7} = 3$.}
\answerTable{$-2$}{$0$}{$1$}{$2$}{$4$}
\taskBlock{Подвійна нерівність --- сесія 2.06, завдання №1}
\zadtask{Розв'яжіть нерівність $-8 \leqslant x - 3 \leqslant 5$.}
\answerTable{$[-11;\, 2]$}{$[-5;\, 8]$}{$[-5;\, 2]$}{$[-8;\, 5]$}{$[-11;\, 8]$}
\solution{Додаємо $3$ до всіх частин: $-8+3\leqslant x\leqslant5+3$, тобто $-5\leqslant x\leqslant8$, або $[-5;\,8]$. Відповідь: \textbf{Б}.}
\zadtask{Розв'яжіть нерівність $-6 \leqslant x - 2 \leqslant 7$.}
\answerTable{$[-8;\, 5]$}{$[-6;\, 7]$}{$[-4;\, 9]$}{$[-4;\, 7]$}{$[-8;\, 9]$}
\zadtask{Розв'яжіть нерівність $-4 \leqslant 1 - x \leqslant 6$.}
\answerTable{$[-5;\, 5]$}{$[-5;\, 7]$}{$[-3;\, 5]$}{$[3;\, 7]$}{$[-7;\, 5]$}
\taskBlock{Логарифмічне рівняння --- сесія 2.06, завдання №10}
\zadtask{Якому проміжку належить корінь рівняння $\dfrac{\log_7(3x-4) - \log_7(x+5)}{2} = 0$?}
\answerTable{$[-1;\, 1)$}{$[1;\, 3)$}{$[3;\, 5)$}{$[5;\, 7)$}{$[7;\, +\infty)$}
\solution{Рівняння рівносильне $\log_7(3x-4)-\log_7(x+5)=0$, тобто $\log_7\dfrac{3x-4}{x+5}=0$. Звідси $3x-4=x+5$, $x=4{,}5\in[3;5)$. Відповідь: \textbf{В}.}
\zadtask{Якому проміжку належить корінь рівняння $\dfrac{\log_3(2x-3) - \log_3(x+4)}{5} = 0$?}
\answerTable{$[-1;\, 1)$}{$[1;\, 3)$}{$[3;\, 5)$}{$[5;\, 7)$}{$[7;\, +\infty)$}
\zadtask{Якому проміжку належить корінь рівняння $\dfrac{\log_5(3x+1)}{\log_5(x+8)} = 1$?}
\answerTable{$[-1;\, 1)$}{$[1;\, 3)$}{$[3;\, 5)$}{$[5;\, 7)$}{$[7;\, +\infty)$}
\taskBlock{Система рівнянь --- сесія 2.06, завдання №13}
\zadtask{Розв'яжіть систему рівнянь $\begin{cases} xy - 2x = 3, \\ xy - 2x + y = 11. \end{cases}$}
\answerTable{$\left(\dfrac{1}{2};\, 8\right)$}{$(2;\, 8)$}{$\left(\dfrac{1}{2};\, -8\right)$}{$\left(-\dfrac{1}{2};\, 8\right)$}{$\left(8;\, \dfrac{1}{2}\right)$}
\solution{Віднявши перше рівняння від другого, дістаємо $y=11-3=8$. Підставляємо у перше: $8x-2x=3$, $6x=3$, $x=\dfrac{1}{2}$. Отже $\left(\dfrac{1}{2};\,8\right)$. Відповідь: \textbf{А}.}
\zadtask{Розв'яжіть систему рівнянь $\begin{cases} xy - 3x = 4, \\ xy - 3x + y = 10. \end{cases}$}
\answerTable{$\left(\dfrac{3}{4};\, 6\right)$}{$\left(\dfrac{4}{3};\, 6\right)$}{$\left(\dfrac{4}{3};\, -6\right)$}{$\left(-\dfrac{4}{3};\, 6\right)$}{$\left(6;\, \dfrac{4}{3}\right)$}
\zadtask{Розв'яжіть систему рівнянь $\begin{cases} xy + 3y = 10, \\ xy + 3y - x = 7. \end{cases}$}
\answerTable{$\left(\dfrac{5}{3};\, 3\right)$}{$\left(3;\, \dfrac{5}{3}\right)$}{$\left(3;\, -\dfrac{5}{3}\right)$}{$\left(-3;\, \dfrac{5}{3}\right)$}{$\left(3;\, \dfrac{3}{5}\right)$}
\taskBlock{Лінійна нерівність --- сесія 3.06, завдання №1}
\zadtask{Розв'яжіть нерівність $2x - 1 \geqslant 0$.}
\answerTableTall{$\left[-\dfrac{1}{2};\, +\infty\right)$}{$\left(-\infty;\, \dfrac{1}{2}\right]$}{$[1;\, +\infty)$}{$\left(-\infty;\, -\dfrac{1}{2}\right]$}{$\left[\dfrac{1}{2};\, +\infty\right)$}
\solution{$2x - 1 \geqslant 0 \Rightarrow 2x \geqslant 1 \Rightarrow x \geqslant \dfrac12$, тобто $\left[\dfrac12;\, +\infty\right)$. Відповідь: \textbf{Д}.}
\zadtask{Розв'яжіть нерівність $3x - 2 \geqslant 0$.}
\answerTableTall{$\left(-\infty;\, \dfrac{2}{3}\right]$}{$\left[\dfrac{2}{3};\, +\infty\right)$}{$\left[-\dfrac{2}{3};\, +\infty\right)$}{$\left(-\infty;\, -\dfrac{2}{3}\right]$}{$\left[\dfrac{3}{2};\, +\infty\right)$}
\zadtask{Розв'яжіть нерівність $5 - 2x > 0$.}
\answerTableTall{$\left(-\infty;\, \dfrac{5}{2}\right)$}{$\left(\dfrac{5}{2};\, +\infty\right)$}{$\left(-\infty;\, -\dfrac{5}{2}\right)$}{$\left(-\dfrac{5}{2};\, +\infty\right)$}{$\left(-\infty;\, \dfrac{5}{2}\right]$}
\taskBlock{Показникове рівняння (проміжок для кореня) --- сесія 3.06, завдання №11}
\zadtask{Якому проміжку належить корінь рівняння $2^{x+1} - 2^x = \dfrac{1}{8}$?}
\answerTable{$(-\infty;\, -5)$}{$[-5;\, -3)$}{$[-3;\, 0)$}{$[0;\, 5)$}{$[5;\, +\infty)$}
\solution{$2^{x+1} - 2^x = 2^x(2-1) = 2^x = \dfrac18 = 2^{-3}$, звідки $x = -3 \in [-3;\, 0)$. Відповідь: \textbf{В}.}
\zadtask{Якому проміжку належить корінь рівняння $3^{x+1} - 3^x = \dfrac{2}{27}$?}
\answerTable{$(-\infty;\, -4)$}{$[-4;\, -2)$}{$[-2;\, 1)$}{$[1;\, 4)$}{$[4;\, +\infty)$}
\zadtask{Якому проміжку належить корінь рівняння $5^{x+1} + 5^x = \dfrac{6}{25}$?}
\answerTable{$(-\infty;\, -3)$}{$[-3;\, 0)$}{$[0;\, 3)$}{$[3;\, 6)$}{$[6;\, +\infty)$}
\taskBlock{Система рівнянь із логарифмом --- сесія 3.06, завдання №15}
\zadtask{Розв'яжіть систему рівнянь $\begin{cases} \log_2 x = 2, \\[0.15cm] \dfrac{y+10}{x+1} = \dfrac{x+2}{2}. \end{cases}$ Якщо $(x_0;\,y_0)$~--- розв'язок системи, то $x_0 + y_0 = {}$}
\answerTable{$18$}{$29$}{$-6$}{$-4$}{$9$}
\solution{Із $\log_2 x = 2$ маємо $x_0 = 4$. Тоді $\dfrac{y+10}{5} = \dfrac{6}{2} = 3$, звідки $y+10 = 15$, $y_0 = 5$. Отже $x_0 + y_0 = 9$. Відповідь: \textbf{Д}.}
\zadtask{Розв'яжіть систему рівнянь $\begin{cases} \log_3 x = 2, \\[0.15cm] \dfrac{y+2}{x+1} = \dfrac{x-4}{2}. \end{cases}$ Якщо $(x_0;\,y_0)$~--- розв'язок системи, то $x_0 + y_0 = {}$}
\answerTable{$11$}{$23$}{$32$}{$9$}{$25$}
\zadtask{Розв'яжіть систему рівнянь $\begin{cases} \log_5 x = 2, \\[0.15cm] \dfrac{y+3}{x-5} = \dfrac{x-15}{4}. \end{cases}$ Якщо $(x_0;\,y_0)$~--- розв'язок системи, то $y_0 - x_0 = {}$}
\answerTable{$22$}{$72$}{$-22$}{$18$}{$28$}
\taskBlock{Система лінійних рівнянь --- сесія 4.06, завдання №4}
\zadtask{Розв'яжіть систему рівнянь $\begin{cases} x + 2y = 8, \\ x - 2y = -4. \end{cases}$}
\answerTable{$(6;\, 1)$}{$(3;\, 2)$}{$(-2;\, 5)$}{$(2;\, 3)$}{$(2;\, 4)$}
\solution{Додаємо рівняння: $2x=4$, $x=2$. Тоді $2+2y=8$, $2y=6$, $y=3$. Отже $(2;\,3)$. Відповідь: \textbf{Г}.}
\zadtask{Розв'яжіть систему рівнянь $\begin{cases} x + 4y = 9, \\ x - 4y = 1. \end{cases}$}
\answerTable{$(5;\, 1)$}{$(1;\, 2)$}{$(5;\, -1)$}{$(4;\, 1)$}{$(1;\, 5)$}
\zadtask{Розв'яжіть систему рівнянь $\begin{cases} y = 2x - 3, \\ 3x + y = 7. \end{cases}$}
\answerTable{$(1;\, 2)$}{$(2;\, 1)$}{$(2;\, -1)$}{$(-2;\, 1)$}{$(3;\, 3)$}
\taskBlock{Логарифмічне рівняння --- сесія 4.06, завдання №7}
\zadtask{Розв'яжіть рівняння $\log_9 (2x - 5) = \dfrac{1}{2}$.}
\answerTable{$-1$}{$-4$}{$43$}{$4$}{$16$}
\solution{За означенням логарифма $2x-5=9^{1/2}=3$, звідки $2x=8$, $x=4$ (ОДЗ $2x-5>0$ виконується). Відповідь: \textbf{Г}.}
\zadtask{Розв'яжіть рівняння $\log_{25} (2x - 1) = \dfrac{1}{2}$.}
\answerTable{$13$}{$3$}{$1$}{$-2$}{$63$}
\zadtask{Розв'яжіть рівняння $\log_x 49 = 2$.}
\answerTable{$7$}{$-7$}{$\pm 7$}{$49$}{$\sqrt{7}$}
\taskBlock{Показникова нерівність --- сесія 4.06, завдання №15}
\zadtask{Розв'яжіть нерівність $3^{x^2 - 2x + 1} \leqslant 1$.}
\answerTable{$1$}{$(-\infty;\, -1]$}{$-1$}{$[1;\, +\infty)$}{$[-1;\, +\infty)$}
\solution{Оскільки $1=3^0$ і основа $3>1$, нерівність рівносильна $x^2-2x+1\leqslant0$, тобто $(x-1)^2\leqslant0$. Це виконується лише при $x=1$. Відповідь: \textbf{А}.}
\zadtask{Розв'яжіть нерівність $5^{x^2 + 4x + 4} \leqslant 1$.}
\answerTable{$2$}{$(-\infty;\, -2]$}{$-2$}{$[2;\, +\infty)$}{$[-2;\, +\infty)$}
\zadtask{Розв'яжіть нерівність $\left(\dfrac{1}{2}\right)^{x^2 - 4x + 4} \geqslant 1$.}
\answerTable{$2$}{$(-\infty;\, 2]$}{$[2;\, +\infty)$}{$-2$}{$(-\infty;\, -2]$}
\taskBlock{Система рівнянь --- сесія 5.06, завдання №4}
\zadtask{Розв'яжіть систему рівнянь $\begin{cases} \dfrac{8}{x} = -2, \\[0.2cm] x + \dfrac{y}{2} = 0. \end{cases}$}
\answerTable{$(-4;\, 8)$}{$(4;\, -8)$}{$(-4;\, -8)$}{$(4;\, 8)$}{$(-2;\, 4)$}
\solution{З першого рівняння $\dfrac{8}{x}=-2\Rightarrow x=-4$. Підставляємо в друге: $-4+\dfrac{y}{2}=0\Rightarrow y=8$. Розв'язок $(-4;\,8)$. Відповідь: \textbf{А}.}
\zadtask{Розв'яжіть систему рівнянь $\begin{cases} \dfrac{6}{x} = -3, \\[0.2cm] x + \dfrac{y}{3} = 0. \end{cases}$}
\answerTable{$(2;\, -6)$}{$(-2;\, 6)$}{$(-2;\, -6)$}{$(2;\, 6)$}{$(-3;\, 9)$}
\zadtask{Розв'яжіть систему рівнянь $\begin{cases} \dfrac{x}{4} = -2, \\[0.2cm] x + \dfrac{y}{4} = 0. \end{cases}$}
\answerTable{$(-8;\, 32)$}{$(8;\, -32)$}{$(-8;\, -32)$}{$(8;\, 32)$}{$(-8;\, 2)$}
\taskBlock{Квадратна нерівність --- сесія 5.06, завдання №11}
\zadtask{Розв'яжіть нерівність $(2x - 3)^2 < 25$.}
\answerTable{$(-1;\, 4)$}{$(-4;\, 1)$}{$(-1;\, +\infty)$}{$(-\infty;\, 4)$}{$(1;\, 4)$}
\solution{$(2x-3)^2<25\iff |2x-3|<5\iff -5<2x-3<5\iff -2<2x<8\iff -1<x<4$. Відповідь: \textbf{А}.}
\zadtask{Розв'яжіть нерівність $(3x - 1)^2 < 16$.}
\answerTable{$\left(-\dfrac{5}{3};\, 1\right)$}{$\left(-1;\, \dfrac{5}{3}\right)$}{$(-1;\, +\infty)$}{$\left(-\infty;\, \dfrac{5}{3}\right)$}{$\left(1;\, \dfrac{5}{3}\right)$}
\zadtask{Розв'яжіть нерівність $(2x + 5)^2 \geqslant 9$.}
\answerTable{$(-4;\, -1)$}{$(-\infty;\, -4] \cup [-1;\, +\infty)$}{$[-4;\, -1]$}{$(-\infty;\, -1]$}{$[-1;\, +\infty)$}
\taskBlock{Показникове рівняння --- сесія 5.06, завдання №13}
\zadtask{Якому проміжку належить корінь рівняння $5^x \cdot 20^x = 1000$?}
\answerTable{$(-\infty;\, 0)$}{$[0;\, 1{,}5)$}{$[1{,}5;\, 2)$}{$[2;\, 4)$}{$[4;\, +\infty)$}
\solution{$5^x\cdot 20^x=(5\cdot 20)^x=100^x=10^{2x}$, а $1000=10^3$. Отже $10^{2x}=10^3\Rightarrow 2x=3\Rightarrow x=1{,}5$, що належить $[1{,}5;\,2)$. Відповідь: \textbf{В}.}
\zadtask{Якому проміжку належить корінь рівняння $2^x \cdot 50^x = 10000$?}
\answerTable{$(-\infty;\, 0)$}{$[0;\, 1{,}5)$}{$[1{,}5;\, 2)$}{$[2;\, 4)$}{$[4;\, +\infty)$}
\zadtask{Якому проміжку належить корінь рівняння $\dfrac{50^x}{5^x} = 1000$?}
\answerTable{$(-\infty;\, 0)$}{$[0;\, 1{,}5)$}{$[1{,}5;\, 2)$}{$[2;\, 4)$}{$[4;\, +\infty)$}
\typeTitle{Завдання з короткою відповіддю}
\taskBlock{Рівняння з параметром (відсутність коренів) --- сесія 23.05, завдання №22}
\zadtask{Знайдіть найбільше значення $a$, за якого рівняння $\dfrac{\sqrt{x^2 - 10x} - \sqrt{36 - 5x}}{x - 2\log_2 a} = 0$ не має коренів. \nmtyear{2026}}
\nmtAnswerBox
\solution{Чисельник нуль при $\sqrt{x^2-10x}=\sqrt{36-5x}$. ОДЗ: $x^2-10x\ge 0$ і $36-5x\ge 0$, тобто $x\le 0$ (бо $x\le 7{,}2$). Піднесемо до квадрата: $x^2-10x=36-5x \Rightarrow x^2-5x-36=0$; $D=25+144=169$, корені $x=\dfrac{5\pm 13}{2}=9$ або $-4$. В ОДЗ лише $x=-4$. Щоб рівняння не мало коренів, цей корінь має обнуляти знаменник: $-4=2\log_2 a$, тобто $\log_2 a=-2$, $a=0{,}25$. Відповідь: \textbf{0,25}.}
\zadtask{Знайдіть найбільше значення $a$, за якого рівняння $\dfrac{\sqrt{x^2 - 14x} - \sqrt{48 - 12x}}{x - 2\log_2 a} = 0$ не має коренів.}
\nmtAnswerBox
\zadtask{Знайдіть найбільше значення $a$, за якого рівняння $\dfrac{\sqrt{x^2 - 2x} - \sqrt{8 - 4x}}{x - 2\log_2 a} = 0$ має рівно один корінь.}
\nmtAnswerBox
\taskBlock{Показникове рівняння з параметром --- сесія 30.05, завдання №22}
\zadtask{Знайдіть \textit{кількість} усіх цілих значень $a$, за кожного з яких рівняння
$$9^x - (a+6) \cdot 3^x + 12a - 2a^2 = 0$$
має два різні корені.
\nmtAnswerBox}
% Розклад: t^2-(a+6)t+(12a-2a^2)=0, D=9(a-2)^2, t1=2a, t2=6-a;
% умови: 2a>0, 6-a>0, 2a≠6-a => a∈(0;6), a≠2 => a∈{1,3,4,5}, кількість 4
\solution{Заміна $t=3^x>0$: $t^2-(a+6)t+(12a-2a^2)=0$. За теоремою Вієта корені $t_1=2a$, $t_2=6-a$ ($t_1t_2=2a(6-a)=12a-2a^2$, $t_1+t_2=a+6$). Для двох різних коренів $x$ потрібні два різні додатні $t$: $2a>0$, $6-a>0$, $2a\neq6-a$, тобто $a\in(0;6)$, $a\neq2$. Цілі: $1,3,4,5$. Відповідь: $4$.}
\zadtask{Знайдіть \textit{кількість} усіх цілих значень $a$, за кожного з яких рівняння
$$9^x - (2a+12) \cdot 3^x + 36a - 3a^2 = 0$$
має два різні корені.
\nmtAnswerBox}
% t1=3a, t2=12-a; умови: a>0, a<12, 3a≠12-a (a≠3) => a∈{1,2,4,...,11}, кількість 10
\zadtask{Знайдіть \textit{кількість} усіх цілих значень $a$ з проміжку $[-5;\, 5]$, за кожного з яких рівняння
$$9^x - (a+4) \cdot 3^x + 4a = 0$$
має \textit{єдиний} корінь.
\nmtAnswerBox}
% Розклад: t^2-(a+4)t+4a=(t-4)(t-a); t=4>0 завжди дає корінь x=log_3 4.
% Єдиний корінь <=> t=a не дає нового кореня: a<=0 або a=4.
% На [-5;5]: a∈{-5,-4,-3,-2,-1,0,4}, кількість 7
\taskBlock{Задача з параметром (два різні корені) --- сесія 1.06, завдання №22}
\zadtask{Знайдіть \textit{суму} всіх цілих значень параметра $a$ з проміжку $(0;\,7)$, за кожного з яких рівняння
$$\dfrac{x^2 + (a-3)x - a + 2}{2x - 3a + 6} = 0$$
має \textit{два різні} корені.}
\nmtAnswerBox
\solution{Чисельник: $x^2+(a-3)x-(a-2)=0$. За теоремою Вієта корені мають суму $3-a$ і добуток $2-a$; підбором (Вієта) це $x_1=1$ і $x_2=2-a$. Для двох різних коренів потрібно $x_1\ne x_2$, тобто $a\ne 1$, і жоден корінь не має обертати знаменник $2x-3a+6$ на нуль; знаменник занулюється коренем $x_2$ при $a=2$. Цілі $a\in(0;7)$: $1,2,3,4,5,6$; вилучаємо $a=1$ і $a=2$. Сума $3+4+5+6=18$. Відповідь: \textbf{$18$}.}
\zadtask{Знайдіть \textit{суму} всіх цілих значень параметра $a$ з проміжку $(0;\,8)$, за кожного з яких рівняння
$$\dfrac{x^2 + (a-4)x - a + 3}{2x - 3a + 9} = 0$$
має \textit{два різні} корені.}
\nmtAnswerBox
\zadtask{Знайдіть \textit{кількість} усіх цілих значень параметра $a$ з проміжку $(0;\,9)$, за кожного з яких рівняння
$$\dfrac{x^2 + (a-5)x - a + 4}{2x - 3a + 12} = 0$$
має \textit{два різні} корені.}
\nmtAnswerBox
\taskBlock{Рівняння з параметром --- сесія 2.06, завдання №22}
\zadtask{Визначте \textit{найменше} значення $a$, за якого рівняння
$$\dfrac{7^{\,3x - 2x^2} - \dfrac{1}{49}}{\sqrt{2x + 5a} - 4} = 0$$
має \textit{два різні} корені.
\nmtAnswerBox
}
\solution{Чисельник $=0$: $7^{3x-2x^2}=7^{-2}$, тобто $3x-2x^2=-2$, або $2x^2-3x-2=0$. За теоремою Вієта корені $x_1=2$, $x_2=-\dfrac{1}{2}$ (бо $x_1+x_2=\dfrac{3}{2}$, $x_1x_2=-1$). Щоб обидва корені були допустимі, потрібно $2x+5a\geqslant0$ у меншому корені $x=-\dfrac{1}{2}$: $-1+5a\geqslant0$, звідки $a\geqslant0{,}2$. Найменше $a=0{,}2$.}
\zadtask{Визначте \textit{найменше} значення $a$, за якого рівняння
$$\dfrac{5^{\,5x - 2x^2} - \dfrac{1}{125}}{\sqrt{3x + 2a} - 3} = 0$$
має \textit{два різні} корені.
\nmtAnswerBox
}
\zadtask{Визначте \textit{найбільше} значення $a$, за якого рівняння
$$\dfrac{3^{\,4x - x^2} - \dfrac{1}{243}}{\sqrt{5x + 2a} - 5} = 0$$
має \textit{рівно один} корінь.
\nmtAnswerBox
}
\taskBlock{Рівняння з параметром і модулем --- сесія 3.06, завдання №22}
\zadtask{Визначте \textit{суму} всіх значень $a$, за кожного з яких рівняння
$$\frac{|4x + 5a - 3| - 3}{(x+2)(x^2 + 9)} = 0$$
має \textit{єдиний} корінь.}
\nmtAnswerBox
\solution{Оскільки $x^2+9>0$, ОДЗ: $x\ne-2$. Рівняння $|4x+5a-3| = 3$ дає $4x+5a-3 = \pm3$, тобто два корені $x_1 = \dfrac{6-5a}{4}$ і $x_2 = \dfrac{-5a}{4}$ ($x_1\ne x_2$). Єдиний корінь буде, коли один із них дорівнює $-2$: $x_1=-2 \Rightarrow a=2{,}8$; $x_2=-2 \Rightarrow a=1{,}6$. Сума: $2{,}8+1{,}6 = 4{,}4$. Відповідь: \textbf{$4{,}4$}.}
\zadtask{Визначте \textit{суму} всіх значень $a$, за кожного з яких рівняння
$$\frac{|3x + 2a - 6| - 6}{(x-1)(x^2 + 4)} = 0$$
має \textit{єдиний} корінь.}
\nmtAnswerBox
\zadtask{Визначте \textit{добуток} усіх значень $a$, за кожного з яких рівняння
$$\frac{|2x + a - 8| - 4}{(x-5)(x^2 + 9)} = 0$$
має \textit{єдиний} корінь.}
\nmtAnswerBox
\taskBlock{Рівняння з параметром --- сесія 4.06, завдання №22}
\zadtask{Визначте \textit{найменше} ціле значення параметра $a$, за якого рівняння
\[ \frac{x + a + \sqrt{x+a} - 12}{2 + \lg^2(2x - a^2)} = 0 \]
має корені.
\nmtAnswerBox}
\solution{Знаменник завжди додатний, тож потрібен нуль чисельника. Нехай $t=\sqrt{x+a}\geqslant0$: $t^2+t-12=0$. За теоремою Вієта $t_1=3$, $t_2=-4$ ($t_1+t_2=-1$, $t_1t_2=-12$); підходить $t=3$, отже $x+a=9$, $x=9-a$. Має існувати $\lg(2x-a^2)$, тобто $2x-a^2>0$: $18-2a-a^2>0$, або $a^2+2a-18<0$. Корені (через дискримінант $D=4+72=76$): $a=-1\pm\sqrt{19}$, тож $-1-\sqrt{19}<a<-1+\sqrt{19}$, тобто приблизно $-5{,}36<a<3{,}36$. Найменше ціле $a=-5$. Відповідь: \textbf{$-5$}.}
\zadtask{Визначте \textit{найменше} ціле значення параметра $a$, за якого рівняння
\[ \frac{x + a + \sqrt{x+a} - 6}{3 + \lg^2(2x - a^2)} = 0 \]
має корені.
\nmtAnswerBox}
\zadtask{Визначте \textit{найбільше} ціле значення параметра $a$, за якого рівняння
\[ \frac{x + a + \sqrt{x+a} - 20}{5 + \lg^2(2x - a^2)} = 0 \]
має корені.
\nmtAnswerBox}
\taskBlock{Рівняння з параметром --- сесія 5.06, завдання №22}
\zadtask{Знайдіть усі значення параметра $a$, за яких рівняння
$$\dfrac{\log_9(2x - 4a + 3) - 2}{\sqrt{a^2 - x^2}} = 0$$
має розв'язок.}
\nmtAnswerBox
\solution{Чисельник $=0$: $\log_9(2x-4a+3)=2\Rightarrow 2x-4a+3=81\Rightarrow x=2a+39$. Знаменник вимагає $a^2-x^2>0$, тобто $-|a|<x<|a|$, що можливо лише при $a<0$. Умова $|2a+39|<-a$ дає систему $a<2a+39<-a$, звідки $-39<a$ і $a<-13$. Відповідь: \textbf{$a\in(-39;\,-13)$}.}
\zadtask{Знайдіть усі значення параметра $a$, за яких рівняння
$$\dfrac{\log_2(2x - 2a + 4) - 4}{\sqrt{a^2 - x^2}} = 0$$
має розв'язок.}
\nmtAnswerBox
\zadtask{Знайдіть усі значення параметра $a$, за яких рівняння
$$\dfrac{\sqrt{2x - 2a + 12} - 4}{\sqrt{a^2 - x^2}} = 0$$
має розв'язок.}
\nmtAnswerBox

\chapterTitle{РОЗДІЛ 3. ФУНКЦІЇ}
\typeTitle{Завдання з вибором однієї відповіді (А--Д)}
\taskBlock{Зсув графіка функції, координатні чверті --- сесія 23.05, завдання №4}
\noindent
\begin{minipage}[c]{0.55\textwidth}
\zadnum На рисунку зображено графік функції $y = f(x)$. У яких координатних чвертях розміщений графік функції $y = f(x) + 3$? \nmtyear{2026}
\end{minipage}%
\hfill
\begin{minipage}[c]{0.43\textwidth}
\centering
\begin{tikzpicture}[scale=0.55, thick]
    \draw[->] (-4,0) -- (4.5,0) node[right] {$x$};
    \draw[->] (0,-3) -- (0,5) node[above] {$y$};
    \node[below] at (-3,-0.1) {\scriptsize $-3$};
    \node[below] at (1,-0.1) {\scriptsize $1$};
    \node[below] at (3,-0.1) {\scriptsize $3$};
    \node[left] at (-0.1,4) {\scriptsize $4$};
    \node[left] at (-0.1,-2) {\scriptsize $-2$};
    \draw[domain=-3:3, smooth, variable=\x, thick]
        plot ({\x},{4 - 0.49*(\x - 0.5)*(\x - 0.5)});
    \fill (-3,-2) circle (2pt);
    \fill (3,0.94) circle (2pt);
    \draw[dashed] (-3,-2) -- (-3,0);
    \draw[dashed] (-3,-2) -- (0,-2);
    \draw[dashed] (3,0.94) -- (3,0);
    \node[above right] at (1.6,2.5) {\scriptsize $f(x)$};
\end{tikzpicture}
\end{minipage}
\par\vspace{0.3cm}
\noindent
\textbf{А)} \ лише у I та II \\[0.15cm]
\textbf{Б)} \ лише у III та IV \\[0.15cm]
\textbf{В)} \ лише у I та III \\[0.15cm]
\textbf{Г)} \ лише у II та IV \\[0.15cm]
\textbf{Д)} \ в усіх чотирьох
\par\vspace{0.2cm}

\solution{Графік $y=f(x)$ має значення від $-2$ до $4$. Зсув на $+3$ угору дає значення від $1$ до $7$~--- усі додатні. Оскільки $x$ набуває як від'ємних, так і додатних значень, графік лежить лише в I та II чвертях. Відповідь: \textbf{А}.}

\noindent
\begin{minipage}[c]{0.55\textwidth}
\zadnum На рисунку зображено графік функції $y = f(x)$. У яких координатних чвертях розміщений графік функції $y = f(x) - 6$?
\end{minipage}%
\hfill
\begin{minipage}[c]{0.43\textwidth}
\centering
\begin{tikzpicture}[scale=0.55, thick]
    \draw[->] (-4,0) -- (4.5,0) node[right] {$x$};
    \draw[->] (0,-1.5) -- (0,6) node[above] {$y$};
    \node[below] at (-3,-0.1) {\scriptsize $-3$};
    \node[below] at (3,-0.1) {\scriptsize $3$};
    \node[left] at (-0.1,5) {\scriptsize $5$};
    \node[left] at (-0.1,1) {\scriptsize $1$};
    \draw[domain=-3:3, smooth, variable=\x, thick]
        plot ({\x},{5 - 0.444*\x*\x});
    \fill (-3,1) circle (2pt);
    \fill (3,1) circle (2pt);
    \draw[dashed] (-3,1) -- (-3,0);
    \draw[dashed] (3,1) -- (3,0);
    \draw[dashed] (-3,1) -- (0,1);
    \node[above right] at (1.4,3.2) {\scriptsize $f(x)$};
\end{tikzpicture}
\end{minipage}
\par\vspace{0.3cm}
\noindent
\textbf{А)} \ лише у I та II \\[0.15cm]
\textbf{Б)} \ лише у III та IV \\[0.15cm]
\textbf{В)} \ лише у I та III \\[0.15cm]
\textbf{Г)} \ лише у II та IV \\[0.15cm]
\textbf{Д)} \ в усіх чотирьох
\par\vspace{0.2cm}

\noindent
\begin{minipage}[c]{0.55\textwidth}
\zadnum На рисунку зображено графік функції $y = f(x)$. У яких координатних чвертях розміщений графік функції $y = f(x) + 6$?
\end{minipage}%
\hfill
\begin{minipage}[c]{0.43\textwidth}
\centering
\begin{tikzpicture}[scale=0.55, thick]
    \draw[->] (-4,0) -- (4.5,0) node[right] {$x$};
    \draw[->] (0,-5) -- (0,1.5) node[above] {$y$};
    \node[below] at (-3,-0.1) {\scriptsize $-3$};
    \node[below] at (3,-0.1) {\scriptsize $3$};
    \node[left] at (-0.1,-4) {\scriptsize $-4$};
    \node[left] at (-0.1,-1) {\scriptsize $-1$};
    \draw[domain=-3:3, smooth, variable=\x, thick]
        plot ({\x},{-4 + 0.333*\x*\x});
    \fill (-3,-1) circle (2pt);
    \fill (3,-1) circle (2pt);
    \draw[dashed] (-3,-1) -- (-3,0);
    \draw[dashed] (3,-1) -- (3,0);
    \draw[dashed] (-3,-1) -- (0,-1);
    \node[below right] at (1.2,-2.6) {\scriptsize $f(x)$};
\end{tikzpicture}
\end{minipage}
\par\vspace{0.3cm}
\noindent
\textbf{А)} \ лише у I та II \\[0.15cm]
\textbf{Б)} \ лише у III та IV \\[0.15cm]
\textbf{В)} \ лише у I та III \\[0.15cm]
\textbf{Г)} \ лише у II та IV \\[0.15cm]
\textbf{Д)} \ в усіх чотирьох
\par\vspace{0.2cm}
\taskBlock{Член послідовності за формулою --- сесія 23.05, завдання №6}
\zadtask{Послідовність задано формулою $a_n = (-2)^n - 7$. Визначте її четвертий член. \nmtyear{2026}}
\answerTable{$-23$}{$23$}{$9$}{$-15$}{$1$}
\solution{Підставимо $n=4$: $a_4=(-2)^4-7=16-7=9$. Відповідь: \textbf{В}.}
\zadtask{Послідовність задано формулою $a_n = (-3)^n + 5$. Визначте її третій член.}
\answerTable{$32$}{$-22$}{$22$}{$-4$}{$14$}
\zadtask{Послідовність задано формулою $a_n = \dfrac{n^2 - 1}{n + 3}$. Визначте її п'ятий член.}
\answerTable{$3$}{$\dfrac{25}{8}$}{$4$}{$\dfrac{8}{3}$}{$2$}
\taskBlock{Паралельність графіків лінійних функцій --- сесія 30.05, завдання №8}
\noindent
\begin{minipage}[c]{0.55\textwidth}
\zadnum На рисунку зображено графік функції $y = -2x$. Графік якої з наведених функцій паралельний графіку цієї функції?
\end{minipage}%
\hfill
\begin{minipage}[c]{0.43\textwidth}
\centering
\begin{tikzpicture}[scale=0.45, thick]
    \draw[gray!30, very thin] (-4.5,-4.5) grid (4.5,4.5);
    \draw[->, thick] (-4.7,0) -- (4.7,0) node[right] {\small $x$};
    \draw[->, thick] (0,-4.7) -- (0,4.7) node[above] {\small $y$};
    \foreach \x in {-4,-3,-2,-1,1,2,3,4} {
        \draw[thick] (\x,-0.1) -- (\x,0.1);
        \node[below] at (\x,-0.1) {\scriptsize $\x$};
    }
    \foreach \y in {-4,-3,-2,-1,1,2,3,4} {
        \draw[thick] (-0.1,\y) -- (0.1,\y);
        \node[left] at (-0.1,\y) {\scriptsize $\y$};
    }
    \node[below left] at (0,0) {\scriptsize $0$};
    \draw[thick, mainGreen, line width=1.2pt] (-2.2,4.4) -- (2.2,-4.4);
    \node[mainGreen, fill=white, inner sep=1pt] at (2.6,-3.5) {\scriptsize $y = -2x$};
\end{tikzpicture}
\end{minipage}
\par\vspace{0.2cm}
\answerTable{$y = 2x + 1$}{$y = 1 - 2x$}{$y = -x - 2$}{$y = \dfrac{x}{2} - 1$}{$y = -\dfrac{1}{2}x$}
\solution{Прямі паралельні, коли кутові коефіцієнти рівні. У $y=-2x$ маємо $k=-2$, тож шукаємо функцію з $k=-2$: це $y=1-2x$. Відповідь: \textbf{Б}.}
\noindent
\begin{minipage}[c]{0.55\textwidth}
\zadnum На рисунку зображено графік функції $y = 3x$. Графік якої з наведених функцій паралельний графіку цієї функції?
\end{minipage}%
\hfill
\begin{minipage}[c]{0.43\textwidth}
\centering
\begin{tikzpicture}[scale=0.45, thick]
    \draw[gray!30, very thin] (-4.5,-4.5) grid (4.5,4.5);
    \draw[->, thick] (-4.7,0) -- (4.7,0) node[right] {\small $x$};
    \draw[->, thick] (0,-4.7) -- (0,4.7) node[above] {\small $y$};
    \foreach \x in {-4,-3,-2,-1,1,2,3,4} {
        \draw[thick] (\x,-0.1) -- (\x,0.1);
        \node[below] at (\x,-0.1) {\scriptsize $\x$};
    }
    \foreach \y in {-4,-3,-2,-1,1,2,3,4} {
        \draw[thick] (-0.1,\y) -- (0.1,\y);
        \node[left] at (-0.1,\y) {\scriptsize $\y$};
    }
    \node[below left] at (0,0) {\scriptsize $0$};
    \draw[thick, mainGreen, line width=1.2pt] (-1.45,-4.35) -- (1.45,4.35);
    \node[mainGreen, fill=white, inner sep=1pt] at (2.6,3.5) {\scriptsize $y = 3x$};
\end{tikzpicture}
\end{minipage}
\par\vspace{0.2cm}
\answerTable{$y = -3x$}{$y = \dfrac{x}{3} + 2$}{$y = 3x - 2$}{$y = -\dfrac{x}{3}$}{$y = x + 3$}
\noindent
\begin{minipage}[c]{0.55\textwidth}
\zadnum На рисунку зображено графік функції $y = \dfrac{x}{2}$. Графік якої з наведених функцій \textit{перпендикулярний} до графіка цієї функції?
\end{minipage}%
\hfill
\begin{minipage}[c]{0.43\textwidth}
\centering
\begin{tikzpicture}[scale=0.45, thick]
    \draw[gray!30, very thin] (-4.5,-4.5) grid (4.5,4.5);
    \draw[->, thick] (-4.7,0) -- (4.7,0) node[right] {\small $x$};
    \draw[->, thick] (0,-4.7) -- (0,4.7) node[above] {\small $y$};
    \foreach \x in {-4,-3,-2,-1,1,2,3,4} {
        \draw[thick] (\x,-0.1) -- (\x,0.1);
        \node[below] at (\x,-0.1) {\scriptsize $\x$};
    }
    \foreach \y in {-4,-3,-2,-1,1,2,3,4} {
        \draw[thick] (-0.1,\y) -- (0.1,\y);
        \node[left] at (-0.1,\y) {\scriptsize $\y$};
    }
    \node[below left] at (0,0) {\scriptsize $0$};
    \draw[thick, mainGreen, line width=1.2pt] (-4.4,-2.2) -- (4.4,2.2);
    \node[mainGreen, fill=white, inner sep=1pt] at (3.3,2.9) {\scriptsize $y = \frac{x}{2}$};
\end{tikzpicture}
\end{minipage}
\par\vspace{0.2cm}
\answerTable{$y = 2x + 1$}{$y = -\dfrac{x}{2}$}{$y = 2 - 2x$}{$y = \dfrac{x}{2} + 3$}{$y = -2x - 1$}
\taskBlock{Порівняння інтегралів --- сесія 30.05, завдання №11}
\zadtask{Укажіть інтеграл, який має \textit{найменше} значення.

\par\vspace{0.2cm}
\begin{flushleft}
\textbf{А} \quad $\displaystyle\int_0^1 \dfrac{1}{4}\,dx$ \\[0.25cm]
\textbf{Б} \quad $\displaystyle\int_0^1 \dfrac{1}{3}\,dx$ \\[0.25cm]
\textbf{В} \quad $\displaystyle\int_0^1 \dfrac{1}{2}\,dx$ \\[0.25cm]
\textbf{Г} \quad $\displaystyle\int_0^1 1\,dx$ \\[0.25cm]
\textbf{Д} \quad $\displaystyle\int_0^1 \sqrt{2}\,dx$
\end{flushleft}
}
\solution{Для сталої $c$ маємо $\int_0^1 c\,dx=c$. Тож значення дорівнюють самим сталим; найменше з $\dfrac14,\dfrac13,\dfrac12,1,\sqrt2$ — це $\dfrac14$. Відповідь: \textbf{А}.}
\zadtask{Укажіть інтеграл, який має \textit{найбільше} значення.

\par\vspace{0.2cm}
\begin{flushleft}
\textbf{А} \quad $\displaystyle\int_0^1 \dfrac{1}{3}\,dx$ \\[0.25cm]
\textbf{Б} \quad $\displaystyle\int_0^1 \dfrac{1}{2}\,dx$ \\[0.25cm]
\textbf{В} \quad $\displaystyle\int_0^1 \sqrt{2}\,dx$ \\[0.25cm]
\textbf{Г} \quad $\displaystyle\int_0^1 1\,dx$ \\[0.25cm]
\textbf{Д} \quad $\displaystyle\int_0^1 \dfrac{2}{5}\,dx$
\end{flushleft}
}
\zadtask{Укажіть інтеграл, який має \textit{найбільше} значення.

\par\vspace{0.2cm}
\begin{flushleft}
\textbf{А} \quad $\displaystyle\int_0^2 0{,}3\,dx$ \\[0.25cm]
\textbf{Б} \quad $\displaystyle\int_0^2 \dfrac{1}{4}\,dx$ \\[0.25cm]
\textbf{В} \quad $\displaystyle\int_0^2 \dfrac{2}{5}\,dx$ \\[0.25cm]
\textbf{Г} \quad $\displaystyle\int_0^2 \dfrac{1}{2}\,dx$ \\[0.25cm]
\textbf{Д} \quad $\displaystyle\int_0^2 1\,dx$
\end{flushleft}
}
\taskBlock{Арифметична прогресія --- сесія 1.06, завдання №5}
\zadtask{В арифметичній прогресії один із членів дорівнює $27$, а різниця прогресії дорівнює $5$. Яке з наведених чисел може бути членом цієї прогресії?}
\answerTable{$49$}{$50$}{$51$}{$52$}{$53$}
\solution{Члени прогресії відрізняються від $27$ на кратне $5$, тобто дають остачу $2$ при діленні на $5$. Серед варіантів лише $52=27+5\cdot 5$. Відповідь: \textbf{Г}.}
\zadtask{В арифметичній прогресії один із членів дорівнює $23$, а різниця прогресії дорівнює $4$. Яке з наведених чисел може бути членом цієї прогресії?}
\answerTable{$45$}{$46$}{$47$}{$48$}{$49$}
\zadtask{У геометричній прогресії один із членів дорівнює $16$, а знаменник прогресії дорівнює $2$. Яке з наведених чисел може бути членом цієї прогресії?}
\answerTable{$24$}{$48$}{$64$}{$96$}{$100$}
\taskBlock{Розпізнавання графіка квадратичної функції --- сесія 1.06, завдання №8}
\zadtask{Графік якої з наведених функцій зображено на рисунку?

\vspace{0.2cm}
\noindent
\begin{minipage}[c]{0.58\textwidth}
\noindent
\textbf{А} \quad $y = (x+1)^2$\\[0.1cm]
\textbf{Б} \quad $y = -(x-1)^2$\\[0.1cm]
\textbf{В} \quad $y = -x^2 + 1$\\[0.1cm]
\textbf{Г} \quad $y = x^2 - 1$\\[0.1cm]
\textbf{Д} \quad $y = -x^2 - 1$
\end{minipage}%
\hfill
\begin{minipage}[c]{0.40\textwidth}
\centering
\begin{tikzpicture}[scale=0.7, thick]
    \draw[->] (-2.3,0) -- (2.3,0) node[right] {\small $x$};
    \draw[->] (0,-2.6) -- (0,1.6) node[above] {\small $y$};
    \draw (1,-0.08) -- (1,0.08) node[above=1pt] {\scriptsize $1$};
    \draw (-0.08,1) -- (0.08,1) node[right=1pt] {\scriptsize $1$};
    \node[below left] at (0,0) {\scriptsize $0$};
    \draw[mainGreen, line width=1.2pt, domain=-1.55:1.55, samples=80] plot (\x, {-(\x)^2 - 1});
\end{tikzpicture}
\end{minipage}
\par\vspace{0.2cm}
}
\solution{Парабола вітками вниз (старший коефіцієнт від'ємний) з вершиною в точці $(0;-1)$. Цьому відповідає $y=-x^2-1$. Відповідь: \textbf{Д}.}
\zadtask{Графік якої з наведених функцій зображено на рисунку?

\vspace{0.2cm}
\noindent
\begin{minipage}[c]{0.58\textwidth}
\noindent
\textbf{А} \quad $y = (x-1)^2$\\[0.1cm]
\textbf{Б} \quad $y = x^2 + 1$\\[0.1cm]
\textbf{В} \quad $y = -x^2 + 1$\\[0.1cm]
\textbf{Г} \quad $y = -(x+1)^2$\\[0.1cm]
\textbf{Д} \quad $y = x^2 - 1$
\end{minipage}%
\hfill
\begin{minipage}[c]{0.40\textwidth}
\centering
\begin{tikzpicture}[scale=0.7, thick]
    \draw[->] (-2.3,0) -- (2.3,0) node[right] {\small $x$};
    \draw[->] (0,-1.6) -- (0,2.0) node[above] {\small $y$};
    \draw (1,-0.08) -- (1,0.08) node[below=1pt] {\scriptsize $1$};
    \draw (-0.08,1) -- (0.08,1) node[left=1pt] {\scriptsize $1$};
    \node[below left] at (0,0) {\scriptsize $0$};
    \draw[mainGreen, line width=1.2pt, domain=-1.55:1.55, samples=80] plot (\x, {-(\x)^2 + 1});
\end{tikzpicture}
\end{minipage}
\par\vspace{0.2cm}
}
\zadtask{Графік якої з наведених функцій зображено на рисунку?

\vspace{0.2cm}
\noindent
\begin{minipage}[c]{0.58\textwidth}
\noindent
\textbf{А} \quad $y = x^2 - 1$\\[0.1cm]
\textbf{Б} \quad $y = (x-1)^2$\\[0.1cm]
\textbf{В} \quad $y = -x^2 + 1$\\[0.1cm]
\textbf{Г} \quad $y = x^2 + 1$\\[0.1cm]
\textbf{Д} \quad $y = -x^2 - 1$
\end{minipage}%
\hfill
\begin{minipage}[c]{0.40\textwidth}
\centering
\begin{tikzpicture}[scale=0.7, thick]
    \draw[->] (-2.3,0) -- (2.3,0) node[right] {\small $x$};
    \draw[->] (0,-1.6) -- (0,2.6) node[above] {\small $y$};
    \draw (1,-0.08) -- (1,0.08) node[below right=0pt] {\scriptsize $1$};
    \draw (-0.08,1) -- (0.08,1) node[left=1pt] {\scriptsize $1$};
    \node[below left] at (0,0) {\scriptsize $0$};
    \draw[mainGreen, line width=1.2pt, domain=-1.8:1.8, samples=80] plot (\x, {(\x)^2 - 1});
\end{tikzpicture}
\end{minipage}
\par\vspace{0.2cm}
}
\taskBlock{Порівняння значень функції --- сесія 2.06, завдання №8}
\zadtask{Задано функцію $f(x) = x^2$. Укажіть правильну нерівність.}
\answerTable{$f(4) < 0$}{$f(-4) < f(4)$}{$f(-1) > f(2)$}{$f(-1) < 0$}{$f(3) < f(4)$}
\solution{Для $f(x)=x^2$ маємо $f(3)=9$, $f(4)=16$, тож $f(3)<f(4)$ -- правильно. Відповідь: \textbf{Д}.}
\zadtask{Задано функцію $f(x) = x^3$. Укажіть правильну нерівність.}
\answerTable{$f(2) < 0$}{$f(-1) > f(1)$}{$f(3) < f(-3)$}{$f(-2) < f(1)$}{$f(0) > f(1)$}
\zadtask{Задано функцію $f(x) = |x|$. Укажіть правильну нерівність.}
\answerTable{$f(-3) < 0$}{$f(-5) > f(2)$}{$f(0) > f(1)$}{$f(4) < f(-4)$}{$f(-2) < f(2)$}
\taskBlock{Сума арифметичної прогресії --- сесія 2.06, завдання №11}
\zadtask{Задано арифметичну прогресію $(a_n)$: $-2;\ -1{,}6;\ -1{,}2;\ \ldots$\ Знайдіть суму перших $21$ членів цієї прогресії.}
\answerTable{$42$}{$-42$}{$21$}{$84$}{$4{,}2$}
\solution{$a_1=-2$, $d=0{,}4$, тож $a_{21}=-2+20\cdot0{,}4=6$. Сума $S_{21}=\dfrac{a_1+a_{21}}{2}\cdot21=\dfrac{-2+6}{2}\cdot21=42$. Відповідь: \textbf{А}.}
\zadtask{Задано арифметичну прогресію $(a_n)$: $-3;\ -2{,}5;\ -2;\ \ldots$\ Знайдіть суму перших $25$ членів цієї прогресії.}
\answerTable{$-75$}{$37{,}5$}{$75$}{$150$}{$7{,}5$}
\zadtask{Задано арифметичну прогресію $(a_n)$: $4;\ 3{,}4;\ 2{,}8;\ \ldots$\ Знайдіть суму перших $11$ членів цієї прогресії.}
\answerTable{$-11$}{$11$}{$22$}{$5{,}5$}{$110$}
\taskBlock{Властивості непарної функції за графіком --- сесія 3.06, завдання №8}
\noindent
\begin{minipage}[c]{0.56\textwidth}
\zadnum У прямокутній системі координат зображено фрагмент графіка \textit{непарної} функції $y = f(x)$ та точки $K$, $L$, $M$, $N$, $P$. Через яку з цих точок проходить графік функції $y = f(x)$?
\end{minipage}%
\hfill
\begin{minipage}[c]{0.40\textwidth}
\centering
\begin{tikzpicture}[scale=0.5, thick]
    \draw[gray!35, very thin] (-3.3,-3.3) grid (3.3,3.3);
    \draw[->] (-3.5,0) -- (3.5,0) node[right, font=\scriptsize] {$x$};
    \draw[->] (0,-3.5) -- (0,3.5) node[above, font=\scriptsize] {$y$};
    \node[below left, font=\tiny] at (0,0) {$0$};
    \node[font=\tiny] at (1,-0.35) {$1$};
    \node[font=\tiny] at (-0.35,1) {$1$};
    % фрагмент у I чверті: відрізок від (1,1) до (2,3) приблизно
    \draw[mainGreen, line width=1.2pt] (1,1) -- (2,3);
    % точки
    \fill (-3,2) circle (2.4pt); \node[left, font=\tiny] at (-3,2) {$L$};
    \fill (-2.6,1.4) circle (2.4pt); \node[left, font=\tiny] at (-2.6,1.4) {$K$};
    \fill (-2,-3) circle (2.4pt); \node[left, font=\tiny] at (-2,-3) {$M$};
    \fill (2,-2.4) circle (2.4pt); \node[right, font=\tiny] at (2,-2.4) {$N$};
    \fill (-3,-3.1) circle (2.4pt); \node[left, font=\tiny] at (-3,-3.1) {$P$};
\end{tikzpicture}
\end{minipage}
\par\vspace{0.2cm}
\answerTable{$K$}{$L$}{$M$}{$N$}{$P$}
\solution{Графік непарної функції симетричний відносно початку координат. Точка, симетрична до точки $(2;\,3)$ цього фрагмента,~--- це $(-2;\,-3)$, тобто точка $M$. Відповідь: \textbf{В} ($M$).}
\noindent
\begin{minipage}[c]{0.56\textwidth}
\zadnum У прямокутній системі координат зображено фрагмент графіка \textit{непарної} функції $y = f(x)$ та точки $K$, $L$, $M$, $N$, $P$. Через яку з цих точок проходить графік функції $y = f(x)$?
\end{minipage}%
\hfill
\begin{minipage}[c]{0.40\textwidth}
\centering
\begin{tikzpicture}[scale=0.5, thick]
    \draw[gray!35, very thin] (-3.3,-3.3) grid (3.3,3.3);
    \draw[->] (-3.5,0) -- (3.5,0) node[right, font=\scriptsize] {$x$};
    \draw[->] (0,-3.5) -- (0,3.5) node[above, font=\scriptsize] {$y$};
    \node[below left, font=\tiny] at (0,0) {$0$};
    \node[font=\tiny] at (1,-0.35) {$1$};
    \node[font=\tiny] at (-0.35,1) {$1$};
    % фрагмент у I чверті: відрізок від (1,3) до (3,1)
    \draw[mainGreen, line width=1.2pt] (1,3) -- (3,1);
    % точки
    \fill (-3,1) circle (2.4pt); \node[left, font=\tiny] at (-3,1) {$K$};
    \fill (-1,-1) circle (2.4pt); \node[below left, font=\tiny] at (-1,-1) {$L$};
    \fill (-2,-2) circle (2.4pt); \node[left, font=\tiny] at (-2,-2) {$M$};
    \fill (2,-2) circle (2.4pt); \node[right, font=\tiny] at (2,-2) {$N$};
    \fill (-3,-3.1) circle (2.4pt); \node[left, font=\tiny] at (-3,-3.1) {$P$};
\end{tikzpicture}
\end{minipage}
\par\vspace{0.2cm}
\answerTable{$K$}{$L$}{$M$}{$N$}{$P$}
\noindent
\begin{minipage}[c]{0.56\textwidth}
\zadnum У прямокутній системі координат зображено фрагмент графіка \textit{парної} функції $y = f(x)$ та точки $K$, $L$, $M$, $N$, $P$. Через яку з цих точок проходить графік функції $y = f(x)$?
\end{minipage}%
\hfill
\begin{minipage}[c]{0.40\textwidth}
\centering
\begin{tikzpicture}[scale=0.5, thick]
    \draw[gray!35, very thin] (-3.3,-3.3) grid (3.3,3.3);
    \draw[->] (-3.5,0) -- (3.5,0) node[right, font=\scriptsize] {$x$};
    \draw[->] (0,-3.5) -- (0,3.5) node[above, font=\scriptsize] {$y$};
    \node[below left, font=\tiny] at (0,0) {$0$};
    \node[font=\tiny] at (1,-0.35) {$1$};
    \node[font=\tiny] at (-0.35,1) {$1$};
    % фрагмент у II чверті: відрізок від (-3,3) до (-1,1)
    \draw[mainGreen, line width=1.2pt] (-3,3) -- (-1,1);
    % точки
    \fill (2,2) circle (2.4pt); \node[right, font=\tiny] at (2,2) {$K$};
    \fill (2,-2) circle (2.4pt); \node[right, font=\tiny] at (2,-2) {$L$};
    \fill (-2,-2) circle (2.4pt); \node[left, font=\tiny] at (-2,-2) {$M$};
    \fill (1,2) circle (2.4pt); \node[above right, font=\tiny] at (1,2) {$N$};
    \fill (3,-3) circle (2.4pt); \node[right, font=\tiny] at (3,-3) {$P$};
\end{tikzpicture}
\end{minipage}
\par\vspace{0.2cm}
\answerTable{$K$}{$L$}{$M$}{$N$}{$P$}
\taskBlock{Первісна функції --- сесія 3.06, завдання №10}
\zadtask{Визначте загальний вигляд первісних для функції $f(x) = 3x^2$.}
\answerTable{$6x^3 + C$}{$6x + C$}{$3x^3 + C$}{$x^2 + C$}{$x^3 + C$}
\solution{Первісна для $f(x) = 3x^2$ має вигляд $F(x) = 3\cdot\dfrac{x^3}{3} + C = x^3 + C$ (перевірка: $(x^3+C)' = 3x^2$). Відповідь: \textbf{Д} ($x^3 + C$).}
\zadtask{Визначте загальний вигляд первісних для функції $f(x) = 4x^3$.}
\answerTable{$12x^2 + C$}{$x^4 + C$}{$4x^4 + C$}{$x^3 + C$}{$12x^4 + C$}
\zadtask{Для функції $f(x) = 6x^2$ визначте первісну $F(x)$, графік якої проходить через точку $(1;\, 5)$.}
\answerTable{$2x^3 + 3$}{$2x^3 + 5$}{$2x^3$}{$x^3 + 4$}{$12x - 7$}
\taskBlock{Точка мінімуму за графіком функції --- сесія 4.06, завдання №8}
\noindent
\begin{minipage}[c]{0.52\textwidth}
\zadnum На рисунку зображено графік функції $y = f(x)$, визначеної на проміжку $[-4;\, 4]$. Укажіть точку мінімуму цієї функції.
\end{minipage}%
\hfill
\begin{minipage}[c]{0.46\textwidth}
\centering
\begin{tikzpicture}[scale=0.55, thick]
    % Сітка
    \draw[gray!40, thin] (-4.5,-4.5) grid (6.5,6.5);

    % Осі
    \draw[->, thick] (-4.5,0) -- (4.8,0) node[below] {$x$};
    \draw[->, thick] (0,-2.5) -- (0,4.8) node[left] {$y$};

    % Підписи осей
    \node[below left] at (0,0) {$0$};
    \node[below] at (1,0) {$1$};
    \node[below] at (-4,0) {$-4$};
    \node[below] at (4,0) {$4$};
    \node[left] at (0,1) {$1$};

    % Графік функції (гладка крива по ключових точках)
    \draw[mainGreen!80!black, very thick] plot[smooth, tension=0.45] coordinates {
        (-4,-2) (-2.3,5) (0,3) (1.5,1) (3,-1) (4,0)
    };

    % Точки на кінцях
    \fill[mainGreen!80!black] (-4,-2) circle (3pt);
    \fill[mainGreen!80!black] (4,0) circle (3pt);

    % Підпис графіка у білому прямокутнику
    \node[fill=white, inner sep=1pt] at (4, 3.8) {$y = f(x)$};
\end{tikzpicture}
\end{minipage}
\par\vspace{0.4cm}
\answerTable{$-2$}{$1$}{$3$}{$2$}{$5$}
\solution{Точка мінімуму~--- абсциса, у якій функція переходить від спадання до зростання (локальна «впадина»). На графіку це відбувається при $x=3$ (там $y=-1$). Відповідь: \textbf{В}.}
\noindent
\begin{minipage}[c]{0.52\textwidth}
\zadnum На рисунку зображено графік функції $y = f(x)$, визначеної на проміжку $[-4;\, 4]$. Укажіть точку максимуму цієї функції.
\end{minipage}%
\hfill
\begin{minipage}[c]{0.46\textwidth}
\centering
\begin{tikzpicture}[scale=0.55, thick]
    % Сітка
    \draw[gray!40, thin] (-4.5,-3.5) grid (6.5,4.5);

    % Осі
    \draw[->, thick] (-4.5,0) -- (4.8,0) node[below] {$x$};
    \draw[->, thick] (0,-3.0) -- (0,4.3) node[left] {$y$};

    % Підписи осей
    \node[below left] at (0,0) {$0$};
    \node[below] at (1,0) {$1$};
    \node[below] at (-4,0) {$-4$};
    \node[below] at (4,0) {$4$};
    \node[left] at (0,1) {$1$};

    % Графік функції
    \draw[mainGreen!80!black, very thick] plot[smooth, tension=0.45] coordinates {
        (-4,1) (-2,-2) (0,1) (1,3) (2.5,1) (4,-1)
    };

    % Точки на кінцях
    \fill[mainGreen!80!black] (-4,1) circle (3pt);
    \fill[mainGreen!80!black] (4,-1) circle (3pt);

    % Підпис графіка у білому прямокутнику
    \node[fill=white, inner sep=1pt] at (4, 3.6) {$y = f(x)$};
\end{tikzpicture}
\end{minipage}
\par\vspace{0.4cm}
\answerTable{$-2$}{$1$}{$3$}{$-1$}{$4$}
\noindent
\begin{minipage}[c]{0.52\textwidth}
\zadnum На рисунку зображено графік функції $y = f(x)$, визначеної на проміжку $[-4;\, 4]$. Укажіть \textit{найбільше значення} цієї функції.
\end{minipage}%
\hfill
\begin{minipage}[c]{0.46\textwidth}
\centering
\begin{tikzpicture}[scale=0.55, thick]
    % Сітка
    \draw[gray!40, thin] (-4.5,-3.5) grid (6.5,4.5);

    % Осі
    \draw[->, thick] (-4.5,0) -- (4.8,0) node[below] {$x$};
    \draw[->, thick] (0,-3.0) -- (0,4.3) node[left] {$y$};

    % Підписи осей
    \node[below left] at (0,0) {$0$};
    \node[below] at (1,0) {$1$};
    \node[below] at (-4,0) {$-4$};
    \node[below] at (4,0) {$4$};
    \node[left] at (0,1) {$1$};

    % Графік функції
    \draw[mainGreen!80!black, very thick] plot[smooth, tension=0.45] coordinates {
        (-4,0) (-2.5,-2) (-1,1) (1,4) (2.5,2) (4,-1)
    };

    % Точки на кінцях
    \fill[mainGreen!80!black] (-4,0) circle (3pt);
    \fill[mainGreen!80!black] (4,-1) circle (3pt);

    % Підпис графіка у білому прямокутнику
    \node[fill=white, inner sep=1pt] at (4, 3.6) {$y = f(x)$};
\end{tikzpicture}
\end{minipage}
\par\vspace{0.4cm}
\answerTable{$-2$}{$1$}{$2$}{$4$}{$3$}
\taskBlock{Геометрична прогресія --- сесія 4.06, завдання №10}
\zadtask{У геометричній прогресії $(b_n)$ з додатними членами відомо, що $\dfrac{b_3}{b_1} = 25$. Обчисліть $\dfrac{b_4}{b_1}$.}
\answerTable{$5$}{$125$}{$525$}{$300$}{$625$}
\solution{$\dfrac{b_3}{b_1}=q^2=25$, а оскільки члени додатні, $q=5$. Тоді $\dfrac{b_4}{b_1}=q^3=125$. Відповідь: \textbf{Б}.}
\zadtask{У геометричній прогресії $(b_n)$ з додатними членами відомо, що $\dfrac{b_3}{b_1} = 9$. Обчисліть $\dfrac{b_4}{b_1}$.}
\answerTable{$3$}{$12$}{$81$}{$27$}{$729$}
\zadtask{У геометричній прогресії $(b_n)$ з додатними членами відомо, що $\dfrac{b_4}{b_2} = 16$. Обчисліть $\dfrac{b_5}{b_1}$.}
\answerTable{$4$}{$64$}{$256$}{$48$}{$32$}
\taskBlock{Графік функції та точка екстремуму зсунутої функції --- сесія 5.06, завдання №8}
\noindent
\begin{minipage}[c]{0.56\textwidth}
\zadnum На рисунку зображено графік функції $y = f(x)$, визначеної на проміжку $[-5;\,5]$. Визначте точку екстремуму функції $y = f(x - 3)$.
\end{minipage}%
\hfill
\begin{minipage}[c]{0.40\textwidth}
\centering
\begin{tikzpicture}[scale=0.45, thick]
    \draw[gray!30, very thin] (-5.3,-2.3) grid (5.3,4.3);
    \draw[->] (-5.5,0) -- (5.6,0) node[right, font=\scriptsize] {$x$};
    \draw[->] (0,-2.5) -- (0,4.5) node[above, font=\scriptsize] {$y$};
    \node[below left, font=\tiny] at (0,0) {$0$};
    \node[font=\tiny] at (1,-0.4) {$1$}; \node[font=\tiny] at (-5,-0.4) {$-5$}; \node[font=\tiny] at (5,-0.4) {$5$};
    \draw[mainGreen, line width=1.2pt]
        (-5,-2) .. controls (-3,0) and (0,3.5) .. (2,4)
        .. controls (3.5,4) and (4.5,2) .. (5,1.5);
    \node[font=\tiny, mainGreen] at (3.5,4) {$y=f(x)$};
\end{tikzpicture}
\end{minipage}
\par\vspace{0.2cm}
\answerTable{$5$}{$1$}{$-1$}{$2$}{$-5$}
\solution{Графік $y=f(x)$ має екстремум (максимум) у точці $x=2$. Заміна $x\to x-3$ зсуває графік праворуч на $3$, тож точка екстремуму $x=2+3=5$. Відповідь: \textbf{А}.}
\noindent
\begin{minipage}[c]{0.56\textwidth}
\zadnum На рисунку зображено графік функції $y = f(x)$, визначеної на проміжку $[-5;\,5]$. Визначте точку екстремуму функції $y = f(x - 2)$.
\end{minipage}%
\hfill
\begin{minipage}[c]{0.40\textwidth}
\centering
\begin{tikzpicture}[scale=0.45, thick]
    \draw[gray!30, very thin] (-5.3,-2.3) grid (5.3,4.3);
    \draw[->] (-5.5,0) -- (5.6,0) node[right, font=\scriptsize] {$x$};
    \draw[->] (0,-2.5) -- (0,4.5) node[above, font=\scriptsize] {$y$};
    \node[below left, font=\tiny] at (0,0) {$0$};
    \node[font=\tiny] at (1,-0.4) {$1$}; \node[font=\tiny] at (-5,-0.4) {$-5$}; \node[font=\tiny] at (5,-0.4) {$5$};
    \draw[mainGreen, line width=1.2pt]
        (-5,-2) .. controls (-3,1) and (0,3.6) .. (1,4)
        .. controls (2.5,4) and (4,1.5) .. (5,0.8);
    \node[font=\tiny, mainGreen] at (3.4,3.6) {$y=f(x)$};
\end{tikzpicture}
\end{minipage}
\par\vspace{0.2cm}
\answerTable{$1$}{$2$}{$-1$}{$3$}{$-3$}
\noindent
\begin{minipage}[c]{0.56\textwidth}
\zadnum На рисунку зображено графік функції $y = f(x)$, визначеної на проміжку $[-5;\,5]$. Визначте точку екстремуму функції $y = f(x) + 3$.
\end{minipage}%
\hfill
\begin{minipage}[c]{0.40\textwidth}
\centering
\begin{tikzpicture}[scale=0.45, thick]
    \draw[gray!30, very thin] (-5.3,-2.3) grid (5.3,4.3);
    \draw[->] (-5.5,0) -- (5.6,0) node[right, font=\scriptsize] {$x$};
    \draw[->] (0,-2.5) -- (0,4.5) node[above, font=\scriptsize] {$y$};
    \node[below left, font=\tiny] at (0,0) {$0$};
    \node[font=\tiny] at (1,-0.4) {$1$}; \node[font=\tiny] at (-5,-0.4) {$-5$}; \node[font=\tiny] at (5,-0.4) {$5$};
    \draw[mainGreen, line width=1.2pt]
        (-5,3.5) .. controls (-3.5,0) and (-2,-2) .. (-1,-2)
        .. controls (1,-2) and (3.5,1.5) .. (5,4);
    \node[font=\tiny, mainGreen] at (-3.6,3.2) {$y=f(x)$};
\end{tikzpicture}
\end{minipage}
\par\vspace{0.2cm}
\answerTable{$-1$}{$2$}{$-4$}{$3$}{$-3$}
\taskBlock{Арифметична прогресія --- сесія 5.06, завдання №15}
\zadtask{В арифметичній прогресії $(a_n)$ відомо, що $a_1 + a_2 = 11{,}2$, а $a_1 + a_4 = 8{,}8$. Визначте різницю $d$ цієї прогресії.}
\answerTable{$-1{,}2$}{$1{,}2$}{$-2{,}4$}{$2{,}4$}{$-0{,}6$}
\solution{$a_1+a_2=2a_1+d=11{,}2$, $a_1+a_4=2a_1+3d=8{,}8$. Віднявши перше від другого: $2d=-2{,}4$, тож $d=-1{,}2$. Відповідь: \textbf{А}.}
\zadtask{В арифметичній прогресії $(a_n)$ відомо, що $a_1 + a_2 = 7{,}4$, а $a_1 + a_4 = 12{,}2$. Визначте різницю $d$ цієї прогресії.}
\answerTable{$-2{,}4$}{$2{,}4$}{$-4{,}8$}{$4{,}8$}{$1{,}2$}
\zadtask{В арифметичній прогресії $(a_n)$ відомо, що $a_3 = 10$, а $a_7 = 2$. Визначте \textit{перший член} $a_1$ цієї прогресії.}
\answerTable{$14$}{$12$}{$6$}{$-14$}{$16$}
\typeTitle{Завдання на встановлення відповідності}
\taskBlock{Відповідність: характеристики функцій та проміжки --- сесія 23.05, завдання №18}
\zadtask{Установіть відповідність між характеристикою функції (1--3) та проміжком (А--Д), якому вона належить. \nmtyear{2026}

\vspace{0.3cm}

\noindent
\begin{minipage}[t]{0.52\textwidth}
\textit{Характеристика} \par \vspace{0.2cm}
\textbf{1} \ \ область визначення функції $y = \sqrt{x - 2}$\\[0.25cm]
\textbf{2} \ \ точка екстремуму функції $y = x^2 + 2x + 3$\\[0.25cm]
\textbf{3} \ \ абсциса точки перетину прямої $y = 6$ з графіком функції $y = x^3$
\end{minipage}
\begin{minipage}[t]{0.24\textwidth}
\textit{Проміжок} \par \vspace{0.2cm}
\textbf{А} \ \ $(-\infty;\, -2)$\\[0.15cm]
\textbf{Б} \ \ $[-2;\, 0)$\\[0.15cm]
\textbf{В} \ \ $[0;\, 1)$\\[0.15cm]
\textbf{Г} \ \ $(1;\, 2)$\\[0.15cm]
\textbf{Д} \ \ $[2;\, +\infty)$
\end{minipage}
\hfill
\begin{minipage}[t]{0.22\textwidth}
\vspace{-0.2cm}
\begin{flushright}\matchingGrid\end{flushright}
\end{minipage}
}
\solution{$1)$ ОДЗ $y=\sqrt{x-2}$: $x\ge 2$, тобто $[2;+\infty)$~(Д). $2)$ Вершина параболи $y=x^2+2x+3$: $x=-\dfrac{2}{2}=-1\in[-2;0)$~(Б). $3)$ $x^3=6 \Rightarrow x=\sqrt[3]{6}\approx 1{,}8\in(1;2)$~(Г). Відповідь: \textbf{1--Д, 2--Б, 3--Г}.}
\zadtask{Установіть відповідність між характеристикою функції (1--3) та проміжком (А--Д), якому вона належить.

\vspace{0.3cm}

\noindent
\begin{minipage}[t]{0.52\textwidth}
\textit{Характеристика} \par \vspace{0.2cm}
\textbf{1} \ \ область визначення функції $y = \sqrt{x - 3}$\\[0.25cm]
\textbf{2} \ \ точка екстремуму функції $y = x^2 + 6x$\\[0.25cm]
\textbf{3} \ \ абсциса точки перетину прямої $y = 15$ з графіком функції $y = x^3$
\end{minipage}
\begin{minipage}[t]{0.24\textwidth}
\textit{Проміжок} \par \vspace{0.2cm}
\textbf{А} \ \ $(-\infty;\, -3)$\\[0.15cm]
\textbf{Б} \ \ $[-3;\, 0)$\\[0.15cm]
\textbf{В} \ \ $[0;\, 2)$\\[0.15cm]
\textbf{Г} \ \ $(2;\, 3)$\\[0.15cm]
\textbf{Д} \ \ $[3;\, +\infty)$
\end{minipage}
\hfill
\begin{minipage}[t]{0.22\textwidth}
\vspace{-0.2cm}
\begin{flushright}\matchingGrid\end{flushright}
\end{minipage}
}
\zadtask{Установіть відповідність між характеристикою функції (1--3) та проміжком (А--Д), якому вона належить.

\vspace{0.3cm}

\noindent
\begin{minipage}[t]{0.52\textwidth}
\textit{Характеристика} \par \vspace{0.2cm}
\textbf{1} \ \ корінь рівняння $2^x = \dfrac{1}{8}$\\[0.25cm]
\textbf{2} \ \ точка мінімуму функції $y = x^2 - x$\\[0.25cm]
\textbf{3} \ \ область визначення функції $y = \sqrt{2x - 8}$
\end{minipage}
\begin{minipage}[t]{0.24\textwidth}
\textit{Проміжок} \par \vspace{0.2cm}
\textbf{А} \ \ $(-\infty;\, -2)$\\[0.15cm]
\textbf{Б} \ \ $[-2;\, 0)$\\[0.15cm]
\textbf{В} \ \ $[0;\, 1)$\\[0.15cm]
\textbf{Г} \ \ $(1;\, 4)$\\[0.15cm]
\textbf{Д} \ \ $[4;\, +\infty)$
\end{minipage}
\hfill
\begin{minipage}[t]{0.22\textwidth}
\vspace{-0.2cm}
\begin{flushright}\matchingGrid\end{flushright}
\end{minipage}
}
\taskBlock{Характеристики функцій (відповідність) --- сесія 30.05, завдання №17}
\zadtask{Установіть відповідність між характеристикою функції (1--3) та проміжком (А--Д), яким є ця характеристика.

\vspace{0.2cm}
\noindent
\begin{minipage}[t]{0.45\textwidth}
\textit{Характеристика функції}\par\vspace{0.15cm}
\textbf{1} \quad область значень функції\\ $y = \sqrt{x+2}$\\[0.25cm]
\textbf{2} \quad область визначення функції \\$y = \lg(x-2)$\\[0.25cm]
\textbf{3} \quad проміжок зростання функції\\ $y = -x^2 + 4x - 7$
\end{minipage}
\hfill
\begin{minipage}[t]{0.3\textwidth}
\textit{Проміжок}\par\vspace{0.15cm}
\textbf{А} \quad $(-\infty;\, +\infty)$\\[0.15cm]
\textbf{Б} \quad $(2;\, +\infty)$\\[0.15cm]
\textbf{В} \quad $[0;\, +\infty)$\\[0.15cm]
\textbf{Г} \quad $(-\infty;\, 2]$\\[0.15cm]
\textbf{Д} \quad $(-\infty;\, 2)$
\end{minipage}
\hfill
\begin{minipage}[t]{0.2\textwidth}
\vspace{0.4cm}
\begin{flushright}\matchingGrid\end{flushright}
\end{minipage}
}
\solution{\textbf{1}: область значень $y=\sqrt{x+2}$ — це $[0;+\infty)$ (В). \textbf{2}: ОДЗ $\lg(x-2)$: $x-2>0$, тобто $(2;+\infty)$ (Б). \textbf{3}: $y=-x^2+4x-7$ — парабола вітками вниз, вершина $x=2$, зростає на $(-\infty;2]$ (Г). Відповідь: \textbf{1~--~В, 2~--~Б, 3~--~Г}.}
\zadtask{Установіть відповідність між характеристикою функції (1--3) та проміжком (А--Д), яким є ця характеристика.

\vspace{0.2cm}
\noindent
\begin{minipage}[t]{0.45\textwidth}
\textit{Характеристика функції}\par\vspace{0.15cm}
\textbf{1} \quad область значень функції\\ $y = \sqrt{x-1}$\\[0.25cm]
\textbf{2} \quad область визначення функції \\$y = \ln(x+3)$\\[0.25cm]
\textbf{3} \quad проміжок зростання функції\\ $y = -x^2 - 6x + 1$
\end{minipage}
\hfill
\begin{minipage}[t]{0.3\textwidth}
\textit{Проміжок}\par\vspace{0.15cm}
\textbf{А} \quad $[0;\, +\infty)$\\[0.15cm]
\textbf{Б} \quad $(-\infty;\, -3]$\\[0.15cm]
\textbf{В} \quad $(-3;\, +\infty)$\\[0.15cm]
\textbf{Г} \quad $(-\infty;\, +\infty)$\\[0.15cm]
\textbf{Д} \quad $[-3;\, +\infty)$
\end{minipage}
\hfill
\begin{minipage}[t]{0.2\textwidth}
\vspace{0.4cm}
\begin{flushright}\matchingGrid\end{flushright}
\end{minipage}
}
\zadtask{Установіть відповідність між характеристикою функції (1--3) та проміжком (А--Д), яким є ця характеристика.

\vspace{0.2cm}
\noindent
\begin{minipage}[t]{0.45\textwidth}
\textit{Характеристика функції}\par\vspace{0.15cm}
\textbf{1} \quad область визначення функції\\ $y = \sqrt{x-5}$\\[0.25cm]
\textbf{2} \quad область визначення функції \\$y = \lg(4-x)$\\[0.25cm]
\textbf{3} \quad проміжок спадання функції\\ $y = -x^2 + 8x - 3$
\end{minipage}
\hfill
\begin{minipage}[t]{0.3\textwidth}
\textit{Проміжок}\par\vspace{0.15cm}
\textbf{А} \quad $[5;\, +\infty)$\\[0.15cm]
\textbf{Б} \quad $(-\infty;\, 4)$\\[0.15cm]
\textbf{В} \quad $(-\infty;\, 4]$\\[0.15cm]
\textbf{Г} \quad $(4;\, +\infty)$\\[0.15cm]
\textbf{Д} \quad $(-\infty;\, +\infty)$
\end{minipage}
\hfill
\begin{minipage}[t]{0.2\textwidth}
\vspace{0.4cm}
\begin{flushright}\matchingGrid\end{flushright}
\end{minipage}
}
\taskBlock{Відповідність: властивості графіків функцій --- сесія 1.06, завдання №17}
\zadtask{Увідповідніть початок речення (1--3) із його закінченням (А--Д) так, щоб утворилося правильне твердження.

\vspace{0.2cm}
\noindent
\begin{minipage}[t]{0.42\textwidth}
\textit{Початок речення}\par\vspace{0.2cm}
\textbf{1} \quad Графік функції $y = \dfrac{2}{x} + 3$\\[0.4cm]
\textbf{2} \quad Графік функції $y = 4^x$\\[0.3cm]
\textbf{3} \quad Графік функції $y = \mathrm{tg}\,x$
\end{minipage}
\hfill
\begin{minipage}[t]{0.42\textwidth}
\textit{Закінчення речення}\par\vspace{0.2cm}
\textbf{А} \quad симетричний відносно початку координат.\\[0.15cm]
\textbf{Б} \quad симетричний відносно осі $y$.\\[0.15cm]
\textbf{В} \quad перетинає лише вісь $x$.\\[0.15cm]
\textbf{Г} \quad перетинає лише вісь $y$.\\[0.15cm]
\textbf{Д} \quad перетинає графік функції $y = x$ лише в одній точці.
\end{minipage}
\hfill
\begin{minipage}[t]{0.14\textwidth}
\vspace{0.4cm}
\begin{flushright}\matchingGrid\end{flushright}
\end{minipage}
}
\solution{\textbf{1)} $y=\dfrac{2}{x}+3$ ($x\ne0$): осі $y$ не перетинає, а вісь $x$ перетинає при $x=-\dfrac{2}{3}$ --- лише вісь $x$ (В); \textbf{2)} $y=4^x>0$: перетинає лише вісь $y$ у точці $(0;1)$ (Г); \textbf{3)} $y=\mathrm{tg}\,x$ --- непарна, симетрична відносно початку координат (А). Відповідь: \textbf{1\,--\,В, 2\,--\,Г, 3\,--\,А}.}
\zadtask{Увідповідніть початок речення (1--3) із його закінченням (А--Д) так, щоб утворилося правильне твердження.

\vspace{0.2cm}
\noindent
\begin{minipage}[t]{0.42\textwidth}
\textit{Початок речення}\par\vspace{0.2cm}
\textbf{1} \quad Графік функції $y = \dfrac{3}{x} - 2$\\[0.4cm]
\textbf{2} \quad Графік функції $y = \left(\dfrac{1}{5}\right)^{x}$\\[0.3cm]
\textbf{3} \quad Графік функції $y = x^3$
\end{minipage}
\hfill
\begin{minipage}[t]{0.42\textwidth}
\textit{Закінчення речення}\par\vspace{0.2cm}
\textbf{А} \quad симетричний відносно початку координат.\\[0.15cm]
\textbf{Б} \quad симетричний відносно осі $y$.\\[0.15cm]
\textbf{В} \quad перетинає лише вісь $x$.\\[0.15cm]
\textbf{Г} \quad перетинає лише вісь $y$.\\[0.15cm]
\textbf{Д} \quad не перетинає жодної з осей координат.
\end{minipage}
\hfill
\begin{minipage}[t]{0.14\textwidth}
\vspace{0.4cm}
\begin{flushright}\matchingGrid\end{flushright}
\end{minipage}
}
\zadtask{Увідповідніть початок речення (1--3) із його закінченням (А--Д) так, щоб утворилося правильне твердження.

\vspace{0.2cm}
\noindent
\begin{minipage}[t]{0.42\textwidth}
\textit{Початок речення}\par\vspace{0.2cm}
\textbf{1} \quad Графік функції $y = -\dfrac{4}{x}$\\[0.4cm]
\textbf{2} \quad Графік функції $y = \log_2 x$\\[0.3cm]
\textbf{3} \quad Графік функції $y = x^2 + 1$
\end{minipage}
\hfill
\begin{minipage}[t]{0.42\textwidth}
\textit{Закінчення речення}\par\vspace{0.2cm}
\textbf{А} \quad симетричний відносно початку координат.\\[0.15cm]
\textbf{Б} \quad симетричний відносно осі $y$.\\[0.15cm]
\textbf{В} \quad перетинає лише вісь $x$.\\[0.15cm]
\textbf{Г} \quad перетинає обидві осі координат.\\[0.15cm]
\textbf{Д} \quad перетинає лише вісь $x$ і вісь $y$ одночасно в початку координат.
\end{minipage}
\hfill
\begin{minipage}[t]{0.14\textwidth}
\vspace{0.4cm}
\begin{flushright}\matchingGrid\end{flushright}
\end{minipage}
}
\taskBlock{Відповідність: функція та координатні чверті --- сесія 2.06, завдання №17}
\zadtask{Установіть відповідність між функцією (1--3) та координатними чвертями (А--Д), у яких розташований її графік.

\vspace{0.3cm}
\noindent
\begin{minipage}[t]{0.4\textwidth}
\textit{Функція}\par\vspace{0.2cm}
\textbf{1} \quad $y = -x - 1$\\[0.4cm]
\textbf{2} \quad $y = \dfrac{1}{x+2}$\\[0.5cm]
\textbf{3} \quad $y = \log_3 x$
\end{minipage}
\hfill
\begin{minipage}[t]{0.32\textwidth}
\textit{Чверті}\par\vspace{0.2cm}
\textbf{А} \quad I і IV\\[0.15cm]
\textbf{Б} \quad I, II і III\\[0.15cm]
\textbf{В} \quad II, III і IV\\[0.15cm]
\textbf{Г} \quad I і III\\[0.15cm]
\textbf{Д} \quad II і IV
\end{minipage}
\hfill
\begin{minipage}[t]{0.2\textwidth}
\vspace{0.5cm}
\begin{flushright}\matchingGrid\end{flushright}
\end{minipage}
}
\solution{$y=-x-1$ -- спадна пряма з від'ємним зсувом, проходить через II, III і IV чверті -- \textbf{В}; $y=\dfrac{1}{x+2}$ -- гіпербола, зміщена ліворуч, її гілки лежать у I, II і III чвертях -- \textbf{Б}; $y=\log_3 x$ зростає при $x>0$, графік у I і IV чвертях -- \textbf{А}. Відповідь: \textbf{1--В, 2--Б, 3--А}.}
\zadtask{Установіть відповідність між функцією (1--3) та координатними чвертями (А--Д), у яких розташований її графік.

\vspace{0.3cm}
\noindent
\begin{minipage}[t]{0.4\textwidth}
\textit{Функція}\par\vspace{0.2cm}
\textbf{1} \quad $y = x + 2$\\[0.4cm]
\textbf{2} \quad $y = -\dfrac{3}{x}$\\[0.5cm]
\textbf{3} \quad $y = \log_{\frac{1}{3}} x$
\end{minipage}
\hfill
\begin{minipage}[t]{0.32\textwidth}
\textit{Чверті}\par\vspace{0.2cm}
\textbf{А} \quad I і IV\\[0.15cm]
\textbf{Б} \quad I, II і III\\[0.15cm]
\textbf{В} \quad II, III і IV\\[0.15cm]
\textbf{Г} \quad I і III\\[0.15cm]
\textbf{Д} \quad II і IV
\end{minipage}
\hfill
\begin{minipage}[t]{0.2\textwidth}
\vspace{0.5cm}
\begin{flushright}\matchingGrid\end{flushright}
\end{minipage}
}
\zadtask{Установіть відповідність між функцією (1--3) та координатними чвертями (А--Д), у яких розташований її графік.

\vspace{0.3cm}
\noindent
\begin{minipage}[t]{0.4\textwidth}
\textit{Функція}\par\vspace{0.2cm}
\textbf{1} \quad $y = 3x$\\[0.4cm]
\textbf{2} \quad $y = -\dfrac{1}{x}$\\[0.5cm]
\textbf{3} \quad $y = -x - 5$
\end{minipage}
\hfill
\begin{minipage}[t]{0.32\textwidth}
\textit{Чверті}\par\vspace{0.2cm}
\textbf{А} \quad I і IV\\[0.15cm]
\textbf{Б} \quad I, II і III\\[0.15cm]
\textbf{В} \quad II, III і IV\\[0.15cm]
\textbf{Г} \quad I і III\\[0.15cm]
\textbf{Д} \quad II і IV
\end{minipage}
\hfill
\begin{minipage}[t]{0.2\textwidth}
\vspace{0.5cm}
\begin{flushright}\matchingGrid\end{flushright}
\end{minipage}
}
\taskBlock{Спільні точки графіків (відповідність) --- сесія 3.06, завдання №17}
\noindent\zadnum До кожного початку речення (1--3) доберіть його закінчення (А--Д) так, щоб утворилося правильне твердження.

\vspace{0.3cm}
\noindent
\begin{minipage}[t]{0.42\textwidth}
\textit{Початок речення}\par\vspace{0.2cm}
\textbf{1} \quad Графіки функцій $y = 2x - 3$ та $y = x + 1$\\[0.3cm]
\textbf{2} \quad Графіки функцій $y = x^2$ та $y = x^2 - 2$\\[0.3cm]
\textbf{3} \quad Графік функції $y = \dfrac{2}{x}$ та графік рівняння $x^2 + y^2 = 4$
\end{minipage}
\hfill
\begin{minipage}[t]{0.42\textwidth}
\textit{Закінчення речення}\par\vspace{0.2cm}
\textbf{А} \quad не мають жодної спільної точки.\\[0.15cm]
\textbf{Б} \quad мають єдину спільну точку.\\[0.15cm]
\textbf{В} \quad мають лише дві спільні точки.\\[0.15cm]
\textbf{Г} \quad мають лише три спільні точки.\\[0.15cm]
\textbf{Д} \quad мають чотири спільні точки.
\end{minipage}
\hfill
\begin{minipage}[t]{0.14\textwidth}
\vspace{0.5cm}
\begin{flushright}\matchingGrid\end{flushright}
\end{minipage}

\solution{\textbf{1}: прямі $y=2x-3$ і $y=x+1$ мають різні кутові коефіцієнти~--- одна спільна точка (Б). \textbf{2}: параболи $y=x^2$ і $y=x^2-2$ зсунуті по вертикалі~--- спільних точок немає (А). \textbf{3}: гіпербола $y=\dfrac{2}{x}$ і коло $x^2+y^2=4$ перетинаються у двох точках (В). Відповідь: \textbf{1~--~Б, 2~--~А, 3~--~В}.}

\vspace{0.5cm}
\noindent\zadnum До кожного початку речення (1--3) доберіть його закінчення (А--Д) так, щоб утворилося правильне твердження.

\vspace{0.3cm}
\noindent
\begin{minipage}[t]{0.42\textwidth}
\textit{Початок речення}\par\vspace{0.2cm}
\textbf{1} \quad Графіки функцій $y = 3x + 1$ та $y = 3x - 2$\\[0.3cm]
\textbf{2} \quad Графіки функцій $y = x^2$ та $y = 4$\\[0.3cm]
\textbf{3} \quad Графіки функцій $y = x^2$ та $y = 2x - 1$
\end{minipage}
\hfill
\begin{minipage}[t]{0.42\textwidth}
\textit{Закінчення речення}\par\vspace{0.2cm}
\textbf{А} \quad не мають жодної спільної точки.\\[0.15cm]
\textbf{Б} \quad мають єдину спільну точку.\\[0.15cm]
\textbf{В} \quad мають лише дві спільні точки.\\[0.15cm]
\textbf{Г} \quad мають лише три спільні точки.\\[0.15cm]
\textbf{Д} \quad мають чотири спільні точки.
\end{minipage}
\hfill
\begin{minipage}[t]{0.14\textwidth}
\vspace{0.5cm}
\begin{flushright}\matchingGrid\end{flushright}
\end{minipage}

\vspace{0.5cm}
\noindent\zadnum До кожного початку речення (1--3) доберіть його закінчення (А--Д) так, щоб утворилося правильне твердження.

\vspace{0.3cm}
\noindent
\begin{minipage}[t]{0.42\textwidth}
\textit{Початок речення}\par\vspace{0.2cm}
\textbf{1} \quad Графіки функцій $y = \dfrac{1}{x}$ та $y = x$\\[0.3cm]
\textbf{2} \quad Графіки функцій $y = \sqrt{x}$ та $y = -1$\\[0.3cm]
\textbf{3} \quad Графік рівняння $x^2 + y^2 = 9$ та графік функції $y = 3$
\end{minipage}
\hfill
\begin{minipage}[t]{0.42\textwidth}
\textit{Закінчення речення}\par\vspace{0.2cm}
\textbf{А} \quad не мають жодної спільної точки.\\[0.15cm]
\textbf{Б} \quad мають єдину спільну точку.\\[0.15cm]
\textbf{В} \quad мають лише дві спільні точки.\\[0.15cm]
\textbf{Г} \quad мають лише три спільні точки.\\[0.15cm]
\textbf{Д} \quad мають чотири спільні точки.
\end{minipage}
\hfill
\begin{minipage}[t]{0.14\textwidth}
\vspace{0.5cm}
\begin{flushright}\matchingGrid\end{flushright}
\end{minipage}
\taskBlock{Найбільше значення функції на проміжку (відповідність) --- сесія 4.06, завдання №17}
\zadtask{Установіть відповідність між функцією, заданою на проміжку (1--3), та її найбільшим значенням (А--Д) на цьому проміжку.

\vspace{0.3cm}
\noindent
\begin{minipage}[t]{0.48\textwidth}
\textit{Функція}\par\vspace{0.2cm}
\textbf{1} \quad $y = \sin 4x + 2,\ x \in (-\infty; +\infty)$\\[0.35cm]
\textbf{2} \quad $y = 4 - x^2,\ x \in [0; 2]$\\[0.35cm]
\textbf{3} \quad $y = 4 + \sqrt{x},\ x \in [0; 1]$
\end{minipage}
\hfill
\begin{minipage}[t]{0.32\textwidth}
\textit{Найбільше значення}\par\vspace{0.2cm}
\textbf{А} \quad $0$\\[0.15cm]
\textbf{Б} \quad $3$\\[0.15cm]
\textbf{В} \quad $4$\\[0.15cm]
\textbf{Г} \quad $5$\\[0.15cm]
\textbf{Д} \quad $6$
\end{minipage}
\hfill
\begin{minipage}[t]{0.16\textwidth}
\vspace{0.4cm}
\begin{flushright}\matchingGrid\end{flushright}
\end{minipage}
}
\solution{1: $\sin4x\leqslant1$, тож найбільше $y=1+2=3$ (Б). 2: $y=4-x^2$ спадає на $[0;2]$, найбільше при $x=0$: $y=4$ (В). 3: $y=4+\sqrt{x}$ зростає на $[0;1]$, найбільше при $x=1$: $y=5$ (Г). Відповідь: \textbf{1~--~Б, 2~--~В, 3~--~Г}.}
\zadtask{Установіть відповідність між функцією, заданою на проміжку (1--3), та її найбільшим значенням (А--Д) на цьому проміжку.

\vspace{0.3cm}
\noindent
\begin{minipage}[t]{0.48\textwidth}
\textit{Функція}\par\vspace{0.2cm}
\textbf{1} \quad $y = \cos 3x + 4,\ x \in (-\infty; +\infty)$\\[0.35cm]
\textbf{2} \quad $y = 9 - x^2,\ x \in [1; 3]$\\[0.35cm]
\textbf{3} \quad $y = 1 + \sqrt{x},\ x \in [0; 4]$
\end{minipage}
\hfill
\begin{minipage}[t]{0.32\textwidth}
\textit{Найбільше значення}\par\vspace{0.2cm}
\textbf{А} \quad $3$\\[0.15cm]
\textbf{Б} \quad $5$\\[0.15cm]
\textbf{В} \quad $8$\\[0.15cm]
\textbf{Г} \quad $9$\\[0.15cm]
\textbf{Д} \quad $10$
\end{minipage}
\hfill
\begin{minipage}[t]{0.16\textwidth}
\vspace{0.4cm}
\begin{flushright}\matchingGrid\end{flushright}
\end{minipage}
}
\zadtask{Установіть відповідність між функцією, заданою на проміжку (1--3), та її \textit{найменшим} значенням (А--Д) на цьому проміжку.

\vspace{0.3cm}
\noindent
\begin{minipage}[t]{0.48\textwidth}
\textit{Функція}\par\vspace{0.2cm}
\textbf{1} \quad $y = 2\sin x + 3,\ x \in (-\infty; +\infty)$\\[0.35cm]
\textbf{2} \quad $y = 6 - x^2,\ x \in [0; 1]$\\[0.35cm]
\textbf{3} \quad $y = \sqrt{x} - 1,\ x \in [1; 4]$
\end{minipage}
\hfill
\begin{minipage}[t]{0.32\textwidth}
\textit{Найменше значення}\par\vspace{0.2cm}
\textbf{А} \quad $1$\\[0.15cm]
\textbf{Б} \quad $0$\\[0.15cm]
\textbf{В} \quad $5$\\[0.15cm]
\textbf{Г} \quad $3$\\[0.15cm]
\textbf{Д} \quad $7$
\end{minipage}
\hfill
\begin{minipage}[t]{0.16\textwidth}
\vspace{0.4cm}
\begin{flushright}\matchingGrid\end{flushright}
\end{minipage}
}
\taskBlock{Властивості функцій (відповідність) --- сесія 5.06, завдання №17}
\zadnum Доберіть до функції (1--3) її властивість (А--Д).

\vspace{0.3cm}
\noindent
\begin{minipage}[t]{0.34\textwidth}
\textit{Функція}\par\vspace{0.2cm}
\textbf{1} \quad $y = x^4 - x^2$\\[0.35cm]
\textbf{2} \quad $y = \log_{0{,}2} x$\\[0.35cm]
\textbf{3} \quad $y = -\dfrac{1}{x}$
\end{minipage}
\hfill
\begin{minipage}[t]{0.52\textwidth}
\textit{Властивість функції}\par\vspace{0.2cm}
\textbf{А} \quad область визначення збігається з областю значень\\[0.12cm]
\textbf{Б} \quad є спадною\\[0.12cm]
\textbf{В} \quad має лише три нулі\\[0.12cm]
\textbf{Г} \quad має лише два нулі\\[0.12cm]
\textbf{Д} \quad графік розташований у I та III чвертях
\end{minipage}
\hfill
\begin{minipage}[t]{0.12\textwidth}
\vspace{0.5cm}
\begin{flushright}\matchingGrid\end{flushright}
\end{minipage}
\par\vspace{0.5cm}
\solution{$1$: $y=x^4-x^2=x^2(x^2-1)$ має нулі $x=0;\,\pm1$, тобто три нулі (В); $2$: $y=\log_{0{,}2}x$ при основі $0<0{,}2<1$ спадна (Б); $3$: $y=-\dfrac1x$ має область визначення $x\ne0$, що збігається з областю значень (А). Відповідь: \textbf{1~--~В, 2~--~Б, 3~--~А}.}
\zadnum Доберіть до функції (1--3) її властивість (А--Д).

\vspace{0.3cm}
\noindent
\begin{minipage}[t]{0.34\textwidth}
\textit{Функція}\par\vspace{0.2cm}
\textbf{1} \quad $y = x^4 - 9x^2$\\[0.35cm]
\textbf{2} \quad $y = 0{,}3^x$\\[0.35cm]
\textbf{3} \quad $y = -\dfrac{5}{x}$
\end{minipage}
\hfill
\begin{minipage}[t]{0.52\textwidth}
\textit{Властивість функції}\par\vspace{0.2cm}
\textbf{А} \quad область визначення збігається з областю значень\\[0.12cm]
\textbf{Б} \quad є спадною\\[0.12cm]
\textbf{В} \quad має лише три нулі\\[0.12cm]
\textbf{Г} \quad має лише два нулі\\[0.12cm]
\textbf{Д} \quad графік розташований у I та III чвертях
\end{minipage}
\hfill
\begin{minipage}[t]{0.12\textwidth}
\vspace{0.5cm}
\begin{flushright}\matchingGrid\end{flushright}
\end{minipage}
\par\vspace{0.5cm}
\zadnum Доберіть до функції (1--3) її властивість (А--Д).

\vspace{0.3cm}
\noindent
\begin{minipage}[t]{0.34\textwidth}
\textit{Функція}\par\vspace{0.2cm}
\textbf{1} \quad $y = \cos x$\\[0.35cm]
\textbf{2} \quad $y = x^3$\\[0.35cm]
\textbf{3} \quad $y = x^2 - 4x$
\end{minipage}
\hfill
\begin{minipage}[t]{0.52\textwidth}
\textit{Властивість функції}\par\vspace{0.2cm}
\textbf{А} \quad область визначення збігається з областю значень\\[0.12cm]
\textbf{Б} \quad є спадною\\[0.12cm]
\textbf{В} \quad має лише три нулі\\[0.12cm]
\textbf{Г} \quad має лише два нулі\\[0.12cm]
\textbf{Д} \quad є парною
\end{minipage}
\hfill
\begin{minipage}[t]{0.12\textwidth}
\vspace{0.5cm}
\begin{flushright}\matchingGrid\end{flushright}
\end{minipage}
\par\vspace{0.5cm}
\typeTitle{Завдання з короткою відповіддю}
\taskBlock{Властивості визначеного інтеграла --- сесія 23.05, завдання №20}
\zadtask{Обчисліть $\displaystyle\int_0^3 4\bigl(f(x) + 2\bigr)\,dx$, якщо $\displaystyle\int_0^3 f(x)\,dx = 10$. \nmtyear{2026}}
\nmtAnswerBox
\solution{$\displaystyle\int_0^3 4(f(x)+2)\,dx=4\int_0^3 f(x)\,dx+4\cdot 2\cdot(3-0)=4\cdot 10+24=64$. Відповідь: \textbf{64}.}
\zadtask{Обчисліть $\displaystyle\int_0^5 3\bigl(f(x) - 1\bigr)\,dx$, якщо $\displaystyle\int_0^5 f(x)\,dx = 8$.}
\nmtAnswerBox
\zadtask{Обчисліть $\displaystyle\int_2^7 f(x)\,dx$, якщо $\displaystyle\int_0^7 f(x)\,dx = 15$ та $\displaystyle\int_0^2 f(x)\,dx = 4$.}
\nmtAnswerBox
\taskBlock{Найменше значення функції на відрізку --- сесія 30.05, завдання №19}
\zadtask{Знайдіть \textit{значення $x_0$}, при якому функція $y = 2x^4 - 9x^2 + 5$ набуває \textit{найменшого} значення на проміжку $[-1;\, 5]$.
\nmtAnswerBox}
\solution{$y'=8x^3-18x=2x(4x^2-9)$; на $[-1;5]$ критичні точки $x=0$ і $x=1{,}5$. Значення: $y(-1)=-2$, $y(0)=5$, $y(1{,}5)=-5{,}125$, $y(5)$ велике. Найменше — при $x_0=1{,}5$. Відповідь: $1{,}5$.}
\zadtask{Знайдіть \textit{значення $x_0$}, при якому функція $y = x^4 - 8x^2 + 3$ набуває \textit{найменшого} значення на проміжку $[-1;\, 3]$.
\nmtAnswerBox}
\zadtask{Знайдіть \textit{значення $x_0$}, при якому функція $y = x^4 - 2x^2 + 7$ набуває \textit{найбільшого} значення на проміжку $[0;\, 5]$.
\nmtAnswerBox}
\taskBlock{Обчислення інтеграла --- сесія 1.06, завдання №19}
\zadtask{Обчисліть інтеграл $\displaystyle\int_{\frac{1}{e}}^{\sqrt{e}} \dfrac{dx}{x}$.}
\nmtAnswerBox
\solution{$\displaystyle\int \dfrac{dx}{x}=\ln x$, тож інтеграл $=\ln\sqrt{e}-\ln\dfrac{1}{e}=\dfrac{1}{2}-(-1)=1{,}5$. Відповідь: \textbf{$1{,}5$}.}
\zadtask{Обчисліть інтеграл $\displaystyle\int_{\frac{1}{e^2}}^{e} \dfrac{dx}{x}$.}
\nmtAnswerBox
\zadtask{Обчисліть інтеграл $\displaystyle\int_{0}^{\ln 3} e^{x}\,dx$.}
\nmtAnswerBox
\taskBlock{Визначений інтеграл --- сесія 2.06, завдання №19}
\zadtask{Обчисліть інтеграл $\displaystyle\int_{1}^{2} \left( \dfrac{2}{x^3} + 4x \right) dx$.
\nmtAnswerBox
}
\solution{Первісна: $\displaystyle\int\!\left(2x^{-3}+4x\right)dx=-\dfrac{1}{x^2}+2x^2$. Підставляємо межі: $\left(-\dfrac{1}{4}+8\right)-\left(-1+2\right)=7{,}75-1=6{,}75$.}
\zadtask{Обчисліть інтеграл $\displaystyle\int_{1}^{2} \left( \dfrac{4}{x^3} + 6x \right) dx$.
\nmtAnswerBox
}
\zadtask{Обчисліть інтеграл $\displaystyle\int_{1}^{3} \left( \dfrac{6}{x^2} + 2x \right) dx$.
\nmtAnswerBox
}
\taskBlock{Критичні точки і точка мінімуму --- сесія 3.06, завдання №19}
\zadtask{Визначте критичні точки функції $f(x) = 4x^3 - 21x^2 - 90x + 19$. У відповідь запишіть ту з них, яка є \textit{точкою мінімуму}.}
\nmtAnswerBox
\solution{$f'(x) = 12x^2 - 42x - 90 = 6(2x^2 - 7x - 15)$. Корені рівняння $2x^2-7x-15=0$ через дискримінант: $D = 49 + 120 = 169$, $\sqrt{D}=13$, $x = \dfrac{7\pm13}{4}$, тобто $x_1 = 5$, $x_2 = -1{,}5$. Похідна змінює знак з $-$ на $+$ у точці $x=5$, тож це точка мінімуму. Відповідь: \textbf{$5$}.}
\zadtask{Визначте критичні точки функції $f(x) = 2x^3 - 3x^2 - 36x + 7$. У відповідь запишіть ту з них, яка є \textit{точкою мінімуму}.}
\nmtAnswerBox
\zadtask{Визначте критичні точки функції $f(x) = 2x^3 + 9x^2 - 24x + 1$. У відповідь запишіть ту з них, яка є \textit{точкою максимуму}.}
\nmtAnswerBox
\taskBlock{Визначений інтеграл --- сесія 4.06, завдання №20}
\zadtask{Обчисліть значення інтеграла
\[ \int_0^3 \frac{x^2 + 3x}{x + 3} \,dx. \]
\nmtAnswerBox}
\solution{$\dfrac{x^2+3x}{x+3}=\dfrac{x(x+3)}{x+3}=x$. Тоді $\displaystyle\int_0^3 x\,dx=\dfrac{x^2}{2}\Big|_0^3=\dfrac{9}{2}=4{,}5$. Відповідь: \textbf{$4{,}5$}.}
\zadtask{Обчисліть значення інтеграла
\[ \int_0^4 \frac{x^2 + 5x}{x + 5} \,dx. \]
\nmtAnswerBox}
\zadtask{Обчисліть значення інтеграла
\[ \int_0^3 \frac{x^2 - 9}{x - 3} \,dx. \]
\nmtAnswerBox}
\taskBlock{Значення похідної в точці --- сесія 5.06, завдання №20}
\zadtask{Знайдіть значення похідної функції $f(x) = \dfrac{\sqrt{x}}{3x - 2}$ у точці $x_0 = 4$.}
\nmtAnswerBox
\solution{За правилом частки $f'(x)=\dfrac{\frac{1}{2\sqrt x}(3x-2)-\sqrt x\cdot3}{(3x-2)^2}$. При $x_0=4$: $\sqrt4=2$, $3\cdot4-2=10$; чисельник $=\frac{1}{4}\cdot10-2\cdot3=-3{,}5$, знаменник $=100$. Отже $f'(4)=-0{,}035$. Відповідь: \textbf{$-0{,}035$}.}
\zadtask{Знайдіть значення похідної функції $f(x) = \dfrac{\sqrt{x}}{x + 6}$ у точці $x_0 = 4$.}
\nmtAnswerBox
\zadtask{Знайдіть значення похідної функції $f(x) = (3x - 2)\sqrt{x}$ у точці $x_0 = 4$.}
\nmtAnswerBox

\chapterTitle{РОЗДІЛ 4. ПЛАНІМЕТРІЯ}
\typeTitle{Завдання з вибором однієї відповіді (А--Д)}
\taskBlock{Кути трикутника --- сесія 23.05, завдання №2}

\noindent
\begin{minipage}[c]{0.60\textwidth}
\zadnum Дано трикутник $ABC$ з кутами $\angle A=30^\circ$ та $\angle B=35^\circ$. Знайдіть градусну міру кута $C$. \nmtyear{2026}
\end{minipage}%
\hfill
\begin{minipage}[c]{0.38\textwidth}
\centering
\begin{tikzpicture}[scale=0.8, thick]
    \coordinate (A) at (0,0);
    \coordinate (C) at (4.5,0);
    \coordinate (B) at (1,2.8);
    \draw (A) -- (B) -- (C) -- cycle;
    \pic [draw, angle radius=0.5cm, "$30^\circ$"{right, shift={(0.15,0.15)}}] {angle = C--A--B};
    \pic [draw, angle radius=0.4cm, "$35^\circ$"{below, shift={(0.2,-0.1)}}] {angle = A--B--C};
    \pic [draw, angle radius=0.5cm, red, "?"{left, shift={(-0.15,0.15)}}] {angle = B--C--A};
    \node[below left] at (A) {$A$};
    \node[above] at (B) {$B$};
    \node[below right] at (C) {$C$};
\end{tikzpicture}
\end{minipage}
\par\vspace{0.2cm}
\answerTable{$75^\circ$}{$100^\circ$}{$115^\circ$}{$120^\circ$}{$30^\circ$}
\solution{Сума кутів трикутника $180^\circ$: $\angle C=180^\circ-30^\circ-35^\circ=115^\circ$. Відповідь: \textbf{В}.}

\vspace{0.3cm}

\noindent
\begin{minipage}[c]{0.60\textwidth}
\zadnum Дано трикутник $MNK$ з кутами $\angle M=40^\circ$ та $\angle N=65^\circ$. Знайдіть градусну міру кута $K$.
\end{minipage}%
\hfill
\begin{minipage}[c]{0.38\textwidth}
\centering
\begin{tikzpicture}[scale=0.8, thick]
    \coordinate (A) at (0,0);
    \coordinate (C) at (4.5,0);
    \coordinate (B) at (1.4,2.6);
    \draw (A) -- (B) -- (C) -- cycle;
    \pic [draw, angle radius=0.5cm, "$40^\circ$"{right, shift={(0.15,0.15)}}] {angle = C--A--B};
    \pic [draw, angle radius=0.4cm, "$65^\circ$"{below, shift={(0.2,-0.1)}}] {angle = A--B--C};
    \pic [draw, angle radius=0.5cm, red, "?"{left, shift={(-0.15,0.15)}}] {angle = B--C--A};
    \node[below left] at (A) {$M$};
    \node[above] at (B) {$N$};
    \node[below right] at (C) {$K$};
\end{tikzpicture}
\end{minipage}
\par\vspace{0.2cm}
\answerTable{$65^\circ$}{$75^\circ$}{$85^\circ$}{$105^\circ$}{$115^\circ$}

\vspace{0.3cm}

\noindent
\begin{minipage}[c]{0.60\textwidth}
\zadnum Дано трикутник $PQR$ з кутами $\angle P=25^\circ$ та $\angle Q=48^\circ$. Знайдіть градусну міру зовнішнього кута трикутника при вершині $R$.
\end{minipage}%
\hfill
\begin{minipage}[c]{0.38\textwidth}
\centering
\begin{tikzpicture}[scale=0.8, thick]
    \coordinate (A) at (0,0);
    \coordinate (C) at (4.5,0);
    \coordinate (B) at (1.2,2.4);
    \coordinate (F) at (6.0,0);
    \draw (A) -- (B) -- (C) -- cycle;
    \draw (C) -- (F);
    \pic [draw, angle radius=0.5cm, "$25^\circ$"{right, shift={(0.15,0.15)}}] {angle = C--A--B};
    \pic [draw, angle radius=0.4cm, "$48^\circ$"{below, shift={(0.2,-0.1)}}] {angle = A--B--C};
    \pic [draw, angle radius=0.5cm, red, "?"{above right, shift={(0.05,0.05)}}] {angle = B--C--F};
    \node[below left] at (A) {$P$};
    \node[above] at (B) {$Q$};
    \node[below] at (C) {$R$};
\end{tikzpicture}
\end{minipage}
\par\vspace{0.2cm}
\answerTable{$97^\circ$}{$103^\circ$}{$107^\circ$}{$117^\circ$}{$73^\circ$}

\vspace{0.3cm}

% ---------------------------------------------------------------------
\taskBlock{Істинність тверджень про коло --- сесія 23.05, завдання №10}

\zadtask{Які з наведених тверджень є істинними? \nmtyear{2026}

\vspace{0.15cm}
\noindent\hspace{1em}I. Пряма, що проходить через центр кола і лежить з ним в одній площині, має з ним дві спільні точки. \\
\noindent\hspace{1em}II. Діаметр кола, перпендикулярний до хорди, проходить через середину цієї хорди. \\
\noindent\hspace{1em}III. У колі можна провести два діаметри, які не перетинаються.}
\answerTable{лише I}{лише II}{лише I і II}{лише II і III}{усі три}
\solution{I~--- істинне (січна через центр перетинає коло у двох точках). II~--- істинне (діаметр, перпендикулярний до хорди, ділить її навпіл). III~--- хибне (будь-які два діаметри перетинаються в центрі). Істинні I і II. Відповідь: \textbf{В}.}

\vspace{0.3cm}

\zadtask{Які з наведених тверджень є істинними?

\vspace{0.15cm}
\noindent\hspace{1em}I. Дотична до кола перпендикулярна до радіуса, проведеного в точку дотику. \\
\noindent\hspace{1em}II. Через будь-які три точки площини можна провести коло. \\
\noindent\hspace{1em}III. Вписаний кут, що спирається на діаметр кола, є прямим.}
\answerTable{лише I}{лише II і III}{лише I і II}{лише I і III}{усі три}

\vspace{0.3cm}

\zadtask{Які з наведених тверджень є істинними?

\vspace{0.15cm}
\noindent\hspace{1em}I. Центральний кут удвічі більший за вписаний кут, що спирається на ту саму дугу. \\
\noindent\hspace{1em}II. Хорди кола, рівновіддалені від його центра, рівні між собою. \\
\noindent\hspace{1em}III. Відрізки дотичних, проведених до кола з однієї точки поза колом, рівні між собою.}
\answerTable{лише I}{лише II}{лише I і III}{лише II і III}{усі три}

\vspace{0.3cm}

% ---------------------------------------------------------------------
\taskBlock{Площа складеної фігури --- сесія 23.05, завдання №13}

\noindent
\begin{minipage}[c]{0.55\textwidth}
\zadnum Для виготовлення картонної моделі трафарету $ABCDEG$ Андрійко наклеїв квадрат зі стороною $6$~см на прямокутник зі сторонами $4$ і $9$~см (див. рисунок). Спільною частиною квадрата та прямокутника є рівнобедрений трикутник, дві вершини якого збігаються з вершинами прямокутника. Обчисліть площу трафарету $ABCDEG$. \nmtyear{2026}
\end{minipage}%
\hfill
\begin{minipage}[c]{0.43\textwidth}
\centering
\begin{tikzpicture}[scale=0.45, thick]
    \coordinate (A) at (-3,0);
    \coordinate (B) at (0,3);
    \coordinate (Rt) at (3,0);
    \coordinate (G) at (0,-3);
    \coordinate (C) at (1,2);
    \coordinate (K) at (1,-2);
    \coordinate (D) at (7,2);
    \coordinate (E) at (7,-2);
    \fill[gray!35] (C) -- (Rt) -- (K) -- cycle;
    \draw (A) -- (B) -- (Rt) -- (G) -- cycle;
    \draw (C) -- (D) -- (E) -- (K);
    \draw (C) -- (K);
    \node[left] at (A) {$A$};
    \node[above] at (B) {$B$};
    \node[above right] at (C) {$C$};
    \node[above right] at (D) {$D$};
    \node[below right] at (E) {$E$};
    \node[below] at (G) {$G$};
\end{tikzpicture}
\end{minipage}
\par\vspace{0.2cm}
\answerTable{$64$ см$^2$}{$68$ см$^2$}{$72$ см$^2$}{$70$ см$^2$}{$76$ см$^2$}
\solution{Площа квадрата $6\cdot 6=36$ см$^2$, площа прямокутника $4\cdot 9=36$ см$^2$. Спільний рівнобедрений трикутник має основу, рівну стороні прямокутника $4$ см, і висоту $2$ см, тож його площа $\dfrac{1}{2}\cdot 4\cdot 2=4$ см$^2$. Площа трафарету $=36+36-4=68$ см$^2$. Відповідь: \textbf{Б}.}

\vspace{0.3cm}

\noindent
\begin{minipage}[c]{0.55\textwidth}
\zadnum Для виготовлення картонної моделі трафарету $ABCDEG$ Оля наклеїла квадрат зі стороною $8$~см на прямокутник зі сторонами $6$ і $10$~см (див. рисунок). Спільною частиною квадрата та прямокутника є рівнобедрений трикутник, дві вершини якого збігаються з вершинами прямокутника. Обчисліть площу трафарету $ABCDEG$.
\end{minipage}%
\hfill
\begin{minipage}[c]{0.43\textwidth}
\centering
\begin{tikzpicture}[scale=0.36, thick]
    \coordinate (A) at (-4,0);
    \coordinate (B) at (0,4);
    \coordinate (Rt) at (4,0);
    \coordinate (G) at (0,-4);
    \coordinate (C) at (1,3);
    \coordinate (K) at (1,-3);
    \coordinate (D) at (8.5,3);
    \coordinate (E) at (8.5,-3);
    \fill[gray!35] (C) -- (Rt) -- (K) -- cycle;
    \draw (A) -- (B) -- (Rt) -- (G) -- cycle;
    \draw (C) -- (D) -- (E) -- (K);
    \draw (C) -- (K);
    \node[left] at (A) {$A$};
    \node[above] at (B) {$B$};
    \node[above right] at (C) {$C$};
    \node[above right] at (D) {$D$};
    \node[below right] at (E) {$E$};
    \node[below] at (G) {$G$};
\end{tikzpicture}
\end{minipage}
\par\vspace{0.2cm}
\answerTable{$124$ см$^2$}{$109$ см$^2$}{$112$ см$^2$}{$115$ см$^2$}{$118$ см$^2$}

\vspace{0.3cm}

\noindent
\begin{minipage}[c]{0.55\textwidth}
\zadnum Для виготовлення картонної моделі трафарету $ABCDE$ Тарас узяв прямокутник $ABCE$ зі сторонами $8$ і $12$~см і вирізав із нього прямокутний трикутник $CDE$ із катетами $6$ і $8$~см (див. рисунок). Обчисліть площу трафарету $ABCDE$.
\end{minipage}%
\hfill
\begin{minipage}[c]{0.43\textwidth}
\centering
\begin{tikzpicture}[scale=0.30, thick]
    \coordinate (A) at (0,0);
    \coordinate (B) at (0,8);
    \coordinate (C) at (12,8);
    \coordinate (E) at (12,0);
    \coordinate (D) at (12,2);
    \coordinate (Dp) at (6,8);
    \fill[gray!35] (Dp) -- (C) -- (D) -- cycle;
    \draw (A) -- (B) -- (Dp) -- (D) -- (E) -- cycle;
    \draw[dashed] (Dp) -- (C) -- (D);
    \node[below left] at (A) {$A$};
    \node[above left] at (B) {$B$};
    \node[above] at (Dp) {$D$};
    \node[right] at (D) {$C$};
    \node[below right] at (E) {$E$};
\end{tikzpicture}
\end{minipage}
\par\vspace{0.2cm}
\answerTable{$72$ см$^2$}{$84$ см$^2$}{$78$ см$^2$}{$90$ см$^2$}{$96$ см$^2$}

\vspace{0.3cm}

% ---------------------------------------------------------------------
\taskBlock{Планіметрія: кути ромба --- сесія 30.05, завдання №3}

\noindent
\begin{minipage}[c]{0.60\textwidth}
\zadnum Тупий кут ромба дорівнює $140^\circ$. Знайдіть кут, який утворює \textit{більша} діагональ цього ромба з його стороною.
\end{minipage}%
\hfill
\begin{minipage}[c]{0.38\textwidth}
\centering
\begin{tikzpicture}[scale=0.7, thick]
    \coordinate (A) at (-2.5, 0);
    \coordinate (B) at (0, 1.2);
    \coordinate (C) at (2.5, 0);
    \coordinate (D) at (0, -1.2);
    \draw[thick] (A) -- (B) -- (C) -- (D) -- cycle;
    \draw[thick] (A) -- (C);
    \node[left] at (A) {$A$};
    \node[above] at (B) {$B$};
    \node[right] at (C) {$C$};
    \node[below] at (D) {$D$};

\end{tikzpicture}
\end{minipage}
\par\vspace{0.2cm}
\answerTable{$20^\circ$}{$40^\circ$}{$50^\circ$}{$70^\circ$}{$110^\circ$}
\solution{Гострий кут ромба $180^\circ-140^\circ=40^\circ$. Більша діагональ сполучає вершини тупих кутів і ділить навпіл гострий кут, тому з стороною утворює $40^\circ:2=20^\circ$. Відповідь: \textbf{А} ($20^\circ$).}

\noindent
\begin{minipage}[c]{0.60\textwidth}
\zadnum Тупий кут ромба дорівнює $110^\circ$. Знайдіть кут, який утворює \textit{більша} діагональ цього ромба з його стороною.
\end{minipage}%
\hfill
\begin{minipage}[c]{0.38\textwidth}
\centering
\begin{tikzpicture}[scale=0.7, thick]
    \coordinate (A) at (-2.2, 0);
    \coordinate (B) at (0, 1.5);
    \coordinate (C) at (2.2, 0);
    \coordinate (D) at (0, -1.5);
    \draw[thick] (A) -- (B) -- (C) -- (D) -- cycle;
    \draw[thick] (A) -- (C);
    \node[left] at (A) {$A$};
    \node[above] at (B) {$B$};
    \node[right] at (C) {$C$};
    \node[below] at (D) {$D$};
\end{tikzpicture}
\end{minipage}
\par\vspace{0.2cm}
\answerTable{$25^\circ$}{$35^\circ$}{$55^\circ$}{$70^\circ$}{$110^\circ$}

\noindent
\begin{minipage}[c]{0.60\textwidth}
\zadnum Тупий кут ромба дорівнює $150^\circ$. Знайдіть кут, який утворює \textit{менша} діагональ цього ромба з його стороною.
\end{minipage}%
\hfill
\begin{minipage}[c]{0.38\textwidth}
\centering
\begin{tikzpicture}[scale=0.7, thick]
    \coordinate (A) at (-2.7, 0);
    \coordinate (B) at (0, 1.0);
    \coordinate (C) at (2.7, 0);
    \coordinate (D) at (0, -1.0);
    \draw[thick] (A) -- (B) -- (C) -- (D) -- cycle;
    \draw[thick] (A) -- (C);
    \node[left] at (A) {$A$};
    \node[above] at (B) {$B$};
    \node[right] at (C) {$C$};
    \node[below] at (D) {$D$};
\end{tikzpicture}
\end{minipage}
\par\vspace{0.2cm}
\answerTable{$15^\circ$}{$30^\circ$}{$60^\circ$}{$75^\circ$}{$105^\circ$}
\taskBlock{Планіметрія: твердження про трикутники --- сесія 30.05, завдання №9}

\zadtask{Розгляньте наведені нижче твердження.

\vspace{0.2cm}
\noindent
\makebox[1.8em][l]{\textbf{I}}\, Найбільший кут прямокутного трикутника дорівнює половині розгорнутого кута.\par\vspace{0.1cm}
\makebox[1.8em][l]{\textbf{II}}\, Сума двох кутів рівностороннього трикутника дорівнює $120^\circ$.\par\vspace{0.1cm}
\makebox[1.8em][l]{\textbf{III}}\, Усі гострі кути гострокутного трикутника більші за $60^\circ$.\par\vspace{0.25cm}

\noindent Які з тверджень є \textit{правильними}?

\par\vspace{0.2cm}
\begin{flushleft}
\textbf{А} \quad лише I \\[0.1cm]
\textbf{Б} \quad лише II \\[0.1cm]
\textbf{В} \quad лише I та II \\[0.1cm]
\textbf{Г} \quad лише II та III \\[0.1cm]
\textbf{Д} \quad I, II та III
\end{flushleft}}
\solution{\textbf{I}: найбільший кут прямокутного трикутника дорівнює $90^\circ$, а половина розгорнутого кута — $180^\circ:2=90^\circ$ (правильно). \textbf{II}: у рівносторонньому трикутнику всі кути по $60^\circ$, сума двох — $120^\circ$ (правильно). \textbf{III}: гострокутний трикутник може мати кут $\le60^\circ$ (напр. $50^\circ,60^\circ,70^\circ$), тож неправильно. Відповідь: \textbf{В} (лише I та II).}
\vspace{0.3cm}

\zadtask{Розгляньте наведені нижче твердження.

\vspace{0.2cm}
\noindent
\makebox[1.8em][l]{\textbf{I}}\, Сума гострих кутів прямокутного трикутника дорівнює $90^\circ$.\par\vspace{0.1cm}
\makebox[1.8em][l]{\textbf{II}}\, У тупокутному трикутнику два кути можуть бути тупими.\par\vspace{0.1cm}
\makebox[1.8em][l]{\textbf{III}}\, Кожен кут рівностороннього трикутника дорівнює третині розгорнутого кута.\par\vspace{0.25cm}

\noindent Які з тверджень є \textit{правильними}?

\par\vspace{0.2cm}
\begin{flushleft}
\textbf{А} \quad лише I \\[0.1cm]
\textbf{Б} \quad лише III \\[0.1cm]
\textbf{В} \quad лише I та III \\[0.1cm]
\textbf{Г} \quad лише II та III \\[0.1cm]
\textbf{Д} \quad I, II та III
\end{flushleft}}
\vspace{0.3cm}

\zadtask{Розгляньте наведені нижче твердження.

\vspace{0.2cm}
\noindent
\makebox[1.8em][l]{\textbf{I}}\, Зовнішній кут трикутника дорівнює сумі двох внутрішніх кутів, не суміжних з ним.\par\vspace{0.1cm}
\makebox[1.8em][l]{\textbf{II}}\, У рівнобедреному трикутнику кути при основі рівні.\par\vspace{0.1cm}
\makebox[1.8em][l]{\textbf{III}}\, У будь-якому трикутнику найменший кут менший за $90^\circ$.\par\vspace{0.25cm}

\noindent Які з тверджень є \textit{правильними}?

\par\vspace{0.2cm}
\begin{flushleft}
\textbf{А} \quad лише I \\[0.1cm]
\textbf{Б} \quad лише II \\[0.1cm]
\textbf{В} \quad лише I та II \\[0.1cm]
\textbf{Г} \quad лише II та III \\[0.1cm]
\textbf{Д} \quad I, II та III
\end{flushleft}}
\vspace{0.3cm}
\taskBlock{Планіметрія: площа складеної фігури --- сесія 30.05, завдання №15}

\noindent
\begin{minipage}[c]{0.51\textwidth}
\zadnum На рисунку зображено схему поверхні декоративного басейну. Воно має форму прямокутника $ABCD$, з кутів якого вирізано чотири однакові чверті кола радіусом $5$~м, а до сторін $BC$ та $AD$ долучено два однакові півкола радіусом $5$~м. Відстань між основами цих півкіл дорівнює $10$~м, а ширина прямокутної частини між дугами вирізів та півкіл становить $10$~м. Знайдіть площу (у м$^2$) поверхні цього басейну.
\end{minipage}%
\hfill
\begin{minipage}[c]{0.48\textwidth}
\centering
\begin{tikzpicture}[scale=0.18, thick]
    \draw[dashed, gray] (0,0) rectangle (40,20);
    \node[below left] at (0,0) {$A$};
    \node[above left] at (0,20) {$B$};
    \node[above right] at (40,20) {$C$};
    \node[below right] at (40,0) {$D$};
    \draw[fill=mainGreen!35, draw=mainGreen!80!black, thick]
        (40,5)
        -- (40,15)
        arc (270:180:5)
        -- (25,20)
        arc (0:180:5)
        -- (5,20)
        arc (360:270:5)
        -- (0,15)
        -- (0,5)
        arc (90:0:5)
        -- (15,0)
        arc (180:360:5)
        -- (35,0)
        arc (180:90:5)
        -- cycle;
    \draw[<->, >=stealth] (-2, 5) -- (-2, 15) node[midway, fill=white, inner sep=1pt] {$10$};
    \draw[<->, >=stealth] (5, 22) -- (15, 22) node[midway, fill=white, inner sep=1pt] {$10$};
    \draw[<->, >=stealth] (25, 22) -- (35, 22) node[midway, fill=white, inner sep=1pt] {$10$};
    \draw[->, >=stealth, gray!90] (20,20) -- (23.5, 23.5) node[right, text=black, font=\scriptsize] {$R=5$};
    \fill[gray] (20,20) circle [radius=0.3];
    \draw[->, >=stealth, gray!90] (0,20) -- (3.5, 16.5) node[right, text=black, font=\scriptsize] {$R=5$};
    \fill[gray] (0,20) circle [radius=0.3];
\end{tikzpicture}
\end{minipage}
\vspace{0.2cm}
\answerTable{$500$}{$50$}{$1000$}{$800$}{$100$}
\solution{Чотири вирізані чверті кола радіусом $5$ разом дають одне коло, а два долучені півкола радіусом $5$ — також одне коло; їх площі взаємно компенсуються. Тому площа фігури дорівнює площі прямокутника $ABCD$ зі сторонами $50$~м і $20$~м: $50\cdot20=1000$ (м$^2$). Відповідь: \textbf{В} ($1000$).}

\noindent
\begin{minipage}[c]{0.51\textwidth}
\zadnum На рисунку зображено схему поверхні декоративного басейну. Воно має форму прямокутника $ABCD$, з кутів якого вирізано чотири однакові чверті кола радіусом $4$~м, а до сторін $BC$ та $AD$ долучено два однакові півкола радіусом $4$~м. Відстань між основами цих півкіл дорівнює $8$~м, а ширина прямокутної частини між дугами вирізів та півкіл становить $8$~м. Знайдіть площу (у м$^2$) поверхні цього басейну.
\end{minipage}%
\hfill
\begin{minipage}[c]{0.48\textwidth}
\centering
\begin{tikzpicture}[scale=0.225, thick]
    \draw[dashed, gray] (0,0) rectangle (32,16);
    \node[below left] at (0,0) {$A$};
    \node[above left] at (0,16) {$B$};
    \node[above right] at (32,16) {$C$};
    \node[below right] at (32,0) {$D$};
    \draw[fill=mainGreen!35, draw=mainGreen!80!black, thick]
        (32,4)
        -- (32,12)
        arc (270:180:4)
        -- (20,16)
        arc (0:180:4)
        -- (4,16)
        arc (360:270:4)
        -- (0,12)
        -- (0,4)
        arc (90:0:4)
        -- (12,0)
        arc (180:360:4)
        -- (28,0)
        arc (180:90:4)
        -- cycle;
    \draw[<->, >=stealth] (-1.6, 4) -- (-1.6, 12) node[midway, fill=white, inner sep=1pt] {$8$};
    \draw[<->, >=stealth] (4, 17.6) -- (12, 17.6) node[midway, fill=white, inner sep=1pt] {$8$};
    \draw[<->, >=stealth] (20, 17.6) -- (28, 17.6) node[midway, fill=white, inner sep=1pt] {$8$};
    \draw[->, >=stealth, gray!90] (16,16) -- (18.8, 18.8) node[right, text=black, font=\scriptsize] {$R=4$};
    \fill[gray] (16,16) circle [radius=0.25];
    \draw[->, >=stealth, gray!90] (0,16) -- (2.8, 13.2) node[right, text=black, font=\scriptsize] {$R=4$};
    \fill[gray] (0,16) circle [radius=0.25];
\end{tikzpicture}
\end{minipage}
\vspace{0.2cm}
\answerTable{$320$}{$256$}{$512$}{$640$}{$1024$}

\noindent
\begin{minipage}[c]{0.51\textwidth}
\zadnum На рисунку зображено схему поверхні декоративного басейну. Воно має форму \textit{квадрата} $ABCD$ зі стороною $24$~м, з кутів якого вирізано чотири однакові чверті кола радіусом $3$~м, а до сторін $BC$ та $AD$ долучено два однакові півкола радіусом $3$~м. Знайдіть площу (у м$^2$) поверхні цього басейну.
\end{minipage}%
\hfill
\begin{minipage}[c]{0.48\textwidth}
\centering
\begin{tikzpicture}[scale=0.16, thick]
    \draw[dashed, gray] (0,0) rectangle (24,24);
    \node[below left] at (0,0) {$A$};
    \node[above left] at (0,24) {$B$};
    \node[above right] at (24,24) {$C$};
    \node[below right] at (24,0) {$D$};
    \draw[fill=mainGreen!35, draw=mainGreen!80!black, thick]
        (24,3)
        -- (24,21)
        arc (270:180:3)
        -- (15,24)
        arc (0:180:3)
        -- (3,24)
        arc (360:270:3)
        -- (0,21)
        -- (0,3)
        arc (90:0:3)
        -- (9,0)
        arc (180:360:3)
        -- (21,0)
        arc (180:90:3)
        -- cycle;
    \draw[->, >=stealth, gray!90] (12,24) -- (14.5, 26.5) node[right, text=black, font=\scriptsize] {$R=3$};
    \fill[gray] (12,24) circle [radius=0.25];
    \draw[->, >=stealth, gray!90] (0,24) -- (2.5, 21.5) node[right, text=black, font=\scriptsize] {$R=3$};
    \fill[gray] (0,24) circle [radius=0.25];
    \draw[<->, >=stealth] (-1.6, 0) -- (-1.6, 24) node[midway, fill=white, inner sep=1pt] {$24$};
\end{tikzpicture}
\end{minipage}
\vspace{0.2cm}
\answerTable{$144$}{$288$}{$432$}{$576$}{$72$}
\taskBlock{Півколо, поділене на рівні частини --- сесія 1.06, завдання №3}

\noindent
\begin{minipage}[c]{0.58\textwidth}
\zadnum Півколо розділене на $5$ рівних частин (див. рисунок). Визначте градусну міру кута $\alpha$.
\end{minipage}%
\hfill
\begin{minipage}[c]{0.40\textwidth}
\centering
\begin{tikzpicture}[scale=1.0, thick]
    \def\R{2}
    \coordinate (O) at (0,0);
    \draw[mainGreen!80!black] (\R,0) arc (0:180:\R);
    \draw[mainGreen!80!black] (-\R,0) -- (\R,0);
    \foreach \a in {36,72,108,144} {
        \draw[mainGreen!80!black] (0,0) -- (\a:\R);
    }
    % точки на двох сусідніх радіусах для дуги кута
    \coordinate (P1) at (36:\R);
    \coordinate (P2) at (72:\R);
    \pic[draw=yearOrange, line width=0.8pt, "$\alpha$", angle radius=0.85cm,
         angle eccentricity=0.65, font=\small] {angle = P1--O--P2};
\end{tikzpicture}
\end{minipage}
\par\vspace{0.2cm}
\answerTable{$24^\circ$}{$32^\circ$}{$36^\circ$}{$42^\circ$}{$48^\circ$}
\solution{Розгорнутий кут півкола $180^\circ$ поділено на $5$ рівних частин, тож кожна частина (і кут $\alpha$ між сусідніми радіусами) дорівнює $\dfrac{180^\circ}{5}=36^\circ$. Відповідь: \textbf{В}.}

\noindent
\begin{minipage}[c]{0.58\textwidth}
\zadnum Півколо розділене на $6$ рівних частин (див. рисунок). Визначте градусну міру кута $\alpha$.
\end{minipage}%
\hfill
\begin{minipage}[c]{0.40\textwidth}
\centering
\begin{tikzpicture}[scale=1.0, thick]
    \def\R{2}
    \coordinate (O) at (0,0);
    \draw[mainGreen!80!black] (\R,0) arc (0:180:\R);
    \draw[mainGreen!80!black] (-\R,0) -- (\R,0);
    \foreach \a in {30,60,90,120,150} {
        \draw[mainGreen!80!black] (0,0) -- (\a:\R);
    }
    \coordinate (P1) at (30:\R);
    \coordinate (P2) at (60:\R);
    \pic[draw=yearOrange, line width=0.8pt, "$\alpha$", angle radius=0.85cm,
         angle eccentricity=0.65, font=\small] {angle = P1--O--P2};
\end{tikzpicture}
\end{minipage}
\par\vspace{0.2cm}
\answerTable{$20^\circ$}{$25^\circ$}{$30^\circ$}{$36^\circ$}{$60^\circ$}

\noindent
\begin{minipage}[c]{0.58\textwidth}
\zadnum Півколо розділене на $6$ рівних частин (див. рисунок). Визначте градусну міру кута $\alpha$, що охоплює \textit{дві} сусідні частини.
\end{minipage}%
\hfill
\begin{minipage}[c]{0.40\textwidth}
\centering
\begin{tikzpicture}[scale=1.0, thick]
    \def\R{2}
    \coordinate (O) at (0,0);
    \draw[mainGreen!80!black] (\R,0) arc (0:180:\R);
    \draw[mainGreen!80!black] (-\R,0) -- (\R,0);
    \foreach \a in {30,60,90,120,150} {
        \draw[mainGreen!80!black] (0,0) -- (\a:\R);
    }
    \coordinate (P1) at (30:\R);
    \coordinate (P2) at (90:\R);
    \pic[draw=yearOrange, line width=0.8pt, "$\alpha$", angle radius=0.95cm,
         angle eccentricity=0.7, font=\small] {angle = P1--O--P2};
\end{tikzpicture}
\end{minipage}
\par\vspace{0.2cm}
\answerTable{$30^\circ$}{$45^\circ$}{$60^\circ$}{$72^\circ$}{$90^\circ$}

% =====================================================================
\taskBlock{Рівнобедрений трикутник, властивості бісектриси (висоти, медіани) --- сесія 1.06, завдання №9}

\zadtask{У рівнобедреному трикутнику $ABC$ ($AB = BC$) $BK$~--- бісектриса кута $B$, точка $K$ належить стороні $AC$. Які з наведених тверджень є \textit{правильними}?

\vspace{0.2cm}
\noindent
\makebox[1.8em][l]{\textbf{I}}\, $AB + BC = AC$.\par\vspace{0.1cm}
\makebox[1.8em][l]{\textbf{II}}\, $BK \perp AC$.\par\vspace{0.1cm}
\makebox[1.8em][l]{\textbf{III}}\, $AK = KC$.}
\answerTable{лише I}{лише I та II}{лише I та III}{лише II та III}{I, II та III}
\solution{У рівнобедреному трикутнику бісектриса, проведена до основи, є також висотою й медіаною, тому $BK\perp AC$ (II) і $AK=KC$ (III). Рівність $AB+BC=AC$ (I) хибна (нерівність трикутника). Правильні II та III. Відповідь: \textbf{Г}.}

\zadtask{У рівнобедреному трикутнику $ABC$ ($AB = BC$) $BM$~--- медіана, точка $M$ належить стороні $AC$. Які з наведених тверджень є \textit{правильними}?

\vspace{0.2cm}
\noindent
\makebox[1.8em][l]{\textbf{I}}\, $BM$~--- бісектриса кута $B$.\par\vspace{0.1cm}
\makebox[1.8em][l]{\textbf{II}}\, $AM = AB$.\par\vspace{0.1cm}
\makebox[1.8em][l]{\textbf{III}}\, $BM \perp AC$.}
\answerTable{лише I}{лише II}{лише I та III}{лише II та III}{I, II та III}

\zadtask{У рівнобедреному трикутнику $ABC$ ($AB = BC$) $BH$~--- висота, точка $H$ належить стороні $AC$. Які з наведених тверджень є \textit{правильними}?

\vspace{0.2cm}
\noindent
\makebox[1.8em][l]{\textbf{I}}\, $AH = HC$.\par\vspace{0.1cm}
\makebox[1.8em][l]{\textbf{II}}\, $\angle ABH = \angle CBH$.\par\vspace{0.1cm}
\makebox[1.8em][l]{\textbf{III}}\, $AB = AC$.}
\answerTable{лише I}{лише I та II}{лише I та III}{лише II та III}{I, II та III}

% =====================================================================
\taskBlock{Квадрат, радіус описаного кола --- сесія 1.06, завдання №10}

\noindent
\begin{minipage}[c]{0.60\textwidth}
\zadnum На рисунку зображено квадрат $ABCD$. На стороні $AD$ вибрано точку $K$ так, що $\angle BKD = 120^\circ$. Визначте радіус кола, описаного навколо цього квадрата, якщо $BK = 12$ см.
\end{minipage}%
\hfill
\begin{minipage}[c]{0.38\textwidth}
\centering
\begin{tikzpicture}[scale=0.85, thick]
    \coordinate (A) at (0,0);
    \coordinate (D) at (3,0);
    \coordinate (C) at (3,3);
    \coordinate (B) at (0,3);
    \coordinate (K) at (1.7,0);
    \draw[thick] (A) -- (B) -- (C) -- (D) -- cycle;
    \draw[thick] (B) -- (K);
    \fill (K) circle (1.8pt);
    \pic[draw=yearOrange, line width=0.8pt, "$120^\circ$", angle radius=0.7cm,
         angle eccentricity=1.45, font=\scriptsize] {angle = D--K--B};
    \node[below left] at (A) {$A$};
    \node[above left] at (B) {$B$};
    \node[above right] at (C) {$C$};
    \node[below right] at (D) {$D$};
    \node[below] at (K) {$K$};
\end{tikzpicture}
\end{minipage}
\par\vspace{0.2cm}
\answerTable{$6\sqrt{2}$ см}{$3\sqrt{3}$ см}{$6\sqrt{3}$ см}{$3\sqrt{2}$ см}{$3\sqrt{6}$ см}
\solution{$\angle BKA=180^\circ-120^\circ=60^\circ$. У прямокутному трикутнику $ABK$ ($\angle A=90^\circ$): $AB=BK\sin 60^\circ=12\cdot\dfrac{\sqrt3}{2}=6\sqrt3$ --- сторона квадрата. Радіус описаного кола $R=\dfrac{a\sqrt2}{2}=\dfrac{6\sqrt3\cdot\sqrt2}{2}=3\sqrt6$ см. Відповідь: \textbf{Д}.}

\noindent
\begin{minipage}[c]{0.60\textwidth}
\zadnum На рисунку зображено квадрат $ABCD$. На стороні $AD$ вибрано точку $K$ так, що $\angle BKD = 135^\circ$. Визначте радіус кола, описаного навколо цього квадрата, якщо $BK = 8$ см.
\end{minipage}%
\hfill
\begin{minipage}[c]{0.38\textwidth}
\centering
\begin{tikzpicture}[scale=0.85, thick]
    \coordinate (A) at (0,0);
    \coordinate (D) at (3,0);
    \coordinate (C) at (3,3);
    \coordinate (B) at (0,3);
    \coordinate (K) at (2.0,0);
    \draw[thick] (A) -- (B) -- (C) -- (D) -- cycle;
    \draw[thick] (B) -- (K);
    \fill (K) circle (1.8pt);
    \pic[draw=yearOrange, line width=0.8pt, "$135^\circ$", angle radius=0.55cm,
         angle eccentricity=1.55, font=\scriptsize] {angle = D--K--B};
    \node[below left] at (A) {$A$};
    \node[above left] at (B) {$B$};
    \node[above right] at (C) {$C$};
    \node[below right] at (D) {$D$};
    \node[below] at (K) {$K$};
\end{tikzpicture}
\end{minipage}
\par\vspace{0.2cm}
\answerTable{$2\sqrt{2}$ см}{$4$ см}{$4\sqrt{2}$ см}{$8$ см}{$8\sqrt{2}$ см}

\noindent
\begin{minipage}[c]{0.60\textwidth}
\zadnum На рисунку зображено квадрат $ABCD$. На стороні $AD$ вибрано точку $K$ так, що $\angle BKD = 150^\circ$. Визначте \textit{площу} цього квадрата, якщо $BK = 10$ см.
\end{minipage}%
\hfill
\begin{minipage}[c]{0.38\textwidth}
\centering
\begin{tikzpicture}[scale=0.85, thick]
    \coordinate (A) at (0,0);
    \coordinate (D) at (3,0);
    \coordinate (C) at (3,3);
    \coordinate (B) at (0,3);
    \coordinate (K) at (2.3,0);
    \draw[thick] (A) -- (B) -- (C) -- (D) -- cycle;
    \draw[thick] (B) -- (K);
    \fill (K) circle (1.8pt);
    \pic[draw=yearOrange, line width=0.8pt, "$150^\circ$", angle radius=0.45cm,
         angle eccentricity=1.7, font=\scriptsize] {angle = D--K--B};
    \node[below left] at (A) {$A$};
    \node[above left] at (B) {$B$};
    \node[above right] at (C) {$C$};
    \node[below right] at (D) {$D$};
    \node[below] at (K) {$K$};
\end{tikzpicture}
\end{minipage}
\par\vspace{0.2cm}
\answerTable{$12{,}5$ см$^2$}{$25$ см$^2$}{$50$ см$^2$}{$25\sqrt{2}$ см$^2$}{$100$ см$^2$}

% =====================================================================
\taskBlock{Площа рівнобічної трапеції --- сесія 2.06, завдання №3}

\noindent
\begin{minipage}[c]{0.54\textwidth}
\zadnum У рівнобічній трапеції $ABCD$ бічна сторона дорівнює $6$, а гострий кут при більшій основі дорівнює $60^\circ$. З вершин меншої основи $B$ і $C$ проведено відрізки до спільної точки $P$ на більшій основі так, що $BP \perp AB$ і $CP \perp CD$ (див. рисунок). Знайдіть площу трапеції.
\end{minipage}%
\hfill
\begin{minipage}[c]{0.44\textwidth}
\centering
\begin{tikzpicture}[scale=0.32, thick]
    \coordinate (A) at (0,0);
    \coordinate (D) at (24,0);
    \coordinate (B) at (3,5.196);
    \coordinate (C) at (21,5.196);
    \coordinate (P) at (12,0);
    \draw[thick] (A) -- (D) -- (C) -- (B) -- cycle;
    \draw[thick, mainGreen] (B) -- (P) -- (C);
    \pic[draw=yearOrange, line width=0.7pt, angle radius=0.3cm] {right angle = A--B--P};
    \pic[draw=yearOrange, line width=0.7pt, angle radius=0.3cm] {right angle = P--C--D};
    \pic[draw=yearOrange, line width=0.7pt, pic text={\scriptsize $60^\circ$}, angle radius=0.3cm, angle eccentricity=1.9] {angle = D--A--B};
    \node[below left] at (A) {$A$};
    \node[above left] at (B) {$B$};
    \node[above right] at (C) {$C$};
    \node[below right] at (D) {$D$};
    \node[below] at (P) {$P$};
    \node[above left, font=\scriptsize] at ($(B)!0.5!(A)$) {$6$};
\end{tikzpicture}
\end{minipage}
\par\vspace{0.2cm}
\answerTable{$63$}{$63\sqrt{3}$}{$126$}{$126\sqrt{3}$}{$252$}
\solution{Висота трапеції $h=AB\sin60^\circ=6\cdot\dfrac{\sqrt3}{2}=3\sqrt3$. У прямокутному $\triangle ABP$ ($BP\perp AB$, $\angle A=60^\circ$) маємо $AP=\dfrac{AB}{\cos60^\circ}=\dfrac{6}{0{,}5}=12$, тож більша основа $AD=2\cdot AP=24$. Проєкція бічної сторони $AB\cos60^\circ=3$, тому менша основа $BC=AD-2\cdot3=18$. Площа $S=\dfrac{AD+BC}{2}\cdot h=\dfrac{24+18}{2}\cdot3\sqrt3=63\sqrt3$. Відповідь: \textbf{Б}.}

\noindent
\begin{minipage}[c]{0.54\textwidth}
\zadnum У рівнобічній трапеції $ABCD$ бічна сторона дорівнює $4$, а гострий кут при більшій основі дорівнює $60^\circ$. З вершин меншої основи $B$ і $C$ проведено відрізки до спільної точки $P$ на більшій основі так, що $BP \perp AB$ і $CP \perp CD$ (див. рисунок). Знайдіть площу трапеції.
\end{minipage}%
\hfill
\begin{minipage}[c]{0.44\textwidth}
\centering
\begin{tikzpicture}[scale=0.45, thick]
    \coordinate (A) at (0,0);
    \coordinate (D) at (16,0);
    \coordinate (B) at (2,3.464);
    \coordinate (C) at (14,3.464);
    \coordinate (P) at (8,0);
    \draw[thick] (A) -- (D) -- (C) -- (B) -- cycle;
    \draw[thick, mainGreen] (B) -- (P) -- (C);
    \pic[draw=yearOrange, line width=0.7pt, angle radius=0.3cm] {right angle = A--B--P};
    \pic[draw=yearOrange, line width=0.7pt, angle radius=0.3cm] {right angle = P--C--D};
    \pic[draw=yearOrange, line width=0.7pt, pic text={\scriptsize $60^\circ$}, angle radius=0.3cm, angle eccentricity=1.9] {angle = D--A--B};
    \node[below left] at (A) {$A$};
    \node[above left] at (B) {$B$};
    \node[above right] at (C) {$C$};
    \node[below right] at (D) {$D$};
    \node[below] at (P) {$P$};
    \node[above left, font=\scriptsize] at ($(B)!0.5!(A)$) {$4$};
\end{tikzpicture}
\end{minipage}
\par\vspace{0.2cm}
\answerTable{$28$}{$14\sqrt{3}$}{$28\sqrt{3}$}{$56\sqrt{3}$}{$56$}

\noindent
\begin{minipage}[c]{0.54\textwidth}
\zadnum У рівнобічній трапеції $ABCD$ бічна сторона дорівнює $8$, а гострий кут при більшій основі дорівнює $45^\circ$. З вершин меншої основи $B$ і $C$ проведено відрізки до спільної точки $P$ на більшій основі так, що $BP \perp AB$ і $CP \perp CD$ (див. рисунок). Знайдіть площу трапеції.
\end{minipage}%
\hfill
\begin{minipage}[c]{0.44\textwidth}
\centering
\begin{tikzpicture}[scale=0.3, thick]
    \coordinate (A) at (0,0);
    \coordinate (D) at (22.627,0);
    \coordinate (B) at (5.657,5.657);
    \coordinate (C) at (16.971,5.657);
    \coordinate (P) at (11.314,0);
    \draw[thick] (A) -- (D) -- (C) -- (B) -- cycle;
    \draw[thick, mainGreen] (B) -- (P) -- (C);
    \pic[draw=yearOrange, line width=0.7pt, angle radius=0.35cm] {right angle = A--B--P};
    \pic[draw=yearOrange, line width=0.7pt, angle radius=0.35cm] {right angle = P--C--D};
    \pic[draw=yearOrange, line width=0.7pt, pic text={\scriptsize $45^\circ$}, angle radius=0.35cm, angle eccentricity=1.9] {angle = D--A--B};
    \node[below left] at (A) {$A$};
    \node[above left] at (B) {$B$};
    \node[above right] at (C) {$C$};
    \node[below right] at (D) {$D$};
    \node[below] at (P) {$P$};
    \node[above left, font=\scriptsize] at ($(B)!0.5!(A)$) {$8$};
\end{tikzpicture}
\end{minipage}
\par\vspace{0.2cm}
\answerTable{$48$}{$96$}{$96\sqrt{2}$}{$192$}{$64\sqrt{2}$}

% =====================================================================
\taskBlock{Кути рівнобедреного трикутника --- сесія 2.06, завдання №4}

\noindent
\begin{minipage}[c]{0.56\textwidth}
\zadnum На рисунку зображено клумбу у формі трикутника $ABC$, у якого $AB = BC$. Знайдіть градусну міру кута $B$, якщо він \textit{удвічі} більший за градусну міру кута $C$.
\end{minipage}%
\hfill
\begin{minipage}[c]{0.40\textwidth}
\centering
\begin{tikzpicture}[scale=0.5, thick]
    \coordinate (A) at (0,0);
    \coordinate (C) at (8,0);
    \coordinate (B) at (4,4);
    \draw[fill=mainGreen!15, draw=mainGreen!80!black, thick] (A) -- (B) -- (C) -- cycle;
    \node[below left] at (A) {$A$};
    \node[above] at (B) {$B$};
    \node[below right] at (C) {$C$};
    \pic[draw=yearOrange, line width=0.7pt, angle radius=0.3cm] {angle = C--A--B};
    \pic[draw=yearOrange, line width=0.7pt, angle radius=0.3cm] {angle = B--C--A};
    \pic[draw=yearOrange, line width=0.7pt, angle radius=0.3cm] {angle = A--B--C};
\end{tikzpicture}
\end{minipage}
\par\vspace{0.2cm}
\answerTable{$60^\circ$}{$30^\circ$}{$45^\circ$}{$90^\circ$}{$120^\circ$}
\solution{Оскільки $AB=BC$, кути при основі рівні: $\angle A=\angle C$. За умовою $\angle B=2\angle C$. Сума кутів: $\angle C+\angle C+2\angle C=180^\circ$, $4\angle C=180^\circ$, $\angle C=45^\circ$, тож $\angle B=90^\circ$. Відповідь: \textbf{Г}.}

\noindent
\begin{minipage}[c]{0.56\textwidth}
\zadnum На рисунку зображено ділянку у формі трикутника $ABC$, у якого $AB = BC$. Знайдіть градусну міру кута $B$, якщо він \textit{утричі} більший за градусну міру кута $C$.
\end{minipage}%
\hfill
\begin{minipage}[c]{0.40\textwidth}
\centering
\begin{tikzpicture}[scale=0.5, thick]
    \coordinate (A) at (0,0);
    \coordinate (C) at (8,0);
    \coordinate (B) at (4,2.9);
    \draw[fill=mainGreen!15, draw=mainGreen!80!black, thick] (A) -- (B) -- (C) -- cycle;
    \node[below left] at (A) {$A$};
    \node[above] at (B) {$B$};
    \node[below right] at (C) {$C$};
    \pic[draw=yearOrange, line width=0.7pt, angle radius=0.3cm] {angle = C--A--B};
    \pic[draw=yearOrange, line width=0.7pt, angle radius=0.3cm] {angle = B--C--A};
    \pic[draw=yearOrange, line width=0.7pt, angle radius=0.3cm] {angle = A--B--C};
\end{tikzpicture}
\end{minipage}
\par\vspace{0.2cm}
\answerTable{$36^\circ$}{$72^\circ$}{$90^\circ$}{$108^\circ$}{$120^\circ$}

\noindent
\begin{minipage}[c]{0.56\textwidth}
\zadnum На рисунку зображено газон у формі трикутника $ABC$, у якого $AB = BC$. Знайдіть градусну міру кута $A$ при основі, якщо кут $B$ \textit{на $30^\circ$ більший} за кут $A$.
\end{minipage}%
\hfill
\begin{minipage}[c]{0.40\textwidth}
\centering
\begin{tikzpicture}[scale=0.5, thick]
    \coordinate (A) at (0,0);
    \coordinate (C) at (8,0);
    \coordinate (B) at (4,4.77);
    \draw[fill=mainGreen!15, draw=mainGreen!80!black, thick] (A) -- (B) -- (C) -- cycle;
    \node[below left] at (A) {$A$};
    \node[above] at (B) {$B$};
    \node[below right] at (C) {$C$};
    \pic[draw=yearOrange, line width=0.7pt, angle radius=0.3cm] {angle = C--A--B};
    \pic[draw=yearOrange, line width=0.7pt, angle radius=0.3cm] {angle = B--C--A};
    \pic[draw=yearOrange, line width=0.7pt, angle radius=0.3cm] {angle = A--B--C};
\end{tikzpicture}
\end{minipage}
\par\vspace{0.2cm}
\answerTable{$40^\circ$}{$50^\circ$}{$65^\circ$}{$80^\circ$}{$100^\circ$}

% =====================================================================
\taskBlock{Твердження про елементи трикутника та чотирикутника --- сесія 2.06, завдання №9}

\zadtask{Які з наведених тверджень є \textit{правильними}?

\vspace{0.1cm}
\begin{enumerate}[label=\textbf{\Roman*}, align=left, leftmargin=2.8em, labelsep=0.6em, itemsep=0.12cm, topsep=0.1cm]
\item У правильному трикутнику всі медіани рівні.
\item Медіана трикутника ділить його на два трикутники з рівними периметрами.
\item Медіана трикутника ділить його на два трикутники з рівними площами.
\end{enumerate}
\vspace{0.15cm}}
\answerTable{лише I}{лише II}{лише I та III}{лише II та III}{I, II та III}
\solution{I -- правильно (у правильному трикутнику всі медіани рівні). II -- хибне: медіана ділить лише сторону навпіл, периметри частин зазвичай різні. III -- правильно: медіана ділить трикутник на два рівновеликі (рівні за площею). Отже правильні I та III. Відповідь: \textbf{В}.}

\zadtask{Які з наведених тверджень є \textit{правильними}?

\vspace{0.1cm}
\begin{enumerate}[label=\textbf{\Roman*}, align=left, leftmargin=2.8em, labelsep=0.6em, itemsep=0.12cm, topsep=0.1cm]
\item У рівносторонньому трикутнику всі висоти рівні.
\item Висота трикутника завжди лежить усередині трикутника.
\item Бісектриса будь-якого трикутника ділить протилежну сторону навпіл.
\end{enumerate}
\vspace{0.15cm}}
\answerTable{лише I}{лише II}{лише I та II}{лише I та III}{I, II та III}

\zadtask{Які з наведених тверджень є \textit{правильними}?

\vspace{0.1cm}
\begin{enumerate}[label=\textbf{\Roman*}, align=left, leftmargin=2.8em, labelsep=0.6em, itemsep=0.12cm, topsep=0.1cm]
\item Діагоналі ромба перпендикулярні.
\item Діагоналі ромба рівні.
\item Діагоналі ромба є бісектрисами його кутів.
\end{enumerate}
\vspace{0.15cm}}
\answerTable{лише I}{лише II}{лише I та III}{лише II та III}{I, II та III}

% =====================================================================
\taskBlock{Планіметрія: кути при перетині прямих --- сесія 3.06, завдання №4}

\noindent
\begin{minipage}[c]{0.58\textwidth}
\zadnum На рисунку схематично зображено підставку для планшета, що стоїть на горизонтальній поверхні (пряма $AD$). Точки $A$, $B$, $C$ і $K$ лежать на одній прямій. Визначте градусну міру кута $BAD$, якщо $\angle BKD = 128^\circ$, $\angle ADK = 68^\circ$.
\end{minipage}%
\hfill
\begin{minipage}[c]{0.40\textwidth}
\centering
\begin{tikzpicture}[scale=0.95, thick]
    \coordinate (A) at (0,0);
    \coordinate (D) at (2.4,0);
    \coordinate (B) at (1.7,3.0);
    \coordinate (K) at (1.25,2.2);
    \draw[thick] (A) -- (B);
    \draw[thick] (D) -- (K);
    \draw[thick] (A) -- (D);
    \node[above] at (B) {\scriptsize $B$};
    \node[below left] at (A) {\scriptsize $A$};
    \node[below right] at (D) {\scriptsize $D$};
    \node[left] at (K) {\scriptsize $K$};
    \pic[draw=yearOrange, line width=0.6pt, angle radius=0.4cm] {angle = D--K--B};
    \pic[draw=mainGreen, line width=0.6pt, angle radius=0.45cm] {angle = K--D--A};
    \pic[draw=yearOrange, line width=0.6pt, angle radius=0.4cm] {angle = D--A--B};
    \node[font=\tiny] at (1.6,2.05) {$128^\circ$};
    \node[font=\tiny] at (1.95,0.27) {$68^\circ$};
    \node[font=\tiny] at (0.42,0.18) {$?$};
\end{tikzpicture}
\end{minipage}
\par\vspace{0.2cm}
\answerTable{$66^\circ$}{$62^\circ$}{$52^\circ$}{$60^\circ$}{$50^\circ$}
\solution{Точки $A$, $B$, $K$ лежать на одній прямій, тому $\angle DKA = 180^\circ - \angle BKD = 180^\circ - 128^\circ = 52^\circ$. У трикутнику $AKD$: $\angle BAD = 180^\circ - \angle ADK - \angle DKA = 180^\circ - 68^\circ - 52^\circ = 60^\circ$. Відповідь: \textbf{Г} ($60^\circ$).}

\noindent
\begin{minipage}[c]{0.58\textwidth}
\zadnum На рисунку схематично зображено підпірку для дзеркала, що стоїть на горизонтальній поверхні (пряма $AD$). Точки $A$, $B$, $C$ і $K$ лежать на одній прямій. Визначте градусну міру кута $BAD$, якщо $\angle BKD = 116^\circ$, $\angle ADK = 54^\circ$.
\end{minipage}%
\hfill
\begin{minipage}[c]{0.40\textwidth}
\centering
\begin{tikzpicture}[scale=0.95, thick]
    \coordinate (A) at (0,0);
    \coordinate (D) at (2.6,0);
    \coordinate (B) at (1.8,2.9);
    \coordinate (K) at (1.3,2.1);
    \draw[thick] (A) -- (B);
    \draw[thick] (D) -- (K);
    \draw[thick] (A) -- (D);
    \node[above] at (B) {\scriptsize $B$};
    \node[below left] at (A) {\scriptsize $A$};
    \node[below right] at (D) {\scriptsize $D$};
    \node[left] at (K) {\scriptsize $K$};
    \pic[draw=yearOrange, line width=0.6pt, angle radius=0.4cm] {angle = D--K--B};
    \pic[draw=mainGreen, line width=0.6pt, angle radius=0.45cm] {angle = K--D--A};
    \pic[draw=yearOrange, line width=0.6pt, angle radius=0.4cm] {angle = D--A--B};
    \node[font=\tiny] at (1.7,1.95) {$116^\circ$};
    \node[font=\tiny] at (2.1,0.27) {$54^\circ$};
    \node[font=\tiny] at (0.45,0.18) {$?$};
\end{tikzpicture}
\end{minipage}
\par\vspace{0.2cm}
\answerTable{$54^\circ$}{$58^\circ$}{$62^\circ$}{$64^\circ$}{$66^\circ$}

\noindent
\begin{minipage}[c]{0.58\textwidth}
\zadnum На рисунку схематично зображено підставку для книжки, що стоїть на горизонтальній поверхні (пряма $AD$). Точки $A$, $B$, $C$ і $K$ лежать на одній прямій. Визначте градусну міру кута $ADK$, якщо $\angle BKD = 140^\circ$, $\angle BAD = 85^\circ$.
\end{minipage}%
\hfill
\begin{minipage}[c]{0.40\textwidth}
\centering
\begin{tikzpicture}[scale=0.95, thick]
    \coordinate (A) at (0,0);
    \coordinate (D) at (2.5,0);
    \coordinate (B) at (1.6,3.0);
    \coordinate (K) at (1.2,2.25);
    \draw[thick] (A) -- (B);
    \draw[thick] (D) -- (K);
    \draw[thick] (A) -- (D);
    \node[above] at (B) {\scriptsize $B$};
    \node[below left] at (A) {\scriptsize $A$};
    \node[below right] at (D) {\scriptsize $D$};
    \node[left] at (K) {\scriptsize $K$};
    \pic[draw=yearOrange, line width=0.6pt, angle radius=0.4cm] {angle = D--K--B};
    \pic[draw=mainGreen, line width=0.6pt, angle radius=0.45cm] {angle = K--D--A};
    \pic[draw=yearOrange, line width=0.6pt, angle radius=0.4cm] {angle = D--A--B};
    \node[font=\tiny] at (1.6,2.05) {$140^\circ$};
    \node[font=\tiny] at (2.0,0.3) {$?$};
    \node[font=\tiny] at (0.42,0.18) {$85^\circ$};
\end{tikzpicture}
\end{minipage}
\par\vspace{0.2cm}
\answerTable{$55^\circ$}{$60^\circ$}{$65^\circ$}{$70^\circ$}{$75^\circ$}
\taskBlock{Планіметрія: властивості кола --- сесія 3.06, завдання №9}

\zadnum Які з наведених тверджень є \textit{правильними}?

\vspace{0.1cm}
\begin{enumerate}[label=\textbf{\Roman*.}, align=left, leftmargin=2.6em, labelsep=0.5em, itemsep=0.1cm, topsep=0.08cm]
\item Будь-яка хорда кола більша за його радіус.
\item Кінці діаметра кола ділять коло навпіл.
\item Рівні хорди кола стягують рівні дуги.
\end{enumerate}
\vspace{0.12cm}

\answerTable{лише II}{лише I та II}{лише III}{лише II та III}{I, II та III}
\solution{I~--- хибне: хорда може бути меншою за радіус. II~--- правильне: кінці діаметра ділять коло на дві рівні півкола. III~--- правильне: рівні хорди стягують рівні дуги. Отже правильні лише II та III. Відповідь: \textbf{Г}.}

\zadnum Які з наведених тверджень є \textit{правильними}?

\vspace{0.1cm}
\begin{enumerate}[label=\textbf{\Roman*.}, align=left, leftmargin=2.6em, labelsep=0.5em, itemsep=0.1cm, topsep=0.08cm]
\item Діаметр є найбільшою хордою кола.
\item Будь-які два радіуси одного кола рівні між собою.
\item Будь-який діаметр кола ділить будь-яку його хорду навпіл.
\end{enumerate}
\vspace{0.12cm}

\answerTable{лише I}{лише I та II}{лише II та III}{лише III}{I, II та III}

\zadnum Які з наведених тверджень є \textit{неправильними}?

\vspace{0.1cm}
\begin{enumerate}[label=\textbf{\Roman*.}, align=left, leftmargin=2.6em, labelsep=0.5em, itemsep=0.1cm, topsep=0.08cm]
\item Через будь-які три точки площини можна провести коло.
\item Центр кола, описаного навколо трикутника, завжди лежить усередині трикутника.
\item Серединний перпендикуляр до хорди проходить через центр кола.
\end{enumerate}
\vspace{0.12cm}

\answerTable{лише I}{лише I та II}{лише II}{лише II та III}{I, II та III}
\taskBlock{Планіметрія: прямокутник і вписане коло --- сесія 3.06, завдання №14}

\noindent
\begin{minipage}[c]{0.56\textwidth}
\zadnum На рисунку зображено прямокутник $ABCD$ та коло з центром $O$, яке дотикається до сторін $AB$, $AD$, а також сторони $BC$ у точці $K$. Визначте площу чотирикутника $OKCD$, якщо $OD = 24$ см, $\angle ADO = 30^\circ$.
\end{minipage}%
\hfill
\begin{minipage}[c]{0.40\textwidth}
\centering
\begin{tikzpicture}[scale=0.42, thick]
    \coordinate (A) at (0,0);
    \coordinate (B) at (0,5);
    \coordinate (Cc) at (8.7,5);
    \coordinate (D) at (8.7,0);
    \coordinate (O) at (2.5,2.5);
    \coordinate (K) at (2.5,5);
    \draw[thick] (A) -- (B) -- (Cc) -- (D) -- cycle;
    \draw[thick, mainGreen] (O) circle (2.5);
    \draw[thick] (O) -- (D);
    \draw[thick] (O) -- (K);
    \fill (O) circle (2.4pt); \node[left, font=\scriptsize] at (O) {$O$};
    \fill (K) circle (2.4pt); \node[above, font=\scriptsize] at (K) {$K$};
    \node[below left, font=\scriptsize] at (A) {$A$};
    \node[above left, font=\scriptsize] at (B) {$B$};
    \node[above right, font=\scriptsize] at (Cc) {$C$};
    \node[below right, font=\scriptsize] at (D) {$D$};
    \node[font=\tiny] at (7.2,0.5) {$30^\circ$};
\end{tikzpicture}
\end{minipage}
\par\vspace{0.2cm}
\answerTable{$288$ см$^2$}{$576$ см$^2$}{$216\sqrt{3}$ см$^2$}{$288\sqrt{3}$ см$^2$}{$576\sqrt{3}$ см$^2$}
\solution{Радіус $r$ дорівнює відстані від $O$ до $AD$: $r = OD\sin30^\circ = 24\cdot\dfrac12 = 12$ см, а $OK = r = 12$. Сторона $CD$ (висота прямокутника) дорівнює діаметру: $CD = 2r = 24$ см. Відстань між паралельними $OK$ і $CD$ дорівнює $OD\cos30^\circ = 24\cdot\dfrac{\sqrt3}{2} = 12\sqrt3$. $OKCD$~--- прямокутна трапеція: $S = \dfrac{OK + CD}{2}\cdot 12\sqrt3 = \dfrac{12+24}{2}\cdot 12\sqrt3 = 216\sqrt3$ см$^2$. Відповідь: \textbf{В}.}

\noindent
\begin{minipage}[c]{0.56\textwidth}
\zadnum На рисунку зображено прямокутник $ABCD$ та коло з центром $O$, яке дотикається до сторін $AB$, $AD$, а також сторони $BC$ у точці $K$. Визначте площу чотирикутника $OKCD$, якщо $OD = 16$ см, $\angle ADO = 30^\circ$.
\end{minipage}%
\hfill
\begin{minipage}[c]{0.40\textwidth}
\centering
\begin{tikzpicture}[scale=0.42, thick]
    \coordinate (A) at (0,0);
    \coordinate (B) at (0,5);
    \coordinate (Cc) at (8.7,5);
    \coordinate (D) at (8.7,0);
    \coordinate (O) at (2.5,2.5);
    \coordinate (K) at (2.5,5);
    \draw[thick] (A) -- (B) -- (Cc) -- (D) -- cycle;
    \draw[thick, mainGreen] (O) circle (2.5);
    \draw[thick] (O) -- (D);
    \draw[thick] (O) -- (K);
    \fill (O) circle (2.4pt); \node[left, font=\scriptsize] at (O) {$O$};
    \fill (K) circle (2.4pt); \node[above, font=\scriptsize] at (K) {$K$};
    \node[below left, font=\scriptsize] at (A) {$A$};
    \node[above left, font=\scriptsize] at (B) {$B$};
    \node[above right, font=\scriptsize] at (Cc) {$C$};
    \node[below right, font=\scriptsize] at (D) {$D$};
    \node[font=\tiny] at (7.2,0.5) {$30^\circ$};
\end{tikzpicture}
\end{minipage}
\par\vspace{0.2cm}
\answerTable{$96$ см$^2$}{$96\sqrt{3}$ см$^2$}{$128\sqrt{3}$ см$^2$}{$192\sqrt{3}$ см$^2$}{$192$ см$^2$}

\noindent
\begin{minipage}[c]{0.56\textwidth}
\zadnum На рисунку зображено прямокутник $ABCD$ та коло з центром $O$, яке дотикається до сторін $AB$, $AD$, а також сторони $BC$ у точці $K$. Визначте площу трикутника $OKC$, якщо $OD = 12\sqrt{2}$ см, $\angle ADO = 45^\circ$.
\end{minipage}%
\hfill
\begin{minipage}[c]{0.40\textwidth}
\centering
\begin{tikzpicture}[scale=0.42, thick]
    \coordinate (A) at (0,0);
    \coordinate (B) at (0,5);
    \coordinate (Cc) at (6.5,5);
    \coordinate (D) at (6.5,0);
    \coordinate (O) at (2.5,2.5);
    \coordinate (K) at (2.5,5);
    \draw[thick] (A) -- (B) -- (Cc) -- (D) -- cycle;
    \draw[thick, mainGreen] (O) circle (2.5);
    \draw[thick] (O) -- (D);
    \draw[thick] (O) -- (K);
    \draw[thick] (O) -- (Cc);
    \fill (O) circle (2.4pt); \node[left, font=\scriptsize] at (O) {$O$};
    \fill (K) circle (2.4pt); \node[above, font=\scriptsize] at (K) {$K$};
    \node[below left, font=\scriptsize] at (A) {$A$};
    \node[above left, font=\scriptsize] at (B) {$B$};
    \node[above right, font=\scriptsize] at (Cc) {$C$};
    \node[below right, font=\scriptsize] at (D) {$D$};
    \node[font=\tiny] at (5.3,0.55) {$45^\circ$};
\end{tikzpicture}
\end{minipage}
\par\vspace{0.2cm}
\answerTable{$72$ см$^2$}{$144$ см$^2$}{$72\sqrt{2}$ см$^2$}{$36$ см$^2$}{$96$ см$^2$}
\taskBlock{Точка на відрізку --- сесія 4.06, завдання №3}
\zadtask{На відрізку $AB$ вибрано точку $K$ так, що відрізок $AK$ на $5$~\textit{см} більший за відрізок $BK$. Знайдіть довжину відрізка $AK$, якщо $AB = 21$~\textit{см}.}
\answerTable{$16$~\textit{см}}{$14$~\textit{см}}{$3$~\textit{см}}{$8$~\textit{см}}{$13$~\textit{см}}
\solution{Нехай $BK=x$, тоді $AK=x+5$, а $AK+BK=AB$: $x+5+x=21$, $2x=16$, $x=8$. Отже $AK=8+5=13$~см. Відповідь: \textbf{Д}.}
\zadtask{На відрізку $AB$ вибрано точку $K$ так, що відрізок $AK$ на $7$~\textit{см} більший за відрізок $BK$. Знайдіть довжину відрізка $AK$, якщо $AB = 25$~\textit{см}.}
\answerTable{$9$~\textit{см}}{$16$~\textit{см}}{$18$~\textit{см}}{$12$~\textit{см}}{$17$~\textit{см}}
\zadtask{На відрізку $MN$ вибрано точку $P$ так, що відрізок $MP$ на $4$~\textit{см} менший за відрізок $PN$. Знайдіть довжину відрізка $PN$, якщо $MN = 18$~\textit{см}.}
\answerTable{$7$~\textit{см}}{$14$~\textit{см}}{$9$~\textit{см}}{$11$~\textit{см}}{$13$~\textit{см}}
\taskBlock{Геометричні твердження (коло і трикутник) --- сесія 4.06, завдання №9}
\zadtask{Які з наведених тверджень є правильними?

\vspace{0.1cm}
\par\noindent\hangindent=2.2em\textbf{I.}\ \ Коло, уписане в рівнобедрений трикутник, дотикається до середини його основи.
\par\noindent\hangindent=2.2em\textbf{II.}\ \ Центр кола, описаного навколо тупокутного трикутника, лежить усередині трикутника.
\par\noindent\hangindent=2.2em\textbf{III.}\ \ Коло, побудоване на гіпотенузі прямокутного трикутника як на діаметрі, проходить через вершину прямого кута.
\vspace{0.15cm}
\answerTable{лише I}{лише I та II}{лише I та III}{лише II та III}{I, II та III}}
\solution{I~--- правильне (точка дотику вписаного кола з основою рівнобедреного трикутника збігається з її серединою). II~--- хибне: центр описаного кола тупокутного трикутника лежить \emph{поза} трикутником. III~--- правильне (вписаний кут, що спирається на діаметр, прямий). Правильні I та III. Відповідь: \textbf{В}.}
\zadtask{Які з наведених тверджень є правильними?

\vspace{0.1cm}
\par\noindent\hangindent=2.2em\textbf{I.}\ \ Центр кола, описаного навколо прямокутного трикутника, лежить на середині гіпотенузи.
\par\noindent\hangindent=2.2em\textbf{II.}\ \ Коло, уписане в трикутник, дотикається до всіх трьох його сторін.
\par\noindent\hangindent=2.2em\textbf{III.}\ \ Центр кола, описаного навколо гострокутного трикутника, лежить поза трикутником.
\vspace{0.15cm}
\answerTable{лише I}{лише I та II}{лише I та III}{лише II та III}{I, II та III}}
\zadtask{Які з наведених тверджень є правильними?

\vspace{0.1cm}
\par\noindent\hangindent=2.2em\textbf{I.}\ \ Центр кола, уписаного в трикутник, є точкою перетину його бісектрис.
\par\noindent\hangindent=2.2em\textbf{II.}\ \ Радіус кола, описаного навколо рівностороннього трикутника, удвічі більший за радіус кола, уписаного в цей трикутник.
\par\noindent\hangindent=2.2em\textbf{III.}\ \ Центр кола, описаного навколо будь-якого трикутника, завжди лежить усередині трикутника.
\vspace{0.15cm}
\answerTable{лише I}{лише I та II}{лише I та III}{лише II та III}{I, II та III}}
\taskBlock{Площа паралелограма --- сесія 4.06, завдання №14}
\noindent
\begin{minipage}[c]{0.54\textwidth}
\zadnum Точка $K$~--- середина висоти $BM$ паралелограма $ABCD$ (див. рисунок). $AB = 8\sqrt{3}$~см, $DK = 12$~см, $\angle ADK = 30^\circ$. Визначте площу паралелограма $ABCD$.
\end{minipage}%
\hfill
\begin{minipage}[c]{0.44\textwidth}
\centering
\begin{tikzpicture}[scale=0.55, thick]
    % Координати вершин паралелограма
    \coordinate (A) at (0,0);
    \coordinate (M) at (2.5,0);
    \coordinate (D) at (7.5,0);
    \coordinate (B) at (2.5,4);
    \coordinate (C) at (10,4);
    \coordinate (K) at (2.5,2); % K - середина BM

    % Паралелограм і відрізки
    \draw[thick] (A) -- (B) -- (C) -- (D) -- cycle;
    \draw[thick] (B) -- (M);
    \draw[thick] (K) -- (D);

    % Точка K
    \fill (K) circle (2.5pt);

    % Позначка прямого кута при основі (з правого боку BM)
    \draw (M) ++(0,0.3) -- ++(0.3,0) -- ++(0,-0.3);

    % Дуга кута 30 градусів та підпис
    \pic[draw, angle radius=0.9cm] {angle = K--D--M};
    \node at ($(D)+(170:2.5)$) {\scriptsize $30^\circ$};

    % Підписи вершин
    \node[below left] at (A) {$A$};
    \node[above left] at (B) {$B$};
    \node[above right] at (C) {$C$};
    \node[below right] at (D) {$D$};
    \node[below] at (M) {$M$};
    \node[left] at (K) {$K$};
\end{tikzpicture}
\end{minipage}
\par\vspace{0.2cm}
\answerTable{$120\sqrt{3}$~см$^2$}{$72\sqrt{3}$~см$^2$}{$96\sqrt{3}$~см$^2$}{$120$~см$^2$}{$60\sqrt{3}$~см$^2$}
\solution{У прямокутному $\triangle DKM$: $KM=DK\sin30^\circ=12\cdot\dfrac{1}{2}=6$, тож висота $BM=2KM=12$, а $DM=DK\cos30^\circ=12\cdot\dfrac{\sqrt3}{2}=6\sqrt3$. У прямокутному $\triangle ABM$: $AM=\sqrt{AB^2-BM^2}=\sqrt{192-144}=4\sqrt3$. Тоді основа $AD=AM+MD=10\sqrt3$, а площа $S=AD\cdot BM=10\sqrt3\cdot12=120\sqrt3$~см$^2$. Відповідь: \textbf{А}.}
\noindent
\begin{minipage}[c]{0.54\textwidth}
\zadnum Точка $K$~--- середина висоти $BM$ паралелограма $ABCD$ (див. рисунок). $AB = 4\sqrt{3}$~см, $DK = 6$~см, $\angle ADK = 30^\circ$. Визначте площу паралелограма $ABCD$.
\end{minipage}%
\hfill
\begin{minipage}[c]{0.44\textwidth}
\centering
\begin{tikzpicture}[scale=0.55, thick]
    \coordinate (A) at (0,0);
    \coordinate (M) at (2.5,0);
    \coordinate (D) at (7.5,0);
    \coordinate (B) at (2.5,4);
    \coordinate (C) at (10,4);
    \coordinate (K) at (2.5,2);

    \draw[thick] (A) -- (B) -- (C) -- (D) -- cycle;
    \draw[thick] (B) -- (M);
    \draw[thick] (K) -- (D);

    \fill (K) circle (2.5pt);

    \draw (M) ++(0,0.3) -- ++(0.3,0) -- ++(0,-0.3);

    \pic[draw, angle radius=0.9cm] {angle = K--D--M};
    \node at ($(D)+(170:2.5)$) {\scriptsize $30^\circ$};

    \node[below left] at (A) {$A$};
    \node[above left] at (B) {$B$};
    \node[above right] at (C) {$C$};
    \node[below right] at (D) {$D$};
    \node[below] at (M) {$M$};
    \node[left] at (K) {$K$};
\end{tikzpicture}
\end{minipage}
\par\vspace{0.2cm}
\answerTable{$24\sqrt{3}$~см$^2$}{$30\sqrt{3}$~см$^2$}{$36\sqrt{3}$~см$^2$}{$30$~см$^2$}{$15\sqrt{3}$~см$^2$}
\noindent
\begin{minipage}[c]{0.54\textwidth}
\zadnum Точка $K$~--- середина висоти $BM$ паралелограма $ABCD$ (див. рисунок). $AB = 10$~см, $DK = 4\sqrt{2}$~см, $\angle ADK = 45^\circ$. Визначте площу паралелограма $ABCD$.
\end{minipage}%
\hfill
\begin{minipage}[c]{0.44\textwidth}
\centering
\begin{tikzpicture}[scale=0.55, thick]
    \coordinate (A) at (0,0);
    \coordinate (M) at (2.5,0);
    \coordinate (D) at (7.5,0);
    \coordinate (B) at (2.5,4);
    \coordinate (C) at (10,4);
    \coordinate (K) at (2.5,2);

    \draw[thick] (A) -- (B) -- (C) -- (D) -- cycle;
    \draw[thick] (B) -- (M);
    \draw[thick] (K) -- (D);

    \fill (K) circle (2.5pt);

    \draw (M) ++(0,0.3) -- ++(0.3,0) -- ++(0,-0.3);

    \pic[draw, angle radius=0.9cm] {angle = K--D--M};
    \node at ($(D)+(170:2.5)$) {\scriptsize $45^\circ$};

    \node[below left] at (A) {$A$};
    \node[above left] at (B) {$B$};
    \node[above right] at (C) {$C$};
    \node[below right] at (D) {$D$};
    \node[below] at (M) {$M$};
    \node[left] at (K) {$K$};
\end{tikzpicture}
\end{minipage}
\par\vspace{0.2cm}
\answerTable{$56$~см$^2$}{$48$~см$^2$}{$80$~см$^2$}{$64$~см$^2$}{$40$~см$^2$}
\taskBlock{Площа кільця (формула) --- сесія 5.06, завдання №3}
\noindent
\begin{minipage}[c]{0.6\textwidth}
\zadnum Плоска сітчаста секція сушарки має вигляд частини круга, яку обмежено колами радіусів $R$ см і $3$ см, $R>3$ (див. рисунок). За якою формулою можна обчислити площу (у см$^2$) $S$ такої секції?
\end{minipage}%
\hfill
\begin{minipage}[c]{0.38\textwidth}
\centering
\begin{tikzpicture}[scale=0.5, thick]
    \filldraw[fill=gray!15, draw=cyan!60!black, thick] (0,0) circle (2);
    \filldraw[fill=white, draw=cyan!60!black, thick] (0,0) circle (0.6);
\end{tikzpicture}
\end{minipage}
\par\vspace{0.1cm}
\begin{flushleft}
\textbf{А}\ \ $S = 2\pi R - 6\pi$ \qquad \textbf{Б}\ \ $S = \pi R^2 - 6\pi$ \qquad \textbf{В}\ \ $S = \pi(R-3)^2$\\[0.15cm]
\textbf{Г}\ \ $S = \pi R^2 + 9\pi$ \qquad \textbf{Д}\ \ $S = \pi R^2 - 9\pi$
\end{flushleft}
\solution{Площа кільця $=$ площа більшого круга мінус площа меншого: $S=\pi R^2-\pi\cdot3^2=\pi R^2-9\pi$. Відповідь: \textbf{Д}.}
\vspace{0.2cm}
\noindent
\begin{minipage}[c]{0.6\textwidth}
\zadnum Металева прокладка має вигляд частини круга, яку обмежено колами радіусів $R$ см і $4$ см, $R>4$ (див. рисунок). За якою формулою можна обчислити площу (у см$^2$) $S$ такої прокладки?
\end{minipage}%
\hfill
\begin{minipage}[c]{0.38\textwidth}
\centering
\begin{tikzpicture}[scale=0.5, thick]
    \filldraw[fill=gray!15, draw=cyan!60!black, thick] (0,0) circle (2);
    \filldraw[fill=white, draw=cyan!60!black, thick] (0,0) circle (0.8);
\end{tikzpicture}
\end{minipage}
\par\vspace{0.1cm}
\begin{flushleft}
\textbf{А}\ \ $S = 2\pi R - 8\pi$ \qquad \textbf{Б}\ \ $S = \pi R^2 - 8\pi$ \qquad \textbf{В}\ \ $S = \pi(R-4)^2$\\[0.15cm]
\textbf{Г}\ \ $S = \pi R^2 - 16\pi$ \qquad \textbf{Д}\ \ $S = \pi R^2 + 16\pi$
\end{flushleft}
\vspace{0.2cm}
\noindent
\begin{minipage}[c]{0.6\textwidth}
\zadnum Гумове ущільнювальне кільце має вигляд частини круга, яку обмежено колами радіусів $R$ см і $5$ см, $R>5$ (див. рисунок). За якою формулою можна обчислити сумарну довжину (у см) $L$ \textit{обох} меж такого кільця?
\end{minipage}%
\hfill
\begin{minipage}[c]{0.38\textwidth}
\centering
\begin{tikzpicture}[scale=0.5, thick]
    \filldraw[fill=gray!15, draw=cyan!60!black, thick] (0,0) circle (2);
    \filldraw[fill=white, draw=cyan!60!black, thick] (0,0) circle (1.0);
\end{tikzpicture}
\end{minipage}
\par\vspace{0.1cm}
\begin{flushleft}
\textbf{А}\ \ $L = 2\pi R + 10\pi$ \qquad \textbf{Б}\ \ $L = 2\pi R - 10\pi$ \qquad \textbf{В}\ \ $L = 2\pi(R-5)$\\[0.15cm]
\textbf{Г}\ \ $L = \pi R^2 + 25\pi$ \qquad \textbf{Д}\ \ $L = \pi R + 5\pi$
\end{flushleft}
\vspace{0.2cm}
\taskBlock{Середня лінія трикутника (твердження) --- сесія 5.06, завдання №9}
\zadnum На сторонах $AB$ та $BC$ довільного трикутника $ABC$ вибрано точки $M$ та $N$ так, що $AM = MB$, $CN = NB$. Які з наведених тверджень є \textit{правильними}?

\vspace{0.1cm}
\begin{enumerate}[label=\textbf{\Roman*.}, align=left, leftmargin=2.6em, labelsep=0.5em, itemsep=0.08cm, topsep=0.06cm]
\item Чотирикутник $AMNC$ є трапецією.
\item Трикутник $MBN$ рівнобедрений.
\item Периметр трикутника $MBN$ дорівнює половині периметра трикутника $ABC$.
\end{enumerate}
\vspace{0.1cm}
\answerTable{лише II}{лише I та II}{лише I та III}{лише II та III}{I, II та III}
\solution{$MN$~--- середня лінія, тож $MN\parallel AC$ і $MN=\tfrac12 AC$: $AMNC$~--- трапеція (I правильне). Тричикутник $MBN$ рівнобедрений лише за умови $AB=BC$, тож для довільного трикутника твердження II хибне. Периметр $MBN=\tfrac12(AB+BC+AC)=\tfrac12 P_{ABC}$ (III правильне). Відповідь: \textbf{В}.}
\zadnum На сторонах $DE$ та $EF$ довільного трикутника $DEF$ вибрано точки $P$ та $Q$ так, що $DP = PE$, $FQ = QE$. Які з наведених тверджень є \textit{правильними}?

\vspace{0.1cm}
\begin{enumerate}[label=\textbf{\Roman*.}, align=left, leftmargin=2.6em, labelsep=0.5em, itemsep=0.08cm, topsep=0.06cm]
\item Відрізок $PQ$ паралельний стороні $DF$.
\item Довжина відрізка $PQ$ удвічі менша за довжину сторони $DF$.
\item Трикутник $PEQ$ подібний до трикутника $DEF$ з коефіцієнтом подібності $\dfrac{1}{3}$.
\end{enumerate}
\vspace{0.1cm}
\answerTable{лише I}{лише I та II}{лише I та III}{лише II та III}{I, II та III}
\zadnum Точки $M$ і $N$~--- середини сторін $AB$ та $BC$ довільного трикутника $ABC$. Які з наведених тверджень є \textit{правильними}?

\vspace{0.1cm}
\begin{enumerate}[label=\textbf{\Roman*.}, align=left, leftmargin=2.6em, labelsep=0.5em, itemsep=0.08cm, topsep=0.06cm]
\item Площа трикутника $MBN$ становить половину площі трикутника $ABC$.
\item Довжина відрізка $MN$ удвічі менша за довжину сторони $AC$.
\item Відрізок $MN$ паралельний стороні $AC$.
\end{enumerate}
\vspace{0.1cm}
\answerTable{лише I}{лише I та II}{лише I та III}{лише II та III}{I, II та III}
\taskBlock{Площа прямокутного трикутника у прямокутнику --- сесія 5.06, завдання №14}
\noindent
\begin{minipage}[c]{0.56\textwidth}
\zadnum У прямокутник $ABCD$ вписано прямокутний трикутник $AKM$ так, як зображено на рисунку. Обчисліть площу трикутника $AKM$, якщо $AK = 12$ см, $AB = 15$ см, $\angle KAD = 30^\circ$.
\end{minipage}%
\hfill
\begin{minipage}[c]{0.40\textwidth}
\centering
\begin{tikzpicture}[scale=0.32, thick]
    \coordinate (A) at (0,0);
    \coordinate (B) at (0,15);
    \coordinate (C) at (11,15);
    \coordinate (D) at (11,0);
    \coordinate (K) at (11,6);
    \coordinate (M) at (5.2,15);
    \draw[thick] (A) -- (B) -- (C) -- (D) -- cycle;
    \draw[thick] (A) -- (K) -- (M) -- cycle;
    \pic[draw=yearOrange, line width=0.6pt, angle radius=0.7cm] {right angle = A--K--M};

    \draw[thick] (K) -- (A) -- (D) -- cycle;
    \pic[draw=yearOrange, line width=0.6pt, angle radius=0.3cm] { angle = D--A--K};
    \node[below left, font=\scriptsize] at (A) {$A$};
    \node[above left, font=\scriptsize] at (B) {$B$};
    \node[above, font=\scriptsize] at (M) {$M$};
    \node[above right, font=\scriptsize] at (C) {$C$};
    \node[right, font=\scriptsize] at (K) {$K$};
    \node[below right, font=\scriptsize] at (D) {$D$};
    \node[font=\tiny] at (1.9,0.5) {$30^\circ$};
\end{tikzpicture}
\end{minipage}
\par\vspace{0.2cm}
\answerTable{$36$ см$^2$}{$96$ см$^2$}{$72$ см$^2$}{$72\sqrt{3}$ см$^2$}{$36\sqrt{3}$ см$^2$}
\solution{Введемо координати $A(0;0)$, $AB$~--- вертикальна сторона. Оскільки $\angle KAD=30^\circ$ і $AK=12$, то $K(6\sqrt3;\,6)$. Точка $M$ лежить на верхній стороні (на висоті $15$); з умови прямого кута при $K$ ($\overrightarrow{KA}\cdot\overrightarrow{KM}=0$) дістаємо $KM=6\sqrt3$ см. Площа $=\tfrac12\cdot AK\cdot KM=\tfrac12\cdot12\cdot6\sqrt3=36\sqrt3$ см$^2$. Відповідь: \textbf{Д}.}
\noindent
\begin{minipage}[c]{0.56\textwidth}
\zadnum У прямокутник $ABCD$ вписано прямокутний трикутник $AKM$ так, як зображено на рисунку. Обчисліть площу трикутника $AKM$, якщо $AK = 8$ см, $AB = 10$ см, $\angle KAD = 30^\circ$.
\end{minipage}%
\hfill
\begin{minipage}[c]{0.40\textwidth}
\centering
\begin{tikzpicture}[scale=0.45, thick]
    \coordinate (A) at (0,0);
    \coordinate (B) at (0,10);
    \coordinate (C) at (7.4,10);
    \coordinate (D) at (7.4,0);
    \coordinate (K) at (7.4,4);
    \coordinate (M) at (3.6,10);
    \draw[thick] (A) -- (B) -- (C) -- (D) -- cycle;
    \draw[thick] (A) -- (K) -- (M) -- cycle;
    \pic[draw=yearOrange, line width=0.6pt, angle radius=0.7cm] {right angle = A--K--M};
    \pic[draw=yearOrange, line width=0.6pt, angle radius=0.3cm] { angle = D--A--K};
    \node[below left, font=\scriptsize] at (A) {$A$};
    \node[above left, font=\scriptsize] at (B) {$B$};
    \node[above, font=\scriptsize] at (M) {$M$};
    \node[above right, font=\scriptsize] at (C) {$C$};
    \node[right, font=\scriptsize] at (K) {$K$};
    \node[below right, font=\scriptsize] at (D) {$D$};
    \node[font=\tiny] at (1.5,0.4) {$30^\circ$};
\end{tikzpicture}
\end{minipage}
\par\vspace{0.2cm}
\answerTable{$16$ см$^2$}{$32$ см$^2$}{$8\sqrt{3}$ см$^2$}{$32\sqrt{3}$ см$^2$}{$16\sqrt{3}$ см$^2$}
\noindent
\begin{minipage}[c]{0.56\textwidth}
\zadnum У прямокутник $ABCD$ вписано прямокутний трикутник $AKM$ так, як зображено на рисунку. Обчисліть довжину катета $KM$ (у см), якщо $AK = 8\sqrt{2}$ см, $AB = 14$ см, $\angle KAD = 45^\circ$.
\end{minipage}%
\hfill
\begin{minipage}[c]{0.40\textwidth}
\centering
\begin{tikzpicture}[scale=0.32, thick]
    \coordinate (A) at (0,0);
    \coordinate (B) at (0,14);
    \coordinate (C) at (10,14);
    \coordinate (D) at (10,0);
    \coordinate (K) at (10,8);
    \coordinate (M) at (2,14);
    \draw[thick] (A) -- (B) -- (C) -- (D) -- cycle;
    \draw[thick] (A) -- (K) -- (M) -- cycle;
    \pic[draw=yearOrange, line width=0.6pt, angle radius=0.7cm] {right angle = A--K--M};
    \pic[draw=yearOrange, line width=0.6pt, angle radius=0.3cm] { angle = D--A--K};
    \node[below left, font=\scriptsize] at (A) {$A$};
    \node[above left, font=\scriptsize] at (B) {$B$};
    \node[above, font=\scriptsize] at (M) {$M$};
    \node[above right, font=\scriptsize] at (C) {$C$};
    \node[right, font=\scriptsize] at (K) {$K$};
    \node[below right, font=\scriptsize] at (D) {$D$};
    \node[font=\tiny] at (2.0,0.55) {$45^\circ$};
\end{tikzpicture}
\end{minipage}
\par\vspace{0.2cm}
\answerTable{$6\sqrt{2}$ см}{$6$ см}{$12$ см}{$3\sqrt{2}$ см}{$12\sqrt{2}$ см}
\typeTitle{Завдання на встановлення відповідності}
\taskBlock{Трапеція з висотою-бісектрисою (відповідність) --- сесія 23.05, завдання №16}

\zadtask{На рисунку зображено трапецію $ABCD$, у якій $\angle CDA = 30^\circ$. Висота $CK$, довжина якої дорівнює $2$, є бісектрисою кута $ACD$. \nmtyear{2026}

\vspace{0.3cm}

\begin{center}
\begin{tikzpicture}[scale=0.9, thick]
    \coordinate (A) at (0,0);
    \coordinate (K) at (4,0);
    \coordinate (D) at (7.46,0);
    \coordinate (C) at (4,2);
    \coordinate (B) at (0,2);
    \draw (A) -- (B) -- (C) -- (D) -- cycle;
    \draw (C) -- (K);
    \draw[dashed] (A) -- (C);
    \draw (0,0) -- (0,0.3) -- (0.3,0.3) -- (0.3,0);
    \draw (4,0) -- (4,0.3) -- (4.3,0.3) -- (4.3,0);
    \pic [draw, angle radius=0.6cm, "$30^\circ$"{above left, shift={(-0.28,-0.15)}}] {angle = C--D--K};
    \pic [draw, angle radius=0.45cm] {angle = K--C--D};
    \pic [draw, angle radius=0.55cm] {angle = A--C--K};
    \node[below left] at (A) {$A$};
    \node[above left] at (B) {$B$};
    \node[above=2pt] at (C) {$C$};
    \node[below right] at (D) {$D$};
    \node[below=2pt] at (K) {$K$};
\end{tikzpicture}
\end{center}

\vspace{0.3cm}

\noindent Установіть відповідність між величиною (1--3) та її значенням (А--Д).

\vspace{0.3cm}

\noindent
\begin{minipage}[t]{0.52\textwidth}
\textit{Величина} \par \vspace{0.2cm}
\textbf{1} \ \ довжина відрізка $CD$\\[0.25cm]
\textbf{2} \ \ довжина відрізка $BC$\\[0.25cm]
\textbf{3} \ \ відстань від точки $B$ до прямої $AC$
\end{minipage}
\begin{minipage}[t]{0.28\textwidth}
\textit{Значення} \par \vspace{0.2cm}
\textbf{А} \ \ $2\sqrt{3}$\\[0.15cm]
\textbf{Б} \ \ $4$\\[0.15cm]
\textbf{В} \ \ $1$\\[0.15cm]
\textbf{Г} \ \ $\sqrt{3}$\\[0.15cm]
\textbf{Д} \ \ $6\sqrt{3}$
\end{minipage}
\hfill
\begin{minipage}[t]{0.22\textwidth}
\vspace{-0.2cm}
\begin{flushright}\matchingGrid\end{flushright}
\end{minipage}
}
\solution{У прямокутному $\triangle CKD$: $\angle CDK=30^\circ$, катет $CK=2$ лежить навпроти $30^\circ$, тож $CD=\dfrac{CK}{\sin 30^\circ}=4$~(Б). $\angle KCD=60^\circ$; оскільки $CK$~--- бісектриса, $\angle ACK=60^\circ$, тоді у прямокутному $\triangle ACK$: $AK=CK\tan 60^\circ=2\sqrt{3}$, а $BC=AK=2\sqrt{3}$~(А). Відстань від $B$ до $AC$ дорівнює $\dfrac{2S_{ABC}}{AC}$; $AC=\dfrac{CK}{\cos 60^\circ}=4$, $S_{ABC}=\dfrac{1}{2}\cdot BC\cdot AB=\dfrac{1}{2}\cdot 2\sqrt{3}\cdot 2=2\sqrt{3}$, тож відстань $=\dfrac{2\cdot 2\sqrt{3}}{4}=\sqrt{3}$~(Г). Відповідь: \textbf{1--Б, 2--А, 3--Г}.}

\vspace{0.5cm}

\zadtask{На рисунку зображено трапецію $ABCD$, у якій $\angle CDA = 45^\circ$. Висота $CK$, довжина якої дорівнює $3$, є бісектрисою кута $ACD$.

\vspace{0.3cm}

\begin{center}
\begin{tikzpicture}[scale=0.9, thick]
    \coordinate (A) at (0,0);
    \coordinate (K) at (3,0);
    \coordinate (D) at (6,0);
    \coordinate (C) at (3,3);
    \coordinate (B) at (0,3);
    \draw (A) -- (B) -- (C) -- (D) -- cycle;
    \draw (C) -- (K);
    \draw[dashed] (A) -- (C);
    \draw (0,0) -- (0,0.3) -- (0.3,0.3) -- (0.3,0);
    \draw (3,0) -- (3,0.3) -- (3.3,0.3) -- (3.3,0);
    \pic [draw, angle radius=0.6cm, "$45^\circ$"{above left, shift={(-0.28,-0.15)}}] {angle = C--D--K};
    \pic [draw, angle radius=0.45cm] {angle = K--C--D};
    \pic [draw, angle radius=0.55cm] {angle = A--C--K};
    \node[below left] at (A) {$A$};
    \node[above left] at (B) {$B$};
    \node[above=2pt] at (C) {$C$};
    \node[below right] at (D) {$D$};
    \node[below=2pt] at (K) {$K$};
\end{tikzpicture}
\end{center}

\vspace{0.3cm}

\noindent Установіть відповідність між величиною (1--3) та її значенням (А--Д).

\vspace{0.3cm}

\noindent
\begin{minipage}[t]{0.52\textwidth}
\textit{Величина} \par \vspace{0.2cm}
\textbf{1} \ \ довжина відрізка $CD$\\[0.25cm]
\textbf{2} \ \ довжина відрізка $BC$\\[0.25cm]
\textbf{3} \ \ відстань від точки $B$ до прямої $AC$
\end{minipage}
\begin{minipage}[t]{0.28\textwidth}
\textit{Значення} \par \vspace{0.2cm}
\textbf{А} \ \ $3$\\[0.15cm]
\textbf{Б} \ \ $\dfrac{3\sqrt{2}}{2}$\\[0.15cm]
\textbf{В} \ \ $6$\\[0.15cm]
\textbf{Г} \ \ $3\sqrt{2}$\\[0.15cm]
\textbf{Д} \ \ $2\sqrt{3}$
\end{minipage}
\hfill
\begin{minipage}[t]{0.22\textwidth}
\vspace{-0.2cm}
\begin{flushright}\matchingGrid\end{flushright}
\end{minipage}
}

\vspace{0.5cm}

\zadtask{На рисунку зображено трапецію $ABCD$, у якій $\angle CDA = 60^\circ$. Висота $CK$, довжина якої дорівнює $\sqrt{3}$, є бісектрисою кута $ACD$.

\vspace{0.3cm}

\begin{center}
\begin{tikzpicture}[scale=1.1, thick]
    \coordinate (A) at (0,0);
    \coordinate (K) at (1.5,0);
    \coordinate (D) at (3.0,0);
    \coordinate (C) at (1.5,1.73);
    \coordinate (B) at (0,1.73);
    \draw (A) -- (B) -- (C) -- (D) -- cycle;
    \draw (C) -- (K);
    \draw[dashed] (A) -- (C);
    \draw (0,0) -- (0,0.25) -- (0.25,0.25) -- (0.25,0);
    \draw (1.5,0) -- (1.5,0.25) -- (1.75,0.25) -- (1.75,0);
    \pic [draw, angle radius=0.5cm, "$60^\circ$"{above left, shift={(-0.28,-0.1)}}] {angle = C--D--K};
    \pic [draw, angle radius=0.4cm] {angle = K--C--D};
    \pic [draw, angle radius=0.5cm] {angle = A--C--K};
    \node[below left] at (A) {$A$};
    \node[above left] at (B) {$B$};
    \node[above=2pt] at (C) {$C$};
    \node[below right] at (D) {$D$};
    \node[below=2pt] at (K) {$K$};
\end{tikzpicture}
\end{center}

\vspace{0.3cm}

\noindent Установіть відповідність між величиною (1--3) та її значенням (А--Д).

\vspace{0.3cm}

\noindent
\begin{minipage}[t]{0.52\textwidth}
\textit{Величина} \par \vspace{0.2cm}
\textbf{1} \ \ площа трикутника $ACD$\\[0.25cm]
\textbf{2} \ \ довжина діагоналі $AC$\\[0.25cm]
\textbf{3} \ \ площа трикутника $CKD$
\end{minipage}
\begin{minipage}[t]{0.28\textwidth}
\textit{Значення} \par \vspace{0.2cm}
\textbf{А} \ \ $1$\\[0.15cm]
\textbf{Б} \ \ $\sqrt{3}$\\[0.15cm]
\textbf{В} \ \ $2$\\[0.15cm]
\textbf{Г} \ \ $\dfrac{\sqrt{3}}{2}$\\[0.15cm]
\textbf{Д} \ \ $2\sqrt{3}$
\end{minipage}
\hfill
\begin{minipage}[t]{0.22\textwidth}
\vspace{-0.2cm}
\begin{flushright}\matchingGrid\end{flushright}
\end{minipage}
}
\taskBlock{Планіметрія: квадрат і трикутник (відповідність) --- сесія 30.05, завдання №16}

\zadtask{На рисунку зображено квадрат $ABCD$ та рівнобедрений прямокутний трикутник $KBC$ (прямий кут при вершині $K$), які мають спільну сторону $BC$. Точка $M$~--- середина сторони $CD$. Довжина ламаної $BADM$ дорівнює $20$. Установіть відповідність між величинами (1--3) та їхніми значеннями (А--Д).

\vspace{0.3cm}
\begin{center}
\begin{tikzpicture}[scale=0.7, thick]
    \coordinate (A) at (0,0);
    \coordinate (B) at (0,4);
    \coordinate (C) at (4,4);
    \coordinate (D) at (4,0);
    \coordinate (K) at (2,6);
    \coordinate (Mid) at (4,2);
    \draw[thick] (A) -- (B) -- (C) -- (D) -- cycle;
    \draw[thick] (B) -- (K) -- (C);
    \fill (Mid) circle (2pt);
    \node[below left] at (A) {$A$};
    \node[above left] at (B) {$B$};
    \node[above right] at (C) {$C$};
    \node[below right] at (D) {$D$};
    \node[above] at (K) {$K$};
    \node[right] at (Mid) {$M$};
\end{tikzpicture}
\end{center}

\vspace{0.2cm}
\noindent
\begin{minipage}[t]{0.45\textwidth}
\textit{Величина}\par\vspace{0.15cm}
\textbf{1} \quad довжина сторони квадрата\\[0.25cm]
\textbf{2} \quad відстань від точки $K$ до прямої $AD$\\[0.25cm]
\textbf{3} \quad довжина відрізка $KM$
\end{minipage}
\hfill
\begin{minipage}[t]{0.3\textwidth}
\textit{Значення}\par\vspace{0.15cm}
\textbf{А} \quad $12$\\[0.15cm]
\textbf{Б} \quad $10$\\[0.15cm]
\textbf{В} \quad $4\sqrt{3}$\\[0.15cm]
\textbf{Г} \quad $4\sqrt{5}$\\[0.15cm]
\textbf{Д} \quad $8$
\end{minipage}
\hfill
\begin{minipage}[t]{0.2\textwidth}
\vspace{0.4cm}
\begin{flushright}\matchingGrid\end{flushright}
\end{minipage}}
\solution{Нехай сторона квадрата $a$. Ламана $BADM=BA+AD+DM=a+a+\dfrac a2=\dfrac{5a}{2}=20$, звідки $a=8$ (\textbf{1}~--~Д). Висота трикутника до $BC$ дорівнює $\dfrac a2=4$, тож відстань від $K$ до $AD$: $a+\dfrac a2=12$ (\textbf{2}~--~А). У координатах $K(4;12)$, $M(8;4)$: $KM=\sqrt{4^2+8^2}=\sqrt{80}=4\sqrt5$ (\textbf{3}~--~Г). Відповідь: \textbf{1~--~Д, 2~--~А, 3~--~Г}.}

\vspace{0.7cm}

\zadtask{На рисунку зображено квадрат $ABCD$ та рівнобедрений прямокутний трикутник $KBC$ (прямий кут при вершині $K$), які мають спільну сторону $BC$. Точка $M$~--- середина сторони $CD$. Довжина ламаної $BADM$ дорівнює $30$. Установіть відповідність між величинами (1--3) та їхніми значеннями (А--Д).

\vspace{0.3cm}
\begin{center}
\begin{tikzpicture}[scale=0.7, thick]
    \coordinate (A) at (0,0);
    \coordinate (B) at (0,4);
    \coordinate (C) at (4,4);
    \coordinate (D) at (4,0);
    \coordinate (K) at (2,6);
    \coordinate (Mid) at (4,2);
    \draw[thick] (A) -- (B) -- (C) -- (D) -- cycle;
    \draw[thick] (B) -- (K) -- (C);
    \fill (Mid) circle (2pt);
    \node[below left] at (A) {$A$};
    \node[above left] at (B) {$B$};
    \node[above right] at (C) {$C$};
    \node[below right] at (D) {$D$};
    \node[above] at (K) {$K$};
    \node[right] at (Mid) {$M$};
\end{tikzpicture}
\end{center}

\vspace{0.2cm}
\noindent
\begin{minipage}[t]{0.45\textwidth}
\textit{Величина}\par\vspace{0.15cm}
\textbf{1} \quad довжина сторони квадрата\\[0.25cm]
\textbf{2} \quad відстань від точки $K$ до прямої $AD$\\[0.25cm]
\textbf{3} \quad довжина відрізка $KM$
\end{minipage}
\hfill
\begin{minipage}[t]{0.3\textwidth}
\textit{Значення}\par\vspace{0.15cm}
\textbf{А} \quad $18$\\[0.15cm]
\textbf{Б} \quad $15$\\[0.15cm]
\textbf{В} \quad $6\sqrt{3}$\\[0.15cm]
\textbf{Г} \quad $6\sqrt{5}$\\[0.15cm]
\textbf{Д} \quad $12$
\end{minipage}
\hfill
\begin{minipage}[t]{0.2\textwidth}
\vspace{0.4cm}
\begin{flushright}\matchingGrid\end{flushright}
\end{minipage}}

\vspace{0.7cm}

\zadtask{На рисунку зображено квадрат $ABCD$ та рівнобедрений прямокутний трикутник $KBC$ (прямий кут при вершині $K$), які мають спільну сторону $BC$. Точка $M$~--- середина сторони $CD$. Довжина ламаної $BADM$ дорівнює $10$. Установіть відповідність між величинами (1--3) та їхніми значеннями (А--Д).

\vspace{0.05cm}

\vspace{0.3cm}
\begin{center}
\begin{tikzpicture}[scale=0.7, thick]
    \coordinate (A) at (0,0);
    \coordinate (B) at (0,4);
    \coordinate (C) at (4,4);
    \coordinate (D) at (4,0);
    \coordinate (K) at (2,6);
    \coordinate (Mid) at (4,2);
    \draw[thick] (A) -- (B) -- (C) -- (D) -- cycle;
    \draw[thick] (B) -- (K) -- (C);
    \fill (Mid) circle (2pt);
    \node[below left] at (A) {$A$};
    \node[above left] at (B) {$B$};
    \node[above right] at (C) {$C$};
    \node[below right] at (D) {$D$};
    \node[above] at (K) {$K$};
    \node[right] at (Mid) {$M$};
\end{tikzpicture}
\end{center}

\vspace{0.2cm}
\noindent
\begin{minipage}[t]{0.45\textwidth}
\textit{Величина}\par\vspace{0.15cm}
\textbf{1} \quad периметр квадрата\\[0.25cm]
\textbf{2} \quad площа трикутника $KBC$\\[0.25cm]
\textbf{3} \quad довжина діагоналі $BD$ квадрата
\end{minipage}
\hfill
\begin{minipage}[t]{0.3\textwidth}
\textit{Значення}\par\vspace{0.15cm}
\textbf{А} \quad $16$\\[0.15cm]
\textbf{Б} \quad $8$\\[0.15cm]
\textbf{В} \quad $4\sqrt{2}$\\[0.15cm]
\textbf{Г} \quad $4$\\[0.15cm]
\textbf{Д} \quad $2\sqrt{5}$
\end{minipage}
\hfill
\begin{minipage}[t]{0.2\textwidth}
\vspace{0.4cm}
\begin{flushright}\matchingGrid\end{flushright}
\end{minipage}}

\vspace{0.7cm}
\taskBlock{Паралелограм, висота й подібність трикутників (відповідність) --- сесія 1.06, завдання №18}

\zadtask{У паралелограмі $ABCD$ на висоті $BM$ вибрано точку $K$ так, що $BK : KM = 4 : 3$ (див. рисунок). $CK = 17$ см, $BK = 8$ см. Прямі $CK$ і $AD$ перетинаються в точці $N$. Установіть відповідність між відрізком (1--3) та його довжиною (А--Д).

\vspace{0.3cm}
\begin{center}
\begin{tikzpicture}[scale=0.8, thick]
    \coordinate (A) at (0,0);
    \coordinate (M) at (1.3,0);
    \coordinate (D) at (5,0);
    \coordinate (B) at (1.3,2.6);
    \coordinate (C) at (5,2.6);
    \coordinate (N) at (-1.1,0);
    \coordinate (K) at (1.3,1.1);
    \draw[thick] (A) -- (B) -- (C) -- (D) -- cycle;
    \draw[thick] (B) -- (M);
    \draw[thick] (N) -- (C);
    \draw[thick] (N) -- (A);
    \draw[thick] (M)++(0,0.28) -- ++(0.28,0) -- ++(0,-0.28);
    \fill (K) circle (1.6pt);
    \fill (N) circle (1.6pt);
    \node[below] at (A) {$A$};
    \node[below] at (M) {$M$};
    \node[below right] at (D) {$D$};
    \node[above left] at (B) {$B$};
    \node[above right] at (C) {$C$};
    \node[below] at (N) {$N$};
    \node[below right] at (K) {$K$};
\end{tikzpicture}
\end{center}

\vspace{0.2cm}
\noindent
\begin{minipage}[t]{0.42\textwidth}
\textit{Відрізок}\par\vspace{0.2cm}
\textbf{1} \quad $BC$\\[0.25cm]
\textbf{2} \quad $BM$\\[0.25cm]
\textbf{3} \quad $NM$
\end{minipage}
\hfill
\begin{minipage}[t]{0.3\textwidth}
\textit{Довжина відрізка}\par\vspace{0.2cm}
\textbf{А} \quad $11{,}25$ см\\[0.15cm]
\textbf{Б} \quad $12$ см\\[0.15cm]
\textbf{В} \quad $14$ см\\[0.15cm]
\textbf{Г} \quad $15$ см\\[0.15cm]
\textbf{Д} \quad $20$ см
\end{minipage}
\hfill
\begin{minipage}[t]{0.2\textwidth}
\vspace{0.4cm}
\begin{flushright}\matchingGrid\end{flushright}
\end{minipage}}
\solution{Оскільки $BM\perp AD$ і $BC\parallel AD$, то $BM\perp BC$, тому в прямокутному трикутнику $BKC$: $BC=\sqrt{CK^2-BK^2}=\sqrt{17^2-8^2}=15$ см (1\,--\,Г). Далі $KM=\dfrac{3}{4}BK=6$, тож $BM=BK+KM=14$ см (2\,--\,В). Трикутники $NMK$ і $CBK$ подібні: $\dfrac{NM}{BC}=\dfrac{KM}{BK}=\dfrac{6}{8}$, звідки $NM=15\cdot\dfrac{6}{8}=11{,}25$ см (3\,--\,А). Відповідь: \textbf{1\,--\,Г, 2\,--\,В, 3\,--\,А}.}

\vspace{0.5cm}

\zadtask{У паралелограмі $ABCD$ на висоті $BM$ вибрано точку $K$ так, що $BK : KM = 3 : 2$ (див. рисунок). $CK = 15$ см, $BK = 9$ см. Прямі $CK$ і $AD$ перетинаються в точці $N$. Установіть відповідність між відрізком (1--3) та його довжиною (А--Д).

\vspace{0.3cm}
\begin{center}
\begin{tikzpicture}[scale=0.8, thick]
    \coordinate (A) at (0,0);
    \coordinate (M) at (1.3,0);
    \coordinate (D) at (5,0);
    \coordinate (B) at (1.3,2.6);
    \coordinate (C) at (5,2.6);
    \coordinate (N) at (-1.2,0);
    \coordinate (K) at (1.3,1.05);
    \draw[thick] (A) -- (B) -- (C) -- (D) -- cycle;
    \draw[thick] (B) -- (M);
    \draw[thick] (N) -- (C);
    \draw[thick] (N) -- (A);
    \draw[thick] (M)++(0,0.28) -- ++(0.28,0) -- ++(0,-0.28);
    \fill (K) circle (1.6pt);
    \fill (N) circle (1.6pt);
    \node[below] at (A) {$A$};
    \node[below] at (M) {$M$};
    \node[below right] at (D) {$D$};
    \node[above left] at (B) {$B$};
    \node[above right] at (C) {$C$};
    \node[below] at (N) {$N$};
    \node[below right] at (K) {$K$};
\end{tikzpicture}
\end{center}

\vspace{0.2cm}
\noindent
\begin{minipage}[t]{0.42\textwidth}
\textit{Відрізок}\par\vspace{0.2cm}
\textbf{1} \quad $BC$\\[0.25cm]
\textbf{2} \quad $BM$\\[0.25cm]
\textbf{3} \quad $NM$
\end{minipage}
\hfill
\begin{minipage}[t]{0.3\textwidth}
\textit{Довжина відрізка}\par\vspace{0.2cm}
\textbf{А} \quad $8$ см\\[0.15cm]
\textbf{Б} \quad $12$ см\\[0.15cm]
\textbf{В} \quad $15$ см\\[0.15cm]
\textbf{Г} \quad $18$ см\\[0.15cm]
\textbf{Д} \quad $20$ см
\end{minipage}
\hfill
\begin{minipage}[t]{0.2\textwidth}
\vspace{0.4cm}
\begin{flushright}\matchingGrid\end{flushright}
\end{minipage}}

\vspace{0.5cm}

\zadtask{У паралелограмі $ABCD$ на висоті $BM$ вибрано точку $K$ так, що $BK : KM = 5 : 4$ (див. рисунок). $CK = 26$ см, $BK = 10$ см. Прямі $CK$ і $AD$ перетинаються в точці $N$. Установіть відповідність між відрізком (1--3) та його довжиною (А--Д).

\vspace{0.3cm}
\begin{center}
\begin{tikzpicture}[scale=0.8, thick]
    \coordinate (A) at (0,0);
    \coordinate (M) at (1.3,0);
    \coordinate (D) at (5,0);
    \coordinate (B) at (1.3,2.6);
    \coordinate (C) at (5,2.6);
    \coordinate (N) at (-1.3,0);
    \coordinate (K) at (1.3,1.15);
    \draw[thick] (A) -- (B) -- (C) -- (D) -- cycle;
    \draw[thick] (B) -- (M);
    \draw[thick] (N) -- (C);
    \draw[thick] (N) -- (A);
    \draw[thick] (M)++(0,0.28) -- ++(0.28,0) -- ++(0,-0.28);
    \fill (K) circle (1.6pt);
    \fill (N) circle (1.6pt);
    \node[below] at (A) {$A$};
    \node[below] at (M) {$M$};
    \node[below right] at (D) {$D$};
    \node[above left] at (B) {$B$};
    \node[above right] at (C) {$C$};
    \node[below] at (N) {$N$};
    \node[below right] at (K) {$K$};
\end{tikzpicture}
\end{center}

\vspace{0.2cm}
\noindent
\begin{minipage}[t]{0.42\textwidth}
\textit{Відрізок}\par\vspace{0.2cm}
\textbf{1} \quad $KM$\\[0.25cm]
\textbf{2} \quad $NM$\\[0.25cm]
\textbf{3} \quad $NK$
\end{minipage}
\hfill
\begin{minipage}[t]{0.3\textwidth}
\textit{Довжина відрізка}\par\vspace{0.2cm}
\textbf{А} \quad $8$ см\\[0.15cm]
\textbf{Б} \quad $19{,}2$ см\\[0.15cm]
\textbf{В} \quad $20{,}8$ см\\[0.15cm]
\textbf{Г} \quad $24$ см\\[0.15cm]
\textbf{Д} \quad $26$ см
\end{minipage}
\hfill
\begin{minipage}[t]{0.2\textwidth}
\vspace{0.4cm}
\begin{flushright}\matchingGrid\end{flushright}
\end{minipage}}
\taskBlock{Прямокутник і коло: відповідність відрізків --- сесія 2.06, завдання №18}

\zadtask{Прямокутник $ABCD$ і коло з центром $O$ розташовані так, що коло дотикається до середин бічних сторін $AB$ і $CD$ та перетинає сторону $BC$ у точках $K$ і $M$ (див. рисунок). Відомо, що $BK = 3$ см, $BM = 27$ см. Установіть відповідність між відрізком (1--3) та його довжиною (А--Д).

\vspace{0.3cm}
\begin{center}
\begin{tikzpicture}[scale=0.32, thick]
    \coordinate (A) at (0,0);
    \coordinate (B) at (0,18);
    \coordinate (C) at (30,18);
    \coordinate (D) at (30,0);
    \coordinate (O) at (15,9);
    \coordinate (K) at (3,18);
    \coordinate (M) at (27,18);
    \draw[thick] (A) -- (B) -- (C) -- (D) -- cycle;
    \draw[thick, mainGreen] (O) circle (15);
    \fill (O) circle (4pt); \node[below right] at (O) {$O$};
    \fill (K) circle (4pt); \node[above] at (K) {$K$};
    \fill (M) circle (4pt); \node[above] at (M) {$M$};
    \node[below left] at (A) {$A$};
    \node[above left] at (B) {$B$};
    \node[above right] at (C) {$C$};
    \node[below right] at (D) {$D$};
\end{tikzpicture}
\end{center}

\vspace{0.2cm}
\noindent
\begin{minipage}[t]{0.42\textwidth}
\textit{Відрізок}\par\vspace{0.2cm}
\textbf{1} \quad $KM$\\[0.25cm]
\textbf{2} \quad $KO$\\[0.25cm]
\textbf{3} \quad $AB$
\end{minipage}
\hfill
\begin{minipage}[t]{0.3\textwidth}
\textit{Довжина}\par\vspace{0.2cm}
\textbf{А} \quad $3$ см\\[0.15cm]
\textbf{Б} \quad $15$ см\\[0.15cm]
\textbf{В} \quad $18$ см\\[0.15cm]
\textbf{Г} \quad $24$ см\\[0.15cm]
\textbf{Д} \quad $30$ см
\end{minipage}
\hfill
\begin{minipage}[t]{0.2\textwidth}
\vspace{0.4cm}
\begin{flushright}\matchingGrid\end{flushright}
\end{minipage}}
\solution{$KM=BM-BK=27-3=24$ см -- \textbf{Г}. Коло дотикається середин $AB$ і $CD$, тож його радіус дорівнює половині сторони $BC$; за симетрією $BC=BK+KM+BK=30$, отже $r=KO=15$ см -- \textbf{Б}. Відстань від $O$ до $BC$ дорівнює $\sqrt{r^2-(KM/2)^2}=\sqrt{225-144}=9$, тому $AB=2\cdot9=18$ см -- \textbf{В}. Відповідь: \textbf{1--Г, 2--Б, 3--В}.}

\zadtask{Прямокутник $ABCD$ і коло з центром $O$ розташовані так, що коло дотикається до середин бічних сторін $AB$ і $CD$ та перетинає сторону $BC$ у точках $K$ і $M$ (див. рисунок). Відомо, що $BK = 2$ см, $BM = 50$ см. Установіть відповідність між відрізком (1--3) та його довжиною (А--Д).

\vspace{0.3cm}
\begin{center}
\begin{tikzpicture}[scale=0.19, thick]
    \coordinate (A) at (0,0);
    \coordinate (B) at (0,20);
    \coordinate (C) at (52,20);
    \coordinate (D) at (52,0);
    \coordinate (O) at (26,10);
    \coordinate (K) at (2,20);
    \coordinate (M) at (50,20);
    \draw[thick] (A) -- (B) -- (C) -- (D) -- cycle;
    \draw[thick, mainGreen] (O) circle (26);
    \fill (O) circle (6pt); \node[below right] at (O) {$O$};
    \fill (K) circle (6pt); \node[above] at (K) {$K$};
    \fill (M) circle (6pt); \node[above] at (M) {$M$};
    \node[below left] at (A) {$A$};
    \node[above left] at (B) {$B$};
    \node[above right] at (C) {$C$};
    \node[below right] at (D) {$D$};
\end{tikzpicture}
\end{center}

\vspace{0.2cm}
\noindent
\begin{minipage}[t]{0.42\textwidth}
\textit{Відрізок}\par\vspace{0.2cm}
\textbf{1} \quad $KM$\\[0.25cm]
\textbf{2} \quad $KO$\\[0.25cm]
\textbf{3} \quad $AB$
\end{minipage}
\hfill
\begin{minipage}[t]{0.3\textwidth}
\textit{Довжина}\par\vspace{0.2cm}
\textbf{А} \quad $20$ см\\[0.15cm]
\textbf{Б} \quad $2$ см\\[0.15cm]
\textbf{В} \quad $48$ см\\[0.15cm]
\textbf{Г} \quad $26$ см\\[0.15cm]
\textbf{Д} \quad $52$ см
\end{minipage}
\hfill
\begin{minipage}[t]{0.2\textwidth}
\vspace{0.4cm}
\begin{flushright}\matchingGrid\end{flushright}
\end{minipage}}

\zadtask{Прямокутник $ABCD$ і коло з центром $O$ розташовані так, що коло дотикається до середин бічних сторін $AB$ і $CD$ та перетинає сторону $BC$ у точках $K$ і $M$ (див. рисунок). Відомо, що $BK = 2$ см, $BM = 18$ см. Установіть відповідність між величиною (1--3) та її значенням (А--Д).

\vspace{0.3cm}
\begin{center}
\begin{tikzpicture}[scale=0.45, thick]
    \coordinate (A) at (0,0);
    \coordinate (B) at (0,12);
    \coordinate (C) at (20,12);
    \coordinate (D) at (20,0);
    \coordinate (O) at (10,6);
    \coordinate (K) at (2,12);
    \coordinate (M) at (18,12);
    \draw[thick] (A) -- (B) -- (C) -- (D) -- cycle;
    \draw[thick, mainGreen] (O) circle (10);
    \fill (O) circle (3pt); \node[below right] at (O) {$O$};
    \fill (K) circle (3pt); \node[above] at (K) {$K$};
    \fill (M) circle (3pt); \node[above] at (M) {$M$};
    \node[below left] at (A) {$A$};
    \node[above left] at (B) {$B$};
    \node[above right] at (C) {$C$};
    \node[below right] at (D) {$D$};
\end{tikzpicture}
\end{center}

\vspace{0.2cm}
\noindent
\begin{minipage}[t]{0.42\textwidth}
\textit{Величина}\par\vspace{0.2cm}
\textbf{1} \quad довжина $KM$\\[0.25cm]
\textbf{2} \quad довжина $BC$\\[0.25cm]
\textbf{3} \quad відстань від $O$ до $BC$
\end{minipage}
\hfill
\begin{minipage}[t]{0.3\textwidth}
\textit{Значення}\par\vspace{0.2cm}
\textbf{А} \quad $6$ см\\[0.15cm]
\textbf{Б} \quad $10$ см\\[0.15cm]
\textbf{В} \quad $12$ см\\[0.15cm]
\textbf{Г} \quad $16$ см\\[0.15cm]
\textbf{Д} \quad $20$ см
\end{minipage}
\hfill
\begin{minipage}[t]{0.2\textwidth}
\vspace{0.4cm}
\begin{flushright}\matchingGrid\end{flushright}
\end{minipage}}
\taskBlock{Планіметрія: прямокутник і рівнобедрений трикутник (відповідність) --- сесія 3.06, завдання №18}

\zadnum На рисунку зображено прямокутник $ABCD$ та рівнобедрений трикутник $DKM$ ($DK = KM$). Пряма, що містить діагональ $AC$ прямокутника, проходить через точку $K$. $AB = 16$ см, $AD = 12$ см, $DM = 2AD$. Установіть відповідність між відрізком (1--3) та його довжиною (А--Д).

\vspace{0.3cm}
\begin{center}
\begin{tikzpicture}[scale=0.4, thick]
    \coordinate (A) at (0,0);
    \coordinate (B) at (0,16);
    \coordinate (Cc) at (12,16);
    \coordinate (D) at (12,0);
    \coordinate (M) at (36,0);
    \coordinate (K) at (24,32);
    \draw[thick] (A) -- (B) -- (Cc) -- (D) -- cycle;
    \draw[thick] (A) -- (M);
    \draw[thick] (D) -- (K) -- (M);
    \draw[thick] (A) -- (K);
    \fill (K) circle (3pt); \node[above, font=\scriptsize] at (K) {$K$};
    \node[below left, font=\scriptsize] at (A) {$A$};
    \node[above left, font=\scriptsize] at (B) {$B$};
    \node[above, font=\scriptsize] at (Cc) {$C$};
    \node[below, font=\scriptsize] at (D) {$D$};
    \node[below right, font=\scriptsize] at (M) {$M$};
\end{tikzpicture}
\end{center}

\vspace{0.2cm}
\noindent
\begin{minipage}[t]{0.46\textwidth}
\textit{Початок речення}\par\vspace{0.2cm}
\textbf{1} \quad Довжина відрізка $AM$ дорівнює\\[0.2cm]
\textbf{2} \quad Відстань від точки $K$ до прямої $AM$ дорівнює\\[0.2cm]
\textbf{3} \quad Довжина відрізка $CK$ дорівнює
\end{minipage}
\hfill
\begin{minipage}[t]{0.28\textwidth}
\textit{Закінчення}\par\vspace{0.2cm}
\textbf{А} \quad $20$ см\\[0.15cm]
\textbf{Б} \quad $24$ см\\[0.15cm]
\textbf{В} \quad $32$ см\\[0.15cm]
\textbf{Г} \quad $36$ см\\[0.15cm]
\textbf{Д} \quad $48$ см
\end{minipage}
\hfill
\begin{minipage}[t]{0.2\textwidth}
\vspace{0.5cm}
\begin{flushright}\matchingGrid\end{flushright}
\end{minipage}

\solution{Уведемо координати: $A(0;0)$, $D(12;0)$, $C(12;16)$. Оскільки $DM = 2AD = 24$ і $M$ лежить на промені $AD$, то $M(36;0)$, тож $AM = 36$ см~--- Г. Трикутник $DKM$ рівнобедрений, тому $K$ над серединою $DM$: $x_K = 24$; пряма $AC$ (нахил $\tfrac43$) дає $y_K = \tfrac43\cdot24 = 32$, отже відстань від $K$ до прямої $AM$ дорівнює $32$ см~--- В. $CK = \sqrt{(24-12)^2 + (32-16)^2} = \sqrt{144+256} = 20$ см~--- А. Відповідь: \textbf{1~--~Г, 2~--~В, 3~--~А}.}

\vspace{0.7cm}

\zadnum На рисунку зображено прямокутник $ABCD$ та рівнобедрений трикутник $DKM$ ($DK = KM$). Пряма, що містить діагональ $AC$ прямокутника, проходить через точку $K$. $AB = 12$ см, $AD = 9$ см, $DM = 2AD$. Установіть відповідність між відрізком (1--3) та його довжиною (А--Д).

\vspace{0.3cm}
\begin{center}
\begin{tikzpicture}[scale=0.53, thick]
    \coordinate (A) at (0,0);
    \coordinate (B) at (0,12);
    \coordinate (Cc) at (9,12);
    \coordinate (D) at (9,0);
    \coordinate (M) at (27,0);
    \coordinate (K) at (18,24);
    \draw[thick] (A) -- (B) -- (Cc) -- (D) -- cycle;
    \draw[thick] (A) -- (M);
    \draw[thick] (D) -- (K) -- (M);
    \draw[thick] (A) -- (K);
    \fill (K) circle (3pt); \node[above, font=\scriptsize] at (K) {$K$};
    \node[below left, font=\scriptsize] at (A) {$A$};
    \node[above left, font=\scriptsize] at (B) {$B$};
    \node[above, font=\scriptsize] at (Cc) {$C$};
    \node[below, font=\scriptsize] at (D) {$D$};
    \node[below right, font=\scriptsize] at (M) {$M$};
\end{tikzpicture}
\end{center}

\vspace{0.2cm}
\noindent
\begin{minipage}[t]{0.46\textwidth}
\textit{Початок речення}\par\vspace{0.2cm}
\textbf{1} \quad Довжина відрізка $AM$ дорівнює\\[0.2cm]
\textbf{2} \quad Відстань від точки $K$ до прямої $AM$ дорівнює\\[0.2cm]
\textbf{3} \quad Довжина відрізка $CK$ дорівнює
\end{minipage}
\hfill
\begin{minipage}[t]{0.28\textwidth}
\textit{Закінчення}\par\vspace{0.2cm}
\textbf{А} \quad $24$ см\\[0.15cm]
\textbf{Б} \quad $15$ см\\[0.15cm]
\textbf{В} \quad $27$ см\\[0.15cm]
\textbf{Г} \quad $30$ см\\[0.15cm]
\textbf{Д} \quad $36$ см
\end{minipage}
\hfill
\begin{minipage}[t]{0.2\textwidth}
\vspace{0.5cm}
\begin{flushright}\matchingGrid\end{flushright}
\end{minipage}

\vspace{0.7cm}

\zadnum На рисунку зображено прямокутник $ABCD$ та рівнобедрений трикутник $DKM$ ($DK = KM$). Пряма, що містить діагональ $AC$ прямокутника, проходить через точку $K$. $AB = 8$ см, $AD = 6$ см, $DM = 2AD$. Установіть відповідність між відрізком (1--3) та його довжиною (А--Д).

\vspace{0.3cm}
\begin{center}
\begin{tikzpicture}[scale=0.8, thick]
    \coordinate (A) at (0,0);
    \coordinate (B) at (0,8);
    \coordinate (Cc) at (6,8);
    \coordinate (D) at (6,0);
    \coordinate (M) at (18,0);
    \coordinate (K) at (12,16);
    \draw[thick] (A) -- (B) -- (Cc) -- (D) -- cycle;
    \draw[thick] (A) -- (M);
    \draw[thick] (D) -- (K) -- (M);
    \draw[thick] (A) -- (K);
    \fill (K) circle (2.4pt); \node[above, font=\scriptsize] at (K) {$K$};
    \node[below left, font=\scriptsize] at (A) {$A$};
    \node[above left, font=\scriptsize] at (B) {$B$};
    \node[above, font=\scriptsize] at (Cc) {$C$};
    \node[below, font=\scriptsize] at (D) {$D$};
    \node[below right, font=\scriptsize] at (M) {$M$};
\end{tikzpicture}
\end{center}

\vspace{0.2cm}
\noindent
\begin{minipage}[t]{0.46\textwidth}
\textit{Початок речення}\par\vspace{0.2cm}
\textbf{1} \quad Довжина відрізка $AK$ дорівнює\\[0.2cm]
\textbf{2} \quad Відстань від точки $K$ до прямої $AB$ дорівнює\\[0.2cm]
\textbf{3} \quad Довжина відрізка $CK$ дорівнює
\end{minipage}
\hfill
\begin{minipage}[t]{0.28\textwidth}
\textit{Закінчення}\par\vspace{0.2cm}
\textbf{А} \quad $16$ см\\[0.15cm]
\textbf{Б} \quad $18$ см\\[0.15cm]
\textbf{В} \quad $10$ см\\[0.15cm]
\textbf{Г} \quad $20$ см\\[0.15cm]
\textbf{Д} \quad $12$ см
\end{minipage}
\hfill
\begin{minipage}[t]{0.2\textwidth}
\vspace{0.5cm}
\begin{flushright}\matchingGrid\end{flushright}
\end{minipage}
\taskBlock{Рівні прямокутні рівнобедрені трикутники (відповідність) --- сесія 4.06, завдання №18}
\zadtask{Прямокутні рівнобедрені трикутники $ABC$ і $KMN$ є рівними (див. рисунок). Точки $K$, $C$, $N$ і $A$ лежать на прямій $AK$, причому $AN = NC = CK$. $AC = 12$~см. Установіть відповідність між відрізком (1--3) та його довжиною (А--Д).

\vspace{0.2cm}
\noindent
\begin{minipage}[c]{0.62\textwidth}
\begin{minipage}[t]{0.55\textwidth}
\textit{Відрізок}\par\vspace{0.2cm}
\textbf{1} \quad $AK$\\[0.25cm]
\textbf{2} \quad $AB$\\[0.25cm]
\textbf{3} \quad $BM$
\end{minipage}%
\hfill
\begin{minipage}[t]{0.42\textwidth}
\textit{Довжина відрізка}\par\vspace{0.2cm}
\textbf{А} \quad $6\sqrt{2}$~см\\[0.15cm]
\textbf{Б} \quad $6\sqrt{5}$~см\\[0.15cm]
\textbf{В} \quad $12\sqrt{2}$~см\\[0.15cm]
\textbf{Г} \quad $18$~см\\[0.15cm]
\textbf{Д} \quad $24$~см
\end{minipage}
\par\vspace{0.3cm}
\matchingGrid
\end{minipage}%
\hfill
\begin{minipage}[c]{0.36\textwidth}
\centering
\begin{tikzpicture}[scale=0.45, thick]
    \coordinate (K) at (-3,3);
    \coordinate (M) at (1,3);
    \coordinate (N) at (1,-1);
    \coordinate (C) at (-1,1);
    \coordinate (B) at (-1,-3);
    \coordinate (A) at (3,-3);

    \draw (K) -- (M) -- (N);
    \draw (C) -- (B) -- (A);
    \draw (K) -- (A);

    \draw (M) ++(-0.4,0) -- ++(0,-0.4) -- ++(0.4,0);
    \draw (B) ++(0,0.4) -- ++(0.4,0) -- ++(0,-0.4);

    \draw (-2-0.15, 2-0.15) -- (-2+0.15, 2+0.15);
    \draw (0-0.15, 0-0.15) -- (0+0.15, 0+0.15);
    \draw (2-0.15, -2-0.15) -- (2+0.15, -2+0.15);

    \fill (C) circle (2pt);
    \fill (N) circle (2pt);

    \node[above left] at (K) {$K$};
    \node[above right] at (M) {$M$};
    \node[left] at (-1.1, 0.9) {$C$};
    \node[right] at (1.1, -0.9) {$N$};
    \node[below left] at (B) {$B$};
    \node[below right] at (A) {$A$};
\end{tikzpicture}
\end{minipage}
}
\solution{$AC$~--- гіпотенуза рівнобедреного прямокутного $\triangle ABC$, тож катети $AB=BC=\dfrac{AC}{\sqrt2}=\dfrac{12}{\sqrt2}=6\sqrt2$~см (А). Оскільки $AN=NC=CK=\dfrac{AC}{2}=6$, маємо $AK=AN+NC+CK=18$~см (Г). Точка $B$ лежить на відстані катета від $C$; $BN$ обчислюється з прямокутного трикутника $BCN$ ($BC=6\sqrt2$, $CN=6$): тоді й $BM=\sqrt{BC^2+CM^2}$ дає $6\sqrt5$~см (Б). Відповідь: \textbf{1~--~Г, 2~--~А, 3~--~Б}.}
\zadtask{Прямокутні рівнобедрені трикутники $ABC$ і $KMN$ є рівними (див. рисунок). Точки $K$, $C$, $N$ і $A$ лежать на прямій $AK$, причому $AN = NC = CK$. $AC = 8$~см. Установіть відповідність між відрізком (1--3) та його довжиною (А--Д).

\vspace{0.2cm}
\noindent
\begin{minipage}[c]{0.62\textwidth}
\begin{minipage}[t]{0.55\textwidth}
\textit{Відрізок}\par\vspace{0.2cm}
\textbf{1} \quad $AK$\\[0.25cm]
\textbf{2} \quad $AB$\\[0.25cm]
\textbf{3} \quad $BM$
\end{minipage}%
\hfill
\begin{minipage}[t]{0.42\textwidth}
\textit{Довжина відрізка}\par\vspace{0.2cm}
\textbf{А} \quad $4\sqrt{2}$~см\\[0.15cm]
\textbf{Б} \quad $4\sqrt{5}$~см\\[0.15cm]
\textbf{В} \quad $8\sqrt{2}$~см\\[0.15cm]
\textbf{Г} \quad $12$~см\\[0.15cm]
\textbf{Д} \quad $16$~см
\end{minipage}
\par\vspace{0.3cm}
\matchingGrid
\end{minipage}%
\hfill
\begin{minipage}[c]{0.36\textwidth}
\centering
\begin{tikzpicture}[scale=0.45, thick]
    \coordinate (K) at (-3,3);
    \coordinate (M) at (1,3);
    \coordinate (N) at (1,-1);
    \coordinate (C) at (-1,1);
    \coordinate (B) at (-1,-3);
    \coordinate (A) at (3,-3);

    \draw (K) -- (M) -- (N);
    \draw (C) -- (B) -- (A);
    \draw (K) -- (A);

    \draw (M) ++(-0.4,0) -- ++(0,-0.4) -- ++(0.4,0);
    \draw (B) ++(0,0.4) -- ++(0.4,0) -- ++(0,-0.4);

    \draw (-2-0.15, 2-0.15) -- (-2+0.15, 2+0.15);
    \draw (0-0.15, 0-0.15) -- (0+0.15, 0+0.15);
    \draw (2-0.15, -2-0.15) -- (2+0.15, -2+0.15);

    \fill (C) circle (2pt);
    \fill (N) circle (2pt);

    \node[above left] at (K) {$K$};
    \node[above right] at (M) {$M$};
    \node[left] at (-1.1, 0.9) {$C$};
    \node[right] at (1.1, -0.9) {$N$};
    \node[below left] at (B) {$B$};
    \node[below right] at (A) {$A$};
\end{tikzpicture}
\end{minipage}
}
\zadtask{Прямокутні рівнобедрені трикутники $ABC$ і $KMN$ є рівними (див. рисунок). Точки $K$, $C$, $N$ і $A$ лежать на прямій $AK$, причому $AN = NC = CK$. $AC = 18$~см. Установіть відповідність між відрізком (1--3) та його довжиною (А--Д).

\vspace{0.2cm}
\noindent
\begin{minipage}[c]{0.62\textwidth}
\begin{minipage}[t]{0.55\textwidth}
\textit{Відрізок}\par\vspace{0.2cm}
\textbf{1} \quad $KN$\\[0.25cm]
\textbf{2} \quad $MN$\\[0.25cm]
\textbf{3} \quad $BN$
\end{minipage}%
\hfill
\begin{minipage}[t]{0.42\textwidth}
\textit{Довжина відрізка}\par\vspace{0.2cm}
\textbf{А} \quad $9$~см\\[0.15cm]
\textbf{Б} \quad $9\sqrt{2}$~см\\[0.15cm]
\textbf{В} \quad $9\sqrt{5}$~см\\[0.15cm]
\textbf{Г} \quad $18$~см\\[0.15cm]
\textbf{Д} \quad $27$~см
\end{minipage}
\par\vspace{0.3cm}
\matchingGrid
\end{minipage}%
\hfill
\begin{minipage}[c]{0.36\textwidth}
\centering
\begin{tikzpicture}[scale=0.45, thick]
    \coordinate (K) at (-3,3);
    \coordinate (M) at (1,3);
    \coordinate (N) at (1,-1);
    \coordinate (C) at (-1,1);
    \coordinate (B) at (-1,-3);
    \coordinate (A) at (3,-3);

    \draw (K) -- (M) -- (N);
    \draw (C) -- (B) -- (A);
    \draw (K) -- (A);

    \draw (M) ++(-0.4,0) -- ++(0,-0.4) -- ++(0.4,0);
    \draw (B) ++(0,0.4) -- ++(0.4,0) -- ++(0,-0.4);

    \draw (-2-0.15, 2-0.15) -- (-2+0.15, 2+0.15);
    \draw (0-0.15, 0-0.15) -- (0+0.15, 0+0.15);
    \draw (2-0.15, -2-0.15) -- (2+0.15, -2+0.15);

    \fill (C) circle (2pt);
    \fill (N) circle (2pt);

    \node[above left] at (K) {$K$};
    \node[above right] at (M) {$M$};
    \node[left] at (-1.1, 0.9) {$C$};
    \node[right] at (1.1, -0.9) {$N$};
    \node[below left] at (B) {$B$};
    \node[below right] at (A) {$A$};
\end{tikzpicture}
\end{minipage}
}
\taskBlock{Кути в комбінації фігур (відповідність) --- сесія 5.06, завдання №18}
\zadnum На рисунку зображено правильний трикутник $ABC$ та прямокутні рівнобедрені трикутники $AKB$ і $BMC$. Установіть відповідність між кутом (1--3) та його градусною мірою (А--Д).

\vspace{0.3cm}
\begin{center}
\begin{tikzpicture}[scale=0.9, thick]
    \coordinate (A) at (0,0);
    \coordinate (C) at (2,0);
    \coordinate (B) at (1,1.732);
    \coordinate (K) at (-0.366,1.366);
    \coordinate (M) at (2.366,1.366);
    \draw[thick] (A) -- (B) -- (C) -- cycle;
    \draw[thick] (A) -- (K) -- (B);
    \draw[thick] (B) -- (M) -- (C);
    \draw[thick] (A) -- (M);
    \draw[thick] (K) -- (B);
    \pic[draw=yearOrange, line width=0.5pt, angle radius=0.3cm] {right angle = A--K--B};
    \pic[draw=yearOrange, line width=0.5pt, angle radius=0.3cm] {right angle = B--M--C};
    \node[below left, font=\scriptsize] at (A) {$A$};
    \node[above, font=\scriptsize] at (B) {$B$};
    \node[below right, font=\scriptsize] at (C) {$C$};
    \node[left, font=\scriptsize] at (K) {$K$};
    \node[right, font=\scriptsize] at (M) {$M$};
\end{tikzpicture}
\end{center}

\vspace{0.2cm}
\noindent
\begin{minipage}[t]{0.32\textwidth}
\textit{Кут}\par\vspace{0.2cm}
\textbf{1} \quad $\angle ACM$\\[0.25cm]
\textbf{2} \quad $\angle KBM$\\[0.25cm]
\textbf{3} \quad $\angle KAM$
\end{minipage}
\hfill
\begin{minipage}[t]{0.34\textwidth}
\textit{Градусна міра}\par\vspace{0.2cm}
\textbf{А} \quad $75^\circ$\\[0.15cm]
\textbf{Б} \quad $90^\circ$\\[0.15cm]
\textbf{В} \quad $105^\circ$\\[0.15cm]
\textbf{Г} \quad $135^\circ$\\[0.15cm]
\textbf{Д} \quad $150^\circ$
\end{minipage}
\hfill
\begin{minipage}[t]{0.2\textwidth}
\vspace{0.5cm}
\begin{flushright}\matchingGrid\end{flushright}
\end{minipage}
\par\vspace{0.5cm}
\solution{Кути правильного трикутника по $60^\circ$, а в прямокутних рівнобедрених трикутниках $AKB$, $BMC$ кути при основі по $45^\circ$. Тоді $1$: $\angle ACM=\angle ACB+\angle BCM=60^\circ+45^\circ=105^\circ$ (В); $2$: $\angle KBM=\angle KBA+\angle ABC+\angle CBM=45^\circ+60^\circ+45^\circ=150^\circ$ (Д); $3$: $\angle KAM=75^\circ$ (А). Відповідь: \textbf{1~--~В, 2~--~Д, 3~--~А}.}
\zadnum На рисунку зображено правильний трикутник $ABC$ та прямокутні рівнобедрені трикутники $AKB$ і $BMC$. Установіть відповідність між кутом (1--3) та його градусною мірою (А--Д).

\vspace{0.3cm}
\begin{center}
\begin{tikzpicture}[scale=0.9, thick]
    \coordinate (A) at (0,0);
    \coordinate (C) at (2,0);
    \coordinate (B) at (1,1.732);
    \coordinate (K) at (-0.366,1.366);
    \coordinate (M) at (2.366,1.366);
    \draw[thick] (A) -- (B) -- (C) -- cycle;
    \draw[thick] (A) -- (K) -- (B);
    \draw[thick] (B) -- (M) -- (C);
    \draw[thick] (K) -- (M);
    \draw[thick] (A) -- (K);
    \pic[draw=yearOrange, line width=0.5pt, angle radius=0.3cm] {right angle = A--K--B};
    \pic[draw=yearOrange, line width=0.5pt, angle radius=0.3cm] {right angle = B--M--C};
    \node[below left, font=\scriptsize] at (A) {$A$};
    \node[above, font=\scriptsize] at (B) {$B$};
    \node[below right, font=\scriptsize] at (C) {$C$};
    \node[left, font=\scriptsize] at (K) {$K$};
    \node[right, font=\scriptsize] at (M) {$M$};
\end{tikzpicture}
\end{center}

\vspace{0.2cm}
\noindent
\begin{minipage}[t]{0.32\textwidth}
\textit{Кут}\par\vspace{0.2cm}
\textbf{1} \quad $\angle MBC$\\[0.25cm]
\textbf{2} \quad $\angle AKM$\\[0.25cm]
\textbf{3} \quad $\angle KAC$
\end{minipage}
\hfill
\begin{minipage}[t]{0.34\textwidth}
\textit{Градусна міра}\par\vspace{0.2cm}
\textbf{А} \quad $45^\circ$\\[0.15cm]
\textbf{Б} \quad $75^\circ$\\[0.15cm]
\textbf{В} \quad $105^\circ$\\[0.15cm]
\textbf{Г} \quad $120^\circ$\\[0.15cm]
\textbf{Д} \quad $135^\circ$
\end{minipage}
\hfill
\begin{minipage}[t]{0.2\textwidth}
\vspace{0.5cm}
\begin{flushright}\matchingGrid\end{flushright}
\end{minipage}
\par\vspace{0.5cm}
\zadnum На рисунку зображено квадрат $ABCD$ та правильний трикутник $ABK$, побудований \textit{усередині} квадрата на стороні $AB$. Установіть відповідність між кутом (1--3) та його градусною мірою (А--Д).

\vspace{0.3cm}
\begin{center}
\begin{tikzpicture}[scale=0.9, thick]
    \coordinate (A) at (0,0);
    \coordinate (B) at (2,0);
    \coordinate (C) at (2,-2);
    \coordinate (D) at (0,-2);
    \coordinate (K) at (1,-1.732);
    \draw[thick] (A) -- (B) -- (C) -- (D) -- cycle;
    \draw[thick] (A) -- (K) -- (B);
    \draw[thick] (K) -- (C);
    \draw[thick] (K) -- (D);
    \node[above left, font=\scriptsize] at (A) {$A$};
    \node[above right, font=\scriptsize] at (B) {$B$};
    \node[below right, font=\scriptsize] at (C) {$C$};
    \node[below left, font=\scriptsize] at (D) {$D$};
    \node[below, font=\scriptsize] at (K) {$K$};
\end{tikzpicture}
\end{center}

\vspace{0.2cm}
\noindent
\begin{minipage}[t]{0.32\textwidth}
\textit{Кут}\par\vspace{0.2cm}
\textbf{1} \quad $\angle KCB$\\[0.25cm]
\textbf{2} \quad $\angle KCD$\\[0.25cm]
\textbf{3} \quad $\angle DKC$
\end{minipage}
\hfill
\begin{minipage}[t]{0.34\textwidth}
\textit{Градусна міра}\par\vspace{0.2cm}
\textbf{А} \quad $15^\circ$\\[0.15cm]
\textbf{Б} \quad $75^\circ$\\[0.15cm]
\textbf{В} \quad $105^\circ$\\[0.15cm]
\textbf{Г} \quad $135^\circ$\\[0.15cm]
\textbf{Д} \quad $150^\circ$
\end{minipage}
\hfill
\begin{minipage}[t]{0.2\textwidth}
\vspace{0.5cm}
\begin{flushright}\matchingGrid\end{flushright}
\end{minipage}
\par\vspace{0.5cm}

\chapterTitle{РОЗДІЛ 5. СТЕРЕОМЕТРІЯ}
\typeTitle{Завдання з вибором однієї відповіді (А--Д)}
\taskBlock{Перетин площин у піраміді --- сесія 23.05, завдання №3}

\noindent
\begin{minipage}[c]{0.58\textwidth}
\zadnum У чотирикутній піраміді $SABCD$ точка $M$ належить ребру $AB$. По якій прямій перетинаються площини $ABS$ і $MCS$? \nmtyear{2026}
\end{minipage}%
\hfill
\begin{minipage}[c]{0.40\textwidth}
\centering
\begin{tikzpicture}[scale=0.85, thick]
    \coordinate (A) at (1.6,0);
    \coordinate (B) at (-0.8,1.0);
    \coordinate (C) at (1.6,1.0);
    \coordinate (D) at (3.4,0.3);
    \coordinate (S) at (1.1,3.2);
    \coordinate (M) at ($(A)!0.55!(B)$);
    \draw (A) -- (S) -- (D) -- cycle;
    \draw (A) -- (D);
    \draw (S) -- (C);
    \draw (C) -- (D);
    \draw[dashed] (B) -- (C);
    \draw (B) -- (S);
    \draw (A) -- (B);
    \fill (M) circle (1.5pt);
    \node[above] at (S) {$S$};
    \node[left] at (B) {$B$};
    \node[above right] at (C) {$C$};
    \node[below] at (A) {$A$};
    \node[right] at (D) {$D$};
    \node[below left] at (M) {$M$};
\end{tikzpicture}
\end{minipage}
\par\vspace{0.2cm}
\answerTable{$MC$}{$AB$}{$SB$}{$BC$}{$SM$}
\solution{Спільними точками площин $ABS$ і $MCS$ є $S$ (вершина) та $M$ (бо $M\in AB\subset ABS$ і $M\in MCS$). Через дві спільні точки площини перетинаються по прямій $SM$. Відповідь: \textbf{Д}.}

\vspace{0.3cm}

\noindent
\begin{minipage}[c]{0.58\textwidth}
\zadnum У трикутній піраміді $SABC$ точка $K$ належить ребру $AC$. По якій прямій перетинаються площини $ACS$ і $KBS$?
\end{minipage}%
\hfill
\begin{minipage}[c]{0.40\textwidth}
\centering
\begin{tikzpicture}[scale=0.85, thick]
    \coordinate (A) at (0,0);
    \coordinate (C) at (3.4,0.4);
    \coordinate (B) at (2.2,1.4);
    \coordinate (S) at (1.6,3.2);
    \coordinate (K) at ($(A)!0.5!(C)$);
    \draw (A) -- (C) -- (S) -- cycle;
    \draw (C) -- (B) -- (S);
    \draw[dashed] (A) -- (B);
    \draw[dashed] (B) -- (K);
    \fill (K) circle (1.5pt);
    \node[above] at (S) {$S$};
    \node[below left] at (A) {$A$};
    \node[right] at (B) {$B$};
    \node[below right] at (C) {$C$};
    \node[below] at (K) {$K$};
\end{tikzpicture}
\end{minipage}
\par\vspace{0.2cm}
\answerTable{$KB$}{$AC$}{$SK$}{$SB$}{$SA$}

\vspace{0.3cm}

\noindent
\begin{minipage}[c]{0.58\textwidth}
\zadnum У чотирикутній піраміді $SABCD$ точка $P$ належить ребру $CD$. По якій прямій перетинаються площини $SCD$ і $PBS$?
\end{minipage}%
\hfill
\begin{minipage}[c]{0.40\textwidth}
\centering
\begin{tikzpicture}[scale=0.85, thick]
    \coordinate (A) at (1.6,0);
    \coordinate (B) at (-0.8,1.0);
    \coordinate (C) at (1.6,1.0);
    \coordinate (D) at (3.4,0.3);
    \coordinate (S) at (1.1,3.2);
    \coordinate (P) at ($(C)!0.5!(D)$);
    \draw (A) -- (S) -- (D) -- cycle;
    \draw (A) -- (D);
    \draw (S) -- (C);
    \draw (C) -- (D);
    \draw[dashed] (B) -- (C);
    \draw (B) -- (S);
    \draw (A) -- (B);
    \draw[dashed] (B) -- (P);
    \fill (P) circle (1.5pt);
    \node[above] at (S) {$S$};
    \node[left] at (B) {$B$};
    \node[above] at (C) {$C$};
    \node[below] at (A) {$A$};
    \node[right] at (D) {$D$};
    \node[below right] at (P) {$P$};
\end{tikzpicture}
\end{minipage}
\par\vspace{0.2cm}
\answerTable{$BP$}{$CD$}{$SC$}{$SP$}{$SD$}

\vspace{0.3cm}

% ---------------------------------------------------------------------
\taskBlock{Модуль вектора у просторі --- сесія 23.05, завдання №8}

\zadtask{Визначте довжину (модуль) вектора $\vec{a}(5;\, 2;\, -1)$. \nmtyear{2026}}
\answerTable{$\sqrt{28}$}{$6$}{$\sqrt{30}$}{$\sqrt{32}$}{$15$}
\solution{$|\vec{a}|=\sqrt{5^2+2^2+(-1)^2}=\sqrt{25+4+1}=\sqrt{30}$. Відповідь: \textbf{В}.}

\vspace{0.3cm}

\zadtask{Визначте довжину (модуль) вектора $\vec{b}(-3;\, 6;\, 2)$.}
\answerTable{$5$}{$\sqrt{47}$}{$11$}{$7$}{$\sqrt{41}$}

\vspace{0.3cm}

\zadtask{Дано точки $A(1;\, -2;\, 3)$ та $B(3;\, 2;\, 1)$. Визначте довжину вектора $\vec{AB}$.}
\answerTable{$2\sqrt{5}$}{$6$}{$2\sqrt{6}$}{$\sqrt{26}$}{$\sqrt{14}$}

\vspace{0.3cm}

% ---------------------------------------------------------------------
\taskBlock{Стереометрія: апофеми піраміди --- сесія 30.05, завдання №2}

\noindent
\begin{minipage}[c]{0.60\textwidth}
\zadnum На рисунку зображено правильну трикутну піраміду $SABC$ з вершиною $S$. Скільки всього апофем можна провести в цій піраміді, якщо у неї $4$ грані?
\end{minipage}%
\hfill
\begin{minipage}[c]{0.38\textwidth}
\centering
\begin{tikzpicture}[scale=0.9, thick]
    \coordinate (S) at (1.2, 3);
    \coordinate (A) at (1.5, 0);
    \coordinate (B) at (3, 1);
    \coordinate (C) at (0, 1);
    \draw[thick] (A) -- (B) -- (S) -- cycle;
    \draw[thick] (S) -- (A);
    \draw[thick, dashed] (C) -- (B);
    \draw[thick] (C) -- (A);
    \draw[thick] (C) -- (S);
    \node[above] at (S) {$S$};
    \node[below] at (A) {$A$};
    \node[right] at (B) {$B$};
    \node[left] at (C) {$C$};
\end{tikzpicture}
\end{minipage}
\par\vspace{0.2cm}
\begin{flushleft}
\textbf{А} \quad жодної \\[0.1cm]
\textbf{Б} \quad лише одну \\[0.1cm]
\textbf{В} \quad лише дві \\[0.1cm]
\textbf{Г} \quad лише три \\[0.1cm]
\textbf{Д} \quad більше трьох
\end{flushleft}
\solution{Апофема — це висота бічної грані, проведена з вершини $S$. Правильна трикутна піраміда має $3$ бічні грані, отже $3$ апофеми. Відповідь: \textbf{Г} (лише три).}
\vspace{0.3cm}

\noindent
\begin{minipage}[c]{0.60\textwidth}
\zadnum На рисунку зображено правильну чотирикутну піраміду $SABCD$ з вершиною $S$. Скільки всього апофем можна провести в цій піраміді, якщо у неї $5$ граней?
\end{minipage}%
\hfill
\begin{minipage}[c]{0.38\textwidth}
\centering
\begin{tikzpicture}[scale=0.9, thick]
    \coordinate (S) at (1.6, 3.2);
    \coordinate (A) at (0.2, 0);
    \coordinate (B) at (3.0, 0);
    \coordinate (C) at (3.6, 1.2);
    \coordinate (D) at (0.8, 1.2);
    \draw[thick] (A) -- (B) -- (S) -- cycle;
    \draw[thick] (S) -- (A);
    \draw[thick] (B) -- (C) -- (S);
    \draw[thick, dashed] (C) -- (D);
    \draw[thick, dashed] (D) -- (A);
    \draw[thick, dashed] (D) -- (S);
    \node[above] at (S) {$S$};
    \node[below left] at (A) {$A$};
    \node[below right] at (B) {$B$};
    \node[right] at (C) {$C$};
    \node[left] at (D) {$D$};
\end{tikzpicture}
\end{minipage}
\par\vspace{0.2cm}
\begin{flushleft}
\textbf{А} \quad жодної \\[0.1cm]
\textbf{Б} \quad лише дві \\[0.1cm]
\textbf{В} \quad лише три \\[0.1cm]
\textbf{Г} \quad лише чотири \\[0.1cm]
\textbf{Д} \quad більше чотирьох
\end{flushleft}
\vspace{0.3cm}

\noindent
\begin{minipage}[c]{0.60\textwidth}
\zadnum На рисунку зображено правильну шестикутну піраміду $SABCDEF$ з вершиною $S$. Скільки всього ребер має ця піраміда, якщо у неї $7$ граней?
\end{minipage}%
\hfill
\begin{minipage}[c]{0.38\textwidth}
\centering
\begin{tikzpicture}[scale=0.85, thick]
    \coordinate (S) at (1.8, 3.4);
    \coordinate (A) at (0.0, 0.4);
    \coordinate (B) at (1.2, 0.0);
    \coordinate (C) at (3.0, 0.2);
    \coordinate (D) at (3.6, 0.9);
    \coordinate (E) at (2.4, 1.3);
    \coordinate (F) at (0.6, 1.1);
    \draw[thick] (A) -- (B) -- (C) -- (D);
    \draw[thick, dashed] (D) -- (E) -- (F) -- (A);
    \draw[thick] (S) -- (A);
    \draw[thick] (S) -- (B);
    \draw[thick] (S) -- (C);
    \draw[thick] (S) -- (D);
    \draw[thick, dashed] (S) -- (E);
    \draw[thick, dashed] (S) -- (F);
    \node[above] at (S) {$S$};
    \node[left] at (A) {$A$};
    \node[below] at (B) {$B$};
    \node[below right] at (C) {$C$};
    \node[right] at (D) {$D$};
    \node[above right] at (E) {$E$};
    \node[above left] at (F) {$F$};
\end{tikzpicture}
\end{minipage}
\par\vspace{0.2cm}
\begin{flushleft}
\textbf{А} \quad лише три \\[0.1cm]
\textbf{Б} \quad лише чотири \\[0.1cm]
\textbf{В} \quad лише п'ять \\[0.1cm]
\textbf{Г} \quad лише шість \\[0.1cm]
\textbf{Д} \quad більше шести
\end{flushleft}
\vspace{0.3cm}
\taskBlock{Стереометрія: сфера в координатах --- сесія 30.05, завдання №10}

\zadtask{У прямокутній системі координат у просторі задано сферу та дві точки $M(-4;\, 5;\, 0)$ і $N(4;\, -3;\, 4)$, які належать цій сфері, причому відстань між ними є \textit{найбільшою} серед усіх можливих відстаней між точками цієї сфери. Знайдіть радіус сфери.}
\answerTable{$4$}{$6$}{$8$}{$12$}{$24$}
\solution{Найбільша відстань між точками сфери — це діаметр, тож $MN$ є діаметром. $MN=\sqrt{(4-(-4))^2+(-3-5)^2+(4-0)^2}=\sqrt{64+64+16}=\sqrt{144}=12$. Радіус $R=\dfrac{12}{2}=6$. Відповідь: \textbf{Б} ($6$).}

\zadtask{У прямокутній системі координат у просторі задано сферу та дві точки $M(3;\, -2;\, 4)$ і $N(3;\, 4;\, -4)$, які належать цій сфері, причому відстань між ними є \textit{найбільшою} серед усіх можливих відстаней між точками цієї сфери. Знайдіть радіус сфери.}
\answerTable{$2{,}5$}{$5$}{$7$}{$10$}{$20$}

\zadtask{У прямокутній системі координат у просторі сферу задано рівнянням $(x-3)^2 + (y+4)^2 + (z-12)^2 = 100$. Знайдіть радіус цієї сфери.}
\answerTable{$5$}{$8$}{$10$}{$20$}{$40$}
\taskBlock{Мимобіжні прямі у призмі --- сесія 1.06, завдання №2}

\noindent
\begin{minipage}[c]{0.58\textwidth}
\zadnum У прямокутній чотирикутній призмі $ABCDA_1B_1C_1D_1$ проведено діагональ $B_1D$. Скільки ребер цієї призми \textit{мимобіжні} з прямою $B_1D$?
\end{minipage}%
\hfill
\begin{minipage}[c]{0.40\textwidth}
\centering
\begin{tikzpicture}[scale=0.6, thick]
    % Координати
    \coordinate (A) at (0,0);
    \coordinate (D) at (3,0);
    \coordinate (B) at (1,1.5); % B ззаду зліва (пунктир)
    \coordinate (C) at (4,1.5); % C ззаду справа
    \coordinate (A1) at (0,4);
    \coordinate (D1) at (3,4);
    \coordinate (B1) at (1,5.5);
    \coordinate (C1) at (4,5.5);
    % Невидимі лінії
    \draw[dashed] (A) -- (B);
    \draw[dashed] (B) -- (C);
    \draw[dashed] (B) -- (B1);
    % Видимі лінії
    \draw[thick] (A) -- (D) -- (C) -- (C1) -- (D1) -- (A1) -- cycle; % Контур
    \draw[thick] (A1) -- (B1) -- (C1); % Верх
    \draw[thick] (D) -- (D1); % Переднє ребро
    \draw[thick] (A) -- (A1); % Переднє ребро
    % Діагональ B1D
    \draw[mainGreen, line width=1.2pt] (B1) -- (D);
    % Підписи
    \node[left] at (A) {$A$};
    \node[right] at (D) {$D$};
    \node[left] at (B) {$B$};
    \node[right] at (C) {$C$};
    \node[left] at (A1) {$A_1$};
    \node[right] at (D1) {$D_1$};
    \node[left] at (B1) {$B_1$};
    \node[right] at (C1) {$C_1$};
\end{tikzpicture}
\end{minipage}
\par\vspace{0.2cm}
\answerTable{$4$}{$5$}{$6$}{$7$}{$8$}
\solution{Призма має $12$ ребер. Із прямою $B_1D$ перетинаються по $3$ ребра біля кожної з вершин $B_1$ і $D$ (разом $6$); решта $6$ ребер не мають із $B_1D$ спільних точок і не паралельні їй, тобто мимобіжні. Відповідь: \textbf{В}.}

\noindent
\begin{minipage}[c]{0.58\textwidth}
\zadnum У правильній трикутній призмі $ABCA_1B_1C_1$ проведено діагональ бічної грані $AB_1$. Скільки ребер цієї призми \textit{мимобіжні} з прямою $AB_1$?
\end{minipage}%
\hfill
\begin{minipage}[c]{0.40\textwidth}
\centering
\begin{tikzpicture}[scale=0.6, thick]
    \coordinate (A) at (0,0);
    \coordinate (B) at (3.4,0);
    \coordinate (C) at (2.2,1.4);
    \coordinate (A1) at (0,4);
    \coordinate (B1) at (3.4,4);
    \coordinate (C1) at (2.2,5.4);
    % Невидимі лінії
    \draw[dashed] (A) -- (C);
    \draw[dashed] (B) -- (C);
    \draw[dashed] (C) -- (C1);
    % Видимі лінії
    \draw[thick] (A) -- (B) -- (B1) -- (A1) -- cycle;
    \draw[thick] (A1) -- (C1) -- (B1);
    % Діагональ AB1
    \draw[mainGreen, line width=1.2pt] (A) -- (B1);
    % Підписи
    \node[below left] at (A) {$A$};
    \node[below right] at (B) {$B$};
    \node[right] at (C) {$C$};
    \node[above left] at (A1) {$A_1$};
    \node[above right] at (B1) {$B_1$};
    \node[above] at (C1) {$C_1$};
\end{tikzpicture}
\end{minipage}
\par\vspace{0.2cm}
\answerTable{$2$}{$3$}{$4$}{$5$}{$6$}

\noindent
\begin{minipage}[c]{0.58\textwidth}
\zadnum У кубі $ABCDA_1B_1C_1D_1$ проведено діагональ основи $AC$. Скільки ребер цього куба \textit{мимобіжні} з прямою $AC$?
\end{minipage}%
\hfill
\begin{minipage}[c]{0.40\textwidth}
\centering
\begin{tikzpicture}[scale=0.6, thick]
    \coordinate (A) at (0,0);
    \coordinate (B) at (3,0);
    \coordinate (C) at (4,1.5);
    \coordinate (D) at (1,1.5);
    \coordinate (A1) at (0,3);
    \coordinate (B1) at (3,3);
    \coordinate (C1) at (4,4.5);
    \coordinate (D1) at (1,4.5);
    % Невидимі лінії
    \draw[dashed] (A) -- (D);
    \draw[dashed] (D) -- (C);
    \draw[dashed] (D) -- (D1);
    % Видимі лінії
    \draw[thick] (A) -- (B) -- (C) -- (C1) -- (B1) -- (A1) -- cycle;
    \draw[thick] (A1) -- (D1) -- (C1);
    \draw[thick] (B) -- (B1);
    % Діагональ AC
    \draw[mainGreen, line width=1.2pt, dash pattern=on 4pt off 2pt] (A) -- (C);
    % Підписи
    \node[below left] at (A) {$A$};
    \node[below right] at (B) {$B$};
    \node[right] at (C) {$C$};
    \node[left] at (D) {$D$};
    \node[left] at (A1) {$A_1$};
    \node[right] at (B1) {$B_1$};
    \node[above right] at (C1) {$C_1$};
    \node[above left] at (D1) {$D_1$};
\end{tikzpicture}
\end{minipage}
\par\vspace{0.2cm}
\answerTable{$4$}{$5$}{$6$}{$7$}{$8$}

% =====================================================================
\taskBlock{Відстань у просторі від початку координат --- сесія 1.06, завдання №15 (відновлено)}
\zadtask{У прямокутній системі координат у просторі задано точку $O(0;\,0;\,0)$. Укажіть з-поміж наведених точку, відстань від якої до точки $O$ \textit{не дорівнює} $10$.}
\answerTable{$(0;\,6;\,8)$}{$(10;\,0;\,0)$}{$(9;\,1;\,0)$}{$(-6;\,0;\,8)$}{$(0;\,-8;\,-6)$}
\solution{Відстань від $O$ до $(x;y;z)$ дорівнює $\sqrt{x^2+y^2+z^2}$. Для чотирьох точок маємо $\sqrt{100}=10$, а для $(9;1;0)$: $\sqrt{81+1}=\sqrt{82}\approx 9{,}06\neq 10$. Відповідь: \textbf{В}.}
\zadtask{У прямокутній системі координат у просторі задано точку $O(0;\,0;\,0)$. Укажіть з-поміж наведених точку, відстань від якої до точки $O$ \textit{не дорівнює} $13$.}
\answerTable{$(12;\,4;\,3)$}{$(13;\,0;\,0)$}{$(0;\,5;\,12)$}{$(8;\,2;\,5)$}{$(3;\,-12;\,-4)$}
\zadtask{У прямокутній системі координат у просторі задано точки $O(0;\,0;\,0)$, $A(1;\,2;\,2)$, $B(2;\,3;\,6)$, $C(6;\,6;\,7)$, $D(4;\,4;\,2)$, $E(0;\,3;\,4)$. Укажіть точку, \textit{найбільш віддалену} від $O$.}
\answerTable{$A$}{$D$}{$C$}{$B$}{$E$}
\taskBlock{Паралельні прямі в кубі --- сесія 2.06, завдання №5}

\noindent
\begin{minipage}[c]{0.56\textwidth}
\zadnum На рисунку зображено куб $ABCDA_1B_1C_1D_1$. Скільки всього є прямих, що проходять через дві вершини куба та є паралельними прямій $AB_1$?

\vspace{0.2cm}
\begin{enumerate}[label=\textbf{\Alph*}, align=left, leftmargin=2.2em, labelsep=0.6em, itemsep=0.08cm, topsep=0.05cm]
\item жодної
\item лише одна
\item лише дві
\item лише три
\item більше як три
\end{enumerate}
\end{minipage}%
\hfill
\begin{minipage}[c]{0.40\textwidth}
\centering
\begin{tikzpicture}[scale=0.7, thick]
    \coordinate (A) at (0,0);
    \coordinate (B) at (1,0.7);
    \coordinate (C) at (3,0.7);
    \coordinate (D) at (2,0);
    \coordinate (A1) at (0,2.4);
    \coordinate (B1) at (1,3.1);
    \coordinate (C1) at (3,3.1);
    \coordinate (D1) at (2,2.4);
    % приховані ребра від B
    \draw[dashed] (B) -- (A);
    \draw[dashed] (B) -- (C);
    \draw[dashed] (B) -- (B1);
    % видимі
    \draw[thick] (A) -- (D) -- (C) -- (C1) -- (D1) -- (A1) -- cycle;
    \draw[thick] (A1) -- (B1) -- (C1);
    \draw[thick] (D) -- (D1);
    \draw[thick] (A) -- (A1);
    % діагональ AB1
    \draw[mainGreen, line width=1.2pt] (A) -- (B1);
    \node[below left] at (A) {$A$};
    \node[below] at (B) {$B$};
    \node[right] at (C) {$C$};
    \node[below right] at (D) {$D$};
    \node[left] at (A1) {$A_1$};
    \node[above left] at (B1) {$B_1$};
    \node[right] at (C1) {$C_1$};
    \node[right] at (D1) {$D_1$};
\end{tikzpicture}
\end{minipage}
\par\vspace{0.3cm}
\solution{$AB_1$ -- діагональ бічної грані. Паралельна їй пряма через дві вершини куба лише одна -- це діагональ протилежної грані $DC_1$. Відповідь: \textbf{Б}.}

\noindent
\begin{minipage}[c]{0.56\textwidth}
\zadnum На рисунку зображено куб $ABCDA_1B_1C_1D_1$. Скільки всього є прямих, що проходять через дві вершини куба та є паралельними прямій $AC_1$?

\vspace{0.2cm}
\begin{enumerate}[label=\textbf{\Alph*}, align=left, leftmargin=2.2em, labelsep=0.6em, itemsep=0.08cm, topsep=0.05cm]
\item жодної
\item лише одна
\item лише дві
\item лише три
\item більше як три
\end{enumerate}
\end{minipage}%
\hfill
\begin{minipage}[c]{0.40\textwidth}
\centering
\begin{tikzpicture}[scale=0.7, thick]
    \coordinate (A) at (0,0);
    \coordinate (B) at (1,0.7);
    \coordinate (C) at (3,0.7);
    \coordinate (D) at (2,0);
    \coordinate (A1) at (0,2.4);
    \coordinate (B1) at (1,3.1);
    \coordinate (C1) at (3,3.1);
    \coordinate (D1) at (2,2.4);
    % приховані ребра від B
    \draw[dashed] (B) -- (A);
    \draw[dashed] (B) -- (C);
    \draw[dashed] (B) -- (B1);
    % видимі
    \draw[thick] (A) -- (D) -- (C) -- (C1) -- (D1) -- (A1) -- cycle;
    \draw[thick] (A1) -- (B1) -- (C1);
    \draw[thick] (D) -- (D1);
    \draw[thick] (A) -- (A1);
    % діагональ AC1
    \draw[mainGreen, line width=1.2pt] (A) -- (C1);
    \node[below left] at (A) {$A$};
    \node[below] at (B) {$B$};
    \node[right] at (C) {$C$};
    \node[below right] at (D) {$D$};
    \node[left] at (A1) {$A_1$};
    \node[above left] at (B1) {$B_1$};
    \node[right] at (C1) {$C_1$};
    \node[right] at (D1) {$D_1$};
\end{tikzpicture}
\end{minipage}
\par\vspace{0.3cm}

\noindent
\begin{minipage}[c]{0.56\textwidth}
\zadnum На рисунку зображено куб $ABCDA_1B_1C_1D_1$. Скільки всього є прямих, що проходять через дві вершини куба та є паралельними прямій $AB$?

\vspace{0.2cm}
\begin{enumerate}[label=\textbf{\Alph*}, align=left, leftmargin=2.2em, labelsep=0.6em, itemsep=0.08cm, topsep=0.05cm]
\item жодної
\item лише одна
\item лише дві
\item лише три
\item більше як три
\end{enumerate}
\end{minipage}%
\hfill
\begin{minipage}[c]{0.40\textwidth}
\centering
\begin{tikzpicture}[scale=0.7, thick]
    \coordinate (A) at (0,0);
    \coordinate (B) at (1,0.7);
    \coordinate (C) at (3,0.7);
    \coordinate (D) at (2,0);
    \coordinate (A1) at (0,2.4);
    \coordinate (B1) at (1,3.1);
    \coordinate (C1) at (3,3.1);
    \coordinate (D1) at (2,2.4);
    % приховані ребра від B
    \draw[dashed] (B) -- (C);
    \draw[dashed] (B) -- (B1);
    % ребро AB (виділене, частково приховане)
    \draw[mainGreen, line width=1.2pt, dashed] (A) -- (B);
    % видимі
    \draw[thick] (A) -- (D) -- (C) -- (C1) -- (D1) -- (A1) -- cycle;
    \draw[thick] (A1) -- (B1) -- (C1);
    \draw[thick] (D) -- (D1);
    \draw[thick] (A) -- (A1);
    \node[below left] at (A) {$A$};
    \node[below] at (B) {$B$};
    \node[right] at (C) {$C$};
    \node[below right] at (D) {$D$};
    \node[left] at (A1) {$A_1$};
    \node[above left] at (B1) {$B_1$};
    \node[right] at (C1) {$C_1$};
    \node[right] at (D1) {$D_1$};
\end{tikzpicture}
\end{minipage}
\par\vspace{0.3cm}

% =====================================================================
\taskBlock{Вектор у просторі: координати точки --- сесія 2.06, завдання №12}

\zadtask{У прямокутній системі координат у просторі задано вектор $\vec{a}\,(5;\, -3;\, 8)$ та точки $K$ і $M(-7;\, 1;\, 4)$. Знайдіть координати точки $K$, якщо $\vec{a} = \overrightarrow{KM}$.}
\answerTable{$(-2;\, -2;\, 12)$}{$(-12;\, 4;\, -4)$}{$(-2;\, 4;\, 4)$}{$(2;\, 2;\, -12)$}{$(12;\, -4;\, 4)$}
\solution{Оскільки $\vec{a}=\overrightarrow{KM}=M-K$, маємо $K=M-\vec{a}=(-7-5;\,1-(-3);\,4-8)=(-12;\,4;\,-4)$. Відповідь: \textbf{Г}.}

\zadtask{У прямокутній системі координат у просторі задано вектор $\vec{a}\,(3;\, -7;\, 2)$ та точки $K$ і $M(-4;\, 5;\, 1)$. Знайдіть координати точки $K$, якщо $\vec{a} = \overrightarrow{KM}$.}
\answerTable{$(-1;\, -2;\, 3)$}{$(7;\, -12;\, 1)$}{$(-7;\, 12;\, -1)$}{$(1;\, 2;\, -3)$}{$(-7;\, -2;\, -1)$}

\zadtask{У прямокутній системі координат у просторі задано вектор $\vec{a}\,(-2;\, 6;\, -5)$ та точки $K(3;\, -1;\, 4)$ і $M$. Знайдіть координати точки $M$, якщо $\vec{a} = \overrightarrow{MK}$.}
\answerTable{$(5;\, -7;\, 9)$}{$(1;\, 5;\, -1)$}{$(-5;\, 7;\, -9)$}{$(-1;\, -5;\, 1)$}{$(1;\, -5;\, -1)$}

% =====================================================================
\taskBlock{Стереометрія: тіла обертання --- сесія 3.06, завдання №6}

\zadnum Конус утворений обертанням

\vspace{0.1cm}
\begin{enumerate}[label=\textbf{\Alph*}, align=left, leftmargin=2.2em, labelsep=0.6em, itemsep=0.06cm, topsep=0.04cm]
\item рівнобедреного трикутника навколо його основи
\item прямокутного трикутника навколо його гіпотенузи
\item рівностороннього трикутника навколо його сторони
\item прямокутного трикутника навколо його катета
\item прямокутника навколо його сторони
\end{enumerate}
\par\vspace{0.25cm}
\solution{Конус дістаємо обертанням прямокутного трикутника навколо одного з його катетів: цей катет стає віссю (висотою), другий катет~--- радіусом основи, гіпотенуза~--- твірною. Відповідь: \textbf{Г}.}

\zadnum Циліндр утворений обертанням

\vspace{0.1cm}
\begin{enumerate}[label=\textbf{\Alph*}, align=left, leftmargin=2.2em, labelsep=0.6em, itemsep=0.06cm, topsep=0.04cm]
\item прямокутного трикутника навколо його катета
\item ромба навколо його сторони
\item прямокутника навколо його сторони
\item прямокутника навколо його діагоналі
\item рівнобедреного трикутника навколо його основи
\end{enumerate}
\par\vspace{0.25cm}

\zadnum Унаслідок обертання прямокутного трикутника навколо одного з його катетів утворюється

\vspace{0.1cm}
\begin{enumerate}[label=\textbf{\Alph*}, align=left, leftmargin=2.2em, labelsep=0.6em, itemsep=0.06cm, topsep=0.04cm]
\item циліндр
\item конус
\item куля
\item зрізаний конус
\item два конуси зі спільною основою
\end{enumerate}
\par\vspace{0.25cm}
\taskBlock{Стереометрія: вектори у просторі --- сесія 3.06, завдання №12}

\zadtask{У прямокутній системі координат у просторі задано вектори $\vec{a}\,(-2;\, 2;\, 1)$ і $\vec{b}\,(4;\, -3;\, 0)$. Визначте модуль вектора $\vec{c} = \vec{a} + \vec{b}$.}
\answerTable{$\sqrt{5}$}{$2$}{$\sqrt{62}$}{$\sqrt{6}$}{$14$}
\solution{$\vec{c} = \vec{a} + \vec{b} = (-2+4;\, 2-3;\, 1+0) = (2;\, -1;\, 1)$. Тоді $|\vec{c}| = \sqrt{2^2 + (-1)^2 + 1^2} = \sqrt{6}$. Відповідь: \textbf{Г} ($\sqrt6$).}

\zadtask{У прямокутній системі координат у просторі задано вектори $\vec{a}\,(3;\, -1;\, 2)$ і $\vec{b}\,(-1;\, 4;\, 2)$. Визначте модуль вектора $\vec{c} = \vec{a} + \vec{b}$.}
\answerTable{$3$}{$\sqrt{29}$}{$\sqrt{21}$}{$29$}{$\sqrt{14}$}

\zadtask{У прямокутній системі координат у просторі задано вектори $\vec{a}\,(2;\, -3;\, 1)$ і $\vec{b}\,(4;\, -1;\, -1)$. Визначте модуль вектора $\vec{c} = \vec{a} - \vec{b}$.}
\answerTable{$2\sqrt{3}$}{$2\sqrt{13}$}{$12$}{$\sqrt{6}$}{$4$}
\taskBlock{Властивості просторових фігур --- сесія 4.06, завдання №5}
\zadtask{Укажіть просторову фігуру, яка не має основи.

\vspace{0.2cm}
\par\noindent\textbf{А}\ \ призма
\par\noindent\textbf{Б}\ \ конус
\par\noindent\textbf{В}\ \ циліндр
\par\noindent\textbf{Г}\ \ куля
\par\noindent\textbf{Д}\ \ піраміда
\vspace{0.3cm}}
\solution{Призма, конус, циліндр і піраміда мають основу (основи). Куля основи не має. Відповідь: \textbf{Г}.}
\zadtask{Укажіть просторову фігуру, яка має дві рівні основи, що є кругами.

\vspace{0.2cm}
\par\noindent\textbf{А}\ \ призма
\par\noindent\textbf{Б}\ \ конус
\par\noindent\textbf{В}\ \ циліндр
\par\noindent\textbf{Г}\ \ куля
\par\noindent\textbf{Д}\ \ піраміда
\vspace{0.3cm}}
\zadtask{Укажіть просторову фігуру, бічна поверхня якої складається лише з трикутників.

\vspace{0.2cm}
\par\noindent\textbf{А}\ \ призма
\par\noindent\textbf{Б}\ \ конус
\par\noindent\textbf{В}\ \ циліндр
\par\noindent\textbf{Г}\ \ куля
\par\noindent\textbf{Д}\ \ піраміда
\vspace{0.3cm}}
\taskBlock{Циліндр у прямокутній системі координат --- сесія 4.06, завдання №12}
\zadtask{У прямокутній системі координат у просторі задано циліндр та точки $O_1(-3;\, 2;\, 1)$ і $O_2(3;\, 4;\, 4)$, що є центрами його основ. Визначте довжину твірної циліндра.}
\answerTable{$7$}{$3\sqrt{5}$}{$\sqrt{13}$}{$\sqrt{61}$}{$3$}
\solution{Твірна дорівнює відстані між центрами основ: $O_1O_2=\sqrt{(3+3)^2+(4-2)^2+(4-1)^2}=\sqrt{36+4+9}=\sqrt{49}=7$. Відповідь: \textbf{А}.}
\zadtask{У прямокутній системі координат у просторі задано циліндр та точки $O_1(1;\, -2;\, 2)$ і $O_2(5;\, 2;\, 0)$, що є центрами його основ. Визначте довжину твірної циліндра.}
\answerTable{$4$}{$6$}{$2\sqrt{5}$}{$\sqrt{34}$}{$2\sqrt{2}$}
\zadtask{У прямокутній системі координат у просторі задано циліндр та точки $O_1(2;\, 0;\, -1)$ і $O_2(6;\, 4;\, 11)$, що є центрами його основ. Визначте координати центра циліндра (середини його осі).}
\answerTable{$(4;\, 2;\, 5)$}{$(8;\, 4;\, 10)$}{$(2;\, 2;\, 6)$}{$(4;\, 2;\, 6)$}{$(3;\, 2;\, 5)$}
\taskBlock{Сфера і дотична площина --- сесія 5.06, завдання №5}
\noindent
\begin{minipage}[c]{0.62\textwidth}
\zadnum Площина $\beta$ дотикається сфери в точці $M$ (див. рисунок). Скільки можна провести площин, кожна з яких перпендикулярна до площини $\beta$ і має зі сферою \textit{лише одну} спільну точку?

\vspace{0.1cm}
\begin{enumerate}[label=\textbf{\Alph*}, align=left, leftmargin=2.2em, labelsep=0.6em, itemsep=0.04cm, topsep=0.03cm]
\item жодної
\item лише одну
\item лише дві
\item лише три
\item безліч
\end{enumerate}
\end{minipage}%
\hfill
\begin{minipage}[c]{0.36\textwidth}
\centering
\begin{tikzpicture}[scale=0.7, thick]
    \draw[fill=cyan!12, draw=cyan!60!black] (0,0.9) circle (0.9);
    \draw[fill=yearOrange!12, draw=gray!60] (-1.5,-0.3) -- (1.5,-0.3) -- (2.1,0.3) -- (-0.9,0.3) -- cycle;
    \fill (0,0) circle (1.5pt); \node[below, font=\scriptsize] at (0.15,-0.05) {$M$};
    \node[font=\scriptsize] at (1.9,0.05) {$\beta$};
\end{tikzpicture}
\end{minipage}
\par\vspace{0.2cm}
\solution{Площина має зі сферою лише одну спільну точку тоді й тільки тоді, коли вона дотична до сфери. Через точку дотику $M$ можна провести безліч площин, перпендикулярних до $\beta$, і кожна з них дотикається сфери в $M$. Отже, таких площин безліч. Відповідь: \textbf{Д}.}
\noindent
\begin{minipage}[c]{0.62\textwidth}
\zadnum Площина $\beta$ дотикається сфери в точці $M$ (див. рисунок). Скільки можна провести площин, кожна з яких паралельна площині $\beta$ і має зі сферою \textit{лише одну} спільну точку?

\vspace{0.1cm}
\begin{enumerate}[label=\textbf{\Alph*}, align=left, leftmargin=2.2em, labelsep=0.6em, itemsep=0.04cm, topsep=0.03cm]
\item жодної
\item лише одну
\item лише дві
\item лише три
\item безліч
\end{enumerate}
\end{minipage}%
\hfill
\begin{minipage}[c]{0.36\textwidth}
\centering
\begin{tikzpicture}[scale=0.7, thick]
    \draw[fill=cyan!12, draw=cyan!60!black] (0,0.9) circle (0.9);
    \draw[fill=yearOrange!12, draw=gray!60] (-1.5,-0.3) -- (1.5,-0.3) -- (2.1,0.3) -- (-0.9,0.3) -- cycle;
    \fill (0,0) circle (1.5pt); \node[below, font=\scriptsize] at (0.15,-0.05) {$M$};
    \node[font=\scriptsize] at (1.9,0.05) {$\beta$};
\end{tikzpicture}
\end{minipage}
\par\vspace{0.2cm}
\noindent
\begin{minipage}[c]{0.62\textwidth}
\zadnum Площина $\beta$ дотикається сфери в точці $M$ (див. рисунок). Скільки можна провести прямих, кожна з яких лежить у площині $\beta$ і має зі сферою \textit{лише одну} спільну точку?

\vspace{0.1cm}
\begin{enumerate}[label=\textbf{\Alph*}, align=left, leftmargin=2.2em, labelsep=0.6em, itemsep=0.04cm, topsep=0.03cm]
\item жодної
\item лише одну
\item лише дві
\item лише три
\item безліч
\end{enumerate}
\end{minipage}%
\hfill
\begin{minipage}[c]{0.36\textwidth}
\centering
\begin{tikzpicture}[scale=0.7, thick]
    \draw[fill=cyan!12, draw=cyan!60!black] (0,0.9) circle (0.9);
    \draw[fill=yearOrange!12, draw=gray!60] (-1.5,-0.3) -- (1.5,-0.3) -- (2.1,0.3) -- (-0.9,0.3) -- cycle;
    \fill (0,0) circle (1.5pt); \node[below, font=\scriptsize] at (0.15,-0.05) {$M$};
    \node[font=\scriptsize] at (1.9,0.05) {$\beta$};
\end{tikzpicture}
\end{minipage}
\par\vspace{0.2cm}
\taskBlock{Довжина вектора у просторі --- сесія 5.06, завдання №7}
\zadtask{У прямокутній системі координат у просторі задано точки $A(1;\,2;\,3)$ і $B(3;\,1;\,1)$. Визначте довжину вектора $\overrightarrow{AB}$.}
\answerTable{$3$}{$\sqrt{6}$}{$9$}{$\sqrt{14}$}{$2\sqrt{3}$}
\solution{$\overrightarrow{AB}=(3-1;\,1-2;\,1-3)=(2;\,-1;\,-2)$. $|\overrightarrow{AB}|=\sqrt{2^2+(-1)^2+(-2)^2}=\sqrt{9}=3$. Відповідь: \textbf{А}.}
\zadtask{У прямокутній системі координат у просторі задано точки $A(1;\,0;\,2)$ і $B(3;\,3;\,8)$. Визначте довжину вектора $\overrightarrow{AB}$.}
\answerTable{$\sqrt{11}$}{$7$}{$49$}{$\sqrt{29}$}{$5$}
\zadtask{У прямокутній системі координат у просторі задано точки $A(1;\,2;\,3)$, $B(4;\,4;\,5)$ і $C(3;\,5;\,9)$. Визначте довжину вектора $\overrightarrow{AB} + \overrightarrow{BC}$.}
\answerTable{$\sqrt{17}$}{$7$}{$49$}{$5$}{$\sqrt{61}$}
\taskBlock{Середина відрізка у просторі --- сесія 5.06, завдання №12}
\zadtask{Визначте координати середини відрізка $AB$, якщо $A(1;\,2;\,3)$ і $B(3;\,1;\,1)$.}
\answerTable{$(-2;\,-1;\,-2)$}{$(1;\,-0{,}5;\,-1)$}{$(2;\,1{,}5;\,2)$}{$(4;\,3;\,1)$}{$(-1;\,-0{,}5;\,-1)$}
\solution{Координати середини~--- середні арифметичні: $\left(\dfrac{1+3}{2};\,\dfrac{2+1}{2};\,\dfrac{3+1}{2}\right)=(2;\,1{,}5;\,2)$. Відповідь: \textbf{В}.}
\zadtask{Визначте координати середини відрізка $AB$, якщо $A(2;\,-3;\,5)$ і $B(4;\,1;\,-1)$.}
\answerTable{$(6;\,-2;\,4)$}{$(1;\,2;\,-3)$}{$(-1;\,-2;\,3)$}{$(3;\,-1;\,2)$}{$(3;\,-2;\,2)$}
\zadtask{Точка $A'$ симетрична точці $A(1;\,4;\,2)$ відносно точки $O(3;\,-1;\,5)$. Визначте координати точки $A'$.}
\answerTable{$(5;\,-6;\,8)$}{$(2;\,1{,}5;\,3{,}5)$}{$(-1;\,9;\,-1)$}{$(4;\,3;\,7)$}{$(2;\,-5;\,3)$}
\typeTitle{Завдання з короткою відповіддю}
\taskBlock{Циліндр і конус зі спільною віссю --- сесія 23.05, завдання №19}

\noindent
\begin{minipage}[c]{0.60\textwidth}
\zadnum На рисунку зображено циліндр і конус, що мають спільну вісь. Радіус основи конуса вдвічі менший за радіус циліндра. Висота конуса дорівнює $15$~см, а твірна конуса~--- $17$~см. Знайдіть об'єм циліндра (у см$^3$), якщо площа його бічної поверхні дорівнює $640\pi$~см$^2$. У відповіді запишіть значення $\dfrac{V}{\pi}$. \nmtyear{2026}
\end{minipage}%
\hfill
\begin{minipage}[c]{0.36\textwidth}
\centering
\begin{tikzpicture}[scale=0.7, thick]
    \draw (0,0) ellipse (1.1 and 0.35);
    \draw (-1.1,0) -- (-1.1,3.2);
    \draw (1.1,0) -- (1.1,3.2);
    \draw (-1.1,3.2) arc (180:360:1.1 and 0.35);
    \draw[dashed] (-1.1,3.2) arc (180:0:1.1 and 0.35);
    \draw (3.7,0) ellipse (0.8 and 0.25);
    \draw (2.9,0) -- (3.7,3.5);
    \draw (4.5,0) -- (3.7,3.5);
\end{tikzpicture}
\end{minipage}
\par\vspace{0.2cm}
\nmtAnswerBox
\solution{Радіус конуса $r=\sqrt{17^2-15^2}=\sqrt{64}=8$ см, тож радіус циліндра $R=2r=16$ см. З бічної поверхні $2\pi R h=640\pi$ маємо $h=\dfrac{640}{2\cdot 16}=20$ см. Об'єм $V=\pi R^2 h=\pi\cdot 256\cdot 20=5120\pi$, тож $\dfrac{V}{\pi}=5120$. Відповідь: \textbf{5120}.}

\vspace{0.3cm}

\noindent
\begin{minipage}[c]{0.60\textwidth}
\zadnum На рисунку зображено циліндр і конус, що мають спільну вісь. Радіус основи конуса вдвічі менший за радіус циліндра. Висота конуса дорівнює $8$~см, а твірна конуса~--- $10$~см. Знайдіть об'єм циліндра (у см$^3$), якщо площа його бічної поверхні дорівнює $360\pi$~см$^2$. У відповіді запишіть значення $\dfrac{V}{\pi}$.
\end{minipage}%
\hfill
\begin{minipage}[c]{0.36\textwidth}
\centering
\begin{tikzpicture}[scale=0.7, thick]
    \draw (0,0) ellipse (1.1 and 0.35);
    \draw (-1.1,0) -- (-1.1,3.0);
    \draw (1.1,0) -- (1.1,3.0);
    \draw (-1.1,3.0) arc (180:360:1.1 and 0.35);
    \draw[dashed] (-1.1,3.0) arc (180:0:1.1 and 0.35);
    \draw (3.7,0) ellipse (0.8 and 0.25);
    \draw (2.9,0) -- (3.7,2.4);
    \draw (4.5,0) -- (3.7,2.4);
\end{tikzpicture}
\end{minipage}
\par\vspace{0.2cm}
\nmtAnswerBox

\vspace{0.3cm}

\noindent
\begin{minipage}[c]{0.60\textwidth}
\zadnum На рисунку зображено циліндр і конус, що мають спільну вісь та рівні радіуси основ. Висота конуса дорівнює $12$~см, а твірна конуса~--- $13$~см. Висота циліндра вдвічі більша за висоту конуса. Знайдіть об'єм циліндра (у см$^3$). У відповіді запишіть значення $\dfrac{V}{\pi}$.
\end{minipage}%
\hfill
\begin{minipage}[c]{0.36\textwidth}
\centering
\begin{tikzpicture}[scale=0.7, thick]
    \draw (0,0) ellipse (1.1 and 0.35);
    \draw (-1.1,0) -- (-1.1,3.1);
    \draw (1.1,0) -- (1.1,3.1);
    \draw (-1.1,3.1) arc (180:360:1.1 and 0.35);
    \draw[dashed] (-1.1,3.1) arc (180:0:1.1 and 0.35);
    \draw (3.7,0) ellipse (0.7 and 0.22);
    \draw (3.0,0) -- (3.7,3.3);
    \draw (4.4,0) -- (3.7,3.3);
\end{tikzpicture}
\end{minipage}
\par\vspace{0.2cm}
\nmtAnswerBox
\taskBlock{Стереометрія: призма і конус --- сесія 30.05, завдання №21}

\zadtask{Пряма чотирикутна призма й конус мають рівні висоти. В основі призми лежить рівнобічна трапеція, у яку можна вписати коло. Висота трапеції дорівнює діаметру основи конуса. Бічна сторона трапеції дорівнює $26$ см, а площа основи призми~--- $624$ см$^2$. Твірна конуса дорівнює $20$ см. Визначте площу (у см$^2$) бічної поверхні призми.

\vspace{0.3cm}
\begin{center}
\begin{tikzpicture}[scale=0.9, thick]
    \coordinate (A) at (0,0);
    \coordinate (B) at (5,0);
    \coordinate (C) at (4,1);
    \coordinate (D) at (1,1);
    \coordinate (A1) at (0,3);
    \coordinate (B1) at (5,3);
    \coordinate (C1) at (4,4);
    \coordinate (D1) at (1,4);

    \draw[thick] (A) -- (A1);
    \draw[thick] (B) -- (B1);
    \draw[thick] (C) -- (C1);
    \draw[thick, dashed] (D) -- (D1);

    \draw[thick] (A1) -- (B1) -- (C1) -- (D1) -- cycle;

    \draw[thick] (A) -- (B);
    \draw[thick] (B) -- (C);
    \draw[thick, dashed] (C) -- (D);
    \draw[thick, dashed] (D) -- (A);

    \draw[thick] ($(B)!0.5!(C) + (-0.1,0)$) -- ($(B)!0.5!(C) + (0.1,0)$);
    \draw[thick] ($(A)!0.5!(D) + (-0.1,0)$) -- ($(A)!0.5!(D) + (0.1,0)$);

    \coordinate (V) at (9, 4);
    \draw[thick] (9, 0) ellipse (1.8 and 0.5);
    \draw[thick] (V) -- (7.2, 0);
    \draw[thick] (V) -- (10.8, 0);
    \draw[thick, dashed] (V) -- (9, 0);
    \draw[thick, dashed] (9, 0) -- (10.8, 0);
    \fill (9, 0) circle (1.5pt);
\end{tikzpicture}
\end{center}
}
\nmtAnswerBox
\solution{У трапецію з вписаним колом сума основ дорівнює сумі бічних сторін: $a+b=2\cdot26=52$. Площа основи $\dfrac{a+b}{2}\cdot h_{\text{тр}}=624\Rightarrow 26\,h_{\text{тр}}=624\Rightarrow h_{\text{тр}}=24$ — це діаметр основи конуса, тож $r=12$. Висота конуса (= висота призми) $h=\sqrt{20^2-12^2}=\sqrt{256}=16$. Периметр трапеції $P=(a+b)+2\cdot26=52+52=104$. Бічна поверхня призми $P\cdot h=104\cdot16=1664$. Відповідь: $1664$.}

\zadtask{Пряма чотирикутна призма й конус мають рівні висоти. В основі призми лежить рівнобічна трапеція, у яку можна вписати коло. Висота трапеції дорівнює діаметру основи конуса. Бічна сторона трапеції дорівнює $13$ см, а площа основи призми~--- $130$ см$^2$. Твірна конуса дорівнює $13$ см. Визначте площу (у см$^2$) бічної поверхні призми.

\vspace{0.3cm}
\begin{center}
\begin{tikzpicture}[scale=0.9, thick]
    \coordinate (A) at (0,0);
    \coordinate (B) at (5,0);
    \coordinate (C) at (4,1);
    \coordinate (D) at (1,1);
    \coordinate (A1) at (0,3);
    \coordinate (B1) at (5,3);
    \coordinate (C1) at (4,4);
    \coordinate (D1) at (1,4);

    \draw[thick] (A) -- (A1);
    \draw[thick] (B) -- (B1);
    \draw[thick] (C) -- (C1);
    \draw[thick, dashed] (D) -- (D1);

    \draw[thick] (A1) -- (B1) -- (C1) -- (D1) -- cycle;

    \draw[thick] (A) -- (B);
    \draw[thick] (B) -- (C);
    \draw[thick, dashed] (C) -- (D);
    \draw[thick, dashed] (D) -- (A);

    \draw[thick] ($(B)!0.5!(C) + (-0.1,0)$) -- ($(B)!0.5!(C) + (0.1,0)$);
    \draw[thick] ($(A)!0.5!(D) + (-0.1,0)$) -- ($(A)!0.5!(D) + (0.1,0)$);

    \coordinate (V) at (9, 4);
    \draw[thick] (9, 0) ellipse (1.8 and 0.5);
    \draw[thick] (V) -- (7.2, 0);
    \draw[thick] (V) -- (10.8, 0);
    \draw[thick, dashed] (V) -- (9, 0);
    \draw[thick, dashed] (9, 0) -- (10.8, 0);
    \fill (9, 0) circle (1.5pt);
\end{tikzpicture}
\end{center}
}
\nmtAnswerBox

\zadtask{Пряма чотирикутна призма й конус мають рівні висоти. В основі призми лежить рівнобічна трапеція, у яку можна вписати коло. Висота трапеції дорівнює діаметру основи конуса. Бічна сторона трапеції дорівнює $15$ см, а площа основи призми~--- $180$ см$^2$. Твірна конуса дорівнює $10$ см. Визначте площу (у см$^2$) \textit{повної} поверхні призми.

\vspace{0.3cm}
\begin{center}
\begin{tikzpicture}[scale=0.9, thick]
    \coordinate (A) at (0,0);
    \coordinate (B) at (5,0);
    \coordinate (C) at (4,1);
    \coordinate (D) at (1,1);
    \coordinate (A1) at (0,3);
    \coordinate (B1) at (5,3);
    \coordinate (C1) at (4,4);
    \coordinate (D1) at (1,4);

    \draw[thick] (A) -- (A1);
    \draw[thick] (B) -- (B1);
    \draw[thick] (C) -- (C1);
    \draw[thick, dashed] (D) -- (D1);

    \draw[thick] (A1) -- (B1) -- (C1) -- (D1) -- cycle;

    \draw[thick] (A) -- (B);
    \draw[thick] (B) -- (C);
    \draw[thick, dashed] (C) -- (D);
    \draw[thick, dashed] (D) -- (A);

    \draw[thick] ($(B)!0.5!(C) + (-0.1,0)$) -- ($(B)!0.5!(C) + (0.1,0)$);
    \draw[thick] ($(A)!0.5!(D) + (-0.1,0)$) -- ($(A)!0.5!(D) + (0.1,0)$);

    \coordinate (V) at (9, 4);
    \draw[thick] (9, 0) ellipse (1.8 and 0.5);
    \draw[thick] (V) -- (7.2, 0);
    \draw[thick] (V) -- (10.8, 0);
    \draw[thick, dashed] (V) -- (9, 0);
    \draw[thick, dashed] (9, 0) -- (10.8, 0);
    \fill (9, 0) circle (1.5pt);
\end{tikzpicture}
\end{center}
}
\nmtAnswerBox
\taskBlock{Конус і правильна трикутна піраміда, об'єм конуса --- сесія 1.06, завдання №21}

\noindent
\begin{minipage}[c]{0.62\textwidth}
\zadnum Конус і правильна трикутна піраміда мають \textit{рівні висоти}. Радіус основи конуса дорівнює радіусу кола, \textit{вписаного} в основу піраміди, і дорівнює $4\sqrt{3}$. Площа бічної поверхні піраміди дорівнює $288$. Знайдіть об'єм конуса. У відповідь запишіть значення $\dfrac{V}{\pi}$.
\end{minipage}%
\hfill
\begin{minipage}[c]{0.36\textwidth}
\centering
\begin{tikzpicture}[scale=0.85, thick]
    % Конус
    \draw[thick] (0,0) ellipse (1.3 and 0.38);
    \coordinate (CV) at (0,3);
    \draw[thick] (CV) -- (-1.3,0);
    \draw[thick] (CV) -- (1.3,0);
    \draw[thick, dashed] (CV) -- (0,0);
    \draw[<-, >=stealth, gray] (0,0) -- (1.3,0);
    \node[below, font=\scriptsize] at (0.7,-0.05) {$R$};
    % Піраміда
    \coordinate (A) at (3.2,-0.35);
    \coordinate (B) at (5.6,-0.35);
    \coordinate (C) at (4.4,0.45);
    \coordinate (S) at (4.4,3);
    \draw[thick] (A) -- (B) -- (S) -- cycle;
    \draw[thick] (A) -- (S);
    \draw[thick, dashed] (A) -- (C) -- (B);
    \draw[thick, dashed] (C) -- (S);
\end{tikzpicture}
\end{minipage}
\nmtAnswerBox
\solution{Для основи піраміди $r=\dfrac{a}{2\sqrt3}=4\sqrt3$, звідки сторона $a=24$. Площа бічної поверхні $S=\dfrac12\cdot 3a\cdot l=288$, тож $l=\dfrac{288}{36}=8$ (апофема). З $l^2=h^2+r^2$: $h^2=64-48=16$, $h=4$. Тоді висота й радіус конуса: $h=4$, $R=4\sqrt3$, а $V=\dfrac13\pi R^2 h=\dfrac13\pi\cdot48\cdot4=64\pi$. Отже $\dfrac{V}{\pi}=64$. Відповідь: \textbf{$64$}.}

\noindent
\begin{minipage}[c]{0.62\textwidth}
\zadnum Конус і правильна трикутна піраміда мають \textit{рівні висоти}. Радіус основи конуса дорівнює радіусу кола, \textit{вписаного} в основу піраміди, і дорівнює $2\sqrt{3}$. Площа бічної поверхні піраміди дорівнює $72$. Знайдіть об'єм конуса. У відповідь запишіть значення $\dfrac{V}{\pi}$.
\end{minipage}%
\hfill
\begin{minipage}[c]{0.36\textwidth}
\centering
\begin{tikzpicture}[scale=0.85, thick]
    % Конус
    \draw[thick] (0,0) ellipse (1.3 and 0.38);
    \coordinate (CV) at (0,3);
    \draw[thick] (CV) -- (-1.3,0);
    \draw[thick] (CV) -- (1.3,0);
    \draw[thick, dashed] (CV) -- (0,0);
    \draw[<-, >=stealth, gray] (0,0) -- (1.3,0);
    \node[below, font=\scriptsize] at (0.7,-0.05) {$R$};
    % Піраміда
    \coordinate (A) at (3.2,-0.35);
    \coordinate (B) at (5.6,-0.35);
    \coordinate (C) at (4.4,0.45);
    \coordinate (S) at (4.4,3);
    \draw[thick] (A) -- (B) -- (S) -- cycle;
    \draw[thick] (A) -- (S);
    \draw[thick, dashed] (A) -- (C) -- (B);
    \draw[thick, dashed] (C) -- (S);
\end{tikzpicture}
\end{minipage}
\nmtAnswerBox

\noindent
\begin{minipage}[c]{0.62\textwidth}
\zadnum Конус і правильна трикутна піраміда мають \textit{рівні висоти}. Радіус основи конуса дорівнює радіусу кола, \textit{описаного} навколо основи піраміди. Радіус кола, \textit{вписаного} в основу піраміди, дорівнює $3\sqrt{3}$, а площа бічної поверхні піраміди дорівнює $162$. Знайдіть об'єм конуса. У відповідь запишіть значення $\dfrac{V}{\pi}$.
\end{minipage}%
\hfill
\begin{minipage}[c]{0.36\textwidth}
\centering
\begin{tikzpicture}[scale=0.85, thick]
    % Конус
    \draw[thick] (0,0) ellipse (1.3 and 0.38);
    \coordinate (CV) at (0,3);
    \draw[thick] (CV) -- (-1.3,0);
    \draw[thick] (CV) -- (1.3,0);
    \draw[thick, dashed] (CV) -- (0,0);
    \draw[<-, >=stealth, gray] (0,0) -- (1.3,0);
    \node[below, font=\scriptsize] at (0.7,-0.05) {$R$};
    % Піраміда
    \coordinate (A) at (3.2,-0.35);
    \coordinate (B) at (5.6,-0.35);
    \coordinate (C) at (4.4,0.45);
    \coordinate (S) at (4.4,3);
    \draw[thick] (A) -- (B) -- (S) -- cycle;
    \draw[thick] (A) -- (S);
    \draw[thick, dashed] (A) -- (C) -- (B);
    \draw[thick, dashed] (C) -- (S);
\end{tikzpicture}
\end{minipage}
\nmtAnswerBox
\taskBlock{Об'єми півсфери та конуса --- сесія 2.06, завдання №21}

\noindent
\begin{minipage}[c]{0.62\textwidth}
\zadnum Радіус основи й твірна конуса відносяться як $3 : 5$, а площа його бічної поверхні дорівнює $135\pi$. Площа повної поверхні півсфери дорівнює $243\pi$. У скільки разів об'єм півсфери більший за об'єм конуса?
\end{minipage}%
\hfill
\begin{minipage}[c]{0.36\textwidth}
\centering
\begin{tikzpicture}[scale=0.8, thick]
    % конус
    \draw[thick] (0,0) ellipse (1.0 and 0.28);
    \draw[thick] (-1.0,0) -- (0,2.4) -- (1.0,0);
    \draw[thick, dashed] (-1.0,0) -- (1.0,0);
    % півсфера
    \begin{scope}[shift={(3.4,0)}]
        \draw[thick] (-1.1,0) arc (180:360:1.1 and 1.1);
        \draw[thick] (-1.1,0) -- (1.1,0);
        \draw[thick, dashed] (0,0) ellipse (1.1 and 0.28);
    \end{scope}
\end{tikzpicture}
\end{minipage}
\nmtAnswerBox
\solution{Конус: $r:l=3:5$, $S_{\text{бік}}=\pi rl=135\pi$. Із $l=\dfrac{5r}{3}$: $\pi r\cdot\dfrac{5r}{3}=135\pi$, тож $r^2=81$, $r=9$, $l=15$, висота $h=\sqrt{l^2-r^2}=\sqrt{225-81}=12$, $V_{\text{кон}}=\dfrac13\pi r^2h=\dfrac13\pi\cdot81\cdot12=324\pi$. Півсфера: $3\pi R^2=243\pi$, $R=9$, $V_{\text{півсф}}=\dfrac23\pi R^3=\dfrac23\pi\cdot729=486\pi$. Відношення $\dfrac{486\pi}{324\pi}=1{,}5$.}

\noindent
\begin{minipage}[c]{0.62\textwidth}
\zadnum Радіус основи й твірна конуса відносяться як $3 : 5$, а площа його бічної поверхні дорівнює $60\pi$. Площа повної поверхні півсфери дорівнює $432\pi$. У скільки разів об'єм півсфери більший за об'єм конуса?
\end{minipage}%
\hfill
\begin{minipage}[c]{0.36\textwidth}
\centering
\begin{tikzpicture}[scale=0.8, thick]
    % конус
    \draw[thick] (0,0) ellipse (1.0 and 0.28);
    \draw[thick] (-1.0,0) -- (0,2.4) -- (1.0,0);
    \draw[thick, dashed] (-1.0,0) -- (1.0,0);
    % півсфера
    \begin{scope}[shift={(3.4,0)}]
        \draw[thick] (-1.3,0) arc (180:360:1.3 and 1.3);
        \draw[thick] (-1.3,0) -- (1.3,0);
        \draw[thick, dashed] (0,0) ellipse (1.3 and 0.32);
    \end{scope}
\end{tikzpicture}
\end{minipage}
\nmtAnswerBox

\noindent
\begin{minipage}[c]{0.62\textwidth}
\zadnum Радіус основи й твірна конуса відносяться як $4 : 5$, а площа його бічної поверхні дорівнює $80\pi$. Площа повної поверхні півсфери дорівнює $48\pi$. У скільки разів об'єм конуса більший за об'єм півсфери?
\end{minipage}%
\hfill
\begin{minipage}[c]{0.36\textwidth}
\centering
\begin{tikzpicture}[scale=0.8, thick]
    % конус
    \draw[thick] (0,0) ellipse (1.2 and 0.3);
    \draw[thick] (-1.2,0) -- (0,2.1) -- (1.2,0);
    \draw[thick, dashed] (-1.2,0) -- (1.2,0);
    % півсфера
    \begin{scope}[shift={(3.4,0)}]
        \draw[thick] (-0.9,0) arc (180:360:0.9 and 0.9);
        \draw[thick] (-0.9,0) -- (0.9,0);
        \draw[thick, dashed] (0,0) ellipse (0.9 and 0.24);
    \end{scope}
\end{tikzpicture}
\end{minipage}
\nmtAnswerBox
\taskBlock{Стереометрія: циліндр і паралелепіпед --- сесія 3.06, завдання №21}

\noindent
\begin{minipage}[c]{0.62\textwidth}
\zadnum На рисунку зображено циліндр і прямокутний паралелепіпед. Діаметр основи циліндра дорівнює діагоналі основи паралелепіпеда. Довжина діагоналі паралелепіпеда дорівнює $39$ см, сторони основи~--- $9$ см і $12$ см. Висота циліндра відноситься до висоти паралелепіпеда як $2 : 3$. Визначте об'єм $V$ циліндра (у см$^3$). У відповідь запишіть значення $\dfrac{V}{\pi}$.
\end{minipage}%
\hfill
\begin{minipage}[c]{0.36\textwidth}
\centering
\begin{tikzpicture}[scale=0.55, thick]
    % циліндр
    \draw[thick] (0,0) ellipse (0.9 and 0.26);
    \draw[thick] (-0.9,0) -- (-0.9,2.6);
    \draw[thick] (0.9,0) -- (0.9,2.6);
    \draw[thick] (0,2.6) ellipse (0.9 and 0.26);
    % паралелепіпед
    \begin{scope}[shift={(3.2,0)}]
        \coordinate (A) at (0,0);
        \coordinate (B) at (2,0);
        \coordinate (C) at (2.7,0.7);
        \coordinate (D) at (0.7,0.7);
        \draw[thick] (A) -- (B) -- (C) -- (D) -- cycle;
        \draw[thick] (A) -- (0,2.6);
        \draw[thick] (B) -- (2,2.6);
        \draw[thick] (C) -- (2.7,3.3);
        \draw[thick, dashed] (D) -- (0.7,3.3);
        \draw[thick] (0,2.6) -- (2,2.6) -- (2.7,3.3) -- (0.7,3.3) -- cycle;
    \end{scope}
\end{tikzpicture}
\end{minipage}
\nmtAnswerBox
\solution{Діагональ основи паралелепіпеда $d = \sqrt{9^2+12^2} = \sqrt{225} = 15$ см~--- це діаметр циліндра, тож $r = 7{,}5$ см. Висота паралелепіпеда $H = \sqrt{39^2 - 15^2} = \sqrt{1296} = 36$ см, а висота циліндра $h = \dfrac23 H = 24$ см. Об'єм $V = \pi r^2 h = \pi\cdot 7{,}5^2\cdot 24 = 1350\pi$, тож $\dfrac{V}{\pi} = 1350$. Відповідь: \textbf{$1350$}.}

\noindent
\begin{minipage}[c]{0.62\textwidth}
\zadnum На рисунку зображено циліндр і прямокутний паралелепіпед. Діаметр основи циліндра дорівнює діагоналі основи паралелепіпеда. Довжина діагоналі паралелепіпеда дорівнює $26$ см, сторони основи~--- $6$ см і $8$ см. Висота циліндра відноситься до висоти паралелепіпеда як $2 : 3$. Визначте об'єм $V$ циліндра (у см$^3$). У відповідь запишіть значення $\dfrac{V}{\pi}$.
\end{minipage}%
\hfill
\begin{minipage}[c]{0.36\textwidth}
\centering
\begin{tikzpicture}[scale=0.55, thick]
    % циліндр
    \draw[thick] (0,0) ellipse (0.9 and 0.26);
    \draw[thick] (-0.9,0) -- (-0.9,2.6);
    \draw[thick] (0.9,0) -- (0.9,2.6);
    \draw[thick] (0,2.6) ellipse (0.9 and 0.26);
    % паралелепіпед
    \begin{scope}[shift={(3.2,0)}]
        \coordinate (A) at (0,0);
        \coordinate (B) at (2,0);
        \coordinate (C) at (2.7,0.7);
        \coordinate (D) at (0.7,0.7);
        \draw[thick] (A) -- (B) -- (C) -- (D) -- cycle;
        \draw[thick] (A) -- (0,2.6);
        \draw[thick] (B) -- (2,2.6);
        \draw[thick] (C) -- (2.7,3.3);
        \draw[thick, dashed] (D) -- (0.7,3.3);
        \draw[thick] (0,2.6) -- (2,2.6) -- (2.7,3.3) -- (0.7,3.3) -- cycle;
    \end{scope}
\end{tikzpicture}
\end{minipage}
\nmtAnswerBox

\noindent
\begin{minipage}[c]{0.62\textwidth}
\zadnum На рисунку зображено циліндр і прямокутний паралелепіпед. Діаметр основи циліндра дорівнює діагоналі основи паралелепіпеда, а \textit{висота циліндра дорівнює діагоналі паралелепіпеда}. Сторони основи паралелепіпеда~--- $12$ см і $16$ см, висота паралелепіпеда~--- $15$ см. Визначте об'єм $V$ циліндра (у см$^3$). У відповідь запишіть значення $\dfrac{V}{\pi}$.
\end{minipage}%
\hfill
\begin{minipage}[c]{0.36\textwidth}
\centering
\begin{tikzpicture}[scale=0.55, thick]
    % циліндр
    \draw[thick] (0,0) ellipse (0.9 and 0.26);
    \draw[thick] (-0.9,0) -- (-0.9,2.6);
    \draw[thick] (0.9,0) -- (0.9,2.6);
    \draw[thick] (0,2.6) ellipse (0.9 and 0.26);
    % паралелепіпед
    \begin{scope}[shift={(3.2,0)}]
        \coordinate (A) at (0,0);
        \coordinate (B) at (2,0);
        \coordinate (C) at (2.7,0.7);
        \coordinate (D) at (0.7,0.7);
        \draw[thick] (A) -- (B) -- (C) -- (D) -- cycle;
        \draw[thick] (A) -- (0,2.6);
        \draw[thick] (B) -- (2,2.6);
        \draw[thick] (C) -- (2.7,3.3);
        \draw[thick, dashed] (D) -- (0.7,3.3);
        \draw[thick] (0,2.6) -- (2,2.6) -- (2.7,3.3) -- (0.7,3.3) -- cycle;
    \end{scope}
\end{tikzpicture}
\end{minipage}
\nmtAnswerBox
\taskBlock{Об'єм призми (призма і піраміда) --- сесія 4.06, завдання №21}
\noindent
\begin{minipage}[c]{0.58\textwidth}
\zadnum На рисунку зображено правильну чотирикутну призму та правильну чотирикутну піраміду. Сторона основи піраміди вдвічі менша від діагоналі основи призми. Усі ребра піраміди рівні й дорівнюють $3\sqrt{2}$~см. Знайдіть об'єм (у см$^3$) призми, якщо сума довжин висот призми та піраміди дорівнює $13$~см.
\par\vspace{0.4cm}
\nmtAnswerBox
\end{minipage}%
\hfill
\begin{minipage}[c]{0.40\textwidth}
\centering
\begin{tikzpicture}[scale=0.65, thick, line join=round]
    % --- Правильна чотирикутна призма ---
    \coordinate (A) at (0,0);
    \coordinate (B) at (1.8,0);
    \coordinate (C) at (2.6, 0.7);
    \coordinate (D) at (0.8, 0.7);

    \coordinate (A1) at (0,3.8);
    \coordinate (B1) at (1.8,3.8);
    \coordinate (C1) at (2.6, 4.5);
    \coordinate (D1) at (0.8, 4.5);

    % Видимі грані
    \draw (A) -- (B) -- (C) -- (C1) -- (D1) -- (A1) -- cycle;
    \draw (B) -- (B1) -- (C1);
    \draw (A1) -- (B1);

    % Невидимі ребра
    \draw[dashed] (A) -- (D) -- (C);
    \draw[dashed] (D) -- (D1);

    % --- Правильна чотирикутна піраміда ---
    \begin{scope}[shift={(3.8,0)}]
        \coordinate (P_A) at (0,0);
        \coordinate (P_B) at (2.2,0);
        \coordinate (P_C) at (3.2, 0.9);
        \coordinate (P_D) at (1.0, 0.9);
        \coordinate (P_S) at (1.6, 4.0);

        % Видимі ребра
        \draw (P_A) -- (P_B) -- (P_C) -- (P_S) -- cycle;
        \draw (P_B) -- (P_S);

        % Невидимі ребра
        \draw[dashed] (P_A) -- (P_D) -- (P_C);
        \draw[dashed] (P_D) -- (P_S);
    \end{scope}
\end{tikzpicture}
\end{minipage}
\par\vspace{0.3cm}
\solution{Усі ребра піраміди $=3\sqrt2$, тож сторона її основи $a_p=3\sqrt2$, а діагональ основи $a_p\sqrt2=6$. Висота піраміди $h_2=\sqrt{(3\sqrt2)^2-3^2}=\sqrt{18-9}=3$. Сторона основи піраміди вдвічі менша від діагоналі основи призми, тож діагональ призми $=6\sqrt2$, а її сторона $a=\dfrac{6\sqrt2}{\sqrt2}=6$, площа основи $S=36$. Висота призми $h_1=13-3=10$, об'єм $V=36\cdot10=360$~см$^3$. Відповідь: \textbf{360}.}
\noindent
\begin{minipage}[c]{0.58\textwidth}
\zadnum На рисунку зображено правильну чотирикутну призму та правильну чотирикутну піраміду. Сторона основи піраміди вдвічі менша від діагоналі основи призми. Усі ребра піраміди рівні й дорівнюють $4\sqrt{2}$~см. Знайдіть об'єм (у см$^3$) призми, якщо сума довжин висот призми та піраміди дорівнює $13$~см.
\par\vspace{0.4cm}
\nmtAnswerBox
\end{minipage}%
\hfill
\begin{minipage}[c]{0.40\textwidth}
\centering
\begin{tikzpicture}[scale=0.65, thick, line join=round]
    \coordinate (A) at (0,0);
    \coordinate (B) at (1.8,0);
    \coordinate (C) at (2.6, 0.7);
    \coordinate (D) at (0.8, 0.7);

    \coordinate (A1) at (0,3.8);
    \coordinate (B1) at (1.8,3.8);
    \coordinate (C1) at (2.6, 4.5);
    \coordinate (D1) at (0.8, 4.5);

    \draw (A) -- (B) -- (C) -- (C1) -- (D1) -- (A1) -- cycle;
    \draw (B) -- (B1) -- (C1);
    \draw (A1) -- (B1);

    \draw[dashed] (A) -- (D) -- (C);
    \draw[dashed] (D) -- (D1);

    \begin{scope}[shift={(3.8,0)}]
        \coordinate (P_A) at (0,0);
        \coordinate (P_B) at (2.2,0);
        \coordinate (P_C) at (3.2, 0.9);
        \coordinate (P_D) at (1.0, 0.9);
        \coordinate (P_S) at (1.6, 4.0);

        \draw (P_A) -- (P_B) -- (P_C) -- (P_S) -- cycle;
        \draw (P_B) -- (P_S);

        \draw[dashed] (P_A) -- (P_D) -- (P_C);
        \draw[dashed] (P_D) -- (P_S);
    \end{scope}
\end{tikzpicture}
\end{minipage}
\noindent
\begin{minipage}[c]{0.58\textwidth}
\zadnum На рисунку зображено правильну чотирикутну призму та правильну чотирикутну піраміду. Сторона основи піраміди дорівнює стороні основи призми. Усі ребра піраміди рівні й дорівнюють $2\sqrt{2}$~см. Знайдіть об'єм (у см$^3$) призми, якщо сума довжин висот призми та піраміди дорівнює $10$~см.
\par\vspace{0.4cm}
\nmtAnswerBox
\end{minipage}%
\hfill
\begin{minipage}[c]{0.40\textwidth}
\centering
\begin{tikzpicture}[scale=0.65, thick, line join=round]
    \coordinate (A) at (0,0);
    \coordinate (B) at (1.8,0);
    \coordinate (C) at (2.6, 0.7);
    \coordinate (D) at (0.8, 0.7);

    \coordinate (A1) at (0,3.8);
    \coordinate (B1) at (1.8,3.8);
    \coordinate (C1) at (2.6, 4.5);
    \coordinate (D1) at (0.8, 4.5);

    \draw (A) -- (B) -- (C) -- (C1) -- (D1) -- (A1) -- cycle;
    \draw (B) -- (B1) -- (C1);
    \draw (A1) -- (B1);

    \draw[dashed] (A) -- (D) -- (C);
    \draw[dashed] (D) -- (D1);

    \begin{scope}[shift={(3.8,0)}]
        \coordinate (P_A) at (0,0);
        \coordinate (P_B) at (2.2,0);
        \coordinate (P_C) at (3.2, 0.9);
        \coordinate (P_D) at (1.0, 0.9);
        \coordinate (P_S) at (1.6, 4.0);

        \draw (P_A) -- (P_B) -- (P_C) -- (P_S) -- cycle;
        \draw (P_B) -- (P_S);

        \draw[dashed] (P_A) -- (P_D) -- (P_C);
        \draw[dashed] (P_D) -- (P_S);
    \end{scope}
\end{tikzpicture}
\end{minipage}
\taskBlock{Призма з ромбом в основі та конус --- сесія 5.06, завдання №21}
\noindent
\begin{minipage}[c]{0.62\textwidth}
\zadnum На рисунку зображено пряму чотирикутну призму та конус. В основі призми лежить ромб, площа якого дорівнює $216$ см$^2$, а сторона~--- $15$ см. Радіус основи конуса дорівнює половині \textit{висоти} ромба. Висота призми дорівнює твірній конуса. Визначте площу \textit{повної} поверхні призми (у см$^2$), якщо площа бічної поверхні конуса дорівнює $144\pi$ см$^2$.
\end{minipage}%
\hfill
\begin{minipage}[c]{0.36\textwidth}
\centering
\begin{tikzpicture}[scale=0.7, thick]
    \coordinate (A) at (0,0);
    \coordinate (B) at (1.4,0.5);
    \coordinate (C) at (2.6,0);
    \coordinate (D) at (1.2,-0.5);
    \draw[thick] (A) -- (B) -- (C) -- (D) -- cycle;
    \draw[thick] (A) -- (0,2.6);
    \draw[thick] (B) -- (1.4,3.1);
    \draw[thick] (C) -- (2.6,2.6);
    \draw[thick, dashed] (D) -- (1.2,2.1);
    \draw[thick] (0,2.6) -- (1.4,3.1) -- (2.6,2.6) -- (1.2,2.1) -- cycle;
    \begin{scope}[shift={(4.0,0)}]
        \draw[thick] (0,0) ellipse (0.7 and 0.2);
        \draw[thick] (-0.7,0) -- (0,2.6) -- (0.7,0);
    \end{scope}
\end{tikzpicture}
\end{minipage}
\nmtAnswerBox
\solution{Висота ромба: $S=a\cdot h\Rightarrow h=\dfrac{216}{15}=14{,}4$. Радіус конуса $r=\dfrac{h}{2}=7{,}2$. З $\pi r l=144\pi$ маємо $l=\dfrac{144}{7{,}2}=20$~--- це й висота призми $H$. Повна поверхня призми $=2S+P\cdot H=2\cdot216+(4\cdot15)\cdot20=432+1200=1632$ см$^2$. Відповідь: \textbf{$1632$}.}
\noindent
\begin{minipage}[c]{0.62\textwidth}
\zadnum На рисунку зображено пряму чотирикутну призму та конус. В основі призми лежить ромб, площа якого дорівнює $192$ см$^2$, а сторона~--- $12$ см. Радіус основи конуса дорівнює половині \textit{висоти} ромба. Висота призми дорівнює твірній конуса. Визначте площу \textit{повної} поверхні призми (у см$^2$), якщо площа бічної поверхні конуса дорівнює $160\pi$ см$^2$.
\end{minipage}%
\hfill
\begin{minipage}[c]{0.36\textwidth}
\centering
\begin{tikzpicture}[scale=0.7, thick]
    \coordinate (A) at (0,0);
    \coordinate (B) at (1.4,0.5);
    \coordinate (C) at (2.6,0);
    \coordinate (D) at (1.2,-0.5);
    \draw[thick] (A) -- (B) -- (C) -- (D) -- cycle;
    \draw[thick] (A) -- (0,2.6);
    \draw[thick] (B) -- (1.4,3.1);
    \draw[thick] (C) -- (2.6,2.6);
    \draw[thick, dashed] (D) -- (1.2,2.1);
    \draw[thick] (0,2.6) -- (1.4,3.1) -- (2.6,2.6) -- (1.2,2.1) -- cycle;
    \begin{scope}[shift={(4.0,0)}]
        \draw[thick] (0,0) ellipse (0.7 and 0.2);
        \draw[thick] (-0.7,0) -- (0,2.6) -- (0.7,0);
    \end{scope}
\end{tikzpicture}
\end{minipage}
\nmtAnswerBox
\noindent
\begin{minipage}[c]{0.62\textwidth}
\zadnum На рисунку зображено пряму чотирикутну призму та конус. В основі призми лежить ромб, площа якого дорівнює $120$ см$^2$, а сторона~--- $10$ см. Радіус основи конуса дорівнює половині \textit{висоти} ромба. Висота призми дорівнює твірній конуса. Визначте \textit{об'єм} призми (у см$^3$), якщо площа бічної поверхні конуса дорівнює $90\pi$ см$^2$.
\end{minipage}%
\hfill
\begin{minipage}[c]{0.36\textwidth}
\centering
\begin{tikzpicture}[scale=0.7, thick]
    \coordinate (A) at (0,0);
    \coordinate (B) at (1.4,0.5);
    \coordinate (C) at (2.6,0);
    \coordinate (D) at (1.2,-0.5);
    \draw[thick] (A) -- (B) -- (C) -- (D) -- cycle;
    \draw[thick] (A) -- (0,2.6);
    \draw[thick] (B) -- (1.4,3.1);
    \draw[thick] (C) -- (2.6,2.6);
    \draw[thick, dashed] (D) -- (1.2,2.1);
    \draw[thick] (0,2.6) -- (1.4,3.1) -- (2.6,2.6) -- (1.2,2.1) -- cycle;
    \begin{scope}[shift={(4.0,0)}]
        \draw[thick] (0,0) ellipse (0.7 and 0.2);
        \draw[thick] (-0.7,0) -- (0,2.6) -- (0.7,0);
    \end{scope}
\end{tikzpicture}
\end{minipage}
\nmtAnswerBox

\chapterTitle{РОЗДІЛ 6. КОМБІНАТОРИКА, ЙМОВІРНІСТЬ, СТАТИСТИКА}
\typeTitle{Завдання з вибором однієї відповіді (А--Д)}
\taskBlock{Комбінаторика за круговою діаграмою --- сесія 23.05, завдання №12}

\noindent
\begin{minipage}[c]{0.55\textwidth}
\zadnum На діаграмі подано інформацію про назви та кількість шкільних гуртків. Макар планує відвідувати два гуртки: один зі спортивних, інший~--- з решти. Скільки різних варіантів вибору має Макар? \nmtyear{2026}
\end{minipage}%
\hfill
\begin{minipage}[c]{0.43\textwidth}
\centering
\begin{tikzpicture}[scale=0.9, thick]
    \draw (0,0) circle (1.7);
    \draw (0,0) -- (1.7,0);
    \draw (0,0) -- ({1.7*cos(77.14)},{1.7*sin(77.14)});
    \draw (0,0) -- (-1.7,0);
    \draw (0,0) -- ({1.7*cos(231.43)},{1.7*sin(231.43)});
    \draw (0,0) -- ({1.7*cos(257.14)},{1.7*sin(257.14)});
    \node at (38.57:1.0) {\small $3$};
    \node at (141.43:1.0) {\small $4$};
    \node at (193:0.9) {\small $1$};
    \node at (244:1.0) {\small $2$};
    \node at (315:1.0) {\small $4$};
    \node at (38.57:2.1) {\scriptsize художні};
    \node at (130:2.2) {\scriptsize танцювальні};
    \node at (200:2.1) {\scriptsize музичні};
    \node[below] at (244:2.0) {\scriptsize STEM};
    \node at (320:2.2) {\scriptsize спортивні};
\end{tikzpicture}
\end{minipage}
\par\vspace{0.2cm}
\answerTable{$36$}{$40$}{$24$}{$14$}{$13$}
\solution{Спортивних гуртків~--- $4$, решти $3+4+1+2=10$. За правилом добутку кількість варіантів $4\cdot 10=40$. Відповідь: \textbf{Б}.}

\vspace{0.3cm}

\noindent
\begin{minipage}[c]{0.55\textwidth}
\zadnum На діаграмі подано інформацію про назви та кількість шкільних гуртків. Соломія планує відвідувати два гуртки: один зі спортивних, інший~--- з решти. Скільки різних варіантів вибору має Соломія?
\end{minipage}%
\hfill
\begin{minipage}[c]{0.43\textwidth}
\centering
\begin{tikzpicture}[scale=0.9, thick]
    \draw (0,0) circle (1.7);
    \draw (0,0) -- (1.7,0);
    \draw (0,0) -- ({1.7*cos(48)},{1.7*sin(48)});
    \draw (0,0) -- ({1.7*cos(144)},{1.7*sin(144)});
    \draw (0,0) -- ({1.7*cos(216)},{1.7*sin(216)});
    \draw (0,0) -- ({1.7*cos(240)},{1.7*sin(240)});
    \node at (24:1.0) {\small $2$};
    \node at (96:1.0) {\small $4$};
    \node at (180:1.0) {\small $3$};
    \node at (228:1.0) {\small $1$};
    \node at (300:1.0) {\small $5$};
    \node at (24:2.2) {\scriptsize художні};
    \node at (96:2.1) {\scriptsize танцювальні};
    \node at (180:2.2) {\scriptsize музичні};
    \node[below] at (228:1.9) {\scriptsize STEM};
    \node at (300:2.2) {\scriptsize спортивні};
\end{tikzpicture}
\end{minipage}
\par\vspace{0.2cm}
\answerTable{$45$}{$15$}{$50$}{$25$}{$10$}

\vspace{0.3cm}

\noindent
\begin{minipage}[c]{0.55\textwidth}
\zadnum На діаграмі подано інформацію про назви та кількість шкільних гуртків. Остап планує відвідувати два гуртки: один із художніх, а інший~--- зі спортивних. Скільки різних варіантів вибору має Остап?
\end{minipage}%
\hfill
\begin{minipage}[c]{0.43\textwidth}
\centering
\begin{tikzpicture}[scale=0.9, thick]
    \draw (0,0) circle (1.7);
    \draw (0,0) -- (1.7,0);
    \draw (0,0) -- ({1.7*cos(120)},{1.7*sin(120)});
    \draw (0,0) -- ({1.7*cos(192)},{1.7*sin(192)});
    \draw (0,0) -- ({1.7*cos(240)},{1.7*sin(240)});
    \draw (0,0) -- ({1.7*cos(288)},{1.7*sin(288)});
    \node at (60:1.0) {\small $5$};
    \node at (156:1.0) {\small $3$};
    \node at (216:1.0) {\small $2$};
    \node at (264:1.0) {\small $2$};
    \node at (324:1.0) {\small $3$};
    \node at (60:2.2) {\scriptsize художні};
    \node at (156:2.3) {\scriptsize танцювальні};
    \node at (216:2.2) {\scriptsize музичні};
    \node[below] at (264:1.9) {\scriptsize STEM};
    \node at (324:2.3) {\scriptsize спортивні};
\end{tikzpicture}
\end{minipage}
\par\vspace{0.2cm}
\answerTable{$12$}{$30$}{$15$}{$36$}{$24$}

\vspace{0.3cm}

% ---------------------------------------------------------------------
\taskBlock{Найбільша різниця за графіками --- сесія 23.05, завдання №14}

\noindent
\begin{minipage}[c]{0.42\textwidth}
\zadnum На рисунку зображено графіки залежності рівня шуму (дБ) від частоти (Гц) для двох джерел. За якої частоти різниця між рівнями шуму джерел є найбільшою? \nmtyear{2026}
\end{minipage}%
\hfill
\begin{minipage}[c]{0.55\textwidth}
\centering
\begin{tikzpicture}
\begin{axis}[
    width=9cm, height=6cm,
    xmode=log, log basis x=10,
    xmin=15, xmax=7000,
    ymin=0, ymax=85,
    ytick={0,20,40,60,80},
    xtick={20,100,200,300,500,1000,2000,5000},
    xticklabels={20,100,200,300,500,1000,2000,5000},
    xlabel={\small частота (Гц)},
    ylabel={\small рівень шуму (дБ)},
    xmajorgrids=true, ymajorgrids=true,
    major grid style={dashed, gray!40},
    axis x line*=bottom, axis y line*=left,
    tick label style={font=\scriptsize},
    log ticks with fixed point,
]
    \addplot[mark=*, thick, black] coordinates {
        (20,60) (100,53) (200,62) (300,60) (500,57) (1000,62) (2000,65) (5000,67)
    };
    \addplot[mark=*, thick, black] coordinates {
        (20,40) (100,30) (200,20) (300,12) (500,16) (1000,20) (2000,25) (5000,40)
    };
\end{axis}
\end{tikzpicture}
\end{minipage}
\par\vspace{0.2cm}
\answerTable{$1000$}{$5000$}{$200$}{$300$}{$2000$}
\solution{Обчислимо різницю рівнів за кожної частоти: при $300$~Гц вона дорівнює $60-12=48$~дБ, що більше за різниці при інших частотах (напр., $200$: $42$; $1000$: $42$). Найбільша різниця~--- за $300$~Гц. Відповідь: \textbf{Г}.}

\vspace{0.3cm}

\noindent
\begin{minipage}[c]{0.42\textwidth}
\zadnum На рисунку зображено графіки залежності рівня сигналу (дБ) від частоти (Гц) для двох антен. За якої частоти різниця між рівнями сигналу антен є найбільшою?
\end{minipage}%
\hfill
\begin{minipage}[c]{0.55\textwidth}
\centering
\begin{tikzpicture}
\begin{axis}[
    width=9cm, height=6cm,
    xmode=log, log basis x=10,
    xmin=15, xmax=7000,
    ymin=0, ymax=85,
    ytick={0,20,40,60,80},
    xtick={20,100,200,300,500,1000,2000,5000},
    xticklabels={20,100,200,300,500,1000,2000,5000},
    xlabel={\small частота (Гц)},
    ylabel={\small рівень сигналу (дБ)},
    xmajorgrids=true, ymajorgrids=true,
    major grid style={dashed, gray!40},
    axis x line*=bottom, axis y line*=left,
    tick label style={font=\scriptsize},
    log ticks with fixed point,
]
    \addplot[mark=*, thick, black] coordinates {
        (20,55) (100,50) (200,58) (300,60) (500,62) (1000,64) (2000,60) (5000,58)
    };
    \addplot[mark=*, thick, black] coordinates {
        (20,35) (100,28) (200,25) (300,22) (500,18) (1000,30) (2000,35) (5000,40)
    };
\end{axis}
\end{tikzpicture}
\end{minipage}
\par\vspace{0.2cm}
\answerTable{$200$}{$300$}{$500$}{$1000$}{$2000$}

\vspace{0.3cm}

\noindent
\begin{minipage}[c]{0.42\textwidth}
\zadnum На рисунку зображено графіки залежності рівня гучності (дБ) від частоти (Гц) для двох динаміків. За якої частоти різниця між рівнями гучності динаміків є найменшою?
\end{minipage}%
\hfill
\begin{minipage}[c]{0.55\textwidth}
\centering
\begin{tikzpicture}
\begin{axis}[
    width=9cm, height=6cm,
    xmode=log, log basis x=10,
    xmin=15, xmax=7000,
    ymin=0, ymax=85,
    ytick={0,20,40,60,80},
    xtick={20,100,200,300,500,1000,2000,5000},
    xticklabels={20,100,200,300,500,1000,2000,5000},
    xlabel={\small частота (Гц)},
    ylabel={\small рівень гучності (дБ)},
    xmajorgrids=true, ymajorgrids=true,
    major grid style={dashed, gray!40},
    axis x line*=bottom, axis y line*=left,
    tick label style={font=\scriptsize},
    log ticks with fixed point,
]
    \addplot[mark=*, thick, black] coordinates {
        (20,50) (100,55) (200,60) (300,58) (500,62) (1000,65) (2000,70) (5000,66)
    };
    \addplot[mark=*, thick, black] coordinates {
        (20,30) (100,32) (200,35) (300,38) (500,45) (1000,58) (2000,40) (5000,42)
    };
\end{axis}
\end{tikzpicture}
\end{minipage}
\par\vspace{0.2cm}
\answerTable{$300$}{$500$}{$1000$}{$2000$}{$5000$}
\taskBlock{Статистика: кругова діаграма --- сесія 30.05, завдання №4}

\noindent
\begin{minipage}[c]{0.60\textwidth}
\zadnum На круговій діаграмі зображено розподіл туристів за країнами походження, що відвідали певний музей за рік. Визначте відсоток туристів, які прибули з \textit{Канади}.
\end{minipage}%
\hfill
\begin{minipage}[c]{0.38\textwidth}
\centering
\begin{tikzpicture}[scale=1.2, thick]
    \def\R{1.7}
    % Європа 50% = 180°, починаємо з 90° (верх)
    \filldraw[fill=mainGreen!50, draw=black] (0,0) -- (90:\R) arc (90:-90:\R) -- cycle;
    \node at (0:0.9) {\scriptsize $50\,\%$};
    \node[right] at (0:\R+0.1) {\scriptsize Європа};
    % США 11,2% = 40,32°
    \filldraw[fill=yearOrange!50, draw=black] (0,0) -- (-90:\R) arc (-90:-130.32:\R) -- cycle;
    \node at (-110:1.0) { \scriptsize $11{,}2\,\%$};
    \node[below] at (-110:\R+0.15) {\scriptsize США};
    % Японія 4,8% = 17,28°
    \filldraw[fill=mainGreen!25, draw=black] (0,0) -- (-130.32:\R) arc (-130.32:-147.6:\R) -- cycle;
    \node[left] at (-139:\R+0.05) {\scriptsize $4{,}8\,\%$ Японія};
    % Австралія 3,5% = 12,6°
    \filldraw[fill=yearOrange!25, draw=black] (0,0) -- (-147.6:\R) arc (-147.6:-160.2:\R) -- cycle;
    \node[left] at (-154:\R+0.05) {\scriptsize $3{,}5\,\%$ Австралія};
    % Канада ? = 109,8°
    \filldraw[fill=gray!30, draw=black] (0,0) -- (-160.2:\R) arc (-160.2:-270:\R) -- cycle;
    \node at (145:1.0) {\scriptsize $?\,\%$};
    \node[left] at (145:\R+0.05) {\scriptsize Канада};
\end{tikzpicture}
\end{minipage}
\par\vspace{0.2cm}
\answerTable{$25\,\%$}{$28{,}5\,\%$}{$30{,}5\,\%$}{$33\,\%$}{$35\,\%$}
\solution{Сума всіх відсотків дорівнює $100\,\%$. Канада: $100-(50+11{,}2+4{,}8+3{,}5)=100-69{,}5=30{,}5\,\%$. Відповідь: \textbf{В} ($30{,}5\,\%$).}

\noindent
\begin{minipage}[c]{0.60\textwidth}
\zadnum На круговій діаграмі зображено розподіл книжок бібліотеки за жанрами. Визначте відсоток книжок \textit{дитячої} літератури.
\end{minipage}%
\hfill
\begin{minipage}[c]{0.38\textwidth}
\centering
\begin{tikzpicture}[scale=1.2, thick]
    \def\R{1.7}
    % Проза 45% = 162°, починаємо з 90° (верх)
    \filldraw[fill=mainGreen!50, draw=black] (0,0) -- (90:\R) arc (90:-72:\R) -- cycle;
    \node at (9:0.9) {\scriptsize $45\,\%$};
    \node[right] at (9:\R+0.1) {\scriptsize Проза};
    % Поезія 12,4% = 44,64°
    \filldraw[fill=yearOrange!50, draw=black] (0,0) -- (-72:\R) arc (-72:-116.64:\R) -- cycle;
    \node at (-94:1.0) {\scriptsize $12{,}4\,\%$};
    \node[below] at (-94:\R+0.15) {\scriptsize Поезія};
    % Фантастика 6,6% = 23,76°
    \filldraw[fill=mainGreen!25, draw=black] (0,0) -- (-116.64:\R) arc (-116.64:-140.4:\R) -- cycle;
    \node[left] at (-128:\R+0.05) {\scriptsize $6{,}6\,\%$ Фантастика};
    % Наукова 8,5% = 30,6°
    \filldraw[fill=yearOrange!25, draw=black] (0,0) -- (-140.4:\R) arc (-140.4:-171:\R) -- cycle;
    \node[left] at (-156:\R+0.05) {\scriptsize $8{,}5\,\%$ Наукова};
    % Дитяча ? = 99°
    \filldraw[fill=gray!30, draw=black] (0,0) -- (-171:\R) arc (-171:-270:\R) -- cycle;
    \node at (140:1.0) {\scriptsize $?\,\%$};
    \node[left] at (140:\R+0.05) {\scriptsize Дитяча};
\end{tikzpicture}
\end{minipage}
\par\vspace{0.2cm}
\answerTable{$24{,}5\,\%$}{$26\,\%$}{$27{,}5\,\%$}{$29\,\%$}{$32{,}5\,\%$}

\noindent
\begin{minipage}[c]{0.60\textwidth}
\zadnum На круговій діаграмі зображено розподіл проданого в кафе морозива за смаками. Визначте \textit{градусну міру} центрального кута сектора, що відповідає смаку \textit{карамелі}.
\end{minipage}%
\hfill
\begin{minipage}[c]{0.38\textwidth}
\centering
\begin{tikzpicture}[scale=1.2, thick]
    \def\R{1.7}
    % Ваніль 40% = 144°, починаємо з 90° (верх)
    \filldraw[fill=mainGreen!50, draw=black] (0,0) -- (90:\R) arc (90:-54:\R) -- cycle;
    \node at (18:0.9) {\scriptsize $40\,\%$};
    \node[right] at (18:\R+0.1) {\scriptsize Ваніль};
    % Шоколад 22,3% = 80,28°
    \filldraw[fill=yearOrange!50, draw=black] (0,0) -- (-54:\R) arc (-54:-134.28:\R) -- cycle;
    \node at (-94:1.0) {\scriptsize $22{,}3\,\%$};
    \node[below] at (-94:\R+0.15) {\scriptsize Шоколад};
    % Полуниця 9,2% = 33,12°
    \filldraw[fill=mainGreen!25, draw=black] (0,0) -- (-134.28:\R) arc (-134.28:-167.4:\R) -- cycle;
    \node[left] at (-151:\R+0.05) {\scriptsize $9{,}2\,\%$ Полуниця};
    % Фісташка 5,5% = 19,8°
    \filldraw[fill=yearOrange!25, draw=black] (0,0) -- (-167.4:\R) arc (-167.4:-187.2:\R) -- cycle;
    \node[left] at (-177:\R+0.05) {\scriptsize $5{,}5\,\%$ Фісташка};
    % Карамель ? = 82,8°
    \filldraw[fill=gray!30, draw=black] (0,0) -- (-187.2:\R) arc (-187.2:-270:\R) -- cycle;
    \node at (131:1.0) {\scriptsize $?\,\%$};
    \node[left] at (131:\R+0.05) {\scriptsize Карамель};
\end{tikzpicture}
\end{minipage}
\par\vspace{0.2cm}
\answerTable{$75{,}6^\circ$}{$80^\circ$}{$82{,}8^\circ$}{$90^\circ$}{$97{,}2^\circ$}
\taskBlock{Комбінаторика: правило добутку --- сесія 30.05, завдання №7}

\zadtask{У філателістичному магазині є $30$ різних марок з історичної серії та $40$ різних марок із серії <<краєвиди>>. Іван хоче придбати по одній марці з кожної серії. Скільки всього є варіантів такого вибору?}
\answerTable{$70$}{$140$}{$600$}{$1200$}{$1400$}
\solution{За правилом добутку кількість варіантів — добуток кількостей у кожній серії: $30\cdot40=1200$. Відповідь: \textbf{Г} ($1200$).}

\zadtask{У чайній крамниці є $25$ різних сортів чаю та $16$ різних видів тістечок. Оксана хоче придбати один сорт чаю та одне тістечко. Скільки всього є варіантів такого вибору?}
\answerTable{$41$}{$82$}{$200$}{$400$}{$800$}

\zadtask{У майстерні є $12$ різних моделей рамок та $35$ різних постерів. Андрій хоче придбати лише \textit{один} предмет~--- \textit{або} одну рамку, \textit{або} один постер. Скільки всього є варіантів такого вибору?}
\answerTable{$47$}{$94$}{$210$}{$420$}{$840$}
\taskBlock{Комбінаторика, кількість способів вибору --- сесія 1.06, завдання №6}

\zadtask{У магазині є $7$ видів подарункових наборів. Скільки всього існує варіантів вибору з них трьох \textit{різних} наборів?}
\answerTable{$21$}{$35$}{$210$}{$343$}{$840$}
\solution{Порядок не важливий, тож це сполучення: $C_7^3=\dfrac{7!}{3!\,4!}=\dfrac{7\cdot6\cdot5}{6}=35$. Відповідь: \textbf{Б}.}

\zadtask{У кондитерській є $8$ видів тістечок. Скільки всього існує варіантів вибору з них трьох \textit{різних} тістечок?}
\answerTable{$24$}{$28$}{$56$}{$336$}{$512$}

\zadtask{У кав'ярні є $6$ сортів морозива. Скількома способами можна скласти стаканчик із трьох \textit{різних} сортів, якщо порядок викладання сортів \textit{має значення}?}
\answerTable{$18$}{$20$}{$120$}{$216$}{$720$}

% =====================================================================
\taskBlock{Ймовірність за діаграмою опитування --- сесія 1.06, завдання №7}

\noindent
\begin{minipage}[c]{0.58\textwidth}
\zadnum На діаграмі відображено результати анкетування $100$ осіб, які навчаються в університеті, щодо відкриття неподалік піцерії. Кожна особа обрала лише одну з трьох відповідей: так, ні, байдуже. Яка ймовірність того, що навмання обрана анкетована особа \textit{підтримала} відкриття піцерії?
\end{minipage}%
\hfill
\begin{minipage}[c]{0.40\textwidth}
\centering
\begin{tikzpicture}[scale=0.75, thick]
    \draw[->] (0,0) -- (0,3.9);
    \draw[->] (0,0) -- (3.7,0);
    \filldraw[fill=mainGreen!50, draw=mainGreen!80!black] (0.3,0) rectangle (1.0,3.2);
    \filldraw[fill=yearOrange!50, draw=yearOrange!80!black] (1.4,0) rectangle (2.1,0.2);
    \filldraw[fill=gray!40, draw=gray!70] (2.5,0) rectangle (3.2,0.6);
    \node[above, font=\scriptsize] at (0.65,3.2) {$80$};
    \node[above, font=\scriptsize] at (1.75,0.2) {$5$};
    \node[above, font=\scriptsize] at (2.85,0.6) {$15$};
    \node[below, font=\scriptsize] at (0.65,0) {так};
    \node[below, font=\scriptsize] at (1.75,0) {ні};
    \node[below, font=\scriptsize] at (2.85,0) {байд.};
\end{tikzpicture}
\end{minipage}
\par\vspace{0.2cm}
\answerTable{$0{,}15$}{$0{,}05$}{$0{,}2$}{$0{,}85$}{$0{,}8$}
\solution{Відкриття піцерії підтримали (відповідь «так») $80$ зі $100$ осіб, тож шукана ймовірність $P=\dfrac{80}{100}=0{,}8$. Відповідь: \textbf{Д}.}

\noindent
\begin{minipage}[c]{0.58\textwidth}
\zadnum На діаграмі відображено результати анкетування $200$ мешканців мікрорайону щодо облаштування нової велодоріжки. Кожна особа обрала лише одну з трьох відповідей: так, ні, байдуже. Яка ймовірність того, що навмання обрана анкетована особа відповіла \textit{«ні»}?
\end{minipage}%
\hfill
\begin{minipage}[c]{0.40\textwidth}
\centering
\begin{tikzpicture}[scale=0.75, thick]
    \draw[->] (0,0) -- (0,3.9);
    \draw[->] (0,0) -- (3.7,0);
    \filldraw[fill=mainGreen!50, draw=mainGreen!80!black] (0.3,0) rectangle (1.0,3.25);
    \filldraw[fill=yearOrange!50, draw=yearOrange!80!black] (1.4,0) rectangle (2.1,1.25);
    \filldraw[fill=gray!40, draw=gray!70] (2.5,0) rectangle (3.2,0.5);
    \node[above, font=\scriptsize] at (0.65,3.25) {$130$};
    \node[above, font=\scriptsize] at (1.75,1.25) {$50$};
    \node[above, font=\scriptsize] at (2.85,0.5) {$20$};
    \node[below, font=\scriptsize] at (0.65,0) {так};
    \node[below, font=\scriptsize] at (1.75,0) {ні};
    \node[below, font=\scriptsize] at (2.85,0) {байд.};
\end{tikzpicture}
\end{minipage}
\par\vspace{0.2cm}
\answerTable{$0{,}1$}{$0{,}25$}{$0{,}5$}{$0{,}65$}{$0{,}75$}

\noindent
\begin{minipage}[c]{0.58\textwidth}
\zadnum На діаграмі відображено результати анкетування $50$ працівників офісу щодо відкриття поруч кав'ярні. Кожна особа обрала лише одну з трьох відповідей: так, ні, байдуже. Яка ймовірність того, що навмання обрана анкетована особа \textit{не підтримала} відкриття кав'ярні (тобто відповіла «ні» або «байдуже»)?
\end{minipage}%
\hfill
\begin{minipage}[c]{0.40\textwidth}
\centering
\begin{tikzpicture}[scale=0.75, thick]
    \draw[->] (0,0) -- (0,3.9);
    \draw[->] (0,0) -- (3.7,0);
    \filldraw[fill=mainGreen!50, draw=mainGreen!80!black] (0.3,0) rectangle (1.0,2.6);
    \filldraw[fill=yearOrange!50, draw=yearOrange!80!black] (1.4,0) rectangle (2.1,3.25);
    \filldraw[fill=gray!40, draw=gray!70] (2.5,0) rectangle (3.2,0.65);
    \node[above, font=\scriptsize] at (0.65,2.6) {$20$};
    \node[above, font=\scriptsize] at (1.75,3.25) {$25$};
    \node[above, font=\scriptsize] at (2.85,0.65) {$5$};
    \node[below, font=\scriptsize] at (0.65,0) {так};
    \node[below, font=\scriptsize] at (1.75,0) {ні};
    \node[below, font=\scriptsize] at (2.85,0) {байд.};
\end{tikzpicture}
\end{minipage}
\par\vspace{0.2cm}
\answerTable{$0{,}1$}{$0{,}5$}{$0{,}6$}{$0{,}7$}{$0{,}9$}
\taskBlock{Читання стовпчастої діаграми --- сесія 2.06, завдання №6}

\zadtask{На діаграмі відображено розподіл кількості працівників фірми за віком. За даними діаграми визначте різницю між кількістю працівників віком від $20$ до $29$ років та кількістю працівників, яким від $40$ до $49$ років.

\vspace{0.2cm}
\begin{center}
\begin{tikzpicture}[scale=0.9, thick]
    % осі
    \draw[->] (0,0) -- (0,4.8) node[above, font=\scriptsize] {Кількість};
    \draw[->] (0,0) -- (7.4,0) node[right, font=\scriptsize] {Вік};
    % горизонтальні лінії сітки + позначки через 2 (0..44)
    \foreach \v in {2,4,6,8,10,12,14,16,18,20,22,24,26,28,30,32,34,36,38,40,42,44} {
        \pgfmathsetmacro\y{\v/10}
        \draw[gray!25, very thin] (0,\y) -- (6.9,\y);
    }
    % підписи кратні 10 жирніше + поділки
    \foreach \v in {0,10,20,30,40} {
        \pgfmathsetmacro\y{\v/10}
        \draw[gray!55] (-0.08,\y) -- (6.9,\y);
        \node[left, font=\scriptsize] at (-0.05,\y) {\v};
    }
    % проміжні поділки через 2 на осі Y (риски)
    \foreach \v in {2,4,6,8,12,14,16,18,22,24,26,28,32,34,36,38,42,44} {
        \pgfmathsetmacro\y{\v/10}
        \draw (-0.06,\y) -- (0.06,\y);
    }
    % бари: 36,40,24,16,4 (масштаб /10)
    \foreach \i/\h/\lab in {0/3.6/{20--29}, 1/4.0/{30--39}, 2/2.4/{40--49}, 3/1.6/{50--59}, 4/0.4/{60--70}} {
        \filldraw[fill=cyan!45, draw=cyan!60!black] ({0.6+\i*1.3},0) rectangle ({1.4+\i*1.3},\h);
        \node[below, font=\scriptsize] at ({1.0+\i*1.3},0) {\lab};
    }
\end{tikzpicture}
\end{center}
\par\vspace{0.2cm}}
\answerTable{$6$}{$10$}{$12$}{$16$}{$20$}
\solution{За діаграмою у віці $20$--$29$ років -- $36$ працівників, у віці $40$--$49$ -- $24$. Різниця $36-24=12$. Відповідь: \textbf{В}.}

\zadtask{На діаграмі відображено розподіл кількості працівників підприємства за віком. За даними діаграми визначте різницю між кількістю працівників віком від $20$ до $29$ років та кількістю працівників, яким від $40$ до $49$ років.

\vspace{0.2cm}
\begin{center}
\begin{tikzpicture}[scale=0.9, thick]
    % осі
    \draw[->] (0,0) -- (0,4.8) node[above, font=\scriptsize] {Кількість};
    \draw[->] (0,0) -- (7.4,0) node[right, font=\scriptsize] {Вік};
    % горизонтальні лінії сітки через 2 (0..44)
    \foreach \v in {2,4,6,8,10,12,14,16,18,20,22,24,26,28,30,32,34,36,38,40,42,44} {
        \pgfmathsetmacro\y{\v/10}
        \draw[gray!25, very thin] (0,\y) -- (6.9,\y);
    }
    % підписи кратні 10
    \foreach \v in {0,10,20,30,40} {
        \pgfmathsetmacro\y{\v/10}
        \draw[gray!55] (-0.08,\y) -- (6.9,\y);
        \node[left, font=\scriptsize] at (-0.05,\y) {\v};
    }
    % проміжні поділки через 2 на осі Y
    \foreach \v in {2,4,6,8,12,14,16,18,22,24,26,28,32,34,36,38,42,44} {
        \pgfmathsetmacro\y{\v/10}
        \draw (-0.06,\y) -- (0.06,\y);
    }
    % бари: 28,38,20,14,6
    \foreach \i/\h/\lab in {0/2.8/{20--29}, 1/3.8/{30--39}, 2/2.0/{40--49}, 3/1.4/{50--59}, 4/0.6/{60--70}} {
        \filldraw[fill=cyan!45, draw=cyan!60!black] ({0.6+\i*1.3},0) rectangle ({1.4+\i*1.3},\h);
        \node[below, font=\scriptsize] at ({1.0+\i*1.3},0) {\lab};
    }
\end{tikzpicture}
\end{center}
\par\vspace{0.2cm}}
\answerTable{$6$}{$8$}{$10$}{$12$}{$18$}

\zadtask{На діаграмі відображено розподіл кількості працівників компанії за віком. За даними діаграми визначте сумарну кількість працівників, яким \textit{не менше} $40$ років.

\vspace{0.2cm}
\begin{center}
\begin{tikzpicture}[scale=0.9, thick]
    % осі
    \draw[->] (0,0) -- (0,4.8) node[above, font=\scriptsize] {Кількість};
    \draw[->] (0,0) -- (7.4,0) node[right, font=\scriptsize] {Вік};
    % горизонтальні лінії сітки через 2 (0..44)
    \foreach \v in {2,4,6,8,10,12,14,16,18,20,22,24,26,28,30,32,34,36,38,40,42,44} {
        \pgfmathsetmacro\y{\v/10}
        \draw[gray!25, very thin] (0,\y) -- (6.9,\y);
    }
    % підписи кратні 10
    \foreach \v in {0,10,20,30,40} {
        \pgfmathsetmacro\y{\v/10}
        \draw[gray!55] (-0.08,\y) -- (6.9,\y);
        \node[left, font=\scriptsize] at (-0.05,\y) {\v};
    }
    % проміжні поділки через 2 на осі Y
    \foreach \v in {2,4,6,8,12,14,16,18,22,24,26,28,32,34,36,38,42,44} {
        \pgfmathsetmacro\y{\v/10}
        \draw (-0.06,\y) -- (0.06,\y);
    }
    % бари: 30,42,26,18,8
    \foreach \i/\h/\lab in {0/3.0/{20--29}, 1/4.2/{30--39}, 2/2.6/{40--49}, 3/1.8/{50--59}, 4/0.8/{60--70}} {
        \filldraw[fill=cyan!45, draw=cyan!60!black] ({0.6+\i*1.3},0) rectangle ({1.4+\i*1.3},\h);
        \node[below, font=\scriptsize] at ({1.0+\i*1.3},0) {\lab};
    }
\end{tikzpicture}
\end{center}
\par\vspace{0.2cm}}
\answerTable{$44$}{$48$}{$52$}{$60$}{$76$}

% =====================================================================
\taskBlock{Комбінаторика: вибір з обов'язковим елементом --- сесія 2.06, завдання №7}

\zadtask{У магазині продають $8$ різних сортів малини, серед яких є сорт <<Альба>>. Покупниця хоче придбати $2$ різні сорти, причому один із них обов'язково має бути <<Альба>>. Скільки в неї є варіантів такого вибору?}
\answerTable{$7$}{$8$}{$14$}{$16$}{$28$}
\solution{Сорт <<Альба>> обов'язково входить, тож залишається вибрати $1$ сорт із решти $8-1=7$. Отже $7$ варіантів. Відповідь: \textbf{А}.}

\zadtask{У розпліднику продають $10$ різних сортів яблунь, серед яких є сорт <<Голден>>. Садівник хоче придбати $2$ різні сорти, причому один із них обов'язково має бути <<Голден>>. Скільки в нього є варіантів такого вибору?}
\answerTable{$9$}{$10$}{$18$}{$20$}{$45$}

\zadtask{У квітковій крамниці продають $12$ різних сортів троянд, серед яких є сорт <<Аваланж>>. Покупець хоче придбати $2$ різні сорти, причому серед них \textit{не повинно бути} сорту <<Аваланж>>. Скільки в нього є варіантів такого вибору?}
\answerTable{$11$}{$22$}{$55$}{$66$}{$110$}

% =====================================================================
\taskBlock{Статистика: читання діаграми --- сесія 3.06, завдання №2}

\noindent
\begin{minipage}[c]{0.56\textwidth}
\zadnum На діаграмі відображено зміну маси (у кг) слона протягом шести місяців (із січня до червня) порівняно з груднем минулого року. Якою стала маса слона в червні, якщо в грудні вона становила $4{,}235$ т?
\end{minipage}%
\hfill
\begin{minipage}[c]{0.42\textwidth}
\centering
\begin{tikzpicture}[scale=1, thick]
    \draw[->] (0,-1.4) -- (0,2.4) node[above, font=\scriptsize] {кг};
    \draw[->] (0,0) -- (7.4,0);
    \foreach \v/\y in {-40/-0.8, 0/0, 40/0.8, 80/1.6} {
        \draw[gray!40] (0,\y) -- (7,\y);
        \node[left, font=\tiny] at (0,\y) {\v};
    }
    % точки зміни: січ -60, лют -40, бер 0, кві 20, тра 80, чер 60
    \foreach \i/\val in {0/-1.2, 1/-0.8, 2/0, 3/0.4, 4/1.6, 5/1.2} {
        \fill[cyan!70!blue] ({0.9+\i*1.1},\val) circle (2.4pt);
    }
    \node[below, font=\tiny, rotate=0] at (0.9,-1.5) {січ};
    \node[below, font=\tiny] at (6.4,-1.5) {чер};
\end{tikzpicture}
\end{minipage}
\par\vspace{0.2cm}
\answerTable{$4305$ кг}{$4295$ кг}{$4235$ кг}{$4275$ кг}{$4215$ кг}
\solution{Маса в грудні $4{,}235$ т $= 4235$ кг. За діаграмою зміна маси у червні порівняно з груднем дорівнює $+60$ кг. Отже маса в червні: $4235 + 60 = 4295$ кг. Відповідь: \textbf{Б}.}

\noindent
\begin{minipage}[c]{0.56\textwidth}
\zadnum На діаграмі відображено зміну маси (у кг) носорога протягом шести місяців (із січня до червня) порівняно з груднем минулого року. Якою стала маса носорога в червні, якщо в грудні вона становила $2{,}105$ т?
\end{minipage}%
\hfill
\begin{minipage}[c]{0.42\textwidth}
\centering
\begin{tikzpicture}[scale=1, thick]
    \draw[->] (0,-1.4) -- (0,2.4) node[above, font=\scriptsize] {кг};
    \draw[->] (0,0) -- (7.4,0);
    \foreach \v/\y in {-40/-0.8, 0/0, 40/0.8, 80/1.6} {
        \draw[gray!40] (0,\y) -- (7,\y);
        \node[left, font=\tiny] at (0,\y) {\v};
    }
    % точки зміни: січ -40, лют -20, бер 0, кві 40, тра 80, чер 40
    \foreach \i/\val in {0/-0.8, 1/-0.4, 2/0, 3/0.8, 4/1.6, 5/0.8} {
        \fill[cyan!70!blue] ({0.9+\i*1.1},\val) circle (2.4pt);
    }
    \node[below, font=\tiny, rotate=0] at (0.9,-1.5) {січ};
    \node[below, font=\tiny] at (6.4,-1.5) {чер};
\end{tikzpicture}
\end{minipage}
\par\vspace{0.2cm}
\answerTable{$2065$ кг}{$2105$ кг}{$2125$ кг}{$2145$ кг}{$2185$ кг}

\noindent
\begin{minipage}[c]{0.56\textwidth}
\zadnum На діаграмі відображено зміну маси (у кг) бегемота протягом шести місяців (із січня до червня) порівняно з груднем минулого року. Якою була \textit{найбільша} маса бегемота за цей період, якщо в грудні вона становила $1{,}760$ т?
\end{minipage}%
\hfill
\begin{minipage}[c]{0.42\textwidth}
\centering
\begin{tikzpicture}[scale=1, thick]
    \draw[->] (0,-1.4) -- (0,2.4) node[above, font=\scriptsize] {кг};
    \draw[->] (0,0) -- (7.4,0);
    \foreach \v/\y in {-40/-0.8, 0/0, 40/0.8, 80/1.6} {
        \draw[gray!40] (0,\y) -- (7,\y);
        \node[left, font=\tiny] at (0,\y) {\v};
    }
    % точки зміни: січ 40, лют 80, бер 20, кві 0, тра -40, чер -20
    \foreach \i/\val in {0/0.8, 1/1.6, 2/0.4, 3/0, 4/-0.8, 5/-0.4} {
        \fill[cyan!70!blue] ({0.9+\i*1.1},\val) circle (2.4pt);
    }
    \node[below, font=\tiny, rotate=0] at (0.9,-1.5) {січ};
    \node[below, font=\tiny] at (6.4,-1.5) {чер};
\end{tikzpicture}
\end{minipage}
\par\vspace{0.2cm}
\answerTable{$1840$ кг}{$1800$ кг}{$1760$ кг}{$1720$ кг}{$1780$ кг}
\taskBlock{Імовірність --- сесія 3.06, завдання №5}

\zadtask{У десятому класі за списком $24$ учні. Імовірність того, що вік навмання вибраного учня менший за $15$ років, дорівнює $\dfrac16$. Скільки всього в класі учнів, вік яких менший за $15$ років?}
\answerTable{$8$}{$4$}{$20$}{$3$}{$6$}
\solution{Кількість учнів молодших за $15$ років дорівнює добутку ймовірності на загальну кількість: $24\cdot\dfrac16 = 4$. Відповідь: \textbf{Б} ($4$).}

\zadtask{В одинадцятому класі за списком $30$ учнів. Імовірність того, що навмання вибраний учень носить окуляри, дорівнює $\dfrac15$. Скільки всього в класі учнів, які носять окуляри?}
\answerTable{$5$}{$10$}{$6$}{$24$}{$3$}

\zadtask{У дев'ятому класі за списком $28$ учнів, із них $4$ учні відвідують спортивну секцію. Яка ймовірність того, що навмання вибраний учень \textit{не} відвідує спортивну секцію?}
\answerTable{$\dfrac17$}{$\dfrac{6}{7}$}{$\dfrac{1}{4}$}{$\dfrac{3}{4}$}{$\dfrac{24}{25}$}
\taskBlock{Читання діаграми (ціни) --- сесія 4.06, завдання №2}
\noindent
\begin{minipage}[t]{0.52\textwidth}
\vspace{0pt}
\zadnum На діаграмі відображено інформацію про ціну пакета цукерок «Чудасія», «Каштанчики», «Шокос», «Смакс» і «Няммі» в деякому магазині. Якщо в Катрусі є лише $200$~грн, то із зазначених цукерок вона може придбати

\vspace{0.3cm}
\par\noindent\textbf{А}\ \ лише «Каштанчики» або «Шокос»
\par\noindent\textbf{Б}\ \ лише «Чудасія»
\par\noindent\textbf{В}\ \ лише «Чудасія», «Смакс» або «Няммі»
\par\noindent\textbf{Г}\ \ будь-який один пакет цукерок
\par\noindent\textbf{Д}\ \ лише «Каштанчики», «Смакс» або «Няммі»
\end{minipage}%
\hfill
\begin{minipage}[t]{0.48\textwidth}
\vspace{0pt}
\centering
\begin{tikzpicture}[scale=0.85, thick]
    % Осі
    \draw (0,0) -- (6.5,0);
    \draw (0,0) -- (0,4.0);

    % Підпис осі Y
    \node[rotate=90, font=\small] at (-1.2, 2.0) {Ціна (грн за пакет)};

    % Сітка та підписи на осі Y
    \foreach \v/\y in {0/0, 50/0.5, 100/1.0, 150/1.5, 200/2.0, 250/2.5, 300/3.0, 350/3.5} {
        \draw[gray!30, thin] (0,\y) -- (6.5,\y);
        \node[left, font=\scriptsize] at (0,\y) {\v};
    }

    % Стовпчики
    \fill[yearOrange!80!yellow] (0.4,0) rectangle (1.2, 1.75);
    \fill[mainGreen] (1.6,0) rectangle (2.4, 2.5);
    \fill[brown!80!black] (2.8,0) rectangle (3.6, 3.25);
    \fill[cyan!90!blue] (4.0,0) rectangle (4.8, 1.9);
    \fill[magenta] (5.2,0) rectangle (6.0, 2.0);

    % Підписи на осі X
    \node[below, rotate=50, anchor=north east, font=\scriptsize, inner sep=2pt] at (1.0, -0.1) {«Чудасія»};
    \node[below, rotate=50, anchor=north east, font=\scriptsize, inner sep=2pt] at (2.2, -0.1) {«Каштанчики»};
    \node[below, rotate=50, anchor=north east, font=\scriptsize, inner sep=2pt] at (3.4, -0.1) {«Шокос»};
    \node[below, rotate=50, anchor=north east, font=\scriptsize, inner sep=2pt] at (4.6, -0.1) {«Смакс»};
    \node[below, rotate=50, anchor=north east, font=\scriptsize, inner sep=2pt] at (5.8, -0.1) {«Няммі»};
\end{tikzpicture}
\end{minipage}
\par\vspace{0.4cm}
\solution{За діаграмою ціни: «Чудасія» $175$, «Каштанчики» $250$, «Шокос» $325$, «Смакс» $190$, «Няммі» $200$~грн. На $200$~грн вистачає лише на ті, що коштують не більше $200$: «Чудасія», «Смакс» або «Няммі». Відповідь: \textbf{В}.}
\noindent
\begin{minipage}[t]{0.52\textwidth}
\vspace{0pt}
\zadnum На діаграмі відображено інформацію про ціну зошитів «Альфа», «Бета», «Гама», «Дельта» і «Омега» в деякій крамниці. Якщо в Марійки є лише $150$~грн, то із зазначених зошитів вона може придбати

\vspace{0.3cm}
\par\noindent\textbf{А}\ \ лише «Бета» або «Дельта»
\par\noindent\textbf{Б}\ \ лише «Гама»
\par\noindent\textbf{В}\ \ будь-який один зошит
\par\noindent\textbf{Г}\ \ лише «Альфа», «Гама» або «Омега»
\par\noindent\textbf{Д}\ \ лише «Альфа», «Дельта» або «Омега»
\end{minipage}%
\hfill
\begin{minipage}[t]{0.48\textwidth}
\vspace{0pt}
\centering
\begin{tikzpicture}[scale=0.85, thick]
    % Осі
    \draw (0,0) -- (6.5,0);
    \draw (0,0) -- (0,4.0);

    % Підпис осі Y
    \node[rotate=90, font=\small] at (-1.2, 2.0) {Ціна (грн за зошит)};

    % Сітка та підписи на осі Y (масштаб: 50 грн = 1 одиниця)
    \foreach \v/\y in {0/0, 25/0.5, 50/1.0, 75/1.5, 100/2.0, 125/2.5, 150/3.0, 175/3.5} {
        \draw[gray!30, thin] (0,\y) -- (6.5,\y);
        \node[left, font=\scriptsize] at (0,\y) {\v};
    }

    % Стовпчики: Альфа 120, Бета 170, Гама 90, Дельта 155, Омега 140
    \fill[yearOrange!80!yellow] (0.4,0) rectangle (1.2, 2.4);
    \fill[mainGreen] (1.6,0) rectangle (2.4, 3.4);
    \fill[brown!80!black] (2.8,0) rectangle (3.6, 1.8);
    \fill[cyan!90!blue] (4.0,0) rectangle (4.8, 3.1);
    \fill[magenta] (5.2,0) rectangle (6.0, 2.8);

    % Підписи на осі X
    \node[below, rotate=50, anchor=north east, font=\scriptsize, inner sep=2pt] at (1.0, -0.1) {«Альфа»};
    \node[below, rotate=50, anchor=north east, font=\scriptsize, inner sep=2pt] at (2.2, -0.1) {«Бета»};
    \node[below, rotate=50, anchor=north east, font=\scriptsize, inner sep=2pt] at (3.4, -0.1) {«Гама»};
    \node[below, rotate=50, anchor=north east, font=\scriptsize, inner sep=2pt] at (4.6, -0.1) {«Дельта»};
    \node[below, rotate=50, anchor=north east, font=\scriptsize, inner sep=2pt] at (5.8, -0.1) {«Омега»};
\end{tikzpicture}
\end{minipage}
\par\vspace{0.4cm}
\noindent
\begin{minipage}[t]{0.52\textwidth}
\vspace{0pt}
\zadnum На діаграмі відображено інформацію про ціну пляшки соку «Яблучний», «Вишневий», «Томатний», «Апельсиновий» і «Мультифрукт» у деякому магазині. Тарас купив по одній пляшці найдешевшого та найдорожчого із зазначених соків. Скільки гривень він заплатив за дві пляшки?

\vspace{0.3cm}
\par\noindent\textbf{А}\ \ $105$~грн
\par\noindent\textbf{Б}\ \ $110$~грн
\par\noindent\textbf{В}\ \ $115$~грн
\par\noindent\textbf{Г}\ \ $125$~грн
\par\noindent\textbf{Д}\ \ $135$~грн
\end{minipage}%
\hfill
\begin{minipage}[t]{0.48\textwidth}
\vspace{0pt}
\centering
\begin{tikzpicture}[scale=0.85, thick]
    % Осі
    \draw (0,0) -- (6.5,0);
    \draw (0,0) -- (0,4.0);

    % Підпис осі Y
    \node[rotate=90, font=\small] at (-1.2, 2.0) {Ціна (грн за пляшку)};

    % Сітка та підписи на осі Y (масштаб: 20 грн = 1 одиниця)
    \foreach \v/\y in {0/0, 10/0.5, 20/1.0, 30/1.5, 40/2.0, 50/2.5, 60/3.0, 70/3.5} {
        \draw[gray!30, thin] (0,\y) -- (6.5,\y);
        \node[left, font=\scriptsize] at (0,\y) {\v};
    }

    % Стовпчики: Яблучний 45, Вишневий 70, Томатний 55, Апельсиновий 65, Мультифрукт 60
    \fill[yearOrange!80!yellow] (0.4,0) rectangle (1.2, 2.25);
    \fill[mainGreen] (1.6,0) rectangle (2.4, 3.5);
    \fill[brown!80!black] (2.8,0) rectangle (3.6, 2.75);
    \fill[cyan!90!blue] (4.0,0) rectangle (4.8, 3.25);
    \fill[magenta] (5.2,0) rectangle (6.0, 3.0);

    % Підписи на осі X
    \node[below, rotate=50, anchor=north east, font=\scriptsize, inner sep=2pt] at (1.0, -0.1) {«Яблучний»};
    \node[below, rotate=50, anchor=north east, font=\scriptsize, inner sep=2pt] at (2.2, -0.1) {«Вишневий»};
    \node[below, rotate=50, anchor=north east, font=\scriptsize, inner sep=2pt] at (3.4, -0.1) {«Томатний»};
    \node[below, rotate=50, anchor=north east, font=\scriptsize, inner sep=2pt] at (4.6, -0.1) {«Апельсиновий»};
    \node[below, rotate=50, anchor=north east, font=\scriptsize, inner sep=2pt] at (5.8, -0.1) {«Мультифрукт»};
\end{tikzpicture}
\end{minipage}
\par\vspace{0.4cm}
\taskBlock{Комбінаторика (перестановки) --- сесія 4.06, завдання №11}
\zadtask{Ді-джей складає трек-лист із $6$ різних музичних композицій. Кожен трек має прозвучати лише один раз. Укажіть вираз, за допомогою якого можна визначити загальну кількість можливих послідовностей звучання цих треків.

\vspace{0.2cm}
\par\noindent\textbf{А}\ \ $6 + 5 + 4 + 3 + 2 + 1$
\par\noindent\textbf{Б}\ \ $6$
\par\noindent\textbf{В}\ \ $36$
\par\noindent\textbf{Г}\ \ $6!$
\par\noindent\textbf{Д}\ \ $6^6$
\vspace{0.4cm}}
\solution{Кількість послідовностей із $6$ різних треків (кожен по разу)~--- це число перестановок $P_6=6!$. Відповідь: \textbf{Г}.}
\zadtask{Бібліотекарка розставляє на полиці $7$ різних книжок в один ряд. Укажіть вираз, за допомогою якого можна визначити загальну кількість можливих способів розташування цих книжок.

\vspace{0.2cm}
\par\noindent\textbf{А}\ \ $7 + 6 + 5 + 4 + 3 + 2 + 1$
\par\noindent\textbf{Б}\ \ $7^7$
\par\noindent\textbf{В}\ \ $49$
\par\noindent\textbf{Г}\ \ $7$
\par\noindent\textbf{Д}\ \ $7!$
\vspace{0.4cm}}
\zadtask{До фіналу конкурсу пройшли $5$ танцювальних пар. Журі визначає, які $3$ з них виступлять у гала-концерті та в якому порядку (по черзі). Укажіть вираз, за допомогою якого можна визначити загальну кількість можливих варіантів такого виступу.

\vspace{0.2cm}
\par\noindent\textbf{А}\ \ $5 \cdot 4 \cdot 3$
\par\noindent\textbf{Б}\ \ $5!$
\par\noindent\textbf{В}\ \ $5^3$
\par\noindent\textbf{Г}\ \ $5 + 4 + 3$
\par\noindent\textbf{Д}\ \ $5^2$
\vspace{0.4cm}}
\taskBlock{Читання діаграми (результати опитування) --- сесія 5.06, завдання №2}
\noindent
\begin{minipage}[c]{0.54\textwidth}
\zadnum На рисунку відображено результати опитування пасажирів щодо якості перевезень (у балах за шкалою $1$--$10$). Визначте кількість опитаних, які поставили оцінку, \textit{вищу за} $8$ балів.
\end{minipage}%
\hfill
\begin{minipage}[c]{0.44\textwidth}
\centering
\begin{tikzpicture}[scale=0.42, thick]
    \draw[->] (0,0) -- (0,9) node[above, font=\scriptsize] {};
    \draw[->] (0,0) -- (7.5,0) node[right, font=\scriptsize] {бали};
    \foreach \v/\yy in {25/1,50/2,75/3,100/4,125/5,150/6,175/7,200/8,225/9} {
        \draw[gray!25] (0,\yy) -- (7,\yy);
    }
    \foreach \v/\yy in {50/2,100/4,150/6,200/8} {\node[left,font=\tiny] at (0,\yy){\v};}
    % точки: 4->75(3), 5->100(4), 6->100(4), 7->200(8), 8->175(7), 9->75(3), 10->25(1)
    \draw[mainGreen,thick] (1,3)--(2,4)--(3,4)--(4,8)--(5,7)--(6,3)--(7,1);
    \foreach \xx/\yy in {1/3,2/4,3/4,4/8,5/7,6/3,7/1} {\fill (\xx,\yy) circle (2.5pt);}
    \foreach \xx/\lab in {1/4,2/5,3/6,4/7,5/8,6/9,7/10} {\node[below,font=\tiny] at (\xx,0){\lab};}
\end{tikzpicture}
\end{minipage}
\par\vspace{0.2cm}
\answerTable{$19$}{$100$}{$200$}{$275$}{$475$}
\solution{Оцінка вища за $8$ балів~--- це $9$ і $10$ балів. За графіком їх поставили $75$ і $25$ опитаних. Разом $75+25=100$. Відповідь: \textbf{Б}.}
\noindent
\begin{minipage}[c]{0.54\textwidth}
\zadnum На рисунку відображено результати опитування відвідувачів кафе щодо якості обслуговування (у балах за шкалою $1$--$10$). Визначте кількість опитаних, які поставили оцінку, \textit{нижчу за} $6$ балів.
\end{minipage}%
\hfill
\begin{minipage}[c]{0.44\textwidth}
\centering
\begin{tikzpicture}[scale=0.42, thick]
    \draw[->] (0,0) -- (0,9) node[above, font=\scriptsize] {};
    \draw[->] (0,0) -- (7.5,0) node[right, font=\scriptsize] {бали};
    \foreach \v/\yy in {25/1,50/2,75/3,100/4,125/5,150/6,175/7,200/8,225/9} {
        \draw[gray!25] (0,\yy) -- (7,\yy);
    }
    \foreach \v/\yy in {50/2,100/4,150/6,200/8} {\node[left,font=\tiny] at (0,\yy){\v};}
    % точки: 4->50(2), 5->100(4), 6->125(5), 7->175(7), 8->150(6), 9->75(3), 10->50(2)
    \draw[mainGreen,thick] (1,2)--(2,4)--(3,5)--(4,7)--(5,6)--(6,3)--(7,2);
    \foreach \xx/\yy in {1/2,2/4,3/5,4/7,5/6,6/3,7/2} {\fill (\xx,\yy) circle (2.5pt);}
    \foreach \xx/\lab in {1/4,2/5,3/6,4/7,5/8,6/9,7/10} {\node[below,font=\tiny] at (\xx,0){\lab};}
\end{tikzpicture}
\end{minipage}
\par\vspace{0.2cm}
\answerTable{$50$}{$100$}{$150$}{$275$}{$425$}
\noindent
\begin{minipage}[c]{0.54\textwidth}
\zadnum На рисунку відображено результати опитування користувачів застосунку щодо зручності сервісу (у балах за шкалою $1$--$10$). Визначте кількість опитаних, які поставили оцінку \textit{від $6$ до $8$ балів включно}.
\end{minipage}%
\hfill
\begin{minipage}[c]{0.44\textwidth}
\centering
\begin{tikzpicture}[scale=0.42, thick]
    \draw[->] (0,0) -- (0,9) node[above, font=\scriptsize] {};
    \draw[->] (0,0) -- (7.5,0) node[right, font=\scriptsize] {бали};
    \foreach \v/\yy in {25/1,50/2,75/3,100/4,125/5,150/6,175/7,200/8,225/9} {
        \draw[gray!25] (0,\yy) -- (7,\yy);
    }
    \foreach \v/\yy in {50/2,100/4,150/6,200/8} {\node[left,font=\tiny] at (0,\yy){\v};}
    % точки: 4->25(1), 5->75(3), 6->150(6), 7->200(8), 8->125(5), 9->75(3), 10->50(2)
    \draw[mainGreen,thick] (1,1)--(2,3)--(3,6)--(4,8)--(5,5)--(6,3)--(7,2);
    \foreach \xx/\yy in {1/1,2/3,3/6,4/8,5/5,6/3,7/2} {\fill (\xx,\yy) circle (2.5pt);}
    \foreach \xx/\lab in {1/4,2/5,3/6,4/7,5/8,6/9,7/10} {\node[below,font=\tiny] at (\xx,0){\lab};}
\end{tikzpicture}
\end{minipage}
\par\vspace{0.2cm}
\answerTable{$350$}{$475$}{$250$}{$700$}{$200$}
\typeTitle{Завдання з короткою відповіддю}
\taskBlock{Інфографіка: вартість таксі з відсотками --- сесія 2.06, завдання №20}

\zadtask{Перед подорожжю турист проаналізував інфографіку й дізнався вартість (у грн) поїздки від аеропорту до центру міста (або у зворотному напрямку) у деяких світових столицях. На час туристичного сезону вартість проїзду в таксі в Осло збільшується на $20\,\%$, а в Рейк'явіку~--- на $30\,\%$ від фіксованої вартості проїзду. Визначте сумарну кількість коштів, яку витратить турист на таксі від аеропорту до центру міста й назад у цих містах, якщо відвідає Осло та Рейк'явік у туристичний сезон.

\vspace{0.3cm}
\begin{center}
\begin{tikzpicture}[scale=1.0]
    % фон-рамка
    \draw[gray!40, rounded corners=3pt] (-0.3,-0.4) rectangle (9.2,5.6);
    % дані: місто / значення / довжина бара (масштаб: 5561 -> 4.2)
    \def\scale{0.000755}
    \foreach \i/\city/\val in {0/{аеропорт Лондона}/2502, 1/{аеропорт Монако}/2502, 2/{аеропорт Рейк'явіка}/2780, 3/{аеропорт Осло}/2780, 4/{аеропорт Токіо}/5561} {
        \pgfmathsetmacro\y{\i*1.05}
        \pgfmathsetmacro\barlen{\val*\scale}
        % прапорець-заглушка
        \filldraw[fill=gray!25, draw=gray!50] (0,\y+0.32) rectangle (0.45,\y+0.62);
        % назва
        \node[right, font=\small] at (0.55,\y+0.47) {\city};
        % бар
        \filldraw[fill=cyan!45, draw=cyan!60!black] (0.55,\y) rectangle (0.55+\barlen,\y+0.28);
        \node[right, font=\scriptsize\bfseries, text=white] at (0.7,\y+0.14) {\val};
    }
    % легенда
    \filldraw[fill=cyan!45, draw=cyan!60!black] (6.0,4.55) rectangle (6.35,4.85);
    \node[right, font=\scriptsize] at (6.4,4.7) {Фіксована вартість};
    \node[right, font=\scriptsize] at (6.4,4.4) {проїзду, \textit{грн}};
\end{tikzpicture}
\end{center}
\nmtAnswerBox}
\solution{За інфографікою фіксована вартість в Осло й Рейк'явіку по $2780$ грн. У сезон: Осло $+20\,\%$ ($\cdot1{,}2$), Рейк'явік $+30\,\%$ ($\cdot1{,}3$); поїздка туди й назад -- подвоюємо: $2\cdot(2780\cdot1{,}2+2780\cdot1{,}3)=2\cdot2780\cdot2{,}5=13900$ грн.}

\zadtask{Перед подорожжю туристка проаналізувала інфографіку й дізналася вартість (у грн) поїздки від аеропорту до центру міста (або у зворотному напрямку) у деяких світових столицях. На час туристичного сезону вартість проїзду в таксі в Осло збільшується на $25\,\%$, а в Рейк'явіку~--- на $15\,\%$ від фіксованої вартості проїзду. Визначте сумарну кількість коштів, яку витратить туристка на таксі від аеропорту до центру міста й назад у цих містах, якщо відвідає Осло та Рейк'явік у туристичний сезон.

\vspace{0.3cm}
\begin{center}
\begin{tikzpicture}[scale=1.0]
    % фон-рамка
    \draw[gray!40, rounded corners=3pt] (-0.3,-0.4) rectangle (9.2,5.6);
    % дані: місто / значення (масштаб: 5800 -> 4.2)
    \def\scale{0.000724}
    \foreach \i/\city/\val in {0/{аеропорт Лондона}/2400, 1/{аеропорт Парижа}/2400, 2/{аеропорт Рейк'явіка}/3200, 3/{аеропорт Осло}/3200, 4/{аеропорт Токіо}/5800} {
        \pgfmathsetmacro\y{\i*1.05}
        \pgfmathsetmacro\barlen{\val*\scale}
        % прапорець-заглушка
        \filldraw[fill=gray!25, draw=gray!50] (0,\y+0.32) rectangle (0.45,\y+0.62);
        % назва
        \node[right, font=\small] at (0.55,\y+0.47) {\city};
        % бар
        \filldraw[fill=cyan!45, draw=cyan!60!black] (0.55,\y) rectangle (0.55+\barlen,\y+0.28);
        \node[right, font=\scriptsize\bfseries, text=white] at (0.7,\y+0.14) {\val};
    }
    % легенда
    \filldraw[fill=cyan!45, draw=cyan!60!black] (6.0,4.55) rectangle (6.35,4.85);
    \node[right, font=\scriptsize] at (6.4,4.7) {Фіксована вартість};
    \node[right, font=\scriptsize] at (6.4,4.4) {проїзду, \textit{грн}};
\end{tikzpicture}
\end{center}
\nmtAnswerBox}

\zadtask{Перед подорожжю турист проаналізував інфографіку й дізнався вартість (у грн) поїздки від аеропорту до центру міста (або у зворотному напрямку) у деяких світових столицях. На час туристичного сезону вартість проїзду в таксі в Осло збільшується на $40\,\%$, а в Рейк'явіку~--- на $10\,\%$ від фіксованої вартості проїзду. Визначте, на скільки гривень \textit{більше} витратить турист на таксі від аеропорту до центру міста й назад в Осло, ніж у Рейк'явіку, у туристичний сезон.

\vspace{0.3cm}
\begin{center}
\begin{tikzpicture}[scale=1.0]
    % фон-рамка
    \draw[gray!40, rounded corners=3pt] (-0.3,-0.4) rectangle (9.2,5.6);
    % дані: місто / значення (масштаб: 5400 -> 4.2)
    \def\scale{0.000778}
    \foreach \i/\city/\val in {0/{аеропорт Лондона}/2200, 1/{аеропорт Монако}/2200, 2/{аеропорт Рейк'явіка}/3000, 3/{аеропорт Осло}/2500, 4/{аеропорт Токіо}/5400} {
        \pgfmathsetmacro\y{\i*1.05}
        \pgfmathsetmacro\barlen{\val*\scale}
        % прапорець-заглушка
        \filldraw[fill=gray!25, draw=gray!50] (0,\y+0.32) rectangle (0.45,\y+0.62);
        % назва
        \node[right, font=\small] at (0.55,\y+0.47) {\city};
        % бар
        \filldraw[fill=cyan!45, draw=cyan!60!black] (0.55,\y) rectangle (0.55+\barlen,\y+0.28);
        \node[right, font=\scriptsize\bfseries, text=white] at (0.7,\y+0.14) {\val};
    }
    % легенда
    \filldraw[fill=cyan!45, draw=cyan!60!black] (6.0,4.55) rectangle (6.35,4.85);
    \node[right, font=\scriptsize] at (6.4,4.7) {Фіксована вартість};
    \node[right, font=\scriptsize] at (6.4,4.4) {проїзду, \textit{грн}};
\end{tikzpicture}
\end{center}
\nmtAnswerBox}

\chapterTitle{ВІДПОВІДІ}
\noindent\small Відповіді наведено у тому ж порядку, що й завдання (наскрізна нумерація).\par\vspace{0.3cm}
\ansTheme{РОЗДІЛ 1. ЧИСЛА І ВИРАЗИ}
\ansType{Завдання з вибором однієї відповіді (А--Д)}
\begin{multicols}{3}\noindent
\ansitem{А}
\ansitem{Б}
\ansitem{Б}
\ansitem{В}
\ansitem{А}
\ansitem{А}
\ansitem{В}
\ansitem{Б}
\ansitem{Б}
\ansitem{А}
\ansitem{Б}
\ansitem{А}
\ansitem{Б}
\ansitem{А}
\ansitem{Б}
\ansitem{Г}
\ansitem{В}
\ansitem{А}
\ansitem{Г}
\ansitem{В}
\ansitem{В}
\ansitem{Д}
\ansitem{Д}
\ansitem{Д}
\ansitem{А}
\ansitem{Б}
\ansitem{А}
\ansitem{А}
\ansitem{Б}
\ansitem{В}
\ansitem{А}
\ansitem{А}
\ansitem{Г}
\ansitem{В}
\ansitem{Б}
\ansitem{Б}
\ansitem{А}
\ansitem{Б}
\ansitem{Б}
\ansitem{Б}
\ansitem{Б}
\ansitem{А}
\ansitem{А}
\ansitem{В}
\ansitem{В}
\ansitem{Б}
\ansitem{В}
\ansitem{Д}
\ansitem{Г}
\ansitem{В}
\ansitem{Б}
\ansitem{А}
\ansitem{А}
\ansitem{А}
\ansitem{А}
\ansitem{Б}
\ansitem{А}
\end{multicols}
\ansType{Завдання на встановлення відповідності}
\begin{multicols}{3}\noindent
\ansitem{1~--~В, 2~--~Б, 3~--~А}
\ansitem{1~--~Б, 2~--~Г, 3~--~А}
\ansitem{1~--~Б, 2~--~В, 3~--~А}
\ansitem{1~--~В, 2~--~Б, 3~--~Д}
\ansitem{1~--~В, 2~--~А, 3~--~Б}
\ansitem{1~--~В, 2~--~Б, 3~--~Д}
\ansitem{1~--~Г, 2~--~А, 3~--~В}
\ansitem{1~--~В, 2~--~Б, 3~--~А}
\ansitem{1~--~В, 2~--~Б, 3~--~А}
\ansitem{1~--~А, 2~--~Г, 3~--~В}
\ansitem{1~--~Б, 2~--~Г, 3~--~В}
\ansitem{1~--~А, 2~--~Д, 3~--~Г}
\ansitem{1~--~Б, 2~--~Г, 3~--~А}
\ansitem{1~--~А, 2~--~В, 3~--~Г}
\ansitem{1~--~В, 2~--~Б, 3~--~А}
\ansitem{1~--~А, 2~--~Д, 3~--~Б}
\ansitem{1~--~Б, 2~--~В, 3~--~А}
\ansitem{1~--~Б, 2~--~А, 3~--~В}
\ansitem{1~--~В, 2~--~Г, 3~--~Б}
\ansitem{1~--~В, 2~--~Г, 3~--~Б}
\ansitem{1~--~В, 2~--~Б, 3~--~А}
\end{multicols}
\ansType{Завдання з короткою відповіддю}
\begin{multicols}{3}\noindent
\ansitem{$90$}
\ansitem{$40$}
\ansitem{$96$}
\ansitem{$128$}
\ansitem{$162$}
\ansitem{$70$}
\ansitem{$51$}
\ansitem{$45$}
\ansitem{$75$}
\ansitem{$13900$}
\ansitem{$11720$}
\ansitem{$6440$}
\ansitem{$2040$}
\ansitem{$1440$}
\ansitem{$450$}
\ansitem{$72$}
\ansitem{$100$}
\ansitem{$49$}
\ansitem{$440$}
\ansitem{$405$}
\ansitem{$560$}
\end{multicols}
\ansTheme{РОЗДІЛ 2. РІВНЯННЯ І НЕРІВНОСТІ}
\ansType{Завдання з вибором однієї відповіді (А--Д)}
\begin{multicols}{3}\noindent
\ansitem{Б}
\ansitem{А}
\ansitem{А}
\ansitem{Б}
\ansitem{В}
\ansitem{В}
\ansitem{Б}
\ansitem{А}
\ansitem{Г}
\ansitem{А}
\ansitem{В}
\ansitem{Г}
\ansitem{Б}
\ansitem{В}
\ansitem{Б}
\ansitem{А}
\ansitem{Б}
\ansitem{А}
\ansitem{А}
\ansitem{Б}
\ansitem{В}
\ansitem{А}
\ansitem{Б}
\ansitem{Б}
\ansitem{А}
\ansitem{В}
\ansitem{А}
\ansitem{А}
\ansitem{Г}
\ansitem{В}
\ansitem{Б}
\ansitem{В}
\ansitem{А}
\ansitem{В}
\ansitem{Д}
\ansitem{В}
\ansitem{А}
\ansitem{Б}
\ansitem{Б}
\ansitem{Д}
\ansitem{Б}
\ansitem{А}
\ansitem{В}
\ansitem{Б}
\ansitem{Б}
\ansitem{Д}
\ansitem{В}
\ansitem{А}
\ansitem{Г}
\ansitem{А}
\ansitem{Б}
\ansitem{Г}
\ansitem{Б}
\ansitem{А}
\ansitem{А}
\ansitem{В}
\ansitem{А}
\ansitem{А}
\ansitem{Б}
\ansitem{А}
\ansitem{А}
\ansitem{Б}
\ansitem{Б}
\ansitem{В}
\ansitem{Г}
\ansitem{Г}
\end{multicols}
\ansType{Завдання з короткою відповіддю}
\begin{multicols}{3}\noindent
\ansitem{$0{,}25$}
\ansitem{$0{,}125$}
\ansitem{$2$}
\ansitem{$4$}
\ansitem{$10$}
\ansitem{$7$}
\ansitem{$18$}
\ansitem{$23$}
\ansitem{$6$}
\ansitem{$0{,}2$}
\ansitem{$0{,}75$}
\ansitem{$15$}
\ansitem{$4{,}4$}
\ansitem{$3$}
\ansitem{$-12$}
\ansitem{$-5$}
\ansitem{$-3$}
\ansitem{$4$}
\ansitem{$a\in(-39;\,-13)$}
\ansitem{$a\in(-\infty;\,-3)$}
\ansitem{$a\in(-\infty;\,-1)$}
\end{multicols}
\ansTheme{РОЗДІЛ 3. ФУНКЦІЇ}
\ansType{Завдання з вибором однієї відповіді (А--Д)}
\begin{multicols}{3}\noindent
\ansitem{А}
\ansitem{Б}
\ansitem{А}
\ansitem{В}
\ansitem{Б}
\ansitem{А}
\ansitem{Б}
\ansitem{В}
\ansitem{Д}
\ansitem{А}
\ansitem{В}
\ansitem{Д}
\ansitem{Г}
\ansitem{В}
\ansitem{В}
\ansitem{Д}
\ansitem{В}
\ansitem{А}
\ansitem{Д}
\ansitem{Г}
\ansitem{Б}
\ansitem{А}
\ansitem{В}
\ansitem{Б}
\ansitem{В}
\ansitem{В}
\ansitem{А}
\ansitem{Д}
\ansitem{Б}
\ansitem{А}
\ansitem{В}
\ansitem{Б}
\ansitem{Г}
\ansitem{Б}
\ansitem{Г}
\ansitem{В}
\ansitem{А}
\ansitem{Г}
\ansitem{А}
\ansitem{А}
\ansitem{Б}
\ansitem{А}
\end{multicols}
\ansType{Завдання на встановлення відповідності}
\begin{multicols}{3}\noindent
\ansitem{1~--~Д, 2~--~Б, 3~--~Г}
\ansitem{1~--~Д, 2~--~Б, 3~--~Г}
\ansitem{1~--~А, 2~--~В, 3~--~Д}
\ansitem{1~--~В, 2~--~Б, 3~--~Г}
\ansitem{1~--~А, 2~--~В, 3~--~Б}
\ansitem{1~--~А, 2~--~Б, 3~--~Г}
\ansitem{1~--~В, 2~--~Г, 3~--~А}
\ansitem{1~--~В, 2~--~Г, 3~--~А}
\ansitem{1~--~А, 2~--~В, 3~--~Б}
\ansitem{1~--~В, 2~--~Б, 3~--~А}
\ansitem{1~--~Б, 2~--~Д, 3~--~А}
\ansitem{1~--~Г, 2~--~Д, 3~--~В}
\ansitem{1~--~Б, 2~--~А, 3~--~В}
\ansitem{1~--~А, 2~--~В, 3~--~Б}
\ansitem{1~--~В, 2~--~А, 3~--~Б}
\ansitem{1~--~Б, 2~--~В, 3~--~Г}
\ansitem{1~--~Б, 2~--~В, 3~--~А}
\ansitem{1~--~А, 2~--~В, 3~--~Б}
\ansitem{1~--~В, 2~--~Б, 3~--~А}
\ansitem{1~--~В, 2~--~Б, 3~--~А}
\ansitem{1~--~Д, 2~--~А, 3~--~Г}
\end{multicols}
\ansType{Завдання з короткою відповіддю}
\begin{multicols}{3}\noindent
\ansitem{$64$}
\ansitem{$9$}
\ansitem{$11$}
\ansitem{$1{,}5$}
\ansitem{$2$}
\ansitem{$5$}
\ansitem{$1{,}5$}
\ansitem{$3$}
\ansitem{$2$}
\ansitem{$6{,}75$}
\ansitem{$10{,}5$}
\ansitem{$12$}
\ansitem{$5$}
\ansitem{$3$}
\ansitem{$-4$}
\ansitem{$4{,}5$}
\ansitem{$8$}
\ansitem{$13{,}5$}
\ansitem{$-0{,}035$}
\ansitem{$0{,}005$}
\ansitem{$8{,}5$}
\end{multicols}
\ansTheme{РОЗДІЛ 4. ПЛАНІМЕТРІЯ}
\ansType{Завдання з вибором однієї відповіді (А--Д)}
\begin{multicols}{3}\noindent
\ansitem{В}
\ansitem{Б}
\ansitem{Д}
\ansitem{В}
\ansitem{Г}
\ansitem{Д}
\ansitem{Б}
\ansitem{Г}
\ansitem{А}
\ansitem{А}
\ansitem{Б}
\ansitem{Г}
\ansitem{В}
\ansitem{В}
\ansitem{Д}
\ansitem{Г}
\ansitem{В}
\ansitem{Г}
\ansitem{В}
\ansitem{В}
\ansitem{В}
\ansitem{Г}
\ansitem{В}
\ansitem{Б}
\ansitem{Д}
\ansitem{Б}
\ansitem{Б}
\ansitem{Б}
\ansitem{В}
\ansitem{Б}
\ansitem{Г}
\ansitem{Г}
\ansitem{Б}
\ansitem{В}
\ansitem{А}
\ansitem{В}
\ansitem{Г}
\ansitem{В}
\ansitem{А}
\ansitem{Г}
\ansitem{Б}
\ansitem{Б}
\ansitem{В}
\ansitem{Б}
\ansitem{А}
\ansitem{Д}
\ansitem{Б}
\ansitem{Г}
\ansitem{В}
\ansitem{Б}
\ansitem{Б}
\ansitem{А}
\ansitem{Б}
\ansitem{В}
\ansitem{Д}
\ansitem{Г}
\ansitem{А}
\ansitem{В}
\ansitem{Б}
\ansitem{Г}
\ansitem{Д}
\ansitem{Д}
\ansitem{А}
\end{multicols}
\ansType{Завдання на встановлення відповідності}
\begin{multicols}{3}\noindent
\ansitem{1~--~Б, 2~--~А, 3~--~Г}
\ansitem{1~--~Г, 2~--~А, 3~--~Б}
\ansitem{1~--~Б, 2~--~В, 3~--~Г}
\ansitem{1~--~Д, 2~--~А, 3~--~Г}
\ansitem{1~--~Д, 2~--~А, 3~--~Г}
\ansitem{1~--~А, 2~--~Г, 3~--~В}
\ansitem{1~--~Г, 2~--~В, 3~--~А}
\ansitem{1~--~Б, 2~--~В, 3~--~А}
\ansitem{1~--~А, 2~--~Б, 3~--~В}
\ansitem{1~--~Г, 2~--~Б, 3~--~В}
\ansitem{1~--~В, 2~--~Г, 3~--~А}
\ansitem{1~--~Г, 2~--~Д, 3~--~А}
\ansitem{1~--~Г, 2~--~В, 3~--~А}
\ansitem{1~--~В, 2~--~А, 3~--~Б}
\ansitem{1~--~Г, 2~--~Д, 3~--~В}
\ansitem{1~--~Г, 2~--~А, 3~--~Б}
\ansitem{1~--~Г, 2~--~А, 3~--~Б}
\ansitem{1~--~Г, 2~--~Б, 3~--~А}
\ansitem{1~--~В, 2~--~Д, 3~--~А}
\ansitem{1~--~А, 2~--~Б, 3~--~В}
\ansitem{1~--~Б, 2~--~А, 3~--~Д}
\end{multicols}
\ansTheme{РОЗДІЛ 5. СТЕРЕОМЕТРІЯ}
\ansType{Завдання з вибором однієї відповіді (А--Д)}
\begin{multicols}{3}\noindent
\ansitem{Д}
\ansitem{В}
\ansitem{Г}
\ansitem{В}
\ansitem{Г}
\ansitem{В}
\ansitem{Г}
\ansitem{Г}
\ansitem{Д}
\ansitem{Б}
\ansitem{Б}
\ansitem{В}
\ansitem{В}
\ansitem{Б}
\ansitem{В}
\ansitem{В}
\ansitem{Г}
\ansitem{В}
\ansitem{Б}
\ansitem{А}
\ansitem{Г}
\ansitem{Б}
\ansitem{В}
\ansitem{А}
\ansitem{Г}
\ansitem{В}
\ansitem{Б}
\ansitem{Г}
\ansitem{Б}
\ansitem{А}
\ansitem{Г}
\ansitem{В}
\ansitem{Д}
\ansitem{А}
\ansitem{Б}
\ansitem{А}
\ansitem{Д}
\ansitem{Б}
\ansitem{Д}
\ansitem{А}
\ansitem{Б}
\ansitem{Б}
\ansitem{В}
\ansitem{Г}
\ansitem{А}
\end{multicols}
\ansType{Завдання з короткою відповіддю}
\begin{multicols}{3}\noindent
\ansitem{$5120$}
\ansitem{$2160$}
\ansitem{$600$}
\ansitem{$1664$}
\ansitem{$624$}
\ansitem{$840$}
\ansitem{$64$}
\ansitem{$8$}
\ansitem{$108$}
\ansitem{$1{,}5$}
\ansitem{$12$}
\ansitem{$3$}
\ansitem{$1350$}
\ansitem{$400$}
\ansitem{$2500$}
\ansitem{$360$}
\ansitem{$576$}
\ansitem{$64$}
\ansitem{$1632$}
\ansitem{$1344$}
\ansitem{$1800$}
\end{multicols}
\ansTheme{РОЗДІЛ 6. КОМБІНАТОРИКА, ЙМОВІРНІСТЬ, СТАТИСТИКА}
\ansType{Завдання з вибором однієї відповіді (А--Д)}
\begin{multicols}{3}\noindent
\ansitem{Б}
\ansitem{В}
\ansitem{В}
\ansitem{Г}
\ansitem{В}
\ansitem{В}
\ansitem{В}
\ansitem{В}
\ansitem{В}
\ansitem{Г}
\ansitem{Г}
\ansitem{А}
\ansitem{Б}
\ansitem{В}
\ansitem{В}
\ansitem{Д}
\ansitem{Б}
\ansitem{В}
\ansitem{В}
\ansitem{Б}
\ansitem{В}
\ansitem{А}
\ansitem{А}
\ansitem{В}
\ansitem{Б}
\ansitem{Г}
\ansitem{А}
\ansitem{Б}
\ansitem{В}
\ansitem{Б}
\ansitem{В}
\ansitem{Г}
\ansitem{В}
\ansitem{Г}
\ansitem{Д}
\ansitem{А}
\ansitem{Б}
\ansitem{В}
\ansitem{Б}
\end{multicols}
\ansType{Завдання з короткою відповіддю}
\begin{multicols}{3}\noindent
\ansitem{$13900$}
\ansitem{$15360$}
\ansitem{$400$}
\end{multicols}

\end{document}
